%% Generated by Sphinx.
\def\sphinxdocclass{report}
\documentclass[letterpaper,10pt,english]{sphinxmanual}
\ifdefined\pdfpxdimen
   \let\sphinxpxdimen\pdfpxdimen\else\newdimen\sphinxpxdimen
\fi \sphinxpxdimen=.75bp\relax
\ifdefined\pdfimageresolution
    \pdfimageresolution= \numexpr \dimexpr1in\relax/\sphinxpxdimen\relax
\fi
%% let collapsible pdf bookmarks panel have high depth per default
\PassOptionsToPackage{bookmarksdepth=5}{hyperref}

\PassOptionsToPackage{booktabs}{sphinx}
\PassOptionsToPackage{colorrows}{sphinx}

\PassOptionsToPackage{warn}{textcomp}
\usepackage[utf8]{inputenc}
\ifdefined\DeclareUnicodeCharacter
% support both utf8 and utf8x syntaxes
  \ifdefined\DeclareUnicodeCharacterAsOptional
    \def\sphinxDUC#1{\DeclareUnicodeCharacter{"#1}}
  \else
    \let\sphinxDUC\DeclareUnicodeCharacter
  \fi
  \sphinxDUC{00A0}{\nobreakspace}
  \sphinxDUC{2500}{\sphinxunichar{2500}}
  \sphinxDUC{2502}{\sphinxunichar{2502}}
  \sphinxDUC{2514}{\sphinxunichar{2514}}
  \sphinxDUC{251C}{\sphinxunichar{251C}}
  \sphinxDUC{2572}{\textbackslash}
\fi
\usepackage{cmap}
\usepackage[T1]{fontenc}
\usepackage{amsmath,amssymb,amstext}
\usepackage{babel}



\usepackage{tgtermes}
\usepackage{tgheros}
\renewcommand{\ttdefault}{txtt}



\usepackage[Bjarne]{fncychap}
\usepackage{sphinx}

\fvset{fontsize=auto}
\usepackage{geometry}


% Include hyperref last.
\usepackage{hyperref}
% Fix anchor placement for figures with captions.
\usepackage{hypcap}% it must be loaded after hyperref.
% Set up styles of URL: it should be placed after hyperref.
\urlstyle{same}


\usepackage{sphinxmessages}
\setcounter{tocdepth}{4}
\setcounter{secnumdepth}{4}


\title{PyGeoX}
\date{Aug 27, 2025}
\release{0.2.0}
\author{Rafael Cabral}
\newcommand{\sphinxlogo}{\vbox{}}
\renewcommand{\releasename}{Release}
\makeindex
\begin{document}

\ifdefined\shorthandoff
  \ifnum\catcode`\=\string=\active\shorthandoff{=}\fi
  \ifnum\catcode`\"=\active\shorthandoff{"}\fi
\fi

\pagestyle{empty}
\sphinxmaketitle
\pagestyle{plain}
\sphinxtableofcontents
\pagestyle{normal}
\phantomsection\label{\detokenize{index::doc}}

\begin{quote}

\sphinxAtStartPar
\sphinxstylestrong{PyGeoX} is an interactive Python library designed to create, manipulate, visualize, and solve geometric problems symbolically and numerically. Ideal for educational purposes, research, and interactive geometry explorations, PyGeoX simplifies geometric constructions and constraint solving with an intuitive, readable Python syntax.


\bigskip\hrule\bigskip



\chapter{Philosophy and Design}
\label{\detokenize{index:philosophy-and-design}}
\sphinxAtStartPar
PyGeoX follows a object\sphinxhyphen{}oriented approach, where geometry becomes intuitive Python code. It supports creating geometric scenes, adding and relating geometric objects, defining precise constraints, and seamlessly solving complex geometry problems.

\sphinxAtStartPar
The philosophy behind PyGeoX revolves around simplicity, clarity, and flexibility:
\begin{itemize}
\item {} 
\sphinxAtStartPar
\sphinxstylestrong{Simple Syntax:} Clearly defined commands mimic geometric language.

\item {} 
\sphinxAtStartPar
\sphinxstylestrong{Automatic Constraint Management:} Constraints are implicitly added when relationships are defined.

\item {} 
\sphinxAtStartPar
\sphinxstylestrong{Rich Property Access:} Directly access properties like areas, lengths, midpoints, and more.

\end{itemize}


\bigskip\hrule\bigskip



\chapter{Getting Started with PyGeoX}
\label{\detokenize{index:getting-started-with-pygeox}}
\sphinxAtStartPar
The main functionalities are described below. Please click \sphinxcode{\sphinxupquote{Examples}} in the left panel to check several use cases.


\section{Create a Scene}
\label{\detokenize{index:create-a-scene}}
\sphinxAtStartPar
Start your geometry exploration by defining a scene:

\begin{sphinxVerbatim}[commandchars=\\\{\}]
\PYG{k+kn}{from} \PYG{n+nn}{pygeox} \PYG{k+kn}{import} \PYG{n}{GeoScene}

\PYG{n}{scene} \PYG{o}{=} \PYG{n}{GeoScene}\PYG{p}{(}\PYG{l+m+mi}{10}\PYG{p}{)}  \PYG{c+c1}{\PYGZsh{} creates a scene with domain from \PYGZhy{}10 to 10}
\end{sphinxVerbatim}


\section{Add Objects}
\label{\detokenize{index:add-objects}}\begin{itemize}
\item {} 
\sphinxAtStartPar
PyGeoX contains \sphinxstylestrong{37 fundamental geometric objects:} Points, Lines, Segments, Rays, Circles, Arcs, Triangles, Quadrilaterals, Regular Polygons, and more.

\item {} 
\sphinxAtStartPar
Please click \sphinxcode{\sphinxupquote{Add geometric objects}} on the left panel to check all objects that can be added.

\item {} 
\sphinxAtStartPar
PyGeoX makes creating geometric figures simple:

\end{itemize}

\begin{sphinxVerbatim}[commandchars=\\\{\}]
\PYG{n}{A}\PYG{p}{,} \PYG{n}{B}\PYG{p}{,} \PYG{n}{C} \PYG{o}{=} \PYG{n}{scene}\PYG{o}{.}\PYG{n}{add}\PYG{o}{.}\PYG{n}{points}\PYG{p}{(}\PYG{p}{[}\PYG{l+s+s1}{\PYGZsq{}}\PYG{l+s+s1}{A}\PYG{l+s+s1}{\PYGZsq{}}\PYG{p}{,} \PYG{l+s+s1}{\PYGZsq{}}\PYG{l+s+s1}{B}\PYG{l+s+s1}{\PYGZsq{}}\PYG{p}{,} \PYG{l+s+s1}{\PYGZsq{}}\PYG{l+s+s1}{C}\PYG{l+s+s1}{\PYGZsq{}}\PYG{p}{]}\PYG{p}{)}
\PYG{n}{triangle} \PYG{o}{=} \PYG{n}{scene}\PYG{o}{.}\PYG{n}{add}\PYG{o}{.}\PYG{n}{triangle}\PYG{p}{(}\PYG{n}{A}\PYG{p}{,} \PYG{n}{B}\PYG{p}{,} \PYG{n}{C}\PYG{p}{)}
\PYG{n}{circle} \PYG{o}{=} \PYG{n}{scene}\PYG{o}{.}\PYG{n}{add}\PYG{o}{.}\PYG{n}{circle}\PYG{p}{(}\PYG{n}{center}\PYG{o}{=}\PYG{n}{A}\PYG{p}{,} \PYG{n}{radius}\PYG{o}{=}\PYG{l+m+mi}{5}\PYG{p}{)}
\end{sphinxVerbatim}
\begin{itemize}
\item {} 
\sphinxAtStartPar
Each object has \sphinxstylestrong{2\sphinxhyphen{}15 detailed properties} such as area, perimeter, midpoint, centroid, circumcircle, incircle, altitudes, and medians.

\item {} 
\sphinxAtStartPar
Example property access:

\end{itemize}

\begin{sphinxVerbatim}[commandchars=\\\{\}]
\PYG{n}{area} \PYG{o}{=} \PYG{n}{triangle}\PYG{o}{.}\PYG{n}{area}
\PYG{n}{length} \PYG{o}{=} \PYG{n}{line\PYGZus{}AB}\PYG{o}{.}\PYG{n}{length}
\PYG{n}{mid\PYGZus{}segment} \PYG{o}{=} \PYG{n}{triangle}\PYG{o}{.}\PYG{n}{midsegment}
\end{sphinxVerbatim}
\begin{itemize}
\item {} 
\sphinxAtStartPar
Please click \sphinxcode{\sphinxupquote{Basic Objects}} and \sphinxcode{\sphinxupquote{Polygons}} on the left panel to see all object properties.

\end{itemize}


\section{Define Relationships}
\label{\detokenize{index:define-relationships}}\begin{itemize}
\item {} 
\sphinxAtStartPar
PyGeoX contains \sphinxstylestrong{27 geometric relationships:} parallel, tangent, similarity, inside/outside containment, etc.

\item {} 
\sphinxAtStartPar
Please click \sphinxcode{\sphinxupquote{Add relationships}} on the left panel to check all relationships that can be added.

\item {} 
\sphinxAtStartPar
Relationships automatically introduce constraints to \sphinxcode{\sphinxupquote{scene}}:

\end{itemize}

\begin{sphinxVerbatim}[commandchars=\\\{\}]
\PYG{c+c1}{\PYGZsh{} Make line AB parallel to line CD}
\PYG{n}{scene}\PYG{o}{.}\PYG{n}{relate}\PYG{o}{.}\PYG{n}{parallel}\PYG{p}{(}\PYG{n}{line\PYGZus{}AB}\PYG{p}{,} \PYG{n}{line\PYGZus{}CD}\PYG{p}{)}

\PYG{c+c1}{\PYGZsh{} Point P lies on circle}
\PYG{n}{scene}\PYG{o}{.}\PYG{n}{relate}\PYG{o}{.}\PYG{n}{point\PYGZus{}lies\PYGZus{}on}\PYG{p}{(}\PYG{n}{P}\PYG{p}{,} \PYG{n}{circle}\PYG{p}{)}
\end{sphinxVerbatim}
\begin{itemize}
\item {} 
\sphinxAtStartPar
This implicitly adds constraints like:

\end{itemize}

\begin{sphinxVerbatim}[commandchars=\\\{\}]
\PYG{c+c1}{\PYGZsh{} Example internal constraint:}
\PYG{c+c1}{\PYGZsh{} (line\PYGZus{}AB.direction\PYGZus{}vector × line\PYGZus{}CD.direction\PYGZus{}vector) = 0}
\end{sphinxVerbatim}


\section{Add Custom Constraints}
\label{\detokenize{index:add-custom-constraints}}\begin{itemize}
\item {} 
\sphinxAtStartPar
Please click \sphinxcode{\sphinxupquote{Add constraints}} on the left panel to check all constraints that can be added.

\item {} 
\sphinxAtStartPar
Define explicit mathematical constraints:

\end{itemize}

\begin{sphinxVerbatim}[commandchars=\\\{\}]
\PYG{k+kn}{import} \PYG{n+nn}{sympy}
\PYG{n}{scene}\PYG{o}{.}\PYG{n}{constraint}\PYG{o}{.}\PYG{n}{eq}\PYG{p}{(}\PYG{n}{triangle}\PYG{o}{.}\PYG{n}{area}\PYG{p}{,} \PYG{n}{sympy}\PYG{o}{.}\PYG{n}{cos}\PYG{p}{(}\PYG{n}{circle}\PYG{o}{.}\PYG{n}{perimeter}\PYG{p}{)} \PYG{o}{/} \PYG{n}{circle}\PYG{o}{.}\PYG{n}{area}\PYG{p}{)}
\PYG{n}{scene}\PYG{o}{.}\PYG{n}{constraint}\PYG{o}{.}\PYG{n}{gt}\PYG{p}{(}\PYG{n}{line\PYGZus{}AB}\PYG{o}{.}\PYG{n}{length}\PYG{p}{,} \PYG{n}{circle}\PYG{o}{.}\PYG{n}{radius}\PYG{p}{)}
\end{sphinxVerbatim}


\bigskip\hrule\bigskip



\chapter{Solving Geometric Problems}
\label{\detokenize{index:solving-geometric-problems}}\begin{itemize}
\item {} 
\sphinxAtStartPar
PyGeoX provides both analytical (symbolic) and numerical solvers.

\end{itemize}

\begin{sphinxVerbatim}[commandchars=\\\{\}]
\PYG{n}{scene}\PYG{o}{.}\PYG{n}{solver}\PYG{o}{.}\PYG{n}{analytical}\PYG{p}{(}\PYG{p}{)}
\PYG{n}{scene}\PYG{o}{.}\PYG{n}{solver}\PYG{o}{.}\PYG{n}{numerical}\PYG{p}{(}\PYG{n}{method}\PYG{o}{=}\PYG{l+s+s2}{\PYGZdq{}}\PYG{l+s+s2}{basinhopping}\PYG{l+s+s2}{\PYGZdq{}}\PYG{p}{)}
\end{sphinxVerbatim}
\begin{itemize}
\item {} 
\sphinxAtStartPar
The numerical solver avoid degenerate solutions (e.g., points too close, radius too small) by adding extra constraints.

\item {} 
\sphinxAtStartPar
Please click \sphinxcode{\sphinxupquote{Solvers}} on the left panel to see all available options.

\end{itemize}


\section{Visualize Results}
\label{\detokenize{index:visualize-results}}
\sphinxAtStartPar
Easily visualize solutions with built\sphinxhyphen{}in plotting:

\begin{sphinxVerbatim}[commandchars=\\\{\}]
\PYG{n}{scene}\PYG{o}{.}\PYG{n}{plot}\PYG{p}{(}\PYG{p}{)}
\end{sphinxVerbatim}


\bigskip\hrule\bigskip



\chapter{Under Construction}
\label{\detokenize{index:under-construction}}

\section{Animation and Interactive Explorations}
\label{\detokenize{index:animation-and-interactive-explorations}}
\sphinxAtStartPar
Future PyGeoX versions will support animations and dynamic geometry, such as moving points:

\begin{sphinxVerbatim}[commandchars=\\\{\}]
\PYG{n}{anim\PYGZus{}lambda} \PYG{o}{=} \PYG{n}{scene}\PYG{o}{.}\PYG{n}{animate}\PYG{o}{.}\PYG{n}{add\PYGZus{}parameter}\PYG{p}{(}\PYG{l+s+s2}{\PYGZdq{}}\PYG{l+s+s2}{lambda}\PYG{l+s+s2}{\PYGZdq{}}\PYG{p}{)}
\PYG{n}{scene}\PYG{o}{.}\PYG{n}{animate}\PYG{o}{.}\PYG{n}{run}\PYG{p}{(}\PYG{n}{frames}\PYG{o}{=}\PYG{l+m+mi}{50}\PYG{p}{)}
\PYG{n}{scene}\PYG{o}{.}\PYG{n}{animate}\PYG{o}{.}\PYG{n}{gif}\PYG{p}{(}\PYG{l+s+s2}{\PYGZdq{}}\PYG{l+s+s2}{animation.gif}\PYG{l+s+s2}{\PYGZdq{}}\PYG{p}{)}
\end{sphinxVerbatim}


\section{Automated Geometry Proofs}
\label{\detokenize{index:automated-geometry-proofs}}
\sphinxAtStartPar
PyGeoX is developing an integrated proving system, enabling automated geometric proofs:

\begin{sphinxVerbatim}[commandchars=\\\{\}]
\PYG{k}{with} \PYG{n}{scene}\PYG{o}{.}\PYG{n}{proving}\PYG{p}{(}\PYG{p}{)} \PYG{k}{as} \PYG{n}{prover}\PYG{p}{:}
    \PYG{n}{scene}\PYG{o}{.}\PYG{n}{relate}\PYG{o}{.}\PYG{n}{collinear}\PYG{p}{(}\PYG{n}{A}\PYG{p}{,} \PYG{n}{B}\PYG{p}{,} \PYG{n}{C}\PYG{p}{)}

\PYG{n}{scene}\PYG{o}{.}\PYG{n}{prove}\PYG{p}{(}\PYG{p}{)}
\end{sphinxVerbatim}

\sphinxstepscope


\chapter{Examples}
\label{\detokenize{examples:examples}}\label{\detokenize{examples::doc}}

\section{Diagram construction + calculation problem}
\label{\detokenize{examples:diagram-construction-calculation-problem}}
\sphinxAtStartPar
Problem: In a regular hexagon ABCDEF, a circle O passing through points E and F is tangent to sides AB and CD at points G and H, respectively, and intersects side DE at point M. Lines GM and FH intersect at point N, the measure of angle GNF is \_\_\_°.

\begin{sphinxVerbatim}[commandchars=\\\{\}]
\PYG{k+kn}{from} \PYG{n+nn}{pygeox} \PYG{k+kn}{import} \PYG{n}{GeoScene}\PYG{p}{,} \PYG{n}{LineSegment}

\PYG{n}{scene} \PYG{o}{=} \PYG{n}{GeoScene}\PYG{p}{(}\PYG{l+m+mi}{10}\PYG{p}{)}

\PYG{n}{A}\PYG{p}{,} \PYG{n}{B}\PYG{p}{,} \PYG{n}{C}\PYG{p}{,} \PYG{n}{D}\PYG{p}{,} \PYG{n}{E}\PYG{p}{,} \PYG{n}{F}\PYG{p}{,} \PYG{n}{G}\PYG{p}{,} \PYG{n}{H}\PYG{p}{,} \PYG{n}{M}\PYG{p}{,} \PYG{n}{N}\PYG{p}{,} \PYG{n}{O} \PYG{o}{=} \PYG{n}{scene}\PYG{o}{.}\PYG{n}{add}\PYG{o}{.}\PYG{n}{points}\PYG{p}{(}\PYG{p}{[}\PYG{l+s+s2}{\PYGZdq{}}\PYG{l+s+s2}{A}\PYG{l+s+s2}{\PYGZdq{}}\PYG{p}{,} \PYG{l+s+s2}{\PYGZdq{}}\PYG{l+s+s2}{B}\PYG{l+s+s2}{\PYGZdq{}}\PYG{p}{,} \PYG{l+s+s2}{\PYGZdq{}}\PYG{l+s+s2}{C}\PYG{l+s+s2}{\PYGZdq{}}\PYG{p}{,} \PYG{l+s+s2}{\PYGZdq{}}\PYG{l+s+s2}{D}\PYG{l+s+s2}{\PYGZdq{}}\PYG{p}{,} \PYG{l+s+s2}{\PYGZdq{}}\PYG{l+s+s2}{E}\PYG{l+s+s2}{\PYGZdq{}}\PYG{p}{,} \PYG{l+s+s2}{\PYGZdq{}}\PYG{l+s+s2}{F}\PYG{l+s+s2}{\PYGZdq{}}\PYG{p}{,} \PYG{l+s+s2}{\PYGZdq{}}\PYG{l+s+s2}{G}\PYG{l+s+s2}{\PYGZdq{}}\PYG{p}{,} \PYG{l+s+s2}{\PYGZdq{}}\PYG{l+s+s2}{H}\PYG{l+s+s2}{\PYGZdq{}}\PYG{p}{,} \PYG{l+s+s2}{\PYGZdq{}}\PYG{l+s+s2}{M}\PYG{l+s+s2}{\PYGZdq{}}\PYG{p}{,} \PYG{l+s+s2}{\PYGZdq{}}\PYG{l+s+s2}{N}\PYG{l+s+s2}{\PYGZdq{}}\PYG{p}{,} \PYG{l+s+s2}{\PYGZdq{}}\PYG{l+s+s2}{O}\PYG{l+s+s2}{\PYGZdq{}}\PYG{p}{]}\PYG{p}{)}

\PYG{c+c1}{\PYGZsh{} In a regular hexagon ABCDEF}
\PYG{n}{hexagon} \PYG{o}{=} \PYG{n}{scene}\PYG{o}{.}\PYG{n}{add}\PYG{o}{.}\PYG{n}{regular\PYGZus{}hexagon}\PYG{p}{(}\PYG{n}{A}\PYG{p}{,} \PYG{n}{B}\PYG{p}{,} \PYG{n}{C}\PYG{p}{,} \PYG{n}{D}\PYG{p}{,} \PYG{n}{E}\PYG{p}{,} \PYG{n}{F}\PYG{p}{)}

\PYG{c+c1}{\PYGZsh{} a circle O passing through points E and F}
\PYG{n}{circle} \PYG{o}{=} \PYG{n}{scene}\PYG{o}{.}\PYG{n}{add}\PYG{o}{.}\PYG{n}{circle}\PYG{p}{(}\PYG{n}{O}\PYG{p}{)}
\PYG{n}{scene}\PYG{o}{.}\PYG{n}{relate}\PYG{o}{.}\PYG{n}{points\PYGZus{}lie\PYGZus{}on}\PYG{p}{(}\PYG{p}{[}\PYG{n}{E}\PYG{p}{,} \PYG{n}{F}\PYG{p}{]}\PYG{p}{,} \PYG{n}{circle}\PYG{p}{)}

\PYG{c+c1}{\PYGZsh{} tangent to sides AB and CD at points G and H}
\PYG{n}{scene}\PYG{o}{.}\PYG{n}{relate}\PYG{o}{.}\PYG{n}{tangent\PYGZus{}to\PYGZus{}circle}\PYG{p}{(}\PYG{n}{LineSegment}\PYG{p}{(}\PYG{n}{A}\PYG{p}{,}\PYG{n}{B}\PYG{p}{)}\PYG{p}{,} \PYG{n}{circle}\PYG{p}{,} \PYG{n}{G}\PYG{p}{)}
\PYG{n}{scene}\PYG{o}{.}\PYG{n}{relate}\PYG{o}{.}\PYG{n}{tangent\PYGZus{}to\PYGZus{}circle}\PYG{p}{(}\PYG{n}{LineSegment}\PYG{p}{(}\PYG{n}{C}\PYG{p}{,}\PYG{n}{D}\PYG{p}{)}\PYG{p}{,} \PYG{n}{circle}\PYG{p}{,} \PYG{n}{H}\PYG{p}{)}

\PYG{c+c1}{\PYGZsh{}\PYGZsh{} intersects side DE at point M}
\PYG{n}{scene}\PYG{o}{.}\PYG{n}{relate}\PYG{o}{.}\PYG{n}{line\PYGZus{}intersects\PYGZus{}circle\PYGZus{}at}\PYG{p}{(}\PYG{n}{LineSegment}\PYG{p}{(}\PYG{n}{D}\PYG{p}{,}\PYG{n}{E}\PYG{p}{)}\PYG{p}{,} \PYG{n}{circle}\PYG{p}{,} \PYG{n}{M}\PYG{p}{)}

\PYG{c+c1}{\PYGZsh{}\PYGZsh{} Lines GM and FH intersect at point N}
\PYG{n}{Line\PYGZus{}GM} \PYG{o}{=} \PYG{n}{scene}\PYG{o}{.}\PYG{n}{add}\PYG{o}{.}\PYG{n}{line\PYGZus{}segment}\PYG{p}{(}\PYG{n}{G}\PYG{p}{,} \PYG{n}{M}\PYG{p}{)}
\PYG{n}{Line\PYGZus{}FH} \PYG{o}{=} \PYG{n}{scene}\PYG{o}{.}\PYG{n}{add}\PYG{o}{.}\PYG{n}{line\PYGZus{}segment}\PYG{p}{(}\PYG{n}{F}\PYG{p}{,} \PYG{n}{H}\PYG{p}{)}
\PYG{n}{scene}\PYG{o}{.}\PYG{n}{relate}\PYG{o}{.}\PYG{n}{lines\PYGZus{}intersect\PYGZus{}at}\PYG{p}{(}\PYG{n}{Line\PYGZus{}GM}\PYG{p}{,} \PYG{n}{Line\PYGZus{}FH}\PYG{p}{,} \PYG{n}{N}\PYG{p}{)}

\PYG{n}{scene}\PYG{o}{.}\PYG{n}{solver}\PYG{o}{.}\PYG{n}{numerical}\PYG{p}{(}\PYG{p}{)}
\PYG{n}{scene}\PYG{o}{.}\PYG{n}{plot}\PYG{p}{(}\PYG{p}{)}
\end{sphinxVerbatim}

\sphinxAtStartPar
\sphinxincludegraphics{{img1}.png}

\sphinxAtStartPar
The function \sphinxcode{\sphinxupquote{scene.solver.numerical()}} finds one of possibly many valid configurations for the diagram. After, we can call any object property to find its numerical value.

\begin{sphinxVerbatim}[commandchars=\\\{\}]
\PYG{c+c1}{\PYGZsh{} the measure of angle GNF is \PYGZus{}\PYGZus{}\PYGZus{}°}
\PYG{n}{scene}\PYG{o}{.}\PYG{n}{angle}\PYG{p}{(}\PYG{n}{G}\PYG{p}{,}\PYG{n}{N}\PYG{p}{,}\PYG{n}{F}\PYG{p}{)}\PYG{o}{.}\PYG{n}{evalf}\PYG{p}{(}\PYG{l+m+mi}{4}\PYG{p}{)} 
\PYG{c+c1}{\PYGZsh{}60}
\end{sphinxVerbatim}


\bigskip\hrule\bigskip



\section{Diagram construction + proof problem (Butterfly theorem)}
\label{\detokenize{examples:diagram-construction-proof-problem-butterfly-theorem}}
\sphinxAtStartPar
Problem: Let M be the midpoint of a chord PQ of a circle, through which two other chords AB and CD are drawn; AD and BC intersect chord PQ at X and Y correspondingly. Prove that M is the midpoint of XY.

\begin{sphinxVerbatim}[commandchars=\\\{\}]
\PYG{k+kn}{from} \PYG{n+nn}{pygeox} \PYG{k+kn}{import} \PYG{n}{GeoScene}\PYG{p}{,} \PYG{n}{LineSegment}


\PYG{n}{scene} \PYG{o}{=} \PYG{n}{GeoScene}\PYG{p}{(}\PYG{l+m+mi}{5}\PYG{p}{)}

\PYG{n}{M}\PYG{p}{,} \PYG{n}{P}\PYG{p}{,} \PYG{n}{Q}\PYG{p}{,} \PYG{n}{A}\PYG{p}{,} \PYG{n}{B}\PYG{p}{,} \PYG{n}{C}\PYG{p}{,} \PYG{n}{D}\PYG{p}{,} \PYG{n}{X}\PYG{p}{,} \PYG{n}{Y}\PYG{p}{,} \PYG{n}{O} \PYG{o}{=} \PYG{n}{scene}\PYG{o}{.}\PYG{n}{add}\PYG{o}{.}\PYG{n}{points}\PYG{p}{(}\PYG{p}{[}\PYG{l+s+s2}{\PYGZdq{}}\PYG{l+s+s2}{M}\PYG{l+s+s2}{\PYGZdq{}}\PYG{p}{,}\PYG{l+s+s2}{\PYGZdq{}}\PYG{l+s+s2}{P}\PYG{l+s+s2}{\PYGZdq{}}\PYG{p}{,}\PYG{l+s+s2}{\PYGZdq{}}\PYG{l+s+s2}{Q}\PYG{l+s+s2}{\PYGZdq{}}\PYG{p}{,}\PYG{l+s+s2}{\PYGZdq{}}\PYG{l+s+s2}{A}\PYG{l+s+s2}{\PYGZdq{}}\PYG{p}{,}\PYG{l+s+s2}{\PYGZdq{}}\PYG{l+s+s2}{B}\PYG{l+s+s2}{\PYGZdq{}}\PYG{p}{,}\PYG{l+s+s2}{\PYGZdq{}}\PYG{l+s+s2}{C}\PYG{l+s+s2}{\PYGZdq{}}\PYG{p}{,}\PYG{l+s+s2}{\PYGZdq{}}\PYG{l+s+s2}{D}\PYG{l+s+s2}{\PYGZdq{}}\PYG{p}{,}\PYG{l+s+s2}{\PYGZdq{}}\PYG{l+s+s2}{X}\PYG{l+s+s2}{\PYGZdq{}}\PYG{p}{,}\PYG{l+s+s2}{\PYGZdq{}}\PYG{l+s+s2}{Y}\PYG{l+s+s2}{\PYGZdq{}}\PYG{p}{,}\PYG{l+s+s2}{\PYGZdq{}}\PYG{l+s+s2}{O}\PYG{l+s+s2}{\PYGZdq{}}\PYG{p}{]}\PYG{p}{)}

\PYG{c+c1}{\PYGZsh{} Let M be the midpoint of a chord PQ of a circle, through which two other chords AB and CD are drawn}
\PYG{n}{circle} \PYG{o}{=} \PYG{n}{scene}\PYG{o}{.}\PYG{n}{add}\PYG{o}{.}\PYG{n}{circle}\PYG{p}{(}\PYG{n}{O}\PYG{p}{)}
\PYG{n}{l\PYGZus{}AB} \PYG{o}{=} \PYG{n}{scene}\PYG{o}{.}\PYG{n}{add}\PYG{o}{.}\PYG{n}{chord}\PYG{p}{(}\PYG{n}{circle}\PYG{p}{,} \PYG{n}{A}\PYG{p}{,} \PYG{n}{B}\PYG{p}{)}
\PYG{n}{l\PYGZus{}CD} \PYG{o}{=} \PYG{n}{scene}\PYG{o}{.}\PYG{n}{add}\PYG{o}{.}\PYG{n}{chord}\PYG{p}{(}\PYG{n}{circle}\PYG{p}{,} \PYG{n}{C}\PYG{p}{,} \PYG{n}{D}\PYG{p}{)}
\PYG{n}{l\PYGZus{}PQ} \PYG{o}{=} \PYG{n}{scene}\PYG{o}{.}\PYG{n}{add}\PYG{o}{.}\PYG{n}{chord}\PYG{p}{(}\PYG{n}{circle}\PYG{p}{,} \PYG{n}{P}\PYG{p}{,} \PYG{n}{Q}\PYG{p}{)}
\PYG{n}{scene}\PYG{o}{.}\PYG{n}{relate}\PYG{o}{.}\PYG{n}{is\PYGZus{}midpoint}\PYG{p}{(}\PYG{n}{M}\PYG{p}{,} \PYG{n}{l\PYGZus{}PQ}\PYG{p}{)}
\PYG{n}{scene}\PYG{o}{.}\PYG{n}{relate}\PYG{o}{.}\PYG{n}{point\PYGZus{}lies\PYGZus{}on}\PYG{p}{(}\PYG{n}{M}\PYG{p}{,} \PYG{n}{l\PYGZus{}CD}\PYG{p}{)}
\PYG{n}{scene}\PYG{o}{.}\PYG{n}{relate}\PYG{o}{.}\PYG{n}{point\PYGZus{}lies\PYGZus{}on}\PYG{p}{(}\PYG{n}{M}\PYG{p}{,} \PYG{n}{l\PYGZus{}AB}\PYG{p}{)}

\PYG{c+c1}{\PYGZsh{} AD and BC intersect chord PQ at X and Y correspondingly}
\PYG{n}{l\PYGZus{}AD} \PYG{o}{=} \PYG{n}{scene}\PYG{o}{.}\PYG{n}{add}\PYG{o}{.}\PYG{n}{line\PYGZus{}segment}\PYG{p}{(}\PYG{n}{A}\PYG{p}{,}\PYG{n}{D}\PYG{p}{)}
\PYG{n}{l\PYGZus{}BC} \PYG{o}{=} \PYG{n}{scene}\PYG{o}{.}\PYG{n}{add}\PYG{o}{.}\PYG{n}{line\PYGZus{}segment}\PYG{p}{(}\PYG{n}{B}\PYG{p}{,}\PYG{n}{C}\PYG{p}{)}
\PYG{n}{scene}\PYG{o}{.}\PYG{n}{relate}\PYG{o}{.}\PYG{n}{lines\PYGZus{}intersect\PYGZus{}at}\PYG{p}{(}\PYG{n}{l\PYGZus{}AD}\PYG{p}{,} \PYG{n}{l\PYGZus{}PQ}\PYG{p}{,} \PYG{n}{X}\PYG{p}{)}
\PYG{n}{scene}\PYG{o}{.}\PYG{n}{relate}\PYG{o}{.}\PYG{n}{lines\PYGZus{}intersect\PYGZus{}at}\PYG{p}{(}\PYG{n}{l\PYGZus{}BC}\PYG{p}{,} \PYG{n}{l\PYGZus{}PQ}\PYG{p}{,} \PYG{n}{Y}\PYG{p}{)}

\PYG{c+c1}{\PYGZsh{} Then M is the midpoint of XY.}
\PYG{k}{with} \PYG{n}{scene}\PYG{o}{.}\PYG{n}{proving}\PYG{p}{(}\PYG{p}{)}\PYG{p}{:}
    \PYG{n}{scene}\PYG{o}{.}\PYG{n}{relate}\PYG{o}{.}\PYG{n}{is\PYGZus{}midpoint}\PYG{p}{(}\PYG{n}{M}\PYG{p}{,} \PYG{n}{LineSegment}\PYG{p}{(}\PYG{n}{X}\PYG{p}{,}\PYG{n}{Y}\PYG{p}{)}\PYG{p}{)}

\PYG{n}{scene}\PYG{o}{.}\PYG{n}{solver}\PYG{o}{.}\PYG{n}{numerical}\PYG{p}{(}\PYG{n}{distance\PYGZus{}penalty}\PYG{o}{=}\PYG{l+m+mf}{0.01}\PYG{p}{)}
\PYG{n}{scene}\PYG{o}{.}\PYG{n}{plot}\PYG{p}{(}\PYG{p}{)}
\end{sphinxVerbatim}

\sphinxAtStartPar
\sphinxincludegraphics{{img2}.png}

\sphinxAtStartPar
Our built\sphinxhyphen{}in prover can prove “M is the midpoint of XY” based on the stated conditions.

\begin{sphinxVerbatim}[commandchars=\\\{\}]
\PYG{n}{scene}\PYG{o}{.}\PYG{n}{prove}\PYG{p}{(}\PYG{p}{)}
\PYG{c+c1}{\PYGZsh{}Conclusion: 2*x\PYGZus{}M \PYGZhy{} x\PYGZus{}X \PYGZhy{} x\PYGZus{}Y = 0}
\PYG{c+c1}{\PYGZsh{}Final Remainder: 0}
\PYG{c+c1}{\PYGZsh{}Verdict: PROVEN TRUE. The conclusion is a consequence of the hypotheses.}

\PYG{c+c1}{\PYGZsh{}Conclusion: 2*y\PYGZus{}M \PYGZhy{} y\PYGZus{}X \PYGZhy{} y\PYGZus{}Y = 0}
\PYG{c+c1}{\PYGZsh{}Final Remainder: 0}
\PYG{c+c1}{\PYGZsh{}Verdict: PROVEN TRUE. The conclusion is a consequence of the hypotheses.}

\end{sphinxVerbatim}


\bigskip\hrule\bigskip



\section{Dynamic geometry with animation: two moving points}
\label{\detokenize{examples:dynamic-geometry-with-animation-two-moving-points}}
\sphinxAtStartPar
Problem: In the rhombus ABCD, draw line segment DE perpendicular to CD, where E is the point of intersection with diagonal AC. Connect BE. Point P is a moving point on segment BE, and construct the symmetry point P’ of point P with respect to line DE. Point Q is a moving point on diagonal AC, connecting PQ and DQ. If AE equals 14 and CE equals 18, then the maximum value of DQ \sphinxhyphen{} P’Q is \_\_\_.

\begin{sphinxVerbatim}[commandchars=\\\{\}]
\PYG{k+kn}{from} \PYG{n+nn}{pygeox} \PYG{k+kn}{import} \PYG{n}{GeoScene}\PYG{p}{,} \PYG{n}{LineSegment}

\PYG{c+c1}{\PYGZsh{} Create a new scene}
\PYG{n}{scene} \PYG{o}{=} \PYG{n}{GeoScene}\PYG{p}{(}\PYG{l+m+mi}{5}\PYG{p}{)}

\PYG{c+c1}{\PYGZsh{} add objects}
\PYG{n}{A}\PYG{p}{,} \PYG{n}{B}\PYG{p}{,} \PYG{n}{C}\PYG{p}{,} \PYG{n}{D}\PYG{p}{,} \PYG{n}{E}\PYG{p}{,} \PYG{n}{P}\PYG{p}{,} \PYG{n}{Q} \PYG{o}{=} \PYG{n}{scene}\PYG{o}{.}\PYG{n}{add}\PYG{o}{.}\PYG{n}{points}\PYG{p}{(}\PYG{p}{[}\PYG{l+s+s2}{\PYGZdq{}}\PYG{l+s+s2}{A}\PYG{l+s+s2}{\PYGZdq{}}\PYG{p}{,}\PYG{l+s+s2}{\PYGZdq{}}\PYG{l+s+s2}{B}\PYG{l+s+s2}{\PYGZdq{}}\PYG{p}{,}\PYG{l+s+s2}{\PYGZdq{}}\PYG{l+s+s2}{C}\PYG{l+s+s2}{\PYGZdq{}}\PYG{p}{,}\PYG{l+s+s2}{\PYGZdq{}}\PYG{l+s+s2}{D}\PYG{l+s+s2}{\PYGZdq{}}\PYG{p}{,} \PYG{l+s+s2}{\PYGZdq{}}\PYG{l+s+s2}{E}\PYG{l+s+s2}{\PYGZdq{}}\PYG{p}{,} \PYG{l+s+s2}{\PYGZdq{}}\PYG{l+s+s2}{P}\PYG{l+s+s2}{\PYGZdq{}}\PYG{p}{,} \PYG{l+s+s2}{\PYGZdq{}}\PYG{l+s+s2}{Q}\PYG{l+s+s2}{\PYGZdq{}}\PYG{p}{]}\PYG{p}{)}

\PYG{c+c1}{\PYGZsh{} In the rhombus ABCD}
\PYG{n}{scene}\PYG{o}{.}\PYG{n}{add}\PYG{o}{.}\PYG{n}{rhombus}\PYG{p}{(}\PYG{n}{A}\PYG{p}{,}\PYG{n}{B}\PYG{p}{,}\PYG{n}{C}\PYG{p}{,}\PYG{n}{D}\PYG{p}{)}

\PYG{c+c1}{\PYGZsh{} draw line segment DE perpendicular to CD}
\PYG{n}{line\PYGZus{}DE} \PYG{o}{=} \PYG{n}{scene}\PYG{o}{.}\PYG{n}{add}\PYG{o}{.}\PYG{n}{line\PYGZus{}segment}\PYG{p}{(}\PYG{n}{D}\PYG{p}{,}\PYG{n}{E}\PYG{p}{)}
\PYG{n}{line\PYGZus{}CD} \PYG{o}{=} \PYG{n}{scene}\PYG{o}{.}\PYG{n}{add}\PYG{o}{.}\PYG{n}{line\PYGZus{}segment}\PYG{p}{(}\PYG{n}{C}\PYG{p}{,}\PYG{n}{D}\PYG{p}{)}
\PYG{n}{scene}\PYG{o}{.}\PYG{n}{relate}\PYG{o}{.}\PYG{n}{perpendicular}\PYG{p}{(}\PYG{n}{line\PYGZus{}DE}\PYG{p}{,} \PYG{n}{line\PYGZus{}CD}\PYG{p}{)}

\PYG{c+c1}{\PYGZsh{}where E is the point of intersection with diagonal AC}
\PYG{n}{line\PYGZus{}AC} \PYG{o}{=} \PYG{n}{scene}\PYG{o}{.}\PYG{n}{add}\PYG{o}{.}\PYG{n}{line\PYGZus{}segment}\PYG{p}{(}\PYG{n}{A}\PYG{p}{,}\PYG{n}{C}\PYG{p}{)}
\PYG{n}{scene}\PYG{o}{.}\PYG{n}{relate}\PYG{o}{.}\PYG{n}{point\PYGZus{}lies\PYGZus{}on}\PYG{p}{(}\PYG{n}{E}\PYG{p}{,} \PYG{n}{line\PYGZus{}AC}\PYG{p}{)}

\PYG{c+c1}{\PYGZsh{}Connect BE}
\PYG{n}{line\PYGZus{}BE} \PYG{o}{=} \PYG{n}{scene}\PYG{o}{.}\PYG{n}{add}\PYG{o}{.}\PYG{n}{line\PYGZus{}segment}\PYG{p}{(}\PYG{n}{B}\PYG{p}{,}\PYG{n}{E}\PYG{p}{)}

\PYG{c+c1}{\PYGZsh{}Point P is a moving point on segment BE}
\PYG{n}{scene}\PYG{o}{.}\PYG{n}{animate}\PYG{o}{.}\PYG{n}{point\PYGZus{}sweep\PYGZus{}line}\PYG{p}{(}\PYG{n}{point} \PYG{o}{=} \PYG{n}{P}\PYG{p}{,} \PYG{n}{start\PYGZus{}point} \PYG{o}{=} \PYG{n}{B}\PYG{p}{,} \PYG{n}{end\PYGZus{}point} \PYG{o}{=} \PYG{n}{E}\PYG{p}{)}

\PYG{c+c1}{\PYGZsh{}construct the symmetry point P\PYGZsq{} of point P with respect to line DE}
\PYG{n}{Ps} \PYG{o}{=} \PYG{n}{scene}\PYG{o}{.}\PYG{n}{add}\PYG{o}{.}\PYG{n}{point}\PYG{p}{(}\PYG{l+s+s2}{\PYGZdq{}}\PYG{l+s+s2}{P}\PYG{l+s+s2}{\PYGZsq{}}\PYG{l+s+s2}{\PYGZdq{}}\PYG{p}{,} \PYG{n}{copy} \PYG{o}{=} \PYG{n}{P}\PYG{o}{.}\PYG{n}{mirror\PYGZus{}across\PYGZus{}line}\PYG{p}{(}\PYG{n}{line\PYGZus{}DE}\PYG{p}{)}\PYG{p}{)}

\PYG{c+c1}{\PYGZsh{}Point Q is a moving point on diagonal AC, connecting PQ and DQ}
\PYG{n}{scene}\PYG{o}{.}\PYG{n}{animate}\PYG{o}{.}\PYG{n}{point\PYGZus{}sweep\PYGZus{}line}\PYG{p}{(}\PYG{n}{point} \PYG{o}{=} \PYG{n}{Q}\PYG{p}{,} \PYG{n}{start\PYGZus{}point} \PYG{o}{=} \PYG{n}{A}\PYG{p}{,} \PYG{n}{end\PYGZus{}point} \PYG{o}{=} \PYG{n}{C}\PYG{p}{)}
\PYG{n}{line\PYGZus{}PQ} \PYG{o}{=} \PYG{n}{scene}\PYG{o}{.}\PYG{n}{add}\PYG{o}{.}\PYG{n}{line\PYGZus{}segment}\PYG{p}{(}\PYG{n}{P}\PYG{p}{,}\PYG{n}{Q}\PYG{p}{)}
\PYG{n}{line\PYGZus{}DQ} \PYG{o}{=} \PYG{n}{scene}\PYG{o}{.}\PYG{n}{add}\PYG{o}{.}\PYG{n}{line\PYGZus{}segment}\PYG{p}{(}\PYG{n}{D}\PYG{p}{,}\PYG{n}{Q}\PYG{p}{)}

\PYG{c+c1}{\PYGZsh{} If AE equals 14 and CE equals 18}
\PYG{n}{scene}\PYG{o}{.}\PYG{n}{constraint}\PYG{o}{.}\PYG{n}{eq}\PYG{p}{(}\PYG{n}{LineSegment}\PYG{p}{(}\PYG{n}{A}\PYG{p}{,}\PYG{n}{E}\PYG{p}{)}\PYG{o}{.}\PYG{n}{length}\PYG{p}{,} \PYG{l+m+mf}{1.4}\PYG{p}{)} \PYG{c+c1}{\PYGZsh{} use small numbers 14 \PYGZhy{}\PYGZgt{} 1.4 }
\PYG{n}{scene}\PYG{o}{.}\PYG{n}{constraint}\PYG{o}{.}\PYG{n}{eq}\PYG{p}{(}\PYG{n}{LineSegment}\PYG{p}{(}\PYG{n}{C}\PYG{p}{,}\PYG{n}{E}\PYG{p}{)}\PYG{o}{.}\PYG{n}{length}\PYG{p}{,} \PYG{l+m+mf}{1.8}\PYG{p}{)} \PYG{c+c1}{\PYGZsh{} use small numbers 18 \PYGZhy{}\PYGZgt{} 1.8 }


\PYG{c+c1}{\PYGZsh{}run animation}
\PYG{n}{scene}\PYG{o}{.}\PYG{n}{animate}\PYG{o}{.}\PYG{n}{run}\PYG{p}{(}\PYG{n}{moving\PYGZus{}points}\PYG{o}{=} \PYG{p}{[}\PYG{n}{P}\PYG{p}{,} \PYG{n}{Ps}\PYG{p}{,} \PYG{n}{Q}\PYG{p}{]}\PYG{p}{,} \PYG{n}{opt\PYGZus{}method} \PYG{o}{=} \PYG{l+s+s2}{\PYGZdq{}}\PYG{l+s+s2}{dual\PYGZus{}annealing}\PYG{l+s+s2}{\PYGZdq{}}\PYG{p}{,} \PYG{n}{frames} \PYG{o}{=} \PYG{l+m+mi}{32}\PYG{p}{)} 
\PYG{c+c1}{\PYGZsh{}when default optimizer fails, use basinhoping \PYGZhy{}\PYGZgt{} slower but it works better}
\PYG{n}{scene}\PYG{o}{.}\PYG{n}{animate}\PYG{o}{.}\PYG{n}{gif}\PYG{p}{(}\PYG{n}{total\PYGZus{}time}\PYG{o}{=}\PYG{l+m+mi}{10}\PYG{p}{)} 
\end{sphinxVerbatim}

\sphinxAtStartPar
%\sphinxincludegraphics{{animation}.gif}

\sphinxstepscope


\chapter{Main geometric scene}
\label{\detokenize{pygeox:module-pygeox.geoscene}}\label{\detokenize{pygeox:main-geometric-scene}}\label{\detokenize{pygeox::doc}}\index{module@\spxentry{module}!pygeox.geoscene@\spxentry{pygeox.geoscene}}\index{pygeox.geoscene@\spxentry{pygeox.geoscene}!module@\spxentry{module}}
\sphinxAtStartPar
GeoScene module: Core framework for geometric scene management, constraint handling,
and problem solving.

\sphinxAtStartPar
This module defines the GeoScene class, which serves as the central manager for
geometric objects (points, lines, circles, polygons, arcs) and constraints. It
provides a unified API to add, retrieve, and manipulate geometric elements, define
symbolic constraints between them, and invoke solvers to find solutions to geometric
problems.

\sphinxAtStartPar
Key features include:
\sphinxhyphen{} Automatic management and storage of geometric objects.
\sphinxhyphen{} Symbolic constraint definition with human\sphinxhyphen{}readable descriptions.
\sphinxhyphen{} Integration of namespaces for relationships, constraint addition, solving,
\begin{quote}

\sphinxAtStartPar
animation, and geometric proofs.
\end{quote}
\begin{itemize}
\item {} 
\sphinxAtStartPar
Sketching capability by auto\sphinxhyphen{}creating points when referenced by name.

\item {} 
\sphinxAtStartPar
Utilities for calculating angles, plotting scenes, and proof management.

\end{itemize}

\sphinxAtStartPar
The GeoScene class acts as the main entry point for users to programmatically
create, analyze, and visualize geometry problems and their solutions.
\index{GeoScene (class in pygeox.geoscene)@\spxentry{GeoScene}\spxextra{class in pygeox.geoscene}}

\begin{fulllineitems}
\phantomsection\label{\detokenize{pygeox:pygeox.geoscene.GeoScene}}
\pysigstartsignatures
\pysiglinewithargsret
{\sphinxbfcode{\sphinxupquote{\DUrole{k}{class}\DUrole{w}{ }}}\sphinxcode{\sphinxupquote{pygeox.geoscene.}}\sphinxbfcode{\sphinxupquote{GeoScene}}}
{\sphinxparam{\DUrole{n}{domain}\DUrole{o}{=}\DUrole{default_value}{None}}}
{}
\pysigstopsignatures
\sphinxAtStartPar
Bases: \sphinxcode{\sphinxupquote{object}}

\sphinxAtStartPar
Manages all geometric objects and constraints, and solves for unknown variables.
This class is the main entry point for defining, manipulating, and solving
geometric problems. It provides sketching capabilities by automatically creating
points when referenced by name.
\begin{quote}\begin{description}
\sphinxlineitem{Parameters}
\sphinxAtStartPar
\sphinxstyleliteralstrong{\sphinxupquote{domain}} (\sphinxstyleliteralemphasis{\sphinxupquote{Any}}\sphinxstyleliteralemphasis{\sphinxupquote{ | }}\sphinxstyleliteralemphasis{\sphinxupquote{None}})

\end{description}\end{quote}
\index{GeoScene.ProvingContext (class in pygeox.geoscene)@\spxentry{GeoScene.ProvingContext}\spxextra{class in pygeox.geoscene}}

\begin{fulllineitems}
\phantomsection\label{\detokenize{pygeox:pygeox.geoscene.GeoScene.ProvingContext}}
\pysigstartsignatures
\pysiglinewithargsret
{\sphinxbfcode{\sphinxupquote{\DUrole{k}{class}\DUrole{w}{ }}}\sphinxbfcode{\sphinxupquote{ProvingContext}}}
{\sphinxparam{\DUrole{n}{scene\_instance}}}
{}
\pysigstopsignatures
\sphinxAtStartPar
Bases: \sphinxcode{\sphinxupquote{object}}

\sphinxAtStartPar
Initialize proving context manager.
\begin{quote}\begin{description}
\sphinxlineitem{Parameters}
\sphinxAtStartPar
\sphinxstyleliteralstrong{\sphinxupquote{scene\_instance}} \textendash{} GeoScene instance.

\end{description}\end{quote}

\end{fulllineitems}

\index{add\_constraint() (pygeox.geoscene.GeoScene method)@\spxentry{add\_constraint()}\spxextra{pygeox.geoscene.GeoScene method}}

\begin{fulllineitems}
\phantomsection\label{\detokenize{pygeox:pygeox.geoscene.GeoScene.add_constraint}}
\pysigstartsignatures
\pysiglinewithargsret
{\sphinxbfcode{\sphinxupquote{add\_constraint}}}
{\sphinxparam{\DUrole{n}{equation}}\sphinxparamcomma \sphinxparam{\DUrole{n}{description}}}
{}
\pysigstopsignatures
\sphinxAtStartPar
Add a symbolic equation constraint with a description.

\sphinxAtStartPar
If in proving mode, stores constraint for proof; else adds to main constraints.
\begin{quote}\begin{description}
\sphinxlineitem{Parameters}\begin{itemize}
\item {} 
\sphinxAtStartPar
\sphinxstyleliteralstrong{\sphinxupquote{equation}} (\sphinxstyleliteralemphasis{\sphinxupquote{Basic}}) \textendash{} Sympy symbolic equation.

\item {} 
\sphinxAtStartPar
\sphinxstyleliteralstrong{\sphinxupquote{description}} (\sphinxstyleliteralemphasis{\sphinxupquote{str}}) \textendash{} Human\sphinxhyphen{}readable description.

\end{itemize}

\sphinxlineitem{Return type}
\sphinxAtStartPar
None

\end{description}\end{quote}

\end{fulllineitems}

\index{add\_mapping\_constraint() (pygeox.geoscene.GeoScene method)@\spxentry{add\_mapping\_constraint()}\spxextra{pygeox.geoscene.GeoScene method}}

\begin{fulllineitems}
\phantomsection\label{\detokenize{pygeox:pygeox.geoscene.GeoScene.add_mapping_constraint}}
\pysigstartsignatures
\pysiglinewithargsret
{\sphinxbfcode{\sphinxupquote{add\_mapping\_constraint}}}
{\sphinxparam{\DUrole{n}{equation}}\sphinxparamcomma \sphinxparam{\DUrole{n}{description}}}
{}
\pysigstopsignatures
\sphinxAtStartPar
Add a mapping constraint equation used for parameter mapping.
\begin{quote}\begin{description}
\sphinxlineitem{Parameters}\begin{itemize}
\item {} 
\sphinxAtStartPar
\sphinxstyleliteralstrong{\sphinxupquote{equation}} (\sphinxstyleliteralemphasis{\sphinxupquote{Basic}}) \textendash{} Sympy symbolic equation.

\item {} 
\sphinxAtStartPar
\sphinxstyleliteralstrong{\sphinxupquote{description}} (\sphinxstyleliteralemphasis{\sphinxupquote{str}}) \textendash{} Human\sphinxhyphen{}readable description.

\end{itemize}

\sphinxlineitem{Return type}
\sphinxAtStartPar
None

\end{description}\end{quote}

\end{fulllineitems}

\index{add\_object() (pygeox.geoscene.GeoScene method)@\spxentry{add\_object()}\spxextra{pygeox.geoscene.GeoScene method}}

\begin{fulllineitems}
\phantomsection\label{\detokenize{pygeox:pygeox.geoscene.GeoScene.add_object}}
\pysigstartsignatures
\pysiglinewithargsret
{\sphinxbfcode{\sphinxupquote{add\_object}}}
{\sphinxparam{\DUrole{n}{obj\_type}}\sphinxparamcomma \sphinxparam{\DUrole{n}{name}}\sphinxparamcomma \sphinxparam{\DUrole{n}{obj}}}
{}
\pysigstopsignatures
\sphinxAtStartPar
Add a geometric object to the scene under the given type and name.
\begin{quote}\begin{description}
\sphinxlineitem{Parameters}\begin{itemize}
\item {} 
\sphinxAtStartPar
\sphinxstyleliteralstrong{\sphinxupquote{obj\_type}} (\sphinxstyleliteralemphasis{\sphinxupquote{str}}) \textendash{} Type of the object (e.g., ‘Point’, ‘Circle’).

\item {} 
\sphinxAtStartPar
\sphinxstyleliteralstrong{\sphinxupquote{name}} (\sphinxstyleliteralemphasis{\sphinxupquote{str}}) \textendash{} Name of the object.

\item {} 
\sphinxAtStartPar
\sphinxstyleliteralstrong{\sphinxupquote{obj}} (\sphinxstyleliteralemphasis{\sphinxupquote{Any}}) \textendash{} The geometric object instance.

\end{itemize}

\sphinxlineitem{Return type}
\sphinxAtStartPar
None

\end{description}\end{quote}

\end{fulllineitems}

\index{add\_parameter() (pygeox.geoscene.GeoScene method)@\spxentry{add\_parameter()}\spxextra{pygeox.geoscene.GeoScene method}}

\begin{fulllineitems}
\phantomsection\label{\detokenize{pygeox:pygeox.geoscene.GeoScene.add_parameter}}
\pysigstartsignatures
\pysiglinewithargsret
{\sphinxbfcode{\sphinxupquote{add\_parameter}}}
{\sphinxparam{\DUrole{n}{symbol}}\sphinxparamcomma \sphinxparam{\DUrole{n}{type}\DUrole{o}{=}\DUrole{default_value}{\textquotesingle{}x\textquotesingle{}}}\sphinxparamcomma \sphinxparam{\DUrole{n}{full}\DUrole{o}{=}\DUrole{default_value}{True}}\sphinxparamcomma \sphinxparam{\DUrole{n}{origin}\DUrole{o}{=}\DUrole{default_value}{\textquotesingle{}Problem\textquotesingle{}}}\sphinxparamcomma \sphinxparam{\DUrole{n}{bound}\DUrole{o}{=}\DUrole{default_value}{None}}}
{}
\pysigstopsignatures
\sphinxAtStartPar
Add a symbolic parameter with metadata to the scene.
\begin{quote}\begin{description}
\sphinxlineitem{Parameters}\begin{itemize}
\item {} 
\sphinxAtStartPar
\sphinxstyleliteralstrong{\sphinxupquote{symbol}} (\sphinxstyleliteralemphasis{\sphinxupquote{Symbol}}) \textendash{} The sympy.Symbol representing the parameter.

\item {} 
\sphinxAtStartPar
\sphinxstyleliteralstrong{\sphinxupquote{type}} (\sphinxstyleliteralemphasis{\sphinxupquote{str}}) \textendash{} Parameter type (default “x”).

\item {} 
\sphinxAtStartPar
\sphinxstyleliteralstrong{\sphinxupquote{full}} (\sphinxstyleliteralemphasis{\sphinxupquote{bool}}) \textendash{} Whether parameter is fully active.

\item {} 
\sphinxAtStartPar
\sphinxstyleliteralstrong{\sphinxupquote{origin}} (\sphinxstyleliteralemphasis{\sphinxupquote{str}}) \textendash{} Origin description.

\item {} 
\sphinxAtStartPar
\sphinxstyleliteralstrong{\sphinxupquote{bound}} \textendash{} Optional numeric bounds {[}min, max{]}.

\end{itemize}

\sphinxlineitem{Returns}
\sphinxAtStartPar
The symbol added.

\end{description}\end{quote}

\end{fulllineitems}

\index{angle() (pygeox.geoscene.GeoScene method)@\spxentry{angle()}\spxextra{pygeox.geoscene.GeoScene method}}

\begin{fulllineitems}
\phantomsection\label{\detokenize{pygeox:pygeox.geoscene.GeoScene.angle}}
\pysigstartsignatures
\pysiglinewithargsret
{\sphinxbfcode{\sphinxupquote{angle}}}
{\sphinxparam{\DUrole{n}{point1}}\sphinxparamcomma \sphinxparam{\DUrole{n}{vertex}}\sphinxparamcomma \sphinxparam{\DUrole{n}{point2}}\sphinxparamcomma \sphinxparam{\DUrole{n}{signed}\DUrole{o}{=}\DUrole{default_value}{False}}}
{}
\pysigstopsignatures
\sphinxAtStartPar
Calculate the angle (degrees between 0 and 180) at vertex between point1 and point2.
\begin{quote}\begin{description}
\sphinxlineitem{Parameters}\begin{itemize}
\item {} 
\sphinxAtStartPar
\sphinxstyleliteralstrong{\sphinxupquote{point1}} (\sphinxstyleliteralemphasis{\sphinxupquote{Point}}) \textendash{} First point defining the angle.

\item {} 
\sphinxAtStartPar
\sphinxstyleliteralstrong{\sphinxupquote{vertex}} (\sphinxstyleliteralemphasis{\sphinxupquote{Point}}) \textendash{} Vertex point of the angle.

\item {} 
\sphinxAtStartPar
\sphinxstyleliteralstrong{\sphinxupquote{point2}} (\sphinxstyleliteralemphasis{\sphinxupquote{Point}}) \textendash{} Second point defining the angle.

\item {} 
\sphinxAtStartPar
\sphinxstyleliteralstrong{\sphinxupquote{signed}} (\sphinxstyleliteralemphasis{\sphinxupquote{bool}}) \textendash{} If True, returns signed angle (CCW positive).

\end{itemize}

\sphinxlineitem{Returns}
\sphinxAtStartPar
Angle in degrees.

\sphinxlineitem{Return type}
\sphinxAtStartPar
int | float | \sphinxstyleemphasis{Expr}

\end{description}\end{quote}

\end{fulllineitems}

\index{get\_all\_objects() (pygeox.geoscene.GeoScene method)@\spxentry{get\_all\_objects()}\spxextra{pygeox.geoscene.GeoScene method}}

\begin{fulllineitems}
\phantomsection\label{\detokenize{pygeox:pygeox.geoscene.GeoScene.get_all_objects}}
\pysigstartsignatures
\pysiglinewithargsret
{\sphinxbfcode{\sphinxupquote{get\_all\_objects}}}
{\sphinxparam{\DUrole{n}{obj\_type}\DUrole{o}{=}\DUrole{default_value}{None}}}
{}
\pysigstopsignatures
\sphinxAtStartPar
Retrieve all objects of a given type or all objects if no type is specified.
\begin{quote}\begin{description}
\sphinxlineitem{Parameters}
\sphinxAtStartPar
\sphinxstyleliteralstrong{\sphinxupquote{obj\_type}} (\sphinxstyleliteralemphasis{\sphinxupquote{str}}\sphinxstyleliteralemphasis{\sphinxupquote{ | }}\sphinxstyleliteralemphasis{\sphinxupquote{None}}) \textendash{} Optional type of objects to retrieve.

\sphinxlineitem{Returns}
\sphinxAtStartPar
Dictionary of object names to objects.

\sphinxlineitem{Return type}
\sphinxAtStartPar
\sphinxstyleemphasis{Dict}{[}str, \sphinxstyleemphasis{Any}{]}

\end{description}\end{quote}

\end{fulllineitems}

\index{get\_object() (pygeox.geoscene.GeoScene method)@\spxentry{get\_object()}\spxextra{pygeox.geoscene.GeoScene method}}

\begin{fulllineitems}
\phantomsection\label{\detokenize{pygeox:pygeox.geoscene.GeoScene.get_object}}
\pysigstartsignatures
\pysiglinewithargsret
{\sphinxbfcode{\sphinxupquote{get\_object}}}
{\sphinxparam{\DUrole{n}{obj\_type}}\sphinxparamcomma \sphinxparam{\DUrole{n}{name}}}
{}
\pysigstopsignatures
\sphinxAtStartPar
Retrieve a specific object by type and name.
\begin{quote}\begin{description}
\sphinxlineitem{Parameters}\begin{itemize}
\item {} 
\sphinxAtStartPar
\sphinxstyleliteralstrong{\sphinxupquote{obj\_type}} (\sphinxstyleliteralemphasis{\sphinxupquote{str}}) \textendash{} Type of the object.

\item {} 
\sphinxAtStartPar
\sphinxstyleliteralstrong{\sphinxupquote{name}} (\sphinxstyleliteralemphasis{\sphinxupquote{str}}) \textendash{} Name of the object.

\end{itemize}

\sphinxlineitem{Returns}
\sphinxAtStartPar
The geometric object or None if not found.

\sphinxlineitem{Return type}
\sphinxAtStartPar
\sphinxstyleemphasis{Any} | None

\end{description}\end{quote}

\end{fulllineitems}

\index{get\_points() (pygeox.geoscene.GeoScene method)@\spxentry{get\_points()}\spxextra{pygeox.geoscene.GeoScene method}}

\begin{fulllineitems}
\phantomsection\label{\detokenize{pygeox:pygeox.geoscene.GeoScene.get_points}}
\pysigstartsignatures
\pysiglinewithargsret
{\sphinxbfcode{\sphinxupquote{get\_points}}}
{}
{}
\pysigstopsignatures
\sphinxAtStartPar
Get dictionary of all points in the scene with their (x, y) coordinates.
\begin{quote}\begin{description}
\sphinxlineitem{Returns}
\sphinxAtStartPar
Dict mapping point names to (x, y) coordinate tuples.

\sphinxlineitem{Return type}
\sphinxAtStartPar
dict

\end{description}\end{quote}

\end{fulllineitems}

\index{plot() (pygeox.geoscene.GeoScene method)@\spxentry{plot()}\spxextra{pygeox.geoscene.GeoScene method}}

\begin{fulllineitems}
\phantomsection\label{\detokenize{pygeox:pygeox.geoscene.GeoScene.plot}}
\pysigstartsignatures
\pysiglinewithargsret
{\sphinxbfcode{\sphinxupquote{plot}}}
{\sphinxparam{\DUrole{n}{show\_names}\DUrole{o}{=}\DUrole{default_value}{True}}\sphinxparamcomma \sphinxparam{\DUrole{n}{show\_axis\_grid}\DUrole{o}{=}\DUrole{default_value}{False}}\sphinxparamcomma \sphinxparam{\DUrole{n}{xlim}\DUrole{o}{=}\DUrole{default_value}{None}}\sphinxparamcomma \sphinxparam{\DUrole{n}{ylim}\DUrole{o}{=}\DUrole{default_value}{None}}\sphinxparamcomma \sphinxparam{\DUrole{n}{return\_fig}\DUrole{o}{=}\DUrole{default_value}{False}}}
{}
\pysigstopsignatures
\sphinxAtStartPar
Plot all geometric objects in the scene.
\begin{quote}\begin{description}
\sphinxlineitem{Parameters}\begin{itemize}
\item {} 
\sphinxAtStartPar
\sphinxstyleliteralstrong{\sphinxupquote{show\_names}} (\sphinxstyleliteralemphasis{\sphinxupquote{bool}}) \textendash{} Show object names on plot.

\item {} 
\sphinxAtStartPar
\sphinxstyleliteralstrong{\sphinxupquote{show\_axis\_grid}} (\sphinxstyleliteralemphasis{\sphinxupquote{bool}}) \textendash{} Show axis grid.

\item {} 
\sphinxAtStartPar
\sphinxstyleliteralstrong{\sphinxupquote{xlim}} (\sphinxstyleliteralemphasis{\sphinxupquote{tuple}}\sphinxstyleliteralemphasis{\sphinxupquote{{[}}}\sphinxstyleliteralemphasis{\sphinxupquote{float}}\sphinxstyleliteralemphasis{\sphinxupquote{, }}\sphinxstyleliteralemphasis{\sphinxupquote{float}}\sphinxstyleliteralemphasis{\sphinxupquote{{]} }}\sphinxstyleliteralemphasis{\sphinxupquote{| }}\sphinxstyleliteralemphasis{\sphinxupquote{None}}) \textendash{} X\sphinxhyphen{}axis limits.

\item {} 
\sphinxAtStartPar
\sphinxstyleliteralstrong{\sphinxupquote{ylim}} (\sphinxstyleliteralemphasis{\sphinxupquote{tuple}}\sphinxstyleliteralemphasis{\sphinxupquote{{[}}}\sphinxstyleliteralemphasis{\sphinxupquote{float}}\sphinxstyleliteralemphasis{\sphinxupquote{, }}\sphinxstyleliteralemphasis{\sphinxupquote{float}}\sphinxstyleliteralemphasis{\sphinxupquote{{]} }}\sphinxstyleliteralemphasis{\sphinxupquote{| }}\sphinxstyleliteralemphasis{\sphinxupquote{None}}) \textendash{} Y\sphinxhyphen{}axis limits.

\item {} 
\sphinxAtStartPar
\sphinxstyleliteralstrong{\sphinxupquote{return\_fig}} (\sphinxstyleliteralemphasis{\sphinxupquote{bool}}) \textendash{} If True, return matplotlib figure object.

\end{itemize}

\sphinxlineitem{Returns}
\sphinxAtStartPar
Matplotlib figure if return\_fig is True; otherwise None.

\end{description}\end{quote}

\end{fulllineitems}

\index{prove() (pygeox.geoscene.GeoScene method)@\spxentry{prove()}\spxextra{pygeox.geoscene.GeoScene method}}

\begin{fulllineitems}
\phantomsection\label{\detokenize{pygeox:pygeox.geoscene.GeoScene.prove}}
\pysigstartsignatures
\pysiglinewithargsret
{\sphinxbfcode{\sphinxupquote{prove}}}
{}
{}
\pysigstopsignatures
\sphinxAtStartPar
Run the prover function on the accumulated proof statements and constraints.
Updates the proof\_statements with prover results.

\end{fulllineitems}

\index{proving() (pygeox.geoscene.GeoScene method)@\spxentry{proving()}\spxextra{pygeox.geoscene.GeoScene method}}

\begin{fulllineitems}
\phantomsection\label{\detokenize{pygeox:pygeox.geoscene.GeoScene.proving}}
\pysigstartsignatures
\pysiglinewithargsret
{\sphinxbfcode{\sphinxupquote{proving}}}
{}
{}
\pysigstopsignatures
\sphinxAtStartPar
Create a context manager for defining proof constraints.
\begin{quote}\begin{description}
\sphinxlineitem{Returns}
\sphinxAtStartPar
ProvingContext instance.

\end{description}\end{quote}

\end{fulllineitems}


\end{fulllineitems}



\chapter{Add geometric objects}
\label{\detokenize{pygeox:module-pygeox.add_objects}}\label{\detokenize{pygeox:add-geometric-objects}}\index{module@\spxentry{module}!pygeox.add\_objects@\spxentry{pygeox.add\_objects}}\index{pygeox.add\_objects@\spxentry{pygeox.add\_objects}!module@\spxentry{module}}
\sphinxAtStartPar
Factory namespace for creating and adding geometric objects to a scene.

\sphinxAtStartPar
This module defines the AddNamespace class, which provides convenient constructors
to create geometric objects such as points, lines, arcs, polygons, and various
special polygons (triangles, quadrilaterals, regular polygons, etc.).

\sphinxAtStartPar
Each method creates the object, registers it in the scene, and adds necessary
constraints automatically (e.g., equal radii for arcs, collinearity for semicircles).

\sphinxAtStartPar
The namespace facilitates easy and consistent object creation for geometric problem
modeling, with support for copying existing objects and symbolic coordinate handling.
\index{AddNamespace (class in pygeox.add\_objects)@\spxentry{AddNamespace}\spxextra{class in pygeox.add\_objects}}

\begin{fulllineitems}
\phantomsection\label{\detokenize{pygeox:pygeox.add_objects.AddNamespace}}
\pysigstartsignatures
\pysiglinewithargsret
{\sphinxbfcode{\sphinxupquote{\DUrole{k}{class}\DUrole{w}{ }}}\sphinxcode{\sphinxupquote{pygeox.add\_objects.}}\sphinxbfcode{\sphinxupquote{AddNamespace}}}
{\sphinxparam{\DUrole{n}{scene}}}
{}
\pysigstopsignatures
\sphinxAtStartPar
Bases: \sphinxcode{\sphinxupquote{object}}

\sphinxAtStartPar
Factory namespace for adding geometric objects to a scene.

\sphinxAtStartPar
This class exposes convenience constructors that create objects,
register them in the scene, and add any required constraints.
\index{acute\_triangle() (pygeox.add\_objects.AddNamespace method)@\spxentry{acute\_triangle()}\spxextra{pygeox.add\_objects.AddNamespace method}}

\begin{fulllineitems}
\phantomsection\label{\detokenize{pygeox:pygeox.add_objects.AddNamespace.acute_triangle}}
\pysigstartsignatures
\pysiglinewithargsret
{\sphinxbfcode{\sphinxupquote{acute\_triangle}}}
{\sphinxparam{\DUrole{n}{p1}\DUrole{o}{=}\DUrole{default_value}{None}}\sphinxparamcomma \sphinxparam{\DUrole{n}{p2}\DUrole{o}{=}\DUrole{default_value}{None}}\sphinxparamcomma \sphinxparam{\DUrole{n}{p3}\DUrole{o}{=}\DUrole{default_value}{None}}\sphinxparamcomma \sphinxparam{\DUrole{n}{name}\DUrole{o}{=}\DUrole{default_value}{None}}\sphinxparamcomma \sphinxparam{\DUrole{n}{copy}\DUrole{o}{=}\DUrole{default_value}{None}}}
{}
\pysigstopsignatures
\sphinxAtStartPar
Add an acute triangle or return a copy.
\begin{quote}\begin{description}
\sphinxlineitem{Parameters}\begin{itemize}
\item {} 
\sphinxAtStartPar
\sphinxstyleliteralstrong{\sphinxupquote{p1}} (\sphinxstyleliteralemphasis{\sphinxupquote{Point}}\sphinxstyleliteralemphasis{\sphinxupquote{ | }}\sphinxstyleliteralemphasis{\sphinxupquote{None}}) \textendash{} Vertex A.

\item {} 
\sphinxAtStartPar
\sphinxstyleliteralstrong{\sphinxupquote{p2}} (\sphinxstyleliteralemphasis{\sphinxupquote{Point}}\sphinxstyleliteralemphasis{\sphinxupquote{ | }}\sphinxstyleliteralemphasis{\sphinxupquote{None}}) \textendash{} Vertex B.

\item {} 
\sphinxAtStartPar
\sphinxstyleliteralstrong{\sphinxupquote{p3}} (\sphinxstyleliteralemphasis{\sphinxupquote{Point}}\sphinxstyleliteralemphasis{\sphinxupquote{ | }}\sphinxstyleliteralemphasis{\sphinxupquote{None}}) \textendash{} Vertex C.

\item {} 
\sphinxAtStartPar
\sphinxstyleliteralstrong{\sphinxupquote{name}} (\sphinxstyleliteralemphasis{\sphinxupquote{str}}\sphinxstyleliteralemphasis{\sphinxupquote{ | }}\sphinxstyleliteralemphasis{\sphinxupquote{None}}) \textendash{} Optional name.

\item {} 
\sphinxAtStartPar
\sphinxstyleliteralstrong{\sphinxupquote{copy}} (\sphinxstyleliteralemphasis{\sphinxupquote{AcuteTriangle}}\sphinxstyleliteralemphasis{\sphinxupquote{ | }}\sphinxstyleliteralemphasis{\sphinxupquote{None}}) \textendash{} Existing triangle to reuse.

\end{itemize}

\sphinxlineitem{Returns}
\sphinxAtStartPar
The triangle instance (copy if provided).

\sphinxlineitem{Return type}
\sphinxAtStartPar
AcuteTriangle

\end{description}\end{quote}
\subsubsection*{Notes}

\sphinxAtStartPar
All internal angles are constrained \textless{} 90°.

\end{fulllineitems}

\index{angle() (pygeox.add\_objects.AddNamespace method)@\spxentry{angle()}\spxextra{pygeox.add\_objects.AddNamespace method}}

\begin{fulllineitems}
\phantomsection\label{\detokenize{pygeox:pygeox.add_objects.AddNamespace.angle}}
\pysigstartsignatures
\pysiglinewithargsret
{\sphinxbfcode{\sphinxupquote{angle}}}
{\sphinxparam{\DUrole{n}{A}}\sphinxparamcomma \sphinxparam{\DUrole{n}{vertex}}\sphinxparamcomma \sphinxparam{\DUrole{n}{B}}\sphinxparamcomma \sphinxparam{\DUrole{n}{plot\_sign}\DUrole{o}{=}\DUrole{default_value}{True}}\sphinxparamcomma \sphinxparam{\DUrole{n}{plot\_text}\DUrole{o}{=}\DUrole{default_value}{False}}\sphinxparamcomma \sphinxparam{\DUrole{n}{name}\DUrole{o}{=}\DUrole{default_value}{None}}}
{}
\pysigstopsignatures
\sphinxAtStartPar
Add an angle object defined by three points: A, vertex, and B.

\sphinxAtStartPar
The angle is formed at the \sphinxtitleref{vertex} point between points \sphinxtitleref{A} and \sphinxtitleref{B}.
\begin{quote}\begin{description}
\sphinxlineitem{Parameters}\begin{itemize}
\item {} 
\sphinxAtStartPar
\sphinxstyleliteralstrong{\sphinxupquote{A}} (\sphinxstyleliteralemphasis{\sphinxupquote{Point}}) \textendash{} One point defining the angle.

\item {} 
\sphinxAtStartPar
\sphinxstyleliteralstrong{\sphinxupquote{vertex}} (\sphinxstyleliteralemphasis{\sphinxupquote{Point}}) \textendash{} The vertex point where the angle is measured.

\item {} 
\sphinxAtStartPar
\sphinxstyleliteralstrong{\sphinxupquote{B}} (\sphinxstyleliteralemphasis{\sphinxupquote{Point}}) \textendash{} The other point defining the angle.

\item {} 
\sphinxAtStartPar
\sphinxstyleliteralstrong{\sphinxupquote{name}} (\sphinxstyleliteralemphasis{\sphinxupquote{Optional}}\sphinxstyleliteralemphasis{\sphinxupquote{{[}}}\sphinxstyleliteralemphasis{\sphinxupquote{str}}\sphinxstyleliteralemphasis{\sphinxupquote{{]}}}\sphinxstyleliteralemphasis{\sphinxupquote{, }}\sphinxstyleliteralemphasis{\sphinxupquote{optional}}) \textendash{} Name for the angle object.

\item {} 
\sphinxAtStartPar
\sphinxstyleliteralstrong{\sphinxupquote{plot\_sign}} (\sphinxstyleliteralemphasis{\sphinxupquote{bool}}\sphinxstyleliteralemphasis{\sphinxupquote{ | }}\sphinxstyleliteralemphasis{\sphinxupquote{None}})

\item {} 
\sphinxAtStartPar
\sphinxstyleliteralstrong{\sphinxupquote{plot\_text}} (\sphinxstyleliteralemphasis{\sphinxupquote{bool}}\sphinxstyleliteralemphasis{\sphinxupquote{ | }}\sphinxstyleliteralemphasis{\sphinxupquote{None}})

\end{itemize}

\sphinxlineitem{Returns}
\sphinxAtStartPar
The measure of the angle at \sphinxtitleref{vertex} between points \sphinxtitleref{A} and \sphinxtitleref{B}.
This is typically a symbolic expression (e.g., a SymPy Expr) if coordinates
are symbolic or unknown.

\end{description}\end{quote}

\end{fulllineitems}

\index{chord() (pygeox.add\_objects.AddNamespace method)@\spxentry{chord()}\spxextra{pygeox.add\_objects.AddNamespace method}}

\begin{fulllineitems}
\phantomsection\label{\detokenize{pygeox:pygeox.add_objects.AddNamespace.chord}}
\pysigstartsignatures
\pysiglinewithargsret
{\sphinxbfcode{\sphinxupquote{chord}}}
{\sphinxparam{\DUrole{n}{circle}}\sphinxparamcomma \sphinxparam{\DUrole{n}{A}}\sphinxparamcomma \sphinxparam{\DUrole{n}{B}}\sphinxparamcomma \sphinxparam{\DUrole{n}{name}\DUrole{o}{=}\DUrole{default_value}{None}}}
{}
\pysigstopsignatures
\sphinxAtStartPar
Add a chord between two points on a circle.

\sphinxAtStartPar
Points \sphinxcode{\sphinxupquote{A}} and \sphinxcode{\sphinxupquote{B}} are constrained to lie on \sphinxcode{\sphinxupquote{circle}}.
\begin{quote}\begin{description}
\sphinxlineitem{Parameters}\begin{itemize}
\item {} 
\sphinxAtStartPar
\sphinxstyleliteralstrong{\sphinxupquote{circle}} (\sphinxstyleliteralemphasis{\sphinxupquote{Circle}}) \textendash{} The circle on which the points lie.

\item {} 
\sphinxAtStartPar
\sphinxstyleliteralstrong{\sphinxupquote{A}} (\sphinxstyleliteralemphasis{\sphinxupquote{Point}}) \textendash{} First point on the circle.

\item {} 
\sphinxAtStartPar
\sphinxstyleliteralstrong{\sphinxupquote{B}} (\sphinxstyleliteralemphasis{\sphinxupquote{Point}}) \textendash{} Second point on the circle.

\item {} 
\sphinxAtStartPar
\sphinxstyleliteralstrong{\sphinxupquote{name}} (\sphinxstyleliteralemphasis{\sphinxupquote{str}}\sphinxstyleliteralemphasis{\sphinxupquote{ | }}\sphinxstyleliteralemphasis{\sphinxupquote{None}}) \textendash{} Optional name. Defaults to \sphinxcode{\sphinxupquote{chord\_A\_B\_on\_\textless{}circle\textgreater{}}}.

\end{itemize}

\sphinxlineitem{Returns}
\sphinxAtStartPar
The created chord, registered in the scene.

\sphinxlineitem{Return type}
\sphinxAtStartPar
LineSegment

\end{description}\end{quote}

\end{fulllineitems}

\index{circle() (pygeox.add\_objects.AddNamespace method)@\spxentry{circle()}\spxextra{pygeox.add\_objects.AddNamespace method}}

\begin{fulllineitems}
\phantomsection\label{\detokenize{pygeox:pygeox.add_objects.AddNamespace.circle}}
\pysigstartsignatures
\pysiglinewithargsret
{\sphinxbfcode{\sphinxupquote{circle}}}
{\sphinxparam{\DUrole{n}{center}\DUrole{o}{=}\DUrole{default_value}{None}}\sphinxparamcomma \sphinxparam{\DUrole{n}{radius}\DUrole{o}{=}\DUrole{default_value}{None}}\sphinxparamcomma \sphinxparam{\DUrole{n}{name}\DUrole{o}{=}\DUrole{default_value}{None}}\sphinxparamcomma \sphinxparam{\DUrole{n}{copy}\DUrole{o}{=}\DUrole{default_value}{None}}}
{}
\pysigstopsignatures
\sphinxAtStartPar
Add a circle or copy an existing one.

\sphinxAtStartPar
If not copying and \sphinxcode{\sphinxupquote{radius}} is \sphinxcode{\sphinxupquote{None}}, a \sphinxcode{\sphinxupquote{sympy.Symbol}}
is created for the radius using the scene’s radius domain.
\begin{quote}\begin{description}
\sphinxlineitem{Parameters}\begin{itemize}
\item {} 
\sphinxAtStartPar
\sphinxstyleliteralstrong{\sphinxupquote{center}} (\sphinxstyleliteralemphasis{\sphinxupquote{Point}}\sphinxstyleliteralemphasis{\sphinxupquote{ | }}\sphinxstyleliteralemphasis{\sphinxupquote{None}}) \textendash{} Center point of the circle.

\item {} 
\sphinxAtStartPar
\sphinxstyleliteralstrong{\sphinxupquote{radius}} (\sphinxstyleliteralemphasis{\sphinxupquote{int}}\sphinxstyleliteralemphasis{\sphinxupquote{ | }}\sphinxstyleliteralemphasis{\sphinxupquote{float}}\sphinxstyleliteralemphasis{\sphinxupquote{ | }}\sphinxstyleliteralemphasis{\sphinxupquote{Expr}}\sphinxstyleliteralemphasis{\sphinxupquote{ | }}\sphinxstyleliteralemphasis{\sphinxupquote{None}}) \textendash{} Radius value or SymPy expression. If \sphinxcode{\sphinxupquote{None}}, a symbol is created.

\item {} 
\sphinxAtStartPar
\sphinxstyleliteralstrong{\sphinxupquote{name}} (\sphinxstyleliteralemphasis{\sphinxupquote{str}}\sphinxstyleliteralemphasis{\sphinxupquote{ | }}\sphinxstyleliteralemphasis{\sphinxupquote{None}}) \textendash{} Optional name for the circle.

\item {} 
\sphinxAtStartPar
\sphinxstyleliteralstrong{\sphinxupquote{copy}} (\sphinxstyleliteralemphasis{\sphinxupquote{Circle}}\sphinxstyleliteralemphasis{\sphinxupquote{ | }}\sphinxstyleliteralemphasis{\sphinxupquote{None}}) \textendash{} Existing circle to duplicate.

\end{itemize}

\sphinxlineitem{Returns}
\sphinxAtStartPar
The created circle, registered in the scene.

\sphinxlineitem{Return type}
\sphinxAtStartPar
Circle

\sphinxlineitem{Raises}
\sphinxAtStartPar
\sphinxstyleliteralstrong{\sphinxupquote{ValueError}} \textendash{} If not copying and \sphinxcode{\sphinxupquote{center}} is not provided.

\end{description}\end{quote}

\end{fulllineitems}

\index{equilateral\_triangle() (pygeox.add\_objects.AddNamespace method)@\spxentry{equilateral\_triangle()}\spxextra{pygeox.add\_objects.AddNamespace method}}

\begin{fulllineitems}
\phantomsection\label{\detokenize{pygeox:pygeox.add_objects.AddNamespace.equilateral_triangle}}
\pysigstartsignatures
\pysiglinewithargsret
{\sphinxbfcode{\sphinxupquote{equilateral\_triangle}}}
{\sphinxparam{\DUrole{n}{p1}\DUrole{o}{=}\DUrole{default_value}{None}}\sphinxparamcomma \sphinxparam{\DUrole{n}{p2}\DUrole{o}{=}\DUrole{default_value}{None}}\sphinxparamcomma \sphinxparam{\DUrole{n}{p3}\DUrole{o}{=}\DUrole{default_value}{None}}\sphinxparamcomma \sphinxparam{\DUrole{n}{name}\DUrole{o}{=}\DUrole{default_value}{None}}\sphinxparamcomma \sphinxparam{\DUrole{n}{copy}\DUrole{o}{=}\DUrole{default_value}{None}}}
{}
\pysigstopsignatures
\sphinxAtStartPar
Add an equilateral triangle or return a copy.
\begin{quote}\begin{description}
\sphinxlineitem{Parameters}\begin{itemize}
\item {} 
\sphinxAtStartPar
\sphinxstyleliteralstrong{\sphinxupquote{p1}} (\sphinxstyleliteralemphasis{\sphinxupquote{Point}}\sphinxstyleliteralemphasis{\sphinxupquote{ | }}\sphinxstyleliteralemphasis{\sphinxupquote{None}}) \textendash{} Vertex A.

\item {} 
\sphinxAtStartPar
\sphinxstyleliteralstrong{\sphinxupquote{p2}} (\sphinxstyleliteralemphasis{\sphinxupquote{Point}}\sphinxstyleliteralemphasis{\sphinxupquote{ | }}\sphinxstyleliteralemphasis{\sphinxupquote{None}}) \textendash{} Vertex B.

\item {} 
\sphinxAtStartPar
\sphinxstyleliteralstrong{\sphinxupquote{p3}} (\sphinxstyleliteralemphasis{\sphinxupquote{Point}}\sphinxstyleliteralemphasis{\sphinxupquote{ | }}\sphinxstyleliteralemphasis{\sphinxupquote{None}}) \textendash{} Vertex C.

\item {} 
\sphinxAtStartPar
\sphinxstyleliteralstrong{\sphinxupquote{name}} (\sphinxstyleliteralemphasis{\sphinxupquote{str}}\sphinxstyleliteralemphasis{\sphinxupquote{ | }}\sphinxstyleliteralemphasis{\sphinxupquote{None}}) \textendash{} Optional name.

\item {} 
\sphinxAtStartPar
\sphinxstyleliteralstrong{\sphinxupquote{copy}} (\sphinxstyleliteralemphasis{\sphinxupquote{EquilateralTriangle}}\sphinxstyleliteralemphasis{\sphinxupquote{ | }}\sphinxstyleliteralemphasis{\sphinxupquote{None}}) \textendash{} Existing triangle to reuse.

\end{itemize}

\sphinxlineitem{Returns}
\sphinxAtStartPar
The triangle instance (copy if provided).

\sphinxlineitem{Return type}
\sphinxAtStartPar
EquilateralTriangle

\end{description}\end{quote}
\subsubsection*{Notes}

\sphinxAtStartPar
Three distinct points are required; all sides are constrained equal.

\end{fulllineitems}

\index{isosceles\_trapezoid() (pygeox.add\_objects.AddNamespace method)@\spxentry{isosceles\_trapezoid()}\spxextra{pygeox.add\_objects.AddNamespace method}}

\begin{fulllineitems}
\phantomsection\label{\detokenize{pygeox:pygeox.add_objects.AddNamespace.isosceles_trapezoid}}
\pysigstartsignatures
\pysiglinewithargsret
{\sphinxbfcode{\sphinxupquote{isosceles\_trapezoid}}}
{\sphinxparam{\DUrole{n}{p1}\DUrole{o}{=}\DUrole{default_value}{None}}\sphinxparamcomma \sphinxparam{\DUrole{n}{p2}\DUrole{o}{=}\DUrole{default_value}{None}}\sphinxparamcomma \sphinxparam{\DUrole{n}{p3}\DUrole{o}{=}\DUrole{default_value}{None}}\sphinxparamcomma \sphinxparam{\DUrole{n}{p4}\DUrole{o}{=}\DUrole{default_value}{None}}\sphinxparamcomma \sphinxparam{\DUrole{n}{name}\DUrole{o}{=}\DUrole{default_value}{None}}\sphinxparamcomma \sphinxparam{\DUrole{n}{force\_convex}\DUrole{o}{=}\DUrole{default_value}{False}}\sphinxparamcomma \sphinxparam{\DUrole{n}{copy}\DUrole{o}{=}\DUrole{default_value}{None}}}
{}
\pysigstopsignatures
\sphinxAtStartPar
Add an isosceles trapezoid or return a copy.

\sphinxAtStartPar
Constrains AB to be parallel to CD and BC = DA. Optional convexity enforcement.
\begin{quote}\begin{description}
\sphinxlineitem{Parameters}\begin{itemize}
\item {} 
\sphinxAtStartPar
\sphinxstyleliteralstrong{\sphinxupquote{p1}} (\sphinxstyleliteralemphasis{\sphinxupquote{Point}}\sphinxstyleliteralemphasis{\sphinxupquote{ | }}\sphinxstyleliteralemphasis{\sphinxupquote{None}}) \textendash{} Vertex A.

\item {} 
\sphinxAtStartPar
\sphinxstyleliteralstrong{\sphinxupquote{p2}} (\sphinxstyleliteralemphasis{\sphinxupquote{Point}}\sphinxstyleliteralemphasis{\sphinxupquote{ | }}\sphinxstyleliteralemphasis{\sphinxupquote{None}}) \textendash{} Vertex B.

\item {} 
\sphinxAtStartPar
\sphinxstyleliteralstrong{\sphinxupquote{p3}} (\sphinxstyleliteralemphasis{\sphinxupquote{Point}}\sphinxstyleliteralemphasis{\sphinxupquote{ | }}\sphinxstyleliteralemphasis{\sphinxupquote{None}}) \textendash{} Vertex C.

\item {} 
\sphinxAtStartPar
\sphinxstyleliteralstrong{\sphinxupquote{p4}} (\sphinxstyleliteralemphasis{\sphinxupquote{Point}}\sphinxstyleliteralemphasis{\sphinxupquote{ | }}\sphinxstyleliteralemphasis{\sphinxupquote{None}}) \textendash{} Vertex D.

\item {} 
\sphinxAtStartPar
\sphinxstyleliteralstrong{\sphinxupquote{name}} (\sphinxstyleliteralemphasis{\sphinxupquote{str}}\sphinxstyleliteralemphasis{\sphinxupquote{ | }}\sphinxstyleliteralemphasis{\sphinxupquote{None}}) \textendash{} Optional name.

\item {} 
\sphinxAtStartPar
\sphinxstyleliteralstrong{\sphinxupquote{force\_convex}} (\sphinxstyleliteralemphasis{\sphinxupquote{bool}}) \textendash{} If \sphinxcode{\sphinxupquote{True}}, enforces convexity and simplicity.

\item {} 
\sphinxAtStartPar
\sphinxstyleliteralstrong{\sphinxupquote{copy}} (\sphinxstyleliteralemphasis{\sphinxupquote{IsoscelesTrapezoid}}\sphinxstyleliteralemphasis{\sphinxupquote{ | }}\sphinxstyleliteralemphasis{\sphinxupquote{None}}) \textendash{} Existing trapezoid to reuse.

\end{itemize}

\sphinxlineitem{Returns}
\sphinxAtStartPar
The quadrilateral instance (copy if provided).

\sphinxlineitem{Return type}
\sphinxAtStartPar
IsoscelesTrapezoid

\end{description}\end{quote}

\end{fulllineitems}

\index{isosceles\_triangle() (pygeox.add\_objects.AddNamespace method)@\spxentry{isosceles\_triangle()}\spxextra{pygeox.add\_objects.AddNamespace method}}

\begin{fulllineitems}
\phantomsection\label{\detokenize{pygeox:pygeox.add_objects.AddNamespace.isosceles_triangle}}
\pysigstartsignatures
\pysiglinewithargsret
{\sphinxbfcode{\sphinxupquote{isosceles\_triangle}}}
{\sphinxparam{\DUrole{n}{p1}\DUrole{o}{=}\DUrole{default_value}{None}}\sphinxparamcomma \sphinxparam{\DUrole{n}{p2}\DUrole{o}{=}\DUrole{default_value}{None}}\sphinxparamcomma \sphinxparam{\DUrole{n}{p3}\DUrole{o}{=}\DUrole{default_value}{None}}\sphinxparamcomma \sphinxparam{\DUrole{n}{name}\DUrole{o}{=}\DUrole{default_value}{None}}\sphinxparamcomma \sphinxparam{\DUrole{n}{copy}\DUrole{o}{=}\DUrole{default_value}{None}}}
{}
\pysigstopsignatures
\sphinxAtStartPar
Add an isosceles triangle or return a copy.
\begin{quote}\begin{description}
\sphinxlineitem{Parameters}\begin{itemize}
\item {} 
\sphinxAtStartPar
\sphinxstyleliteralstrong{\sphinxupquote{p1}} (\sphinxstyleliteralemphasis{\sphinxupquote{Point}}\sphinxstyleliteralemphasis{\sphinxupquote{ | }}\sphinxstyleliteralemphasis{\sphinxupquote{None}}) \textendash{} Vertex A.

\item {} 
\sphinxAtStartPar
\sphinxstyleliteralstrong{\sphinxupquote{p2}} (\sphinxstyleliteralemphasis{\sphinxupquote{Point}}\sphinxstyleliteralemphasis{\sphinxupquote{ | }}\sphinxstyleliteralemphasis{\sphinxupquote{None}}) \textendash{} Vertex B (apex for equal sides AB and BC).

\item {} 
\sphinxAtStartPar
\sphinxstyleliteralstrong{\sphinxupquote{p3}} (\sphinxstyleliteralemphasis{\sphinxupquote{Point}}\sphinxstyleliteralemphasis{\sphinxupquote{ | }}\sphinxstyleliteralemphasis{\sphinxupquote{None}}) \textendash{} Vertex C.

\item {} 
\sphinxAtStartPar
\sphinxstyleliteralstrong{\sphinxupquote{name}} (\sphinxstyleliteralemphasis{\sphinxupquote{str}}\sphinxstyleliteralemphasis{\sphinxupquote{ | }}\sphinxstyleliteralemphasis{\sphinxupquote{None}}) \textendash{} Optional name.

\item {} 
\sphinxAtStartPar
\sphinxstyleliteralstrong{\sphinxupquote{copy}} (\sphinxstyleliteralemphasis{\sphinxupquote{IsoscelesTriangle}}\sphinxstyleliteralemphasis{\sphinxupquote{ | }}\sphinxstyleliteralemphasis{\sphinxupquote{None}}) \textendash{} Existing triangle to reuse.

\end{itemize}

\sphinxlineitem{Returns}
\sphinxAtStartPar
The triangle instance (copy if provided).

\sphinxlineitem{Return type}
\sphinxAtStartPar
IsoscelesTriangle

\end{description}\end{quote}
\subsubsection*{Notes}

\sphinxAtStartPar
Sides AB and BC are constrained equal (order matters).

\end{fulllineitems}

\index{kite() (pygeox.add\_objects.AddNamespace method)@\spxentry{kite()}\spxextra{pygeox.add\_objects.AddNamespace method}}

\begin{fulllineitems}
\phantomsection\label{\detokenize{pygeox:pygeox.add_objects.AddNamespace.kite}}
\pysigstartsignatures
\pysiglinewithargsret
{\sphinxbfcode{\sphinxupquote{kite}}}
{\sphinxparam{\DUrole{n}{p1}\DUrole{o}{=}\DUrole{default_value}{None}}\sphinxparamcomma \sphinxparam{\DUrole{n}{p2}\DUrole{o}{=}\DUrole{default_value}{None}}\sphinxparamcomma \sphinxparam{\DUrole{n}{p3}\DUrole{o}{=}\DUrole{default_value}{None}}\sphinxparamcomma \sphinxparam{\DUrole{n}{p4}\DUrole{o}{=}\DUrole{default_value}{None}}\sphinxparamcomma \sphinxparam{\DUrole{n}{name}\DUrole{o}{=}\DUrole{default_value}{None}}\sphinxparamcomma \sphinxparam{\DUrole{n}{force\_convex}\DUrole{o}{=}\DUrole{default_value}{False}}\sphinxparamcomma \sphinxparam{\DUrole{n}{copy}\DUrole{o}{=}\DUrole{default_value}{None}}}
{}
\pysigstopsignatures
\sphinxAtStartPar
Add a kite or return a copy.

\sphinxAtStartPar
Constrains AB = BC and CD = DA. Optional convexity enforcement.
\begin{quote}\begin{description}
\sphinxlineitem{Parameters}\begin{itemize}
\item {} 
\sphinxAtStartPar
\sphinxstyleliteralstrong{\sphinxupquote{p1}} (\sphinxstyleliteralemphasis{\sphinxupquote{Point}}\sphinxstyleliteralemphasis{\sphinxupquote{ | }}\sphinxstyleliteralemphasis{\sphinxupquote{None}}) \textendash{} Vertex A.

\item {} 
\sphinxAtStartPar
\sphinxstyleliteralstrong{\sphinxupquote{p2}} (\sphinxstyleliteralemphasis{\sphinxupquote{Point}}\sphinxstyleliteralemphasis{\sphinxupquote{ | }}\sphinxstyleliteralemphasis{\sphinxupquote{None}}) \textendash{} Vertex B.

\item {} 
\sphinxAtStartPar
\sphinxstyleliteralstrong{\sphinxupquote{p3}} (\sphinxstyleliteralemphasis{\sphinxupquote{Point}}\sphinxstyleliteralemphasis{\sphinxupquote{ | }}\sphinxstyleliteralemphasis{\sphinxupquote{None}}) \textendash{} Vertex C.

\item {} 
\sphinxAtStartPar
\sphinxstyleliteralstrong{\sphinxupquote{p4}} (\sphinxstyleliteralemphasis{\sphinxupquote{Point}}\sphinxstyleliteralemphasis{\sphinxupquote{ | }}\sphinxstyleliteralemphasis{\sphinxupquote{None}}) \textendash{} Vertex D.

\item {} 
\sphinxAtStartPar
\sphinxstyleliteralstrong{\sphinxupquote{name}} (\sphinxstyleliteralemphasis{\sphinxupquote{str}}\sphinxstyleliteralemphasis{\sphinxupquote{ | }}\sphinxstyleliteralemphasis{\sphinxupquote{None}}) \textendash{} Optional name.

\item {} 
\sphinxAtStartPar
\sphinxstyleliteralstrong{\sphinxupquote{force\_convex}} (\sphinxstyleliteralemphasis{\sphinxupquote{bool}}) \textendash{} If \sphinxcode{\sphinxupquote{True}}, enforces convexity and simplicity.

\item {} 
\sphinxAtStartPar
\sphinxstyleliteralstrong{\sphinxupquote{copy}} (\sphinxstyleliteralemphasis{\sphinxupquote{Kite}}\sphinxstyleliteralemphasis{\sphinxupquote{ | }}\sphinxstyleliteralemphasis{\sphinxupquote{None}}) \textendash{} Existing kite to reuse.

\end{itemize}

\sphinxlineitem{Returns}
\sphinxAtStartPar
The quadrilateral instance (copy if provided).

\sphinxlineitem{Return type}
\sphinxAtStartPar
Kite

\end{description}\end{quote}

\end{fulllineitems}

\index{line() (pygeox.add\_objects.AddNamespace method)@\spxentry{line()}\spxextra{pygeox.add\_objects.AddNamespace method}}

\begin{fulllineitems}
\phantomsection\label{\detokenize{pygeox:pygeox.add_objects.AddNamespace.line}}
\pysigstartsignatures
\pysiglinewithargsret
{\sphinxbfcode{\sphinxupquote{line}}}
{\sphinxparam{\DUrole{n}{point1}\DUrole{o}{=}\DUrole{default_value}{None}}\sphinxparamcomma \sphinxparam{\DUrole{n}{point2}\DUrole{o}{=}\DUrole{default_value}{None}}\sphinxparamcomma \sphinxparam{\DUrole{n}{name}\DUrole{o}{=}\DUrole{default_value}{None}}\sphinxparamcomma \sphinxparam{\DUrole{n}{copy}\DUrole{o}{=}\DUrole{default_value}{None}}}
{}
\pysigstopsignatures
\sphinxAtStartPar
Add a line or copy an existing one.

\sphinxAtStartPar
If \sphinxcode{\sphinxupquote{copy}} is provided and \sphinxcode{\sphinxupquote{point1}}/\sphinxcode{\sphinxupquote{point2}} are given,
their coordinates are overwritten to match the copy. A default name
is inferred if none is given.
\begin{quote}\begin{description}
\sphinxlineitem{Parameters}\begin{itemize}
\item {} 
\sphinxAtStartPar
\sphinxstyleliteralstrong{\sphinxupquote{point1}} (\sphinxstyleliteralemphasis{\sphinxupquote{Point}}\sphinxstyleliteralemphasis{\sphinxupquote{ | }}\sphinxstyleliteralemphasis{\sphinxupquote{None}}) \textendash{} First point on the line.

\item {} 
\sphinxAtStartPar
\sphinxstyleliteralstrong{\sphinxupquote{point2}} (\sphinxstyleliteralemphasis{\sphinxupquote{Point}}\sphinxstyleliteralemphasis{\sphinxupquote{ | }}\sphinxstyleliteralemphasis{\sphinxupquote{None}}) \textendash{} Second point on the line.

\item {} 
\sphinxAtStartPar
\sphinxstyleliteralstrong{\sphinxupquote{name}} (\sphinxstyleliteralemphasis{\sphinxupquote{str}}\sphinxstyleliteralemphasis{\sphinxupquote{ | }}\sphinxstyleliteralemphasis{\sphinxupquote{None}}) \textendash{} Optional name for the line.

\item {} 
\sphinxAtStartPar
\sphinxstyleliteralstrong{\sphinxupquote{copy}} (\sphinxstyleliteralemphasis{\sphinxupquote{Line}}\sphinxstyleliteralemphasis{\sphinxupquote{ | }}\sphinxstyleliteralemphasis{\sphinxupquote{None}}) \textendash{} Existing line to duplicate.

\end{itemize}

\sphinxlineitem{Returns}
\sphinxAtStartPar
The created line, registered in the scene.

\sphinxlineitem{Return type}
\sphinxAtStartPar
Line

\sphinxlineitem{Raises}
\sphinxAtStartPar
\sphinxstyleliteralstrong{\sphinxupquote{ValueError}} \textendash{} If not copying and fewer than two points are provided.

\end{description}\end{quote}

\end{fulllineitems}

\index{line\_segment() (pygeox.add\_objects.AddNamespace method)@\spxentry{line\_segment()}\spxextra{pygeox.add\_objects.AddNamespace method}}

\begin{fulllineitems}
\phantomsection\label{\detokenize{pygeox:pygeox.add_objects.AddNamespace.line_segment}}
\pysigstartsignatures
\pysiglinewithargsret
{\sphinxbfcode{\sphinxupquote{line\_segment}}}
{\sphinxparam{\DUrole{n}{point1}\DUrole{o}{=}\DUrole{default_value}{None}}\sphinxparamcomma \sphinxparam{\DUrole{n}{point2}\DUrole{o}{=}\DUrole{default_value}{None}}\sphinxparamcomma \sphinxparam{\DUrole{n}{name}\DUrole{o}{=}\DUrole{default_value}{None}}\sphinxparamcomma \sphinxparam{\DUrole{n}{copy}\DUrole{o}{=}\DUrole{default_value}{None}}}
{}
\pysigstopsignatures
\sphinxAtStartPar
Add a line segment or copy an existing one.
\begin{quote}\begin{description}
\sphinxlineitem{Parameters}\begin{itemize}
\item {} 
\sphinxAtStartPar
\sphinxstyleliteralstrong{\sphinxupquote{point1}} (\sphinxstyleliteralemphasis{\sphinxupquote{Point}}\sphinxstyleliteralemphasis{\sphinxupquote{ | }}\sphinxstyleliteralemphasis{\sphinxupquote{None}}) \textendash{} First endpoint.

\item {} 
\sphinxAtStartPar
\sphinxstyleliteralstrong{\sphinxupquote{point2}} (\sphinxstyleliteralemphasis{\sphinxupquote{Point}}\sphinxstyleliteralemphasis{\sphinxupquote{ | }}\sphinxstyleliteralemphasis{\sphinxupquote{None}}) \textendash{} Second endpoint.

\item {} 
\sphinxAtStartPar
\sphinxstyleliteralstrong{\sphinxupquote{name}} (\sphinxstyleliteralemphasis{\sphinxupquote{str}}\sphinxstyleliteralemphasis{\sphinxupquote{ | }}\sphinxstyleliteralemphasis{\sphinxupquote{None}}) \textendash{} Optional name for the segment.

\item {} 
\sphinxAtStartPar
\sphinxstyleliteralstrong{\sphinxupquote{copy}} (\sphinxstyleliteralemphasis{\sphinxupquote{LineSegment}}\sphinxstyleliteralemphasis{\sphinxupquote{ | }}\sphinxstyleliteralemphasis{\sphinxupquote{None}}) \textendash{} Existing segment to duplicate.

\end{itemize}

\sphinxlineitem{Returns}
\sphinxAtStartPar
The created segment, registered in the scene.

\sphinxlineitem{Return type}
\sphinxAtStartPar
LineSegment

\sphinxlineitem{Raises}
\sphinxAtStartPar
\sphinxstyleliteralstrong{\sphinxupquote{ValueError}} \textendash{} If not copying and fewer than two points are provided.

\end{description}\end{quote}

\end{fulllineitems}

\index{major\_arc() (pygeox.add\_objects.AddNamespace method)@\spxentry{major\_arc()}\spxextra{pygeox.add\_objects.AddNamespace method}}

\begin{fulllineitems}
\phantomsection\label{\detokenize{pygeox:pygeox.add_objects.AddNamespace.major_arc}}
\pysigstartsignatures
\pysiglinewithargsret
{\sphinxbfcode{\sphinxupquote{major\_arc}}}
{\sphinxparam{\DUrole{n}{center}}\sphinxparamcomma \sphinxparam{\DUrole{n}{start\_point}}\sphinxparamcomma \sphinxparam{\DUrole{n}{end\_point}}\sphinxparamcomma \sphinxparam{\DUrole{n}{name}\DUrole{o}{=}\DUrole{default_value}{None}}\sphinxparamcomma \sphinxparam{\DUrole{n}{on\_circle}\DUrole{o}{=}\DUrole{default_value}{None}}}
{}
\pysigstopsignatures
\sphinxAtStartPar
Create a \sphinxcode{\sphinxupquote{MajorArc}} and register it in the scene.

\sphinxAtStartPar
Adds equal\sphinxhyphen{}radius constraints and optionally enforces that the arc
lies on a given circle (same center and radius).
\begin{quote}\begin{description}
\sphinxlineitem{Parameters}\begin{itemize}
\item {} 
\sphinxAtStartPar
\sphinxstyleliteralstrong{\sphinxupquote{center}} (\sphinxstyleliteralemphasis{\sphinxupquote{Point}}) \textendash{} Center of the circle.

\item {} 
\sphinxAtStartPar
\sphinxstyleliteralstrong{\sphinxupquote{start\_point}} (\sphinxstyleliteralemphasis{\sphinxupquote{Point}}) \textendash{} Start point of the arc.

\item {} 
\sphinxAtStartPar
\sphinxstyleliteralstrong{\sphinxupquote{end\_point}} (\sphinxstyleliteralemphasis{\sphinxupquote{Point}}) \textendash{} End point of the arc.

\item {} 
\sphinxAtStartPar
\sphinxstyleliteralstrong{\sphinxupquote{name}} (\sphinxstyleliteralemphasis{\sphinxupquote{str}}\sphinxstyleliteralemphasis{\sphinxupquote{ | }}\sphinxstyleliteralemphasis{\sphinxupquote{None}}) \textendash{} Optional name. If \sphinxcode{\sphinxupquote{None}}, one is generated.

\item {} 
\sphinxAtStartPar
\sphinxstyleliteralstrong{\sphinxupquote{on\_circle}} (\sphinxstyleliteralemphasis{\sphinxupquote{Circle}}\sphinxstyleliteralemphasis{\sphinxupquote{ | }}\sphinxstyleliteralemphasis{\sphinxupquote{None}}) \textendash{} If provided, enforces arc radius equals the circle’s radius.

\end{itemize}

\sphinxlineitem{Returns}
\sphinxAtStartPar
The created arc.

\sphinxlineitem{Return type}
\sphinxAtStartPar
MajorArc

\end{description}\end{quote}

\end{fulllineitems}

\index{minor\_arc() (pygeox.add\_objects.AddNamespace method)@\spxentry{minor\_arc()}\spxextra{pygeox.add\_objects.AddNamespace method}}

\begin{fulllineitems}
\phantomsection\label{\detokenize{pygeox:pygeox.add_objects.AddNamespace.minor_arc}}
\pysigstartsignatures
\pysiglinewithargsret
{\sphinxbfcode{\sphinxupquote{minor\_arc}}}
{\sphinxparam{\DUrole{n}{center}}\sphinxparamcomma \sphinxparam{\DUrole{n}{start\_point}}\sphinxparamcomma \sphinxparam{\DUrole{n}{end\_point}}\sphinxparamcomma \sphinxparam{\DUrole{n}{name}\DUrole{o}{=}\DUrole{default_value}{None}}\sphinxparamcomma \sphinxparam{\DUrole{n}{on\_circle}\DUrole{o}{=}\DUrole{default_value}{None}}}
{}
\pysigstopsignatures
\sphinxAtStartPar
Create a \sphinxcode{\sphinxupquote{MinorArc}} and register it in the scene.

\sphinxAtStartPar
Adds equal\sphinxhyphen{}radius constraints and optionally enforces that the arc
lies on a given circle (same center and radius).
\begin{quote}\begin{description}
\sphinxlineitem{Parameters}\begin{itemize}
\item {} 
\sphinxAtStartPar
\sphinxstyleliteralstrong{\sphinxupquote{center}} (\sphinxstyleliteralemphasis{\sphinxupquote{Point}}) \textendash{} Center of the circle.

\item {} 
\sphinxAtStartPar
\sphinxstyleliteralstrong{\sphinxupquote{start\_point}} (\sphinxstyleliteralemphasis{\sphinxupquote{Point}}) \textendash{} Start point of the arc.

\item {} 
\sphinxAtStartPar
\sphinxstyleliteralstrong{\sphinxupquote{end\_point}} (\sphinxstyleliteralemphasis{\sphinxupquote{Point}}) \textendash{} End point of the arc.

\item {} 
\sphinxAtStartPar
\sphinxstyleliteralstrong{\sphinxupquote{name}} (\sphinxstyleliteralemphasis{\sphinxupquote{str}}\sphinxstyleliteralemphasis{\sphinxupquote{ | }}\sphinxstyleliteralemphasis{\sphinxupquote{None}}) \textendash{} Optional name. If \sphinxcode{\sphinxupquote{None}}, one is generated.

\item {} 
\sphinxAtStartPar
\sphinxstyleliteralstrong{\sphinxupquote{on\_circle}} (\sphinxstyleliteralemphasis{\sphinxupquote{Circle}}\sphinxstyleliteralemphasis{\sphinxupquote{ | }}\sphinxstyleliteralemphasis{\sphinxupquote{None}}) \textendash{} If provided, enforces arc radius equals the circle’s radius.

\end{itemize}

\sphinxlineitem{Returns}
\sphinxAtStartPar
The created arc.

\sphinxlineitem{Return type}
\sphinxAtStartPar
MinorArc

\end{description}\end{quote}

\end{fulllineitems}

\index{obtuse\_triangle() (pygeox.add\_objects.AddNamespace method)@\spxentry{obtuse\_triangle()}\spxextra{pygeox.add\_objects.AddNamespace method}}

\begin{fulllineitems}
\phantomsection\label{\detokenize{pygeox:pygeox.add_objects.AddNamespace.obtuse_triangle}}
\pysigstartsignatures
\pysiglinewithargsret
{\sphinxbfcode{\sphinxupquote{obtuse\_triangle}}}
{\sphinxparam{\DUrole{n}{p1}\DUrole{o}{=}\DUrole{default_value}{None}}\sphinxparamcomma \sphinxparam{\DUrole{n}{p2}\DUrole{o}{=}\DUrole{default_value}{None}}\sphinxparamcomma \sphinxparam{\DUrole{n}{p3}\DUrole{o}{=}\DUrole{default_value}{None}}\sphinxparamcomma \sphinxparam{\DUrole{n}{name}\DUrole{o}{=}\DUrole{default_value}{None}}\sphinxparamcomma \sphinxparam{\DUrole{n}{copy}\DUrole{o}{=}\DUrole{default_value}{None}}}
{}
\pysigstopsignatures
\sphinxAtStartPar
Add an obtuse triangle or return a copy.
\begin{quote}\begin{description}
\sphinxlineitem{Parameters}\begin{itemize}
\item {} 
\sphinxAtStartPar
\sphinxstyleliteralstrong{\sphinxupquote{p1}} (\sphinxstyleliteralemphasis{\sphinxupquote{Point}}\sphinxstyleliteralemphasis{\sphinxupquote{ | }}\sphinxstyleliteralemphasis{\sphinxupquote{None}}) \textendash{} Vertex A.

\item {} 
\sphinxAtStartPar
\sphinxstyleliteralstrong{\sphinxupquote{p2}} (\sphinxstyleliteralemphasis{\sphinxupquote{Point}}\sphinxstyleliteralemphasis{\sphinxupquote{ | }}\sphinxstyleliteralemphasis{\sphinxupquote{None}}) \textendash{} Vertex B.

\item {} 
\sphinxAtStartPar
\sphinxstyleliteralstrong{\sphinxupquote{p3}} (\sphinxstyleliteralemphasis{\sphinxupquote{Point}}\sphinxstyleliteralemphasis{\sphinxupquote{ | }}\sphinxstyleliteralemphasis{\sphinxupquote{None}}) \textendash{} Vertex C.

\item {} 
\sphinxAtStartPar
\sphinxstyleliteralstrong{\sphinxupquote{name}} (\sphinxstyleliteralemphasis{\sphinxupquote{str}}\sphinxstyleliteralemphasis{\sphinxupquote{ | }}\sphinxstyleliteralemphasis{\sphinxupquote{None}}) \textendash{} Optional name.

\item {} 
\sphinxAtStartPar
\sphinxstyleliteralstrong{\sphinxupquote{copy}} (\sphinxstyleliteralemphasis{\sphinxupquote{ObtuseTriangle}}\sphinxstyleliteralemphasis{\sphinxupquote{ | }}\sphinxstyleliteralemphasis{\sphinxupquote{None}}) \textendash{} Existing triangle to reuse.

\end{itemize}

\sphinxlineitem{Returns}
\sphinxAtStartPar
The triangle instance (copy if provided).

\sphinxlineitem{Return type}
\sphinxAtStartPar
ObtuseTriangle

\end{description}\end{quote}
\subsubsection*{Notes}

\sphinxAtStartPar
At least one internal angle is constrained \textgreater{} 90°.

\end{fulllineitems}

\index{parallelogram() (pygeox.add\_objects.AddNamespace method)@\spxentry{parallelogram()}\spxextra{pygeox.add\_objects.AddNamespace method}}

\begin{fulllineitems}
\phantomsection\label{\detokenize{pygeox:pygeox.add_objects.AddNamespace.parallelogram}}
\pysigstartsignatures
\pysiglinewithargsret
{\sphinxbfcode{\sphinxupquote{parallelogram}}}
{\sphinxparam{\DUrole{n}{p1}\DUrole{o}{=}\DUrole{default_value}{None}}\sphinxparamcomma \sphinxparam{\DUrole{n}{p2}\DUrole{o}{=}\DUrole{default_value}{None}}\sphinxparamcomma \sphinxparam{\DUrole{n}{p3}\DUrole{o}{=}\DUrole{default_value}{None}}\sphinxparamcomma \sphinxparam{\DUrole{n}{p4}\DUrole{o}{=}\DUrole{default_value}{None}}\sphinxparamcomma \sphinxparam{\DUrole{n}{name}\DUrole{o}{=}\DUrole{default_value}{None}}\sphinxparamcomma \sphinxparam{\DUrole{n}{force\_convex}\DUrole{o}{=}\DUrole{default_value}{False}}\sphinxparamcomma \sphinxparam{\DUrole{n}{copy}\DUrole{o}{=}\DUrole{default_value}{None}}}
{}
\pysigstopsignatures
\sphinxAtStartPar
Add a parallelogram or return a copy.

\sphinxAtStartPar
Opposite sides are constrained parallel. Optional convexity enforcement.
\begin{quote}\begin{description}
\sphinxlineitem{Parameters}\begin{itemize}
\item {} 
\sphinxAtStartPar
\sphinxstyleliteralstrong{\sphinxupquote{p1}} (\sphinxstyleliteralemphasis{\sphinxupquote{Point}}\sphinxstyleliteralemphasis{\sphinxupquote{ | }}\sphinxstyleliteralemphasis{\sphinxupquote{None}}) \textendash{} Vertex A.

\item {} 
\sphinxAtStartPar
\sphinxstyleliteralstrong{\sphinxupquote{p2}} (\sphinxstyleliteralemphasis{\sphinxupquote{Point}}\sphinxstyleliteralemphasis{\sphinxupquote{ | }}\sphinxstyleliteralemphasis{\sphinxupquote{None}}) \textendash{} Vertex B.

\item {} 
\sphinxAtStartPar
\sphinxstyleliteralstrong{\sphinxupquote{p3}} (\sphinxstyleliteralemphasis{\sphinxupquote{Point}}\sphinxstyleliteralemphasis{\sphinxupquote{ | }}\sphinxstyleliteralemphasis{\sphinxupquote{None}}) \textendash{} Vertex C.

\item {} 
\sphinxAtStartPar
\sphinxstyleliteralstrong{\sphinxupquote{p4}} (\sphinxstyleliteralemphasis{\sphinxupquote{Point}}\sphinxstyleliteralemphasis{\sphinxupquote{ | }}\sphinxstyleliteralemphasis{\sphinxupquote{None}}) \textendash{} Vertex D.

\item {} 
\sphinxAtStartPar
\sphinxstyleliteralstrong{\sphinxupquote{name}} (\sphinxstyleliteralemphasis{\sphinxupquote{str}}\sphinxstyleliteralemphasis{\sphinxupquote{ | }}\sphinxstyleliteralemphasis{\sphinxupquote{None}}) \textendash{} Optional name.

\item {} 
\sphinxAtStartPar
\sphinxstyleliteralstrong{\sphinxupquote{force\_convex}} (\sphinxstyleliteralemphasis{\sphinxupquote{bool}}) \textendash{} If \sphinxcode{\sphinxupquote{True}}, enforces convexity and simplicity.

\item {} 
\sphinxAtStartPar
\sphinxstyleliteralstrong{\sphinxupquote{copy}} (\sphinxstyleliteralemphasis{\sphinxupquote{Parallelogram}}\sphinxstyleliteralemphasis{\sphinxupquote{ | }}\sphinxstyleliteralemphasis{\sphinxupquote{None}}) \textendash{} Existing parallelogram to reuse.

\end{itemize}

\sphinxlineitem{Returns}
\sphinxAtStartPar
The quadrilateral instance (copy if provided).

\sphinxlineitem{Return type}
\sphinxAtStartPar
Parallelogram

\end{description}\end{quote}

\end{fulllineitems}

\index{point() (pygeox.add\_objects.AddNamespace method)@\spxentry{point()}\spxextra{pygeox.add\_objects.AddNamespace method}}

\begin{fulllineitems}
\phantomsection\label{\detokenize{pygeox:pygeox.add_objects.AddNamespace.point}}
\pysigstartsignatures
\pysiglinewithargsret
{\sphinxbfcode{\sphinxupquote{point}}}
{\sphinxparam{\DUrole{n}{name}\DUrole{o}{=}\DUrole{default_value}{None}}\sphinxparamcomma \sphinxparam{\DUrole{n}{x}\DUrole{o}{=}\DUrole{default_value}{None}}\sphinxparamcomma \sphinxparam{\DUrole{n}{y}\DUrole{o}{=}\DUrole{default_value}{None}}\sphinxparamcomma \sphinxparam{\DUrole{n}{copy}\DUrole{o}{=}\DUrole{default_value}{None}}}
{}
\pysigstopsignatures
\sphinxAtStartPar
Add a point or copy an existing one.

\sphinxAtStartPar
If \sphinxcode{\sphinxupquote{copy}} is provided, its coordinates and name are used unless
overridden. If \sphinxcode{\sphinxupquote{x}} or \sphinxcode{\sphinxupquote{y}} is \sphinxcode{\sphinxupquote{None}} (and not copying), a
\sphinxcode{\sphinxupquote{sympy.Symbol}} is created for that coordinate, using any
domain set on the scene.
\begin{quote}\begin{description}
\sphinxlineitem{Parameters}\begin{itemize}
\item {} 
\sphinxAtStartPar
\sphinxstyleliteralstrong{\sphinxupquote{name}} (\sphinxstyleliteralemphasis{\sphinxupquote{str}}) \textendash{} Optional name for the point.

\item {} 
\sphinxAtStartPar
\sphinxstyleliteralstrong{\sphinxupquote{x}} (\sphinxstyleliteralemphasis{\sphinxupquote{int}}\sphinxstyleliteralemphasis{\sphinxupquote{ | }}\sphinxstyleliteralemphasis{\sphinxupquote{float}}\sphinxstyleliteralemphasis{\sphinxupquote{ | }}\sphinxstyleliteralemphasis{\sphinxupquote{Expr}}\sphinxstyleliteralemphasis{\sphinxupquote{ | }}\sphinxstyleliteralemphasis{\sphinxupquote{None}}) \textendash{} X coordinate (number or SymPy expression). If \sphinxcode{\sphinxupquote{None}}, a symbol is created.

\item {} 
\sphinxAtStartPar
\sphinxstyleliteralstrong{\sphinxupquote{y}} (\sphinxstyleliteralemphasis{\sphinxupquote{int}}\sphinxstyleliteralemphasis{\sphinxupquote{ | }}\sphinxstyleliteralemphasis{\sphinxupquote{float}}\sphinxstyleliteralemphasis{\sphinxupquote{ | }}\sphinxstyleliteralemphasis{\sphinxupquote{Expr}}\sphinxstyleliteralemphasis{\sphinxupquote{ | }}\sphinxstyleliteralemphasis{\sphinxupquote{None}}) \textendash{} Y coordinate (number or SymPy expression). If \sphinxcode{\sphinxupquote{None}}, a symbol is created.

\item {} 
\sphinxAtStartPar
\sphinxstyleliteralstrong{\sphinxupquote{copy}} (\sphinxstyleliteralemphasis{\sphinxupquote{Point}}\sphinxstyleliteralemphasis{\sphinxupquote{ | }}\sphinxstyleliteralemphasis{\sphinxupquote{None}}) \textendash{} Existing point to duplicate.

\end{itemize}

\sphinxlineitem{Returns}
\sphinxAtStartPar
The created point, registered in the scene.

\sphinxlineitem{Return type}
\sphinxAtStartPar
Point

\end{description}\end{quote}

\end{fulllineitems}

\index{points() (pygeox.add\_objects.AddNamespace method)@\spxentry{points()}\spxextra{pygeox.add\_objects.AddNamespace method}}

\begin{fulllineitems}
\phantomsection\label{\detokenize{pygeox:pygeox.add_objects.AddNamespace.points}}
\pysigstartsignatures
\pysiglinewithargsret
{\sphinxbfcode{\sphinxupquote{points}}}
{\sphinxparam{\DUrole{n}{names}}}
{}
\pysigstopsignatures
\sphinxAtStartPar
Add multiple points by name.
\begin{quote}\begin{description}
\sphinxlineitem{Parameters}
\sphinxAtStartPar
\sphinxstyleliteralstrong{\sphinxupquote{names}} (\sphinxstyleliteralemphasis{\sphinxupquote{list}}\sphinxstyleliteralemphasis{\sphinxupquote{{[}}}\sphinxstyleliteralemphasis{\sphinxupquote{str}}\sphinxstyleliteralemphasis{\sphinxupquote{{]}}}) \textendash{} List of point names.

\sphinxlineitem{Returns}
\sphinxAtStartPar
The created points.

\sphinxlineitem{Return type}
\sphinxAtStartPar
list{[}Point{]}

\end{description}\end{quote}

\end{fulllineitems}

\index{polygon() (pygeox.add\_objects.AddNamespace method)@\spxentry{polygon()}\spxextra{pygeox.add\_objects.AddNamespace method}}

\begin{fulllineitems}
\phantomsection\label{\detokenize{pygeox:pygeox.add_objects.AddNamespace.polygon}}
\pysigstartsignatures
\pysiglinewithargsret
{\sphinxbfcode{\sphinxupquote{polygon}}}
{\sphinxparam{\DUrole{o}{*}\DUrole{n}{points}}\sphinxparamcomma \sphinxparam{\DUrole{n}{name}\DUrole{o}{=}\DUrole{default_value}{None}}\sphinxparamcomma \sphinxparam{\DUrole{n}{copy}\DUrole{o}{=}\DUrole{default_value}{None}}}
{}
\pysigstopsignatures
\sphinxAtStartPar
Add a general polygon or return a copy.

\sphinxAtStartPar
Assumptions \& constraints (when creating):
at least three points, given in order. If \sphinxcode{\sphinxupquote{copy}} is provided and the
number of points matches, provided points are overwritten to match.
\begin{quote}\begin{description}
\sphinxlineitem{Parameters}\begin{itemize}
\item {} 
\sphinxAtStartPar
\sphinxstyleliteralstrong{\sphinxupquote{*points}} (\sphinxstyleliteralemphasis{\sphinxupquote{Point}}) \textendash{} Vertices in order.

\item {} 
\sphinxAtStartPar
\sphinxstyleliteralstrong{\sphinxupquote{name}} (\sphinxstyleliteralemphasis{\sphinxupquote{str}}\sphinxstyleliteralemphasis{\sphinxupquote{ | }}\sphinxstyleliteralemphasis{\sphinxupquote{None}}) \textendash{} Optional name.

\item {} 
\sphinxAtStartPar
\sphinxstyleliteralstrong{\sphinxupquote{copy}} (\sphinxstyleliteralemphasis{\sphinxupquote{Polygon}}\sphinxstyleliteralemphasis{\sphinxupquote{ | }}\sphinxstyleliteralemphasis{\sphinxupquote{None}}) \textendash{} Existing polygon to reuse.

\end{itemize}

\sphinxlineitem{Returns}
\sphinxAtStartPar
The polygon instance (copy if provided).

\sphinxlineitem{Return type}
\sphinxAtStartPar
Polygon

\end{description}\end{quote}

\end{fulllineitems}

\index{quadrilateral() (pygeox.add\_objects.AddNamespace method)@\spxentry{quadrilateral()}\spxextra{pygeox.add\_objects.AddNamespace method}}

\begin{fulllineitems}
\phantomsection\label{\detokenize{pygeox:pygeox.add_objects.AddNamespace.quadrilateral}}
\pysigstartsignatures
\pysiglinewithargsret
{\sphinxbfcode{\sphinxupquote{quadrilateral}}}
{\sphinxparam{\DUrole{n}{p1}\DUrole{o}{=}\DUrole{default_value}{None}}\sphinxparamcomma \sphinxparam{\DUrole{n}{p2}\DUrole{o}{=}\DUrole{default_value}{None}}\sphinxparamcomma \sphinxparam{\DUrole{n}{p3}\DUrole{o}{=}\DUrole{default_value}{None}}\sphinxparamcomma \sphinxparam{\DUrole{n}{p4}\DUrole{o}{=}\DUrole{default_value}{None}}\sphinxparamcomma \sphinxparam{\DUrole{n}{name}\DUrole{o}{=}\DUrole{default_value}{None}}\sphinxparamcomma \sphinxparam{\DUrole{n}{force\_convex}\DUrole{o}{=}\DUrole{default_value}{False}}\sphinxparamcomma \sphinxparam{\DUrole{n}{copy}\DUrole{o}{=}\DUrole{default_value}{None}}}
{}
\pysigstopsignatures
\sphinxAtStartPar
Add a quadrilateral or return a copy.
\begin{quote}\begin{description}
\sphinxlineitem{Parameters}\begin{itemize}
\item {} 
\sphinxAtStartPar
\sphinxstyleliteralstrong{\sphinxupquote{p1}} (\sphinxstyleliteralemphasis{\sphinxupquote{Point}}\sphinxstyleliteralemphasis{\sphinxupquote{ | }}\sphinxstyleliteralemphasis{\sphinxupquote{None}}) \textendash{} Vertex A.

\item {} 
\sphinxAtStartPar
\sphinxstyleliteralstrong{\sphinxupquote{p2}} (\sphinxstyleliteralemphasis{\sphinxupquote{Point}}\sphinxstyleliteralemphasis{\sphinxupquote{ | }}\sphinxstyleliteralemphasis{\sphinxupquote{None}}) \textendash{} Vertex B.

\item {} 
\sphinxAtStartPar
\sphinxstyleliteralstrong{\sphinxupquote{p3}} (\sphinxstyleliteralemphasis{\sphinxupquote{Point}}\sphinxstyleliteralemphasis{\sphinxupquote{ | }}\sphinxstyleliteralemphasis{\sphinxupquote{None}}) \textendash{} Vertex C.

\item {} 
\sphinxAtStartPar
\sphinxstyleliteralstrong{\sphinxupquote{p4}} (\sphinxstyleliteralemphasis{\sphinxupquote{Point}}\sphinxstyleliteralemphasis{\sphinxupquote{ | }}\sphinxstyleliteralemphasis{\sphinxupquote{None}}) \textendash{} Vertex D.

\item {} 
\sphinxAtStartPar
\sphinxstyleliteralstrong{\sphinxupquote{name}} (\sphinxstyleliteralemphasis{\sphinxupquote{str}}\sphinxstyleliteralemphasis{\sphinxupquote{ | }}\sphinxstyleliteralemphasis{\sphinxupquote{None}}) \textendash{} Optional name.

\item {} 
\sphinxAtStartPar
\sphinxstyleliteralstrong{\sphinxupquote{force\_convex}} (\sphinxstyleliteralemphasis{\sphinxupquote{bool}}) \textendash{} If \sphinxcode{\sphinxupquote{True}}, enforces convexity and simplicity.

\item {} 
\sphinxAtStartPar
\sphinxstyleliteralstrong{\sphinxupquote{copy}} (\sphinxstyleliteralemphasis{\sphinxupquote{Quadrilateral}}\sphinxstyleliteralemphasis{\sphinxupquote{ | }}\sphinxstyleliteralemphasis{\sphinxupquote{None}}) \textendash{} Existing quadrilateral to reuse.

\end{itemize}

\sphinxlineitem{Returns}
\sphinxAtStartPar
The quadrilateral instance (copy if provided).

\sphinxlineitem{Return type}
\sphinxAtStartPar
Quadrilateral

\end{description}\end{quote}

\end{fulllineitems}

\index{ray() (pygeox.add\_objects.AddNamespace method)@\spxentry{ray()}\spxextra{pygeox.add\_objects.AddNamespace method}}

\begin{fulllineitems}
\phantomsection\label{\detokenize{pygeox:pygeox.add_objects.AddNamespace.ray}}
\pysigstartsignatures
\pysiglinewithargsret
{\sphinxbfcode{\sphinxupquote{ray}}}
{\sphinxparam{\DUrole{n}{start\_point}\DUrole{o}{=}\DUrole{default_value}{None}}\sphinxparamcomma \sphinxparam{\DUrole{n}{direction\_point}\DUrole{o}{=}\DUrole{default_value}{None}}\sphinxparamcomma \sphinxparam{\DUrole{n}{name}\DUrole{o}{=}\DUrole{default_value}{None}}\sphinxparamcomma \sphinxparam{\DUrole{n}{copy}\DUrole{o}{=}\DUrole{default_value}{None}}}
{}
\pysigstopsignatures
\sphinxAtStartPar
Add a ray or copy an existing one.
\begin{quote}\begin{description}
\sphinxlineitem{Parameters}\begin{itemize}
\item {} 
\sphinxAtStartPar
\sphinxstyleliteralstrong{\sphinxupquote{start\_point}} (\sphinxstyleliteralemphasis{\sphinxupquote{Point}}\sphinxstyleliteralemphasis{\sphinxupquote{ | }}\sphinxstyleliteralemphasis{\sphinxupquote{None}}) \textendash{} Ray origin.

\item {} 
\sphinxAtStartPar
\sphinxstyleliteralstrong{\sphinxupquote{direction\_point}} (\sphinxstyleliteralemphasis{\sphinxupquote{Point}}\sphinxstyleliteralemphasis{\sphinxupquote{ | }}\sphinxstyleliteralemphasis{\sphinxupquote{None}}) \textendash{} Second point indicating direction.

\item {} 
\sphinxAtStartPar
\sphinxstyleliteralstrong{\sphinxupquote{name}} (\sphinxstyleliteralemphasis{\sphinxupquote{str}}\sphinxstyleliteralemphasis{\sphinxupquote{ | }}\sphinxstyleliteralemphasis{\sphinxupquote{None}}) \textendash{} Optional name for the ray.

\item {} 
\sphinxAtStartPar
\sphinxstyleliteralstrong{\sphinxupquote{copy}} (\sphinxstyleliteralemphasis{\sphinxupquote{Ray}}\sphinxstyleliteralemphasis{\sphinxupquote{ | }}\sphinxstyleliteralemphasis{\sphinxupquote{None}}) \textendash{} Existing ray to duplicate.

\end{itemize}

\sphinxlineitem{Returns}
\sphinxAtStartPar
The created ray, registered in the scene.

\sphinxlineitem{Return type}
\sphinxAtStartPar
Ray

\sphinxlineitem{Raises}
\sphinxAtStartPar
\sphinxstyleliteralstrong{\sphinxupquote{ValueError}} \textendash{} If not copying and a start/direction point is missing.

\end{description}\end{quote}

\end{fulllineitems}

\index{rectangle() (pygeox.add\_objects.AddNamespace method)@\spxentry{rectangle()}\spxextra{pygeox.add\_objects.AddNamespace method}}

\begin{fulllineitems}
\phantomsection\label{\detokenize{pygeox:pygeox.add_objects.AddNamespace.rectangle}}
\pysigstartsignatures
\pysiglinewithargsret
{\sphinxbfcode{\sphinxupquote{rectangle}}}
{\sphinxparam{\DUrole{n}{p1}\DUrole{o}{=}\DUrole{default_value}{None}}\sphinxparamcomma \sphinxparam{\DUrole{n}{p2}\DUrole{o}{=}\DUrole{default_value}{None}}\sphinxparamcomma \sphinxparam{\DUrole{n}{p3}\DUrole{o}{=}\DUrole{default_value}{None}}\sphinxparamcomma \sphinxparam{\DUrole{n}{p4}\DUrole{o}{=}\DUrole{default_value}{None}}\sphinxparamcomma \sphinxparam{\DUrole{n}{name}\DUrole{o}{=}\DUrole{default_value}{None}}\sphinxparamcomma \sphinxparam{\DUrole{n}{force\_convex}\DUrole{o}{=}\DUrole{default_value}{False}}\sphinxparamcomma \sphinxparam{\DUrole{n}{copy}\DUrole{o}{=}\DUrole{default_value}{None}}}
{}
\pysigstopsignatures
\sphinxAtStartPar
Add a rectangle or return a copy.

\sphinxAtStartPar
All angles are constrained 90°; opposite sides parallel. Optional
convexity enforcement.
\begin{quote}\begin{description}
\sphinxlineitem{Parameters}\begin{itemize}
\item {} 
\sphinxAtStartPar
\sphinxstyleliteralstrong{\sphinxupquote{p1}} (\sphinxstyleliteralemphasis{\sphinxupquote{Point}}\sphinxstyleliteralemphasis{\sphinxupquote{ | }}\sphinxstyleliteralemphasis{\sphinxupquote{None}}) \textendash{} Vertex A.

\item {} 
\sphinxAtStartPar
\sphinxstyleliteralstrong{\sphinxupquote{p2}} (\sphinxstyleliteralemphasis{\sphinxupquote{Point}}\sphinxstyleliteralemphasis{\sphinxupquote{ | }}\sphinxstyleliteralemphasis{\sphinxupquote{None}}) \textendash{} Vertex B.

\item {} 
\sphinxAtStartPar
\sphinxstyleliteralstrong{\sphinxupquote{p3}} (\sphinxstyleliteralemphasis{\sphinxupquote{Point}}\sphinxstyleliteralemphasis{\sphinxupquote{ | }}\sphinxstyleliteralemphasis{\sphinxupquote{None}}) \textendash{} Vertex C.

\item {} 
\sphinxAtStartPar
\sphinxstyleliteralstrong{\sphinxupquote{p4}} (\sphinxstyleliteralemphasis{\sphinxupquote{Point}}\sphinxstyleliteralemphasis{\sphinxupquote{ | }}\sphinxstyleliteralemphasis{\sphinxupquote{None}}) \textendash{} Vertex D.

\item {} 
\sphinxAtStartPar
\sphinxstyleliteralstrong{\sphinxupquote{name}} (\sphinxstyleliteralemphasis{\sphinxupquote{str}}\sphinxstyleliteralemphasis{\sphinxupquote{ | }}\sphinxstyleliteralemphasis{\sphinxupquote{None}}) \textendash{} Optional name.

\item {} 
\sphinxAtStartPar
\sphinxstyleliteralstrong{\sphinxupquote{force\_convex}} (\sphinxstyleliteralemphasis{\sphinxupquote{bool}}) \textendash{} If \sphinxcode{\sphinxupquote{True}}, enforces convexity and simplicity.

\item {} 
\sphinxAtStartPar
\sphinxstyleliteralstrong{\sphinxupquote{copy}} (\sphinxstyleliteralemphasis{\sphinxupquote{Rectangle}}\sphinxstyleliteralemphasis{\sphinxupquote{ | }}\sphinxstyleliteralemphasis{\sphinxupquote{None}}) \textendash{} Existing rectangle to reuse.

\end{itemize}

\sphinxlineitem{Returns}
\sphinxAtStartPar
The quadrilateral instance (copy if provided).

\sphinxlineitem{Return type}
\sphinxAtStartPar
Rectangle

\end{description}\end{quote}

\end{fulllineitems}

\index{regular\_heptagon() (pygeox.add\_objects.AddNamespace method)@\spxentry{regular\_heptagon()}\spxextra{pygeox.add\_objects.AddNamespace method}}

\begin{fulllineitems}
\phantomsection\label{\detokenize{pygeox:pygeox.add_objects.AddNamespace.regular_heptagon}}
\pysigstartsignatures
\pysiglinewithargsret
{\sphinxbfcode{\sphinxupquote{regular\_heptagon}}}
{\sphinxparam{\DUrole{o}{*}\DUrole{n}{points}}\sphinxparamcomma \sphinxparam{\DUrole{n}{name}\DUrole{o}{=}\DUrole{default_value}{None}}\sphinxparamcomma \sphinxparam{\DUrole{n}{copy}\DUrole{o}{=}\DUrole{default_value}{None}}}
{}
\pysigstopsignatures
\sphinxAtStartPar
Add a regular heptagon or return a copy.

\sphinxAtStartPar
Assumptions \& constraints: exactly seven ordered points; all sides and
angles equal.
\begin{quote}\begin{description}
\sphinxlineitem{Parameters}\begin{itemize}
\item {} 
\sphinxAtStartPar
\sphinxstyleliteralstrong{\sphinxupquote{*points}} (\sphinxstyleliteralemphasis{\sphinxupquote{Point}}) \textendash{} Seven vertices in order.

\item {} 
\sphinxAtStartPar
\sphinxstyleliteralstrong{\sphinxupquote{name}} (\sphinxstyleliteralemphasis{\sphinxupquote{str}}\sphinxstyleliteralemphasis{\sphinxupquote{ | }}\sphinxstyleliteralemphasis{\sphinxupquote{None}}) \textendash{} Optional name.

\item {} 
\sphinxAtStartPar
\sphinxstyleliteralstrong{\sphinxupquote{copy}} (\sphinxstyleliteralemphasis{\sphinxupquote{RegularHeptagon}}\sphinxstyleliteralemphasis{\sphinxupquote{ | }}\sphinxstyleliteralemphasis{\sphinxupquote{None}}) \textendash{} Existing heptagon to reuse.

\end{itemize}

\sphinxlineitem{Returns}
\sphinxAtStartPar
The polygon instance (copy if provided).

\sphinxlineitem{Return type}
\sphinxAtStartPar
RegularHeptagon

\end{description}\end{quote}

\end{fulllineitems}

\index{regular\_hexagon() (pygeox.add\_objects.AddNamespace method)@\spxentry{regular\_hexagon()}\spxextra{pygeox.add\_objects.AddNamespace method}}

\begin{fulllineitems}
\phantomsection\label{\detokenize{pygeox:pygeox.add_objects.AddNamespace.regular_hexagon}}
\pysigstartsignatures
\pysiglinewithargsret
{\sphinxbfcode{\sphinxupquote{regular\_hexagon}}}
{\sphinxparam{\DUrole{o}{*}\DUrole{n}{points}}\sphinxparamcomma \sphinxparam{\DUrole{n}{name}\DUrole{o}{=}\DUrole{default_value}{None}}\sphinxparamcomma \sphinxparam{\DUrole{n}{copy}\DUrole{o}{=}\DUrole{default_value}{None}}}
{}
\pysigstopsignatures
\sphinxAtStartPar
Add a regular hexagon or return a copy.

\sphinxAtStartPar
Assumptions \& constraints: exactly six ordered points; all sides and
angles equal; inscribed in a circle.
\begin{quote}\begin{description}
\sphinxlineitem{Parameters}\begin{itemize}
\item {} 
\sphinxAtStartPar
\sphinxstyleliteralstrong{\sphinxupquote{*points}} (\sphinxstyleliteralemphasis{\sphinxupquote{Point}}) \textendash{} Six vertices in order.

\item {} 
\sphinxAtStartPar
\sphinxstyleliteralstrong{\sphinxupquote{name}} (\sphinxstyleliteralemphasis{\sphinxupquote{str}}\sphinxstyleliteralemphasis{\sphinxupquote{ | }}\sphinxstyleliteralemphasis{\sphinxupquote{None}}) \textendash{} Optional name.

\item {} 
\sphinxAtStartPar
\sphinxstyleliteralstrong{\sphinxupquote{copy}} (\sphinxstyleliteralemphasis{\sphinxupquote{RegularHexagon}}\sphinxstyleliteralemphasis{\sphinxupquote{ | }}\sphinxstyleliteralemphasis{\sphinxupquote{None}}) \textendash{} Existing hexagon to reuse.

\end{itemize}

\sphinxlineitem{Returns}
\sphinxAtStartPar
The polygon instance (copy if provided).

\sphinxlineitem{Return type}
\sphinxAtStartPar
RegularHexagon

\end{description}\end{quote}

\end{fulllineitems}

\index{regular\_octagon() (pygeox.add\_objects.AddNamespace method)@\spxentry{regular\_octagon()}\spxextra{pygeox.add\_objects.AddNamespace method}}

\begin{fulllineitems}
\phantomsection\label{\detokenize{pygeox:pygeox.add_objects.AddNamespace.regular_octagon}}
\pysigstartsignatures
\pysiglinewithargsret
{\sphinxbfcode{\sphinxupquote{regular\_octagon}}}
{\sphinxparam{\DUrole{o}{*}\DUrole{n}{points}}\sphinxparamcomma \sphinxparam{\DUrole{n}{name}\DUrole{o}{=}\DUrole{default_value}{None}}\sphinxparamcomma \sphinxparam{\DUrole{n}{copy}\DUrole{o}{=}\DUrole{default_value}{None}}}
{}
\pysigstopsignatures
\sphinxAtStartPar
Add a regular octagon or return a copy.

\sphinxAtStartPar
Assumptions \& constraints: exactly eight ordered points; all sides and
angles equal.
\begin{quote}\begin{description}
\sphinxlineitem{Parameters}\begin{itemize}
\item {} 
\sphinxAtStartPar
\sphinxstyleliteralstrong{\sphinxupquote{*points}} (\sphinxstyleliteralemphasis{\sphinxupquote{Point}}) \textendash{} Eight vertices in order.

\item {} 
\sphinxAtStartPar
\sphinxstyleliteralstrong{\sphinxupquote{name}} (\sphinxstyleliteralemphasis{\sphinxupquote{str}}\sphinxstyleliteralemphasis{\sphinxupquote{ | }}\sphinxstyleliteralemphasis{\sphinxupquote{None}}) \textendash{} Optional name.

\item {} 
\sphinxAtStartPar
\sphinxstyleliteralstrong{\sphinxupquote{copy}} (\sphinxstyleliteralemphasis{\sphinxupquote{RegularOctagon}}\sphinxstyleliteralemphasis{\sphinxupquote{ | }}\sphinxstyleliteralemphasis{\sphinxupquote{None}}) \textendash{} Existing octagon to reuse.

\end{itemize}

\sphinxlineitem{Returns}
\sphinxAtStartPar
The polygon instance (copy if provided).

\sphinxlineitem{Return type}
\sphinxAtStartPar
RegularOctagon

\end{description}\end{quote}

\end{fulllineitems}

\index{regular\_pentagon() (pygeox.add\_objects.AddNamespace method)@\spxentry{regular\_pentagon()}\spxextra{pygeox.add\_objects.AddNamespace method}}

\begin{fulllineitems}
\phantomsection\label{\detokenize{pygeox:pygeox.add_objects.AddNamespace.regular_pentagon}}
\pysigstartsignatures
\pysiglinewithargsret
{\sphinxbfcode{\sphinxupquote{regular\_pentagon}}}
{\sphinxparam{\DUrole{o}{*}\DUrole{n}{points}}\sphinxparamcomma \sphinxparam{\DUrole{n}{name}\DUrole{o}{=}\DUrole{default_value}{None}}\sphinxparamcomma \sphinxparam{\DUrole{n}{copy}\DUrole{o}{=}\DUrole{default_value}{None}}}
{}
\pysigstopsignatures
\sphinxAtStartPar
Add a regular pentagon or return a copy.

\sphinxAtStartPar
Assumptions \& constraints: exactly five ordered points; all sides and
angles equal; inscribed in a circle.
\begin{quote}\begin{description}
\sphinxlineitem{Parameters}\begin{itemize}
\item {} 
\sphinxAtStartPar
\sphinxstyleliteralstrong{\sphinxupquote{*points}} (\sphinxstyleliteralemphasis{\sphinxupquote{Point}}) \textendash{} Five vertices in order.

\item {} 
\sphinxAtStartPar
\sphinxstyleliteralstrong{\sphinxupquote{name}} (\sphinxstyleliteralemphasis{\sphinxupquote{str}}\sphinxstyleliteralemphasis{\sphinxupquote{ | }}\sphinxstyleliteralemphasis{\sphinxupquote{None}}) \textendash{} Optional name.

\item {} 
\sphinxAtStartPar
\sphinxstyleliteralstrong{\sphinxupquote{copy}} (\sphinxstyleliteralemphasis{\sphinxupquote{RegularPentagon}}\sphinxstyleliteralemphasis{\sphinxupquote{ | }}\sphinxstyleliteralemphasis{\sphinxupquote{None}}) \textendash{} Existing pentagon to reuse.

\end{itemize}

\sphinxlineitem{Returns}
\sphinxAtStartPar
The polygon instance (copy if provided).

\sphinxlineitem{Return type}
\sphinxAtStartPar
RegularPentagon

\end{description}\end{quote}

\end{fulllineitems}

\index{regular\_polygon() (pygeox.add\_objects.AddNamespace method)@\spxentry{regular\_polygon()}\spxextra{pygeox.add\_objects.AddNamespace method}}

\begin{fulllineitems}
\phantomsection\label{\detokenize{pygeox:pygeox.add_objects.AddNamespace.regular_polygon}}
\pysigstartsignatures
\pysiglinewithargsret
{\sphinxbfcode{\sphinxupquote{regular\_polygon}}}
{\sphinxparam{\DUrole{o}{*}\DUrole{n}{points}}\sphinxparamcomma \sphinxparam{\DUrole{n}{name}\DUrole{o}{=}\DUrole{default_value}{None}}\sphinxparamcomma \sphinxparam{\DUrole{n}{copy}\DUrole{o}{=}\DUrole{default_value}{None}}}
{}
\pysigstopsignatures
\sphinxAtStartPar
Add a regular polygon or return a copy.

\sphinxAtStartPar
Assumptions \& constraints: at least three ordered points; all sides and
angles are constrained equal; polygon is inscribed in a circle.
\begin{quote}\begin{description}
\sphinxlineitem{Parameters}\begin{itemize}
\item {} 
\sphinxAtStartPar
\sphinxstyleliteralstrong{\sphinxupquote{*points}} (\sphinxstyleliteralemphasis{\sphinxupquote{Point}}) \textendash{} Vertices in order.

\item {} 
\sphinxAtStartPar
\sphinxstyleliteralstrong{\sphinxupquote{name}} (\sphinxstyleliteralemphasis{\sphinxupquote{str}}\sphinxstyleliteralemphasis{\sphinxupquote{ | }}\sphinxstyleliteralemphasis{\sphinxupquote{None}}) \textendash{} Optional name.

\item {} 
\sphinxAtStartPar
\sphinxstyleliteralstrong{\sphinxupquote{copy}} (\sphinxstyleliteralemphasis{\sphinxupquote{RegularPolygon}}\sphinxstyleliteralemphasis{\sphinxupquote{ | }}\sphinxstyleliteralemphasis{\sphinxupquote{None}}) \textendash{} Existing regular polygon to reuse.

\end{itemize}

\sphinxlineitem{Returns}
\sphinxAtStartPar
The polygon instance (copy if provided).

\sphinxlineitem{Return type}
\sphinxAtStartPar
RegularPolygon

\end{description}\end{quote}

\end{fulllineitems}

\index{rhomboid() (pygeox.add\_objects.AddNamespace method)@\spxentry{rhomboid()}\spxextra{pygeox.add\_objects.AddNamespace method}}

\begin{fulllineitems}
\phantomsection\label{\detokenize{pygeox:pygeox.add_objects.AddNamespace.rhomboid}}
\pysigstartsignatures
\pysiglinewithargsret
{\sphinxbfcode{\sphinxupquote{rhomboid}}}
{\sphinxparam{\DUrole{n}{p1}\DUrole{o}{=}\DUrole{default_value}{None}}\sphinxparamcomma \sphinxparam{\DUrole{n}{p2}\DUrole{o}{=}\DUrole{default_value}{None}}\sphinxparamcomma \sphinxparam{\DUrole{n}{p3}\DUrole{o}{=}\DUrole{default_value}{None}}\sphinxparamcomma \sphinxparam{\DUrole{n}{p4}\DUrole{o}{=}\DUrole{default_value}{None}}\sphinxparamcomma \sphinxparam{\DUrole{n}{name}\DUrole{o}{=}\DUrole{default_value}{None}}\sphinxparamcomma \sphinxparam{\DUrole{n}{force\_convex}\DUrole{o}{=}\DUrole{default_value}{False}}\sphinxparamcomma \sphinxparam{\DUrole{n}{copy}\DUrole{o}{=}\DUrole{default_value}{None}}}
{}
\pysigstopsignatures
\sphinxAtStartPar
Add a rhomboid or return a copy.

\sphinxAtStartPar
Opposite sides parallel; adjacent sides unequal; no right\sphinxhyphen{}angle
constraint. Optional convexity enforcement.
\begin{quote}\begin{description}
\sphinxlineitem{Parameters}\begin{itemize}
\item {} 
\sphinxAtStartPar
\sphinxstyleliteralstrong{\sphinxupquote{p1}} (\sphinxstyleliteralemphasis{\sphinxupquote{Point}}\sphinxstyleliteralemphasis{\sphinxupquote{ | }}\sphinxstyleliteralemphasis{\sphinxupquote{None}}) \textendash{} Vertex A.

\item {} 
\sphinxAtStartPar
\sphinxstyleliteralstrong{\sphinxupquote{p2}} (\sphinxstyleliteralemphasis{\sphinxupquote{Point}}\sphinxstyleliteralemphasis{\sphinxupquote{ | }}\sphinxstyleliteralemphasis{\sphinxupquote{None}}) \textendash{} Vertex B.

\item {} 
\sphinxAtStartPar
\sphinxstyleliteralstrong{\sphinxupquote{p3}} (\sphinxstyleliteralemphasis{\sphinxupquote{Point}}\sphinxstyleliteralemphasis{\sphinxupquote{ | }}\sphinxstyleliteralemphasis{\sphinxupquote{None}}) \textendash{} Vertex C.

\item {} 
\sphinxAtStartPar
\sphinxstyleliteralstrong{\sphinxupquote{p4}} (\sphinxstyleliteralemphasis{\sphinxupquote{Point}}\sphinxstyleliteralemphasis{\sphinxupquote{ | }}\sphinxstyleliteralemphasis{\sphinxupquote{None}}) \textendash{} Vertex D.

\item {} 
\sphinxAtStartPar
\sphinxstyleliteralstrong{\sphinxupquote{name}} (\sphinxstyleliteralemphasis{\sphinxupquote{str}}\sphinxstyleliteralemphasis{\sphinxupquote{ | }}\sphinxstyleliteralemphasis{\sphinxupquote{None}}) \textendash{} Optional name.

\item {} 
\sphinxAtStartPar
\sphinxstyleliteralstrong{\sphinxupquote{force\_convex}} (\sphinxstyleliteralemphasis{\sphinxupquote{bool}}) \textendash{} If \sphinxcode{\sphinxupquote{True}}, enforces convexity and simplicity.

\item {} 
\sphinxAtStartPar
\sphinxstyleliteralstrong{\sphinxupquote{copy}} (\sphinxstyleliteralemphasis{\sphinxupquote{Rhomboid}}\sphinxstyleliteralemphasis{\sphinxupquote{ | }}\sphinxstyleliteralemphasis{\sphinxupquote{None}}) \textendash{} Existing rhomboid to reuse.

\end{itemize}

\sphinxlineitem{Returns}
\sphinxAtStartPar
The quadrilateral instance (copy if provided).

\sphinxlineitem{Return type}
\sphinxAtStartPar
Rhomboid

\end{description}\end{quote}

\end{fulllineitems}

\index{rhombus() (pygeox.add\_objects.AddNamespace method)@\spxentry{rhombus()}\spxextra{pygeox.add\_objects.AddNamespace method}}

\begin{fulllineitems}
\phantomsection\label{\detokenize{pygeox:pygeox.add_objects.AddNamespace.rhombus}}
\pysigstartsignatures
\pysiglinewithargsret
{\sphinxbfcode{\sphinxupquote{rhombus}}}
{\sphinxparam{\DUrole{n}{p1}\DUrole{o}{=}\DUrole{default_value}{None}}\sphinxparamcomma \sphinxparam{\DUrole{n}{p2}\DUrole{o}{=}\DUrole{default_value}{None}}\sphinxparamcomma \sphinxparam{\DUrole{n}{p3}\DUrole{o}{=}\DUrole{default_value}{None}}\sphinxparamcomma \sphinxparam{\DUrole{n}{p4}\DUrole{o}{=}\DUrole{default_value}{None}}\sphinxparamcomma \sphinxparam{\DUrole{n}{name}\DUrole{o}{=}\DUrole{default_value}{None}}\sphinxparamcomma \sphinxparam{\DUrole{n}{force\_convex}\DUrole{o}{=}\DUrole{default_value}{False}}\sphinxparamcomma \sphinxparam{\DUrole{n}{copy}\DUrole{o}{=}\DUrole{default_value}{None}}}
{}
\pysigstopsignatures
\sphinxAtStartPar
Add a rhombus or return a copy.

\sphinxAtStartPar
All sides equal; opposite sides parallel; angles not necessarily 90°.
Optional convexity enforcement.
\begin{quote}\begin{description}
\sphinxlineitem{Parameters}\begin{itemize}
\item {} 
\sphinxAtStartPar
\sphinxstyleliteralstrong{\sphinxupquote{p1}} (\sphinxstyleliteralemphasis{\sphinxupquote{Point}}\sphinxstyleliteralemphasis{\sphinxupquote{ | }}\sphinxstyleliteralemphasis{\sphinxupquote{None}}) \textendash{} Vertex A.

\item {} 
\sphinxAtStartPar
\sphinxstyleliteralstrong{\sphinxupquote{p2}} (\sphinxstyleliteralemphasis{\sphinxupquote{Point}}\sphinxstyleliteralemphasis{\sphinxupquote{ | }}\sphinxstyleliteralemphasis{\sphinxupquote{None}}) \textendash{} Vertex B.

\item {} 
\sphinxAtStartPar
\sphinxstyleliteralstrong{\sphinxupquote{p3}} (\sphinxstyleliteralemphasis{\sphinxupquote{Point}}\sphinxstyleliteralemphasis{\sphinxupquote{ | }}\sphinxstyleliteralemphasis{\sphinxupquote{None}}) \textendash{} Vertex C.

\item {} 
\sphinxAtStartPar
\sphinxstyleliteralstrong{\sphinxupquote{p4}} (\sphinxstyleliteralemphasis{\sphinxupquote{Point}}\sphinxstyleliteralemphasis{\sphinxupquote{ | }}\sphinxstyleliteralemphasis{\sphinxupquote{None}}) \textendash{} Vertex D.

\item {} 
\sphinxAtStartPar
\sphinxstyleliteralstrong{\sphinxupquote{name}} (\sphinxstyleliteralemphasis{\sphinxupquote{str}}\sphinxstyleliteralemphasis{\sphinxupquote{ | }}\sphinxstyleliteralemphasis{\sphinxupquote{None}}) \textendash{} Optional name.

\item {} 
\sphinxAtStartPar
\sphinxstyleliteralstrong{\sphinxupquote{force\_convex}} (\sphinxstyleliteralemphasis{\sphinxupquote{bool}}) \textendash{} If \sphinxcode{\sphinxupquote{True}}, enforces convexity and simplicity.

\item {} 
\sphinxAtStartPar
\sphinxstyleliteralstrong{\sphinxupquote{copy}} (\sphinxstyleliteralemphasis{\sphinxupquote{Rhombus}}\sphinxstyleliteralemphasis{\sphinxupquote{ | }}\sphinxstyleliteralemphasis{\sphinxupquote{None}}) \textendash{} Existing rhombus to reuse.

\end{itemize}

\sphinxlineitem{Returns}
\sphinxAtStartPar
The quadrilateral instance (copy if provided).

\sphinxlineitem{Return type}
\sphinxAtStartPar
Rhombus

\end{description}\end{quote}

\end{fulllineitems}

\index{right\_isosceles\_triangle() (pygeox.add\_objects.AddNamespace method)@\spxentry{right\_isosceles\_triangle()}\spxextra{pygeox.add\_objects.AddNamespace method}}

\begin{fulllineitems}
\phantomsection\label{\detokenize{pygeox:pygeox.add_objects.AddNamespace.right_isosceles_triangle}}
\pysigstartsignatures
\pysiglinewithargsret
{\sphinxbfcode{\sphinxupquote{right\_isosceles\_triangle}}}
{\sphinxparam{\DUrole{n}{p1}\DUrole{o}{=}\DUrole{default_value}{None}}\sphinxparamcomma \sphinxparam{\DUrole{n}{p2}\DUrole{o}{=}\DUrole{default_value}{None}}\sphinxparamcomma \sphinxparam{\DUrole{n}{p3}\DUrole{o}{=}\DUrole{default_value}{None}}\sphinxparamcomma \sphinxparam{\DUrole{n}{name}\DUrole{o}{=}\DUrole{default_value}{None}}\sphinxparamcomma \sphinxparam{\DUrole{n}{copy}\DUrole{o}{=}\DUrole{default_value}{None}}}
{}
\pysigstopsignatures
\sphinxAtStartPar
Add a right isosceles triangle or return a copy.
\begin{quote}\begin{description}
\sphinxlineitem{Parameters}\begin{itemize}
\item {} 
\sphinxAtStartPar
\sphinxstyleliteralstrong{\sphinxupquote{p1}} (\sphinxstyleliteralemphasis{\sphinxupquote{Point}}\sphinxstyleliteralemphasis{\sphinxupquote{ | }}\sphinxstyleliteralemphasis{\sphinxupquote{None}}) \textendash{} Vertex A.

\item {} 
\sphinxAtStartPar
\sphinxstyleliteralstrong{\sphinxupquote{p2}} (\sphinxstyleliteralemphasis{\sphinxupquote{Point}}\sphinxstyleliteralemphasis{\sphinxupquote{ | }}\sphinxstyleliteralemphasis{\sphinxupquote{None}}) \textendash{} Vertex B (right angle at B).

\item {} 
\sphinxAtStartPar
\sphinxstyleliteralstrong{\sphinxupquote{p3}} (\sphinxstyleliteralemphasis{\sphinxupquote{Point}}\sphinxstyleliteralemphasis{\sphinxupquote{ | }}\sphinxstyleliteralemphasis{\sphinxupquote{None}}) \textendash{} Vertex C.

\item {} 
\sphinxAtStartPar
\sphinxstyleliteralstrong{\sphinxupquote{name}} (\sphinxstyleliteralemphasis{\sphinxupquote{str}}\sphinxstyleliteralemphasis{\sphinxupquote{ | }}\sphinxstyleliteralemphasis{\sphinxupquote{None}}) \textendash{} Optional name.

\item {} 
\sphinxAtStartPar
\sphinxstyleliteralstrong{\sphinxupquote{copy}} (\sphinxstyleliteralemphasis{\sphinxupquote{RightIsoscelesTriangle}}\sphinxstyleliteralemphasis{\sphinxupquote{ | }}\sphinxstyleliteralemphasis{\sphinxupquote{None}}) \textendash{} Existing triangle to reuse.

\end{itemize}

\sphinxlineitem{Returns}
\sphinxAtStartPar
The triangle instance (copy if provided).

\sphinxlineitem{Return type}
\sphinxAtStartPar
RightIsoscelesTriangle

\end{description}\end{quote}
\subsubsection*{Notes}

\sphinxAtStartPar
Sides AB and BC are constrained equal (order matters).

\end{fulllineitems}

\index{right\_trapezoid() (pygeox.add\_objects.AddNamespace method)@\spxentry{right\_trapezoid()}\spxextra{pygeox.add\_objects.AddNamespace method}}

\begin{fulllineitems}
\phantomsection\label{\detokenize{pygeox:pygeox.add_objects.AddNamespace.right_trapezoid}}
\pysigstartsignatures
\pysiglinewithargsret
{\sphinxbfcode{\sphinxupquote{right\_trapezoid}}}
{\sphinxparam{\DUrole{n}{p1}\DUrole{o}{=}\DUrole{default_value}{None}}\sphinxparamcomma \sphinxparam{\DUrole{n}{p2}\DUrole{o}{=}\DUrole{default_value}{None}}\sphinxparamcomma \sphinxparam{\DUrole{n}{p3}\DUrole{o}{=}\DUrole{default_value}{None}}\sphinxparamcomma \sphinxparam{\DUrole{n}{p4}\DUrole{o}{=}\DUrole{default_value}{None}}\sphinxparamcomma \sphinxparam{\DUrole{n}{name}\DUrole{o}{=}\DUrole{default_value}{None}}\sphinxparamcomma \sphinxparam{\DUrole{n}{force\_convex}\DUrole{o}{=}\DUrole{default_value}{False}}\sphinxparamcomma \sphinxparam{\DUrole{n}{copy}\DUrole{o}{=}\DUrole{default_value}{None}}}
{}
\pysigstopsignatures
\sphinxAtStartPar
Add a right trapezoid or return a copy.

\sphinxAtStartPar
Constrains AB to be parallel to CD and angles A, D = 90°. Optional convexity enforcement.
\begin{quote}\begin{description}
\sphinxlineitem{Parameters}\begin{itemize}
\item {} 
\sphinxAtStartPar
\sphinxstyleliteralstrong{\sphinxupquote{p1}} (\sphinxstyleliteralemphasis{\sphinxupquote{Point}}\sphinxstyleliteralemphasis{\sphinxupquote{ | }}\sphinxstyleliteralemphasis{\sphinxupquote{None}}) \textendash{} Vertex A.

\item {} 
\sphinxAtStartPar
\sphinxstyleliteralstrong{\sphinxupquote{p2}} (\sphinxstyleliteralemphasis{\sphinxupquote{Point}}\sphinxstyleliteralemphasis{\sphinxupquote{ | }}\sphinxstyleliteralemphasis{\sphinxupquote{None}}) \textendash{} Vertex B.

\item {} 
\sphinxAtStartPar
\sphinxstyleliteralstrong{\sphinxupquote{p3}} (\sphinxstyleliteralemphasis{\sphinxupquote{Point}}\sphinxstyleliteralemphasis{\sphinxupquote{ | }}\sphinxstyleliteralemphasis{\sphinxupquote{None}}) \textendash{} Vertex C.

\item {} 
\sphinxAtStartPar
\sphinxstyleliteralstrong{\sphinxupquote{p4}} (\sphinxstyleliteralemphasis{\sphinxupquote{Point}}\sphinxstyleliteralemphasis{\sphinxupquote{ | }}\sphinxstyleliteralemphasis{\sphinxupquote{None}}) \textendash{} Vertex D.

\item {} 
\sphinxAtStartPar
\sphinxstyleliteralstrong{\sphinxupquote{name}} (\sphinxstyleliteralemphasis{\sphinxupquote{str}}\sphinxstyleliteralemphasis{\sphinxupquote{ | }}\sphinxstyleliteralemphasis{\sphinxupquote{None}}) \textendash{} Optional name.

\item {} 
\sphinxAtStartPar
\sphinxstyleliteralstrong{\sphinxupquote{force\_convex}} (\sphinxstyleliteralemphasis{\sphinxupquote{bool}}) \textendash{} If \sphinxcode{\sphinxupquote{True}}, enforces convexity and simplicity.

\item {} 
\sphinxAtStartPar
\sphinxstyleliteralstrong{\sphinxupquote{copy}} (\sphinxstyleliteralemphasis{\sphinxupquote{RightTrapezoid}}\sphinxstyleliteralemphasis{\sphinxupquote{ | }}\sphinxstyleliteralemphasis{\sphinxupquote{None}}) \textendash{} Existing trapezoid to reuse.

\end{itemize}

\sphinxlineitem{Returns}
\sphinxAtStartPar
The quadrilateral instance (copy if provided).

\sphinxlineitem{Return type}
\sphinxAtStartPar
RightTrapezoid

\end{description}\end{quote}

\end{fulllineitems}

\index{right\_triangle() (pygeox.add\_objects.AddNamespace method)@\spxentry{right\_triangle()}\spxextra{pygeox.add\_objects.AddNamespace method}}

\begin{fulllineitems}
\phantomsection\label{\detokenize{pygeox:pygeox.add_objects.AddNamespace.right_triangle}}
\pysigstartsignatures
\pysiglinewithargsret
{\sphinxbfcode{\sphinxupquote{right\_triangle}}}
{\sphinxparam{\DUrole{n}{p1}\DUrole{o}{=}\DUrole{default_value}{None}}\sphinxparamcomma \sphinxparam{\DUrole{n}{p2}\DUrole{o}{=}\DUrole{default_value}{None}}\sphinxparamcomma \sphinxparam{\DUrole{n}{p3}\DUrole{o}{=}\DUrole{default_value}{None}}\sphinxparamcomma \sphinxparam{\DUrole{n}{name}\DUrole{o}{=}\DUrole{default_value}{None}}\sphinxparamcomma \sphinxparam{\DUrole{n}{copy}\DUrole{o}{=}\DUrole{default_value}{None}}}
{}
\pysigstopsignatures
\sphinxAtStartPar
Add a right triangle or return a copy.
\begin{quote}\begin{description}
\sphinxlineitem{Parameters}\begin{itemize}
\item {} 
\sphinxAtStartPar
\sphinxstyleliteralstrong{\sphinxupquote{p1}} (\sphinxstyleliteralemphasis{\sphinxupquote{Point}}\sphinxstyleliteralemphasis{\sphinxupquote{ | }}\sphinxstyleliteralemphasis{\sphinxupquote{None}}) \textendash{} Vertex A.

\item {} 
\sphinxAtStartPar
\sphinxstyleliteralstrong{\sphinxupquote{p2}} (\sphinxstyleliteralemphasis{\sphinxupquote{Point}}\sphinxstyleliteralemphasis{\sphinxupquote{ | }}\sphinxstyleliteralemphasis{\sphinxupquote{None}}) \textendash{} Vertex B (right angle at B).

\item {} 
\sphinxAtStartPar
\sphinxstyleliteralstrong{\sphinxupquote{p3}} (\sphinxstyleliteralemphasis{\sphinxupquote{Point}}\sphinxstyleliteralemphasis{\sphinxupquote{ | }}\sphinxstyleliteralemphasis{\sphinxupquote{None}}) \textendash{} Vertex C.

\item {} 
\sphinxAtStartPar
\sphinxstyleliteralstrong{\sphinxupquote{name}} (\sphinxstyleliteralemphasis{\sphinxupquote{str}}\sphinxstyleliteralemphasis{\sphinxupquote{ | }}\sphinxstyleliteralemphasis{\sphinxupquote{None}}) \textendash{} Optional name.

\item {} 
\sphinxAtStartPar
\sphinxstyleliteralstrong{\sphinxupquote{copy}} (\sphinxstyleliteralemphasis{\sphinxupquote{RightTriangle}}\sphinxstyleliteralemphasis{\sphinxupquote{ | }}\sphinxstyleliteralemphasis{\sphinxupquote{None}}) \textendash{} Existing triangle to reuse.

\end{itemize}

\sphinxlineitem{Returns}
\sphinxAtStartPar
The triangle instance (copy if provided).

\sphinxlineitem{Return type}
\sphinxAtStartPar
RightTriangle

\end{description}\end{quote}

\end{fulllineitems}

\index{scalene\_triangle() (pygeox.add\_objects.AddNamespace method)@\spxentry{scalene\_triangle()}\spxextra{pygeox.add\_objects.AddNamespace method}}

\begin{fulllineitems}
\phantomsection\label{\detokenize{pygeox:pygeox.add_objects.AddNamespace.scalene_triangle}}
\pysigstartsignatures
\pysiglinewithargsret
{\sphinxbfcode{\sphinxupquote{scalene\_triangle}}}
{\sphinxparam{\DUrole{n}{p1}\DUrole{o}{=}\DUrole{default_value}{None}}\sphinxparamcomma \sphinxparam{\DUrole{n}{p2}\DUrole{o}{=}\DUrole{default_value}{None}}\sphinxparamcomma \sphinxparam{\DUrole{n}{p3}\DUrole{o}{=}\DUrole{default_value}{None}}\sphinxparamcomma \sphinxparam{\DUrole{n}{name}\DUrole{o}{=}\DUrole{default_value}{None}}\sphinxparamcomma \sphinxparam{\DUrole{n}{copy}\DUrole{o}{=}\DUrole{default_value}{None}}}
{}
\pysigstopsignatures
\sphinxAtStartPar
Add a scalene triangle or return a copy.
\begin{quote}\begin{description}
\sphinxlineitem{Parameters}\begin{itemize}
\item {} 
\sphinxAtStartPar
\sphinxstyleliteralstrong{\sphinxupquote{p1}} (\sphinxstyleliteralemphasis{\sphinxupquote{Point}}\sphinxstyleliteralemphasis{\sphinxupquote{ | }}\sphinxstyleliteralemphasis{\sphinxupquote{None}}) \textendash{} Vertex A.

\item {} 
\sphinxAtStartPar
\sphinxstyleliteralstrong{\sphinxupquote{p2}} (\sphinxstyleliteralemphasis{\sphinxupquote{Point}}\sphinxstyleliteralemphasis{\sphinxupquote{ | }}\sphinxstyleliteralemphasis{\sphinxupquote{None}}) \textendash{} Vertex B.

\item {} 
\sphinxAtStartPar
\sphinxstyleliteralstrong{\sphinxupquote{p3}} (\sphinxstyleliteralemphasis{\sphinxupquote{Point}}\sphinxstyleliteralemphasis{\sphinxupquote{ | }}\sphinxstyleliteralemphasis{\sphinxupquote{None}}) \textendash{} Vertex C.

\item {} 
\sphinxAtStartPar
\sphinxstyleliteralstrong{\sphinxupquote{name}} (\sphinxstyleliteralemphasis{\sphinxupquote{str}}\sphinxstyleliteralemphasis{\sphinxupquote{ | }}\sphinxstyleliteralemphasis{\sphinxupquote{None}}) \textendash{} Optional name.

\item {} 
\sphinxAtStartPar
\sphinxstyleliteralstrong{\sphinxupquote{copy}} (\sphinxstyleliteralemphasis{\sphinxupquote{ScaleneTriangle}}\sphinxstyleliteralemphasis{\sphinxupquote{ | }}\sphinxstyleliteralemphasis{\sphinxupquote{None}}) \textendash{} Existing triangle to reuse.

\end{itemize}

\sphinxlineitem{Returns}
\sphinxAtStartPar
The triangle instance (copy if provided).

\sphinxlineitem{Return type}
\sphinxAtStartPar
ScaleneTriangle

\end{description}\end{quote}
\subsubsection*{Notes}

\sphinxAtStartPar
All side lengths are constrained to be distinct.

\end{fulllineitems}

\index{semicircle() (pygeox.add\_objects.AddNamespace method)@\spxentry{semicircle()}\spxextra{pygeox.add\_objects.AddNamespace method}}

\begin{fulllineitems}
\phantomsection\label{\detokenize{pygeox:pygeox.add_objects.AddNamespace.semicircle}}
\pysigstartsignatures
\pysiglinewithargsret
{\sphinxbfcode{\sphinxupquote{semicircle}}}
{\sphinxparam{\DUrole{n}{center}}\sphinxparamcomma \sphinxparam{\DUrole{n}{A}}\sphinxparamcomma \sphinxparam{\DUrole{n}{B}}\sphinxparamcomma \sphinxparam{\DUrole{n}{name}\DUrole{o}{=}\DUrole{default_value}{None}}}
{}
\pysigstopsignatures
\sphinxAtStartPar
Add a semicircle (180° major arc) from A to B with a given center.

\sphinxAtStartPar
Enforces collinearity of \sphinxcode{\sphinxupquote{center}}, \sphinxcode{\sphinxupquote{A}}, and \sphinxcode{\sphinxupquote{B}} (direction not forced).
\begin{quote}\begin{description}
\sphinxlineitem{Parameters}\begin{itemize}
\item {} 
\sphinxAtStartPar
\sphinxstyleliteralstrong{\sphinxupquote{center}} (\sphinxstyleliteralemphasis{\sphinxupquote{Point}}) \textendash{} Center of the circle.

\item {} 
\sphinxAtStartPar
\sphinxstyleliteralstrong{\sphinxupquote{A}} (\sphinxstyleliteralemphasis{\sphinxupquote{Point}}) \textendash{} First endpoint on the circle.

\item {} 
\sphinxAtStartPar
\sphinxstyleliteralstrong{\sphinxupquote{B}} (\sphinxstyleliteralemphasis{\sphinxupquote{Point}}) \textendash{} Second endpoint on the circle.

\item {} 
\sphinxAtStartPar
\sphinxstyleliteralstrong{\sphinxupquote{name}} (\sphinxstyleliteralemphasis{\sphinxupquote{str}}\sphinxstyleliteralemphasis{\sphinxupquote{ | }}\sphinxstyleliteralemphasis{\sphinxupquote{None}}) \textendash{} Optional name. Defaults to \sphinxcode{\sphinxupquote{semicircle\_A\_B\_center\_\textless{}center\textgreater{}}}.

\end{itemize}

\sphinxlineitem{Returns}
\sphinxAtStartPar
The created semicircle as a major arc.

\sphinxlineitem{Return type}
\sphinxAtStartPar
MajorArc

\end{description}\end{quote}

\end{fulllineitems}

\index{square() (pygeox.add\_objects.AddNamespace method)@\spxentry{square()}\spxextra{pygeox.add\_objects.AddNamespace method}}

\begin{fulllineitems}
\phantomsection\label{\detokenize{pygeox:pygeox.add_objects.AddNamespace.square}}
\pysigstartsignatures
\pysiglinewithargsret
{\sphinxbfcode{\sphinxupquote{square}}}
{\sphinxparam{\DUrole{n}{p1}\DUrole{o}{=}\DUrole{default_value}{None}}\sphinxparamcomma \sphinxparam{\DUrole{n}{p2}\DUrole{o}{=}\DUrole{default_value}{None}}\sphinxparamcomma \sphinxparam{\DUrole{n}{p3}\DUrole{o}{=}\DUrole{default_value}{None}}\sphinxparamcomma \sphinxparam{\DUrole{n}{p4}\DUrole{o}{=}\DUrole{default_value}{None}}\sphinxparamcomma \sphinxparam{\DUrole{n}{name}\DUrole{o}{=}\DUrole{default_value}{None}}\sphinxparamcomma \sphinxparam{\DUrole{n}{copy}\DUrole{o}{=}\DUrole{default_value}{None}}}
{}
\pysigstopsignatures
\sphinxAtStartPar
Add a square or return a copy.

\sphinxAtStartPar
All sides equal, all angles 90°, opposite sides parallel.
\begin{quote}\begin{description}
\sphinxlineitem{Parameters}\begin{itemize}
\item {} 
\sphinxAtStartPar
\sphinxstyleliteralstrong{\sphinxupquote{p1}} (\sphinxstyleliteralemphasis{\sphinxupquote{Point}}\sphinxstyleliteralemphasis{\sphinxupquote{ | }}\sphinxstyleliteralemphasis{\sphinxupquote{None}}) \textendash{} Vertex A.

\item {} 
\sphinxAtStartPar
\sphinxstyleliteralstrong{\sphinxupquote{p2}} (\sphinxstyleliteralemphasis{\sphinxupquote{Point}}\sphinxstyleliteralemphasis{\sphinxupquote{ | }}\sphinxstyleliteralemphasis{\sphinxupquote{None}}) \textendash{} Vertex B.

\item {} 
\sphinxAtStartPar
\sphinxstyleliteralstrong{\sphinxupquote{p3}} (\sphinxstyleliteralemphasis{\sphinxupquote{Point}}\sphinxstyleliteralemphasis{\sphinxupquote{ | }}\sphinxstyleliteralemphasis{\sphinxupquote{None}}) \textendash{} Vertex C.

\item {} 
\sphinxAtStartPar
\sphinxstyleliteralstrong{\sphinxupquote{p4}} (\sphinxstyleliteralemphasis{\sphinxupquote{Point}}\sphinxstyleliteralemphasis{\sphinxupquote{ | }}\sphinxstyleliteralemphasis{\sphinxupquote{None}}) \textendash{} Vertex D.

\item {} 
\sphinxAtStartPar
\sphinxstyleliteralstrong{\sphinxupquote{name}} (\sphinxstyleliteralemphasis{\sphinxupquote{str}}\sphinxstyleliteralemphasis{\sphinxupquote{ | }}\sphinxstyleliteralemphasis{\sphinxupquote{None}}) \textendash{} Optional name.

\item {} 
\sphinxAtStartPar
\sphinxstyleliteralstrong{\sphinxupquote{copy}} (\sphinxstyleliteralemphasis{\sphinxupquote{Square}}\sphinxstyleliteralemphasis{\sphinxupquote{ | }}\sphinxstyleliteralemphasis{\sphinxupquote{None}}) \textendash{} Existing square to reuse.

\end{itemize}

\sphinxlineitem{Returns}
\sphinxAtStartPar
The quadrilateral instance (copy if provided).

\sphinxlineitem{Return type}
\sphinxAtStartPar
Square

\end{description}\end{quote}

\end{fulllineitems}

\index{trapezoid() (pygeox.add\_objects.AddNamespace method)@\spxentry{trapezoid()}\spxextra{pygeox.add\_objects.AddNamespace method}}

\begin{fulllineitems}
\phantomsection\label{\detokenize{pygeox:pygeox.add_objects.AddNamespace.trapezoid}}
\pysigstartsignatures
\pysiglinewithargsret
{\sphinxbfcode{\sphinxupquote{trapezoid}}}
{\sphinxparam{\DUrole{n}{p1}\DUrole{o}{=}\DUrole{default_value}{None}}\sphinxparamcomma \sphinxparam{\DUrole{n}{p2}\DUrole{o}{=}\DUrole{default_value}{None}}\sphinxparamcomma \sphinxparam{\DUrole{n}{p3}\DUrole{o}{=}\DUrole{default_value}{None}}\sphinxparamcomma \sphinxparam{\DUrole{n}{p4}\DUrole{o}{=}\DUrole{default_value}{None}}\sphinxparamcomma \sphinxparam{\DUrole{n}{name}\DUrole{o}{=}\DUrole{default_value}{None}}\sphinxparamcomma \sphinxparam{\DUrole{n}{force\_convex}\DUrole{o}{=}\DUrole{default_value}{False}}\sphinxparamcomma \sphinxparam{\DUrole{n}{copy}\DUrole{o}{=}\DUrole{default_value}{None}}}
{}
\pysigstopsignatures
\sphinxAtStartPar
Add a trapezoid or return a copy.

\sphinxAtStartPar
Constrains AB to be parallel to CD. Optional convexity enforcement.
\begin{quote}\begin{description}
\sphinxlineitem{Parameters}\begin{itemize}
\item {} 
\sphinxAtStartPar
\sphinxstyleliteralstrong{\sphinxupquote{p1}} (\sphinxstyleliteralemphasis{\sphinxupquote{Point}}\sphinxstyleliteralemphasis{\sphinxupquote{ | }}\sphinxstyleliteralemphasis{\sphinxupquote{None}}) \textendash{} Vertex A.

\item {} 
\sphinxAtStartPar
\sphinxstyleliteralstrong{\sphinxupquote{p2}} (\sphinxstyleliteralemphasis{\sphinxupquote{Point}}\sphinxstyleliteralemphasis{\sphinxupquote{ | }}\sphinxstyleliteralemphasis{\sphinxupquote{None}}) \textendash{} Vertex B.

\item {} 
\sphinxAtStartPar
\sphinxstyleliteralstrong{\sphinxupquote{p3}} (\sphinxstyleliteralemphasis{\sphinxupquote{Point}}\sphinxstyleliteralemphasis{\sphinxupquote{ | }}\sphinxstyleliteralemphasis{\sphinxupquote{None}}) \textendash{} Vertex C.

\item {} 
\sphinxAtStartPar
\sphinxstyleliteralstrong{\sphinxupquote{p4}} (\sphinxstyleliteralemphasis{\sphinxupquote{Point}}\sphinxstyleliteralemphasis{\sphinxupquote{ | }}\sphinxstyleliteralemphasis{\sphinxupquote{None}}) \textendash{} Vertex D.

\item {} 
\sphinxAtStartPar
\sphinxstyleliteralstrong{\sphinxupquote{name}} (\sphinxstyleliteralemphasis{\sphinxupquote{str}}\sphinxstyleliteralemphasis{\sphinxupquote{ | }}\sphinxstyleliteralemphasis{\sphinxupquote{None}}) \textendash{} Optional name.

\item {} 
\sphinxAtStartPar
\sphinxstyleliteralstrong{\sphinxupquote{force\_convex}} (\sphinxstyleliteralemphasis{\sphinxupquote{bool}}) \textendash{} If \sphinxcode{\sphinxupquote{True}}, enforces convexity and simplicity.

\item {} 
\sphinxAtStartPar
\sphinxstyleliteralstrong{\sphinxupquote{copy}} (\sphinxstyleliteralemphasis{\sphinxupquote{Trapezoid}}\sphinxstyleliteralemphasis{\sphinxupquote{ | }}\sphinxstyleliteralemphasis{\sphinxupquote{None}}) \textendash{} Existing trapezoid to reuse.

\end{itemize}

\sphinxlineitem{Returns}
\sphinxAtStartPar
The quadrilateral instance (copy if provided).

\sphinxlineitem{Return type}
\sphinxAtStartPar
Trapezoid

\end{description}\end{quote}

\end{fulllineitems}

\index{triangle() (pygeox.add\_objects.AddNamespace method)@\spxentry{triangle()}\spxextra{pygeox.add\_objects.AddNamespace method}}

\begin{fulllineitems}
\phantomsection\label{\detokenize{pygeox:pygeox.add_objects.AddNamespace.triangle}}
\pysigstartsignatures
\pysiglinewithargsret
{\sphinxbfcode{\sphinxupquote{triangle}}}
{\sphinxparam{\DUrole{n}{p1}\DUrole{o}{=}\DUrole{default_value}{None}}\sphinxparamcomma \sphinxparam{\DUrole{n}{p2}\DUrole{o}{=}\DUrole{default_value}{None}}\sphinxparamcomma \sphinxparam{\DUrole{n}{p3}\DUrole{o}{=}\DUrole{default_value}{None}}\sphinxparamcomma \sphinxparam{\DUrole{n}{name}\DUrole{o}{=}\DUrole{default_value}{None}}\sphinxparamcomma \sphinxparam{\DUrole{n}{copy}\DUrole{o}{=}\DUrole{default_value}{None}}}
{}
\pysigstopsignatures
\sphinxAtStartPar
Add a general triangle or return a copy.
\begin{quote}\begin{description}
\sphinxlineitem{Parameters}\begin{itemize}
\item {} 
\sphinxAtStartPar
\sphinxstyleliteralstrong{\sphinxupquote{p1}} (\sphinxstyleliteralemphasis{\sphinxupquote{Point}}\sphinxstyleliteralemphasis{\sphinxupquote{ | }}\sphinxstyleliteralemphasis{\sphinxupquote{None}}) \textendash{} Vertex A.

\item {} 
\sphinxAtStartPar
\sphinxstyleliteralstrong{\sphinxupquote{p2}} (\sphinxstyleliteralemphasis{\sphinxupquote{Point}}\sphinxstyleliteralemphasis{\sphinxupquote{ | }}\sphinxstyleliteralemphasis{\sphinxupquote{None}}) \textendash{} Vertex B.

\item {} 
\sphinxAtStartPar
\sphinxstyleliteralstrong{\sphinxupquote{p3}} (\sphinxstyleliteralemphasis{\sphinxupquote{Point}}\sphinxstyleliteralemphasis{\sphinxupquote{ | }}\sphinxstyleliteralemphasis{\sphinxupquote{None}}) \textendash{} Vertex C.

\item {} 
\sphinxAtStartPar
\sphinxstyleliteralstrong{\sphinxupquote{name}} (\sphinxstyleliteralemphasis{\sphinxupquote{str}}\sphinxstyleliteralemphasis{\sphinxupquote{ | }}\sphinxstyleliteralemphasis{\sphinxupquote{None}}) \textendash{} Optional name.

\item {} 
\sphinxAtStartPar
\sphinxstyleliteralstrong{\sphinxupquote{copy}} (\sphinxstyleliteralemphasis{\sphinxupquote{Triangle}}\sphinxstyleliteralemphasis{\sphinxupquote{ | }}\sphinxstyleliteralemphasis{\sphinxupquote{None}}) \textendash{} Existing triangle to reuse.

\end{itemize}

\sphinxlineitem{Returns}
\sphinxAtStartPar
The triangle instance (copy if provided).

\sphinxlineitem{Return type}
\sphinxAtStartPar
Triangle

\end{description}\end{quote}
\subsubsection*{Notes}

\sphinxAtStartPar
Three distinct points are required when creating.

\end{fulllineitems}


\end{fulllineitems}



\chapter{Add relationships}
\label{\detokenize{pygeox:module-pygeox.relationship}}\label{\detokenize{pygeox:add-relationships}}\index{module@\spxentry{module}!pygeox.relationship@\spxentry{pygeox.relationship}}\index{pygeox.relationship@\spxentry{pygeox.relationship}!module@\spxentry{module}}
\sphinxAtStartPar
RelationshipNamespace class provides geometric relationship constraints
between various geometric objects such as points, lines, circles, arcs, and polygons.
These constraints can be used within a scene to enforce geometric properties like
parallelism, perpendicularity, collinearity, congruence, similarity, and containment.

\sphinxAtStartPar
Each method adds symbolic constraints to the scene’s constraint solver to represent
the desired geometric relationships.
\index{RelationshipNamespace (class in pygeox.relationship)@\spxentry{RelationshipNamespace}\spxextra{class in pygeox.relationship}}

\begin{fulllineitems}
\phantomsection\label{\detokenize{pygeox:pygeox.relationship.RelationshipNamespace}}
\pysigstartsignatures
\pysiglinewithargsret
{\sphinxbfcode{\sphinxupquote{\DUrole{k}{class}\DUrole{w}{ }}}\sphinxcode{\sphinxupquote{pygeox.relationship.}}\sphinxbfcode{\sphinxupquote{RelationshipNamespace}}}
{\sphinxparam{\DUrole{n}{scene}}}
{}
\pysigstopsignatures
\sphinxAtStartPar
Bases: \sphinxcode{\sphinxupquote{object}}

\sphinxAtStartPar
Namespace class for defining geometric relationship constraints
in a scene.
\index{scene (pygeox.relationship.RelationshipNamespace attribute)@\spxentry{scene}\spxextra{pygeox.relationship.RelationshipNamespace attribute}}

\begin{fulllineitems}
\phantomsection\label{\detokenize{pygeox:pygeox.relationship.RelationshipNamespace.scene}}
\pysigstartsignatures
\pysigline
{\sphinxbfcode{\sphinxupquote{scene}}}
\pysigstopsignatures
\sphinxAtStartPar
The geometric scene containing objects and constraint solvers.

\end{fulllineitems}

\index{acute\_angle() (pygeox.relationship.RelationshipNamespace method)@\spxentry{acute\_angle()}\spxextra{pygeox.relationship.RelationshipNamespace method}}

\begin{fulllineitems}
\phantomsection\label{\detokenize{pygeox:pygeox.relationship.RelationshipNamespace.acute_angle}}
\pysigstartsignatures
\pysiglinewithargsret
{\sphinxbfcode{\sphinxupquote{acute\_angle}}}
{\sphinxparam{\DUrole{o}{*}\DUrole{n}{args}}\sphinxparamcomma \sphinxparam{\DUrole{o}{**}\DUrole{n}{kwargs}}}
{}
\pysigstopsignatures
\end{fulllineitems}

\index{angle\_bisector() (pygeox.relationship.RelationshipNamespace method)@\spxentry{angle\_bisector()}\spxextra{pygeox.relationship.RelationshipNamespace method}}

\begin{fulllineitems}
\phantomsection\label{\detokenize{pygeox:pygeox.relationship.RelationshipNamespace.angle_bisector}}
\pysigstartsignatures
\pysiglinewithargsret
{\sphinxbfcode{\sphinxupquote{angle\_bisector}}}
{\sphinxparam{\DUrole{o}{*}\DUrole{n}{args}}\sphinxparamcomma \sphinxparam{\DUrole{o}{**}\DUrole{n}{kwargs}}}
{}
\pysigstopsignatures
\end{fulllineitems}

\index{collinear() (pygeox.relationship.RelationshipNamespace method)@\spxentry{collinear()}\spxextra{pygeox.relationship.RelationshipNamespace method}}

\begin{fulllineitems}
\phantomsection\label{\detokenize{pygeox:pygeox.relationship.RelationshipNamespace.collinear}}
\pysigstartsignatures
\pysiglinewithargsret
{\sphinxbfcode{\sphinxupquote{collinear}}}
{\sphinxparam{\DUrole{o}{*}\DUrole{n}{args}}\sphinxparamcomma \sphinxparam{\DUrole{o}{**}\DUrole{n}{kwargs}}}
{}
\pysigstopsignatures
\end{fulllineitems}

\index{congruent() (pygeox.relationship.RelationshipNamespace method)@\spxentry{congruent()}\spxextra{pygeox.relationship.RelationshipNamespace method}}

\begin{fulllineitems}
\phantomsection\label{\detokenize{pygeox:pygeox.relationship.RelationshipNamespace.congruent}}
\pysigstartsignatures
\pysiglinewithargsret
{\sphinxbfcode{\sphinxupquote{congruent}}}
{\sphinxparam{\DUrole{o}{*}\DUrole{n}{args}}\sphinxparamcomma \sphinxparam{\DUrole{o}{**}\DUrole{n}{kwargs}}}
{}
\pysigstopsignatures
\end{fulllineitems}

\index{inside() (pygeox.relationship.RelationshipNamespace method)@\spxentry{inside()}\spxextra{pygeox.relationship.RelationshipNamespace method}}

\begin{fulllineitems}
\phantomsection\label{\detokenize{pygeox:pygeox.relationship.RelationshipNamespace.inside}}
\pysigstartsignatures
\pysiglinewithargsret
{\sphinxbfcode{\sphinxupquote{inside}}}
{\sphinxparam{\DUrole{o}{*}\DUrole{n}{args}}\sphinxparamcomma \sphinxparam{\DUrole{o}{**}\DUrole{n}{kwargs}}}
{}
\pysigstopsignatures
\end{fulllineitems}

\index{is\_chord() (pygeox.relationship.RelationshipNamespace method)@\spxentry{is\_chord()}\spxextra{pygeox.relationship.RelationshipNamespace method}}

\begin{fulllineitems}
\phantomsection\label{\detokenize{pygeox:pygeox.relationship.RelationshipNamespace.is_chord}}
\pysigstartsignatures
\pysiglinewithargsret
{\sphinxbfcode{\sphinxupquote{is\_chord}}}
{\sphinxparam{\DUrole{o}{*}\DUrole{n}{args}}\sphinxparamcomma \sphinxparam{\DUrole{o}{**}\DUrole{n}{kwargs}}}
{}
\pysigstopsignatures
\end{fulllineitems}

\index{is\_diameter() (pygeox.relationship.RelationshipNamespace method)@\spxentry{is\_diameter()}\spxextra{pygeox.relationship.RelationshipNamespace method}}

\begin{fulllineitems}
\phantomsection\label{\detokenize{pygeox:pygeox.relationship.RelationshipNamespace.is_diameter}}
\pysigstartsignatures
\pysiglinewithargsret
{\sphinxbfcode{\sphinxupquote{is\_diameter}}}
{\sphinxparam{\DUrole{o}{*}\DUrole{n}{args}}\sphinxparamcomma \sphinxparam{\DUrole{o}{**}\DUrole{n}{kwargs}}}
{}
\pysigstopsignatures
\end{fulllineitems}

\index{is\_midpoint() (pygeox.relationship.RelationshipNamespace method)@\spxentry{is\_midpoint()}\spxextra{pygeox.relationship.RelationshipNamespace method}}

\begin{fulllineitems}
\phantomsection\label{\detokenize{pygeox:pygeox.relationship.RelationshipNamespace.is_midpoint}}
\pysigstartsignatures
\pysiglinewithargsret
{\sphinxbfcode{\sphinxupquote{is\_midpoint}}}
{\sphinxparam{\DUrole{o}{*}\DUrole{n}{args}}\sphinxparamcomma \sphinxparam{\DUrole{o}{**}\DUrole{n}{kwargs}}}
{}
\pysigstopsignatures
\end{fulllineitems}

\index{is\_radius() (pygeox.relationship.RelationshipNamespace method)@\spxentry{is\_radius()}\spxextra{pygeox.relationship.RelationshipNamespace method}}

\begin{fulllineitems}
\phantomsection\label{\detokenize{pygeox:pygeox.relationship.RelationshipNamespace.is_radius}}
\pysigstartsignatures
\pysiglinewithargsret
{\sphinxbfcode{\sphinxupquote{is\_radius}}}
{\sphinxparam{\DUrole{o}{*}\DUrole{n}{args}}\sphinxparamcomma \sphinxparam{\DUrole{o}{**}\DUrole{n}{kwargs}}}
{}
\pysigstopsignatures
\end{fulllineitems}

\index{line\_extension\_intersects\_circle\_at() (pygeox.relationship.RelationshipNamespace method)@\spxentry{line\_extension\_intersects\_circle\_at()}\spxextra{pygeox.relationship.RelationshipNamespace method}}

\begin{fulllineitems}
\phantomsection\label{\detokenize{pygeox:pygeox.relationship.RelationshipNamespace.line_extension_intersects_circle_at}}
\pysigstartsignatures
\pysiglinewithargsret
{\sphinxbfcode{\sphinxupquote{line\_extension\_intersects\_circle\_at}}}
{\sphinxparam{\DUrole{o}{*}\DUrole{n}{args}}\sphinxparamcomma \sphinxparam{\DUrole{o}{**}\DUrole{n}{kwargs}}}
{}
\pysigstopsignatures
\end{fulllineitems}

\index{line\_extensions\_intersect\_at() (pygeox.relationship.RelationshipNamespace method)@\spxentry{line\_extensions\_intersect\_at()}\spxextra{pygeox.relationship.RelationshipNamespace method}}

\begin{fulllineitems}
\phantomsection\label{\detokenize{pygeox:pygeox.relationship.RelationshipNamespace.line_extensions_intersect_at}}
\pysigstartsignatures
\pysiglinewithargsret
{\sphinxbfcode{\sphinxupquote{line\_extensions\_intersect\_at}}}
{\sphinxparam{\DUrole{o}{*}\DUrole{n}{args}}\sphinxparamcomma \sphinxparam{\DUrole{o}{**}\DUrole{n}{kwargs}}}
{}
\pysigstopsignatures
\end{fulllineitems}

\index{line\_intersects\_circle\_at() (pygeox.relationship.RelationshipNamespace method)@\spxentry{line\_intersects\_circle\_at()}\spxextra{pygeox.relationship.RelationshipNamespace method}}

\begin{fulllineitems}
\phantomsection\label{\detokenize{pygeox:pygeox.relationship.RelationshipNamespace.line_intersects_circle_at}}
\pysigstartsignatures
\pysiglinewithargsret
{\sphinxbfcode{\sphinxupquote{line\_intersects\_circle\_at}}}
{\sphinxparam{\DUrole{o}{*}\DUrole{n}{args}}\sphinxparamcomma \sphinxparam{\DUrole{o}{**}\DUrole{n}{kwargs}}}
{}
\pysigstopsignatures
\end{fulllineitems}

\index{lines\_intersect\_at() (pygeox.relationship.RelationshipNamespace method)@\spxentry{lines\_intersect\_at()}\spxextra{pygeox.relationship.RelationshipNamespace method}}

\begin{fulllineitems}
\phantomsection\label{\detokenize{pygeox:pygeox.relationship.RelationshipNamespace.lines_intersect_at}}
\pysigstartsignatures
\pysiglinewithargsret
{\sphinxbfcode{\sphinxupquote{lines\_intersect\_at}}}
{\sphinxparam{\DUrole{o}{*}\DUrole{n}{args}}\sphinxparamcomma \sphinxparam{\DUrole{o}{**}\DUrole{n}{kwargs}}}
{}
\pysigstopsignatures
\end{fulllineitems}

\index{mirror\_across\_line() (pygeox.relationship.RelationshipNamespace method)@\spxentry{mirror\_across\_line()}\spxextra{pygeox.relationship.RelationshipNamespace method}}

\begin{fulllineitems}
\phantomsection\label{\detokenize{pygeox:pygeox.relationship.RelationshipNamespace.mirror_across_line}}
\pysigstartsignatures
\pysiglinewithargsret
{\sphinxbfcode{\sphinxupquote{mirror\_across\_line}}}
{\sphinxparam{\DUrole{o}{*}\DUrole{n}{args}}\sphinxparamcomma \sphinxparam{\DUrole{o}{**}\DUrole{n}{kwargs}}}
{}
\pysigstopsignatures
\end{fulllineitems}

\index{obtuse\_angle() (pygeox.relationship.RelationshipNamespace method)@\spxentry{obtuse\_angle()}\spxextra{pygeox.relationship.RelationshipNamespace method}}

\begin{fulllineitems}
\phantomsection\label{\detokenize{pygeox:pygeox.relationship.RelationshipNamespace.obtuse_angle}}
\pysigstartsignatures
\pysiglinewithargsret
{\sphinxbfcode{\sphinxupquote{obtuse\_angle}}}
{\sphinxparam{\DUrole{o}{*}\DUrole{n}{args}}\sphinxparamcomma \sphinxparam{\DUrole{o}{**}\DUrole{n}{kwargs}}}
{}
\pysigstopsignatures
\end{fulllineitems}

\index{outside() (pygeox.relationship.RelationshipNamespace method)@\spxentry{outside()}\spxextra{pygeox.relationship.RelationshipNamespace method}}

\begin{fulllineitems}
\phantomsection\label{\detokenize{pygeox:pygeox.relationship.RelationshipNamespace.outside}}
\pysigstartsignatures
\pysiglinewithargsret
{\sphinxbfcode{\sphinxupquote{outside}}}
{\sphinxparam{\DUrole{o}{*}\DUrole{n}{args}}\sphinxparamcomma \sphinxparam{\DUrole{o}{**}\DUrole{n}{kwargs}}}
{}
\pysigstopsignatures
\end{fulllineitems}

\index{parallel() (pygeox.relationship.RelationshipNamespace method)@\spxentry{parallel()}\spxextra{pygeox.relationship.RelationshipNamespace method}}

\begin{fulllineitems}
\phantomsection\label{\detokenize{pygeox:pygeox.relationship.RelationshipNamespace.parallel}}
\pysigstartsignatures
\pysiglinewithargsret
{\sphinxbfcode{\sphinxupquote{parallel}}}
{\sphinxparam{\DUrole{o}{*}\DUrole{n}{args}}\sphinxparamcomma \sphinxparam{\DUrole{o}{**}\DUrole{n}{kwargs}}}
{}
\pysigstopsignatures
\end{fulllineitems}

\index{perpendicular() (pygeox.relationship.RelationshipNamespace method)@\spxentry{perpendicular()}\spxextra{pygeox.relationship.RelationshipNamespace method}}

\begin{fulllineitems}
\phantomsection\label{\detokenize{pygeox:pygeox.relationship.RelationshipNamespace.perpendicular}}
\pysigstartsignatures
\pysiglinewithargsret
{\sphinxbfcode{\sphinxupquote{perpendicular}}}
{\sphinxparam{\DUrole{o}{*}\DUrole{n}{args}}\sphinxparamcomma \sphinxparam{\DUrole{o}{**}\DUrole{n}{kwargs}}}
{}
\pysigstopsignatures
\end{fulllineitems}

\index{perpendicular\_bisector\_at() (pygeox.relationship.RelationshipNamespace method)@\spxentry{perpendicular\_bisector\_at()}\spxextra{pygeox.relationship.RelationshipNamespace method}}

\begin{fulllineitems}
\phantomsection\label{\detokenize{pygeox:pygeox.relationship.RelationshipNamespace.perpendicular_bisector_at}}
\pysigstartsignatures
\pysiglinewithargsret
{\sphinxbfcode{\sphinxupquote{perpendicular\_bisector\_at}}}
{\sphinxparam{\DUrole{o}{*}\DUrole{n}{args}}\sphinxparamcomma \sphinxparam{\DUrole{o}{**}\DUrole{n}{kwargs}}}
{}
\pysigstopsignatures
\end{fulllineitems}

\index{point\_lies\_on() (pygeox.relationship.RelationshipNamespace method)@\spxentry{point\_lies\_on()}\spxextra{pygeox.relationship.RelationshipNamespace method}}

\begin{fulllineitems}
\phantomsection\label{\detokenize{pygeox:pygeox.relationship.RelationshipNamespace.point_lies_on}}
\pysigstartsignatures
\pysiglinewithargsret
{\sphinxbfcode{\sphinxupquote{point\_lies\_on}}}
{\sphinxparam{\DUrole{o}{*}\DUrole{n}{args}}\sphinxparamcomma \sphinxparam{\DUrole{o}{**}\DUrole{n}{kwargs}}}
{}
\pysigstopsignatures
\end{fulllineitems}

\index{points\_lie\_on() (pygeox.relationship.RelationshipNamespace method)@\spxentry{points\_lie\_on()}\spxextra{pygeox.relationship.RelationshipNamespace method}}

\begin{fulllineitems}
\phantomsection\label{\detokenize{pygeox:pygeox.relationship.RelationshipNamespace.points_lie_on}}
\pysigstartsignatures
\pysiglinewithargsret
{\sphinxbfcode{\sphinxupquote{points\_lie\_on}}}
{\sphinxparam{\DUrole{o}{*}\DUrole{n}{args}}\sphinxparamcomma \sphinxparam{\DUrole{o}{**}\DUrole{n}{kwargs}}}
{}
\pysigstopsignatures
\end{fulllineitems}

\index{right\_angle() (pygeox.relationship.RelationshipNamespace method)@\spxentry{right\_angle()}\spxextra{pygeox.relationship.RelationshipNamespace method}}

\begin{fulllineitems}
\phantomsection\label{\detokenize{pygeox:pygeox.relationship.RelationshipNamespace.right_angle}}
\pysigstartsignatures
\pysiglinewithargsret
{\sphinxbfcode{\sphinxupquote{right\_angle}}}
{\sphinxparam{\DUrole{o}{*}\DUrole{n}{args}}\sphinxparamcomma \sphinxparam{\DUrole{o}{**}\DUrole{n}{kwargs}}}
{}
\pysigstopsignatures
\end{fulllineitems}

\index{rotation\_around\_point() (pygeox.relationship.RelationshipNamespace method)@\spxentry{rotation\_around\_point()}\spxextra{pygeox.relationship.RelationshipNamespace method}}

\begin{fulllineitems}
\phantomsection\label{\detokenize{pygeox:pygeox.relationship.RelationshipNamespace.rotation_around_point}}
\pysigstartsignatures
\pysiglinewithargsret
{\sphinxbfcode{\sphinxupquote{rotation\_around\_point}}}
{\sphinxparam{\DUrole{o}{*}\DUrole{n}{args}}\sphinxparamcomma \sphinxparam{\DUrole{o}{**}\DUrole{n}{kwargs}}}
{}
\pysigstopsignatures
\end{fulllineitems}

\index{scale() (pygeox.relationship.RelationshipNamespace method)@\spxentry{scale()}\spxextra{pygeox.relationship.RelationshipNamespace method}}

\begin{fulllineitems}
\phantomsection\label{\detokenize{pygeox:pygeox.relationship.RelationshipNamespace.scale}}
\pysigstartsignatures
\pysiglinewithargsret
{\sphinxbfcode{\sphinxupquote{scale}}}
{\sphinxparam{\DUrole{o}{*}\DUrole{n}{args}}\sphinxparamcomma \sphinxparam{\DUrole{o}{**}\DUrole{n}{kwargs}}}
{}
\pysigstopsignatures
\end{fulllineitems}

\index{similar() (pygeox.relationship.RelationshipNamespace method)@\spxentry{similar()}\spxextra{pygeox.relationship.RelationshipNamespace method}}

\begin{fulllineitems}
\phantomsection\label{\detokenize{pygeox:pygeox.relationship.RelationshipNamespace.similar}}
\pysigstartsignatures
\pysiglinewithargsret
{\sphinxbfcode{\sphinxupquote{similar}}}
{\sphinxparam{\DUrole{o}{*}\DUrole{n}{args}}\sphinxparamcomma \sphinxparam{\DUrole{o}{**}\DUrole{n}{kwargs}}}
{}
\pysigstopsignatures
\end{fulllineitems}

\index{tangent\_to\_circle() (pygeox.relationship.RelationshipNamespace method)@\spxentry{tangent\_to\_circle()}\spxextra{pygeox.relationship.RelationshipNamespace method}}

\begin{fulllineitems}
\phantomsection\label{\detokenize{pygeox:pygeox.relationship.RelationshipNamespace.tangent_to_circle}}
\pysigstartsignatures
\pysiglinewithargsret
{\sphinxbfcode{\sphinxupquote{tangent\_to\_circle}}}
{\sphinxparam{\DUrole{o}{*}\DUrole{n}{args}}\sphinxparamcomma \sphinxparam{\DUrole{o}{**}\DUrole{n}{kwargs}}}
{}
\pysigstopsignatures
\end{fulllineitems}

\index{translation() (pygeox.relationship.RelationshipNamespace method)@\spxentry{translation()}\spxextra{pygeox.relationship.RelationshipNamespace method}}

\begin{fulllineitems}
\phantomsection\label{\detokenize{pygeox:pygeox.relationship.RelationshipNamespace.translation}}
\pysigstartsignatures
\pysiglinewithargsret
{\sphinxbfcode{\sphinxupquote{translation}}}
{\sphinxparam{\DUrole{o}{*}\DUrole{n}{args}}\sphinxparamcomma \sphinxparam{\DUrole{o}{**}\DUrole{n}{kwargs}}}
{}
\pysigstopsignatures
\end{fulllineitems}


\end{fulllineitems}



\chapter{Add constraints}
\label{\detokenize{pygeox:module-pygeox.constraint}}\label{\detokenize{pygeox:add-constraints}}\index{module@\spxentry{module}!pygeox.constraint@\spxentry{pygeox.constraint}}\index{pygeox.constraint@\spxentry{pygeox.constraint}!module@\spxentry{module}}\index{ConstraintNamespace (class in pygeox.constraint)@\spxentry{ConstraintNamespace}\spxextra{class in pygeox.constraint}}

\begin{fulllineitems}
\phantomsection\label{\detokenize{pygeox:pygeox.constraint.ConstraintNamespace}}
\pysigstartsignatures
\pysiglinewithargsret
{\sphinxbfcode{\sphinxupquote{\DUrole{k}{class}\DUrole{w}{ }}}\sphinxcode{\sphinxupquote{pygeox.constraint.}}\sphinxbfcode{\sphinxupquote{ConstraintNamespace}}}
{\sphinxparam{\DUrole{n}{scene}}}
{}
\pysigstopsignatures
\sphinxAtStartPar
Bases: \sphinxcode{\sphinxupquote{object}}
\index{eq() (pygeox.constraint.ConstraintNamespace method)@\spxentry{eq()}\spxextra{pygeox.constraint.ConstraintNamespace method}}

\begin{fulllineitems}
\phantomsection\label{\detokenize{pygeox:pygeox.constraint.ConstraintNamespace.eq}}
\pysigstartsignatures
\pysiglinewithargsret
{\sphinxbfcode{\sphinxupquote{eq}}}
{\sphinxparam{\DUrole{n}{a}}\sphinxparamcomma \sphinxparam{\DUrole{n}{b}}\sphinxparamcomma \sphinxparam{\DUrole{n}{description}\DUrole{o}{=}\DUrole{default_value}{None}}}
{}
\pysigstopsignatures
\sphinxAtStartPar
Adds an equality constraint (a == b) to the scene.
If ‘a’ is a square root expression, it squares both sides.
\begin{quote}\begin{description}
\sphinxlineitem{Parameters}\begin{itemize}
\item {} 
\sphinxAtStartPar
\sphinxstyleliteralstrong{\sphinxupquote{a}} (\sphinxstyleliteralemphasis{\sphinxupquote{Expr}}\sphinxstyleliteralemphasis{\sphinxupquote{ | }}\sphinxstyleliteralemphasis{\sphinxupquote{int}}\sphinxstyleliteralemphasis{\sphinxupquote{ | }}\sphinxstyleliteralemphasis{\sphinxupquote{float}})

\item {} 
\sphinxAtStartPar
\sphinxstyleliteralstrong{\sphinxupquote{b}} (\sphinxstyleliteralemphasis{\sphinxupquote{Expr}}\sphinxstyleliteralemphasis{\sphinxupquote{ | }}\sphinxstyleliteralemphasis{\sphinxupquote{int}}\sphinxstyleliteralemphasis{\sphinxupquote{ | }}\sphinxstyleliteralemphasis{\sphinxupquote{float}})

\item {} 
\sphinxAtStartPar
\sphinxstyleliteralstrong{\sphinxupquote{description}} (\sphinxstyleliteralemphasis{\sphinxupquote{str}}\sphinxstyleliteralemphasis{\sphinxupquote{ | }}\sphinxstyleliteralemphasis{\sphinxupquote{None}})

\end{itemize}

\sphinxlineitem{Return type}
\sphinxAtStartPar
\sphinxstyleemphasis{Equality} | None

\end{description}\end{quote}

\end{fulllineitems}

\index{geq() (pygeox.constraint.ConstraintNamespace method)@\spxentry{geq()}\spxextra{pygeox.constraint.ConstraintNamespace method}}

\begin{fulllineitems}
\phantomsection\label{\detokenize{pygeox:pygeox.constraint.ConstraintNamespace.geq}}
\pysigstartsignatures
\pysiglinewithargsret
{\sphinxbfcode{\sphinxupquote{geq}}}
{\sphinxparam{\DUrole{n}{a}}\sphinxparamcomma \sphinxparam{\DUrole{n}{b}}\sphinxparamcomma \sphinxparam{\DUrole{n}{description}\DUrole{o}{=}\DUrole{default_value}{None}}}
{}
\pysigstopsignatures
\sphinxAtStartPar
Adds a ‘greater than or equal to’ constraint (a \textgreater{}= b) to the scene.
\begin{quote}\begin{description}
\sphinxlineitem{Parameters}\begin{itemize}
\item {} 
\sphinxAtStartPar
\sphinxstyleliteralstrong{\sphinxupquote{a}} (\sphinxstyleliteralemphasis{\sphinxupquote{Expr}}\sphinxstyleliteralemphasis{\sphinxupquote{ | }}\sphinxstyleliteralemphasis{\sphinxupquote{int}}\sphinxstyleliteralemphasis{\sphinxupquote{ | }}\sphinxstyleliteralemphasis{\sphinxupquote{float}})

\item {} 
\sphinxAtStartPar
\sphinxstyleliteralstrong{\sphinxupquote{b}} (\sphinxstyleliteralemphasis{\sphinxupquote{Expr}}\sphinxstyleliteralemphasis{\sphinxupquote{ | }}\sphinxstyleliteralemphasis{\sphinxupquote{int}}\sphinxstyleliteralemphasis{\sphinxupquote{ | }}\sphinxstyleliteralemphasis{\sphinxupquote{float}})

\item {} 
\sphinxAtStartPar
\sphinxstyleliteralstrong{\sphinxupquote{description}} (\sphinxstyleliteralemphasis{\sphinxupquote{str}}\sphinxstyleliteralemphasis{\sphinxupquote{ | }}\sphinxstyleliteralemphasis{\sphinxupquote{None}})

\end{itemize}

\sphinxlineitem{Return type}
\sphinxAtStartPar
\sphinxstyleemphasis{GreaterThan} | None

\end{description}\end{quote}

\end{fulllineitems}

\index{gt() (pygeox.constraint.ConstraintNamespace method)@\spxentry{gt()}\spxextra{pygeox.constraint.ConstraintNamespace method}}

\begin{fulllineitems}
\phantomsection\label{\detokenize{pygeox:pygeox.constraint.ConstraintNamespace.gt}}
\pysigstartsignatures
\pysiglinewithargsret
{\sphinxbfcode{\sphinxupquote{gt}}}
{\sphinxparam{\DUrole{n}{a}}\sphinxparamcomma \sphinxparam{\DUrole{n}{b}}\sphinxparamcomma \sphinxparam{\DUrole{n}{description}\DUrole{o}{=}\DUrole{default_value}{None}}}
{}
\pysigstopsignatures
\sphinxAtStartPar
Adds a ‘greater than’ constraint (a \textgreater{} b) to the scene.
\begin{quote}\begin{description}
\sphinxlineitem{Parameters}\begin{itemize}
\item {} 
\sphinxAtStartPar
\sphinxstyleliteralstrong{\sphinxupquote{a}} (\sphinxstyleliteralemphasis{\sphinxupquote{Expr}}\sphinxstyleliteralemphasis{\sphinxupquote{ | }}\sphinxstyleliteralemphasis{\sphinxupquote{int}}\sphinxstyleliteralemphasis{\sphinxupquote{ | }}\sphinxstyleliteralemphasis{\sphinxupquote{float}})

\item {} 
\sphinxAtStartPar
\sphinxstyleliteralstrong{\sphinxupquote{b}} (\sphinxstyleliteralemphasis{\sphinxupquote{Expr}}\sphinxstyleliteralemphasis{\sphinxupquote{ | }}\sphinxstyleliteralemphasis{\sphinxupquote{int}}\sphinxstyleliteralemphasis{\sphinxupquote{ | }}\sphinxstyleliteralemphasis{\sphinxupquote{float}})

\item {} 
\sphinxAtStartPar
\sphinxstyleliteralstrong{\sphinxupquote{description}} (\sphinxstyleliteralemphasis{\sphinxupquote{str}}\sphinxstyleliteralemphasis{\sphinxupquote{ | }}\sphinxstyleliteralemphasis{\sphinxupquote{None}})

\end{itemize}

\sphinxlineitem{Return type}
\sphinxAtStartPar
\sphinxstyleemphasis{StrictGreaterThan} | None

\end{description}\end{quote}

\end{fulllineitems}

\index{leq() (pygeox.constraint.ConstraintNamespace method)@\spxentry{leq()}\spxextra{pygeox.constraint.ConstraintNamespace method}}

\begin{fulllineitems}
\phantomsection\label{\detokenize{pygeox:pygeox.constraint.ConstraintNamespace.leq}}
\pysigstartsignatures
\pysiglinewithargsret
{\sphinxbfcode{\sphinxupquote{leq}}}
{\sphinxparam{\DUrole{n}{a}}\sphinxparamcomma \sphinxparam{\DUrole{n}{b}}\sphinxparamcomma \sphinxparam{\DUrole{n}{description}\DUrole{o}{=}\DUrole{default_value}{None}}}
{}
\pysigstopsignatures
\sphinxAtStartPar
Adds a ‘less than or equal to’ constraint (a \textless{}= b) to the scene.
\begin{quote}\begin{description}
\sphinxlineitem{Parameters}\begin{itemize}
\item {} 
\sphinxAtStartPar
\sphinxstyleliteralstrong{\sphinxupquote{a}} (\sphinxstyleliteralemphasis{\sphinxupquote{Expr}}\sphinxstyleliteralemphasis{\sphinxupquote{ | }}\sphinxstyleliteralemphasis{\sphinxupquote{int}}\sphinxstyleliteralemphasis{\sphinxupquote{ | }}\sphinxstyleliteralemphasis{\sphinxupquote{float}})

\item {} 
\sphinxAtStartPar
\sphinxstyleliteralstrong{\sphinxupquote{b}} (\sphinxstyleliteralemphasis{\sphinxupquote{Expr}}\sphinxstyleliteralemphasis{\sphinxupquote{ | }}\sphinxstyleliteralemphasis{\sphinxupquote{int}}\sphinxstyleliteralemphasis{\sphinxupquote{ | }}\sphinxstyleliteralemphasis{\sphinxupquote{float}})

\item {} 
\sphinxAtStartPar
\sphinxstyleliteralstrong{\sphinxupquote{description}} (\sphinxstyleliteralemphasis{\sphinxupquote{str}}\sphinxstyleliteralemphasis{\sphinxupquote{ | }}\sphinxstyleliteralemphasis{\sphinxupquote{None}})

\end{itemize}

\sphinxlineitem{Return type}
\sphinxAtStartPar
\sphinxstyleemphasis{LessThan} | None

\end{description}\end{quote}

\end{fulllineitems}

\index{lt() (pygeox.constraint.ConstraintNamespace method)@\spxentry{lt()}\spxextra{pygeox.constraint.ConstraintNamespace method}}

\begin{fulllineitems}
\phantomsection\label{\detokenize{pygeox:pygeox.constraint.ConstraintNamespace.lt}}
\pysigstartsignatures
\pysiglinewithargsret
{\sphinxbfcode{\sphinxupquote{lt}}}
{\sphinxparam{\DUrole{n}{a}}\sphinxparamcomma \sphinxparam{\DUrole{n}{b}}\sphinxparamcomma \sphinxparam{\DUrole{n}{description}\DUrole{o}{=}\DUrole{default_value}{None}}}
{}
\pysigstopsignatures
\sphinxAtStartPar
Adds a ‘less than’ constraint (a \textless{} b) to the scene.
\begin{quote}\begin{description}
\sphinxlineitem{Parameters}\begin{itemize}
\item {} 
\sphinxAtStartPar
\sphinxstyleliteralstrong{\sphinxupquote{a}} (\sphinxstyleliteralemphasis{\sphinxupquote{Expr}}\sphinxstyleliteralemphasis{\sphinxupquote{ | }}\sphinxstyleliteralemphasis{\sphinxupquote{int}}\sphinxstyleliteralemphasis{\sphinxupquote{ | }}\sphinxstyleliteralemphasis{\sphinxupquote{float}})

\item {} 
\sphinxAtStartPar
\sphinxstyleliteralstrong{\sphinxupquote{b}} (\sphinxstyleliteralemphasis{\sphinxupquote{Expr}}\sphinxstyleliteralemphasis{\sphinxupquote{ | }}\sphinxstyleliteralemphasis{\sphinxupquote{int}}\sphinxstyleliteralemphasis{\sphinxupquote{ | }}\sphinxstyleliteralemphasis{\sphinxupquote{float}})

\item {} 
\sphinxAtStartPar
\sphinxstyleliteralstrong{\sphinxupquote{description}} (\sphinxstyleliteralemphasis{\sphinxupquote{str}}\sphinxstyleliteralemphasis{\sphinxupquote{ | }}\sphinxstyleliteralemphasis{\sphinxupquote{None}})

\end{itemize}

\sphinxlineitem{Return type}
\sphinxAtStartPar
\sphinxstyleemphasis{StrictLessThan} | None

\end{description}\end{quote}

\end{fulllineitems}

\index{neq() (pygeox.constraint.ConstraintNamespace method)@\spxentry{neq()}\spxextra{pygeox.constraint.ConstraintNamespace method}}

\begin{fulllineitems}
\phantomsection\label{\detokenize{pygeox:pygeox.constraint.ConstraintNamespace.neq}}
\pysigstartsignatures
\pysiglinewithargsret
{\sphinxbfcode{\sphinxupquote{neq}}}
{\sphinxparam{\DUrole{n}{a}}\sphinxparamcomma \sphinxparam{\DUrole{n}{b}}\sphinxparamcomma \sphinxparam{\DUrole{n}{description}\DUrole{o}{=}\DUrole{default_value}{None}}}
{}
\pysigstopsignatures
\sphinxAtStartPar
Adds an inequality constraint (a != b) to the scene.
\begin{quote}\begin{description}
\sphinxlineitem{Parameters}\begin{itemize}
\item {} 
\sphinxAtStartPar
\sphinxstyleliteralstrong{\sphinxupquote{a}} (\sphinxstyleliteralemphasis{\sphinxupquote{Expr}}\sphinxstyleliteralemphasis{\sphinxupquote{ | }}\sphinxstyleliteralemphasis{\sphinxupquote{int}}\sphinxstyleliteralemphasis{\sphinxupquote{ | }}\sphinxstyleliteralemphasis{\sphinxupquote{float}})

\item {} 
\sphinxAtStartPar
\sphinxstyleliteralstrong{\sphinxupquote{b}} (\sphinxstyleliteralemphasis{\sphinxupquote{Expr}}\sphinxstyleliteralemphasis{\sphinxupquote{ | }}\sphinxstyleliteralemphasis{\sphinxupquote{int}}\sphinxstyleliteralemphasis{\sphinxupquote{ | }}\sphinxstyleliteralemphasis{\sphinxupquote{float}})

\item {} 
\sphinxAtStartPar
\sphinxstyleliteralstrong{\sphinxupquote{description}} (\sphinxstyleliteralemphasis{\sphinxupquote{str}}\sphinxstyleliteralemphasis{\sphinxupquote{ | }}\sphinxstyleliteralemphasis{\sphinxupquote{None}})

\end{itemize}

\sphinxlineitem{Return type}
\sphinxAtStartPar
\sphinxstyleemphasis{Unequality} | None

\end{description}\end{quote}

\end{fulllineitems}


\end{fulllineitems}



\chapter{Solvers}
\label{\detokenize{pygeox:module-pygeox.solver}}\label{\detokenize{pygeox:solvers}}\index{module@\spxentry{module}!pygeox.solver@\spxentry{pygeox.solver}}\index{pygeox.solver@\spxentry{pygeox.solver}!module@\spxentry{module}}\index{SolverNamespace (class in pygeox.solver)@\spxentry{SolverNamespace}\spxextra{class in pygeox.solver}}

\begin{fulllineitems}
\phantomsection\label{\detokenize{pygeox:pygeox.solver.SolverNamespace}}
\pysigstartsignatures
\pysiglinewithargsret
{\sphinxbfcode{\sphinxupquote{\DUrole{k}{class}\DUrole{w}{ }}}\sphinxcode{\sphinxupquote{pygeox.solver.}}\sphinxbfcode{\sphinxupquote{SolverNamespace}}}
{\sphinxparam{\DUrole{n}{scene}}}
{}
\pysigstopsignatures
\sphinxAtStartPar
Bases: \sphinxcode{\sphinxupquote{object}}
\index{analytical() (pygeox.solver.SolverNamespace method)@\spxentry{analytical()}\spxextra{pygeox.solver.SolverNamespace method}}

\begin{fulllineitems}
\phantomsection\label{\detokenize{pygeox:pygeox.solver.SolverNamespace.analytical}}
\pysigstartsignatures
\pysiglinewithargsret
{\sphinxbfcode{\sphinxupquote{analytical}}}
{}
{}
\pysigstopsignatures
\sphinxAtStartPar
Attempt to solve the constraint system symbolically (analytically).
\begin{description}
\sphinxlineitem{Steps:}\begin{itemize}
\item {} 
\sphinxAtStartPar
Substitute mapping constraints.

\item {} 
\sphinxAtStartPar
Separate equations and inequalities.

\item {} 
\sphinxAtStartPar
Solve equalities using SymPy.

\item {} 
\sphinxAtStartPar
Filter and validate solutions.

\item {} 
\sphinxAtStartPar
Return unique solution or multiple candidates.

\end{itemize}

\end{description}
\begin{quote}\begin{description}
\sphinxlineitem{Raises}
\sphinxAtStartPar
\sphinxstyleliteralstrong{\sphinxupquote{ValueError}} \textendash{} If no solution or invalid solutions found.

\sphinxlineitem{Returns}
\sphinxAtStartPar
Solution mapping symbols to values.

\sphinxlineitem{Return type}
\sphinxAtStartPar
dict or list of dicts

\end{description}\end{quote}

\end{fulllineitems}

\index{get\_objective\_function() (pygeox.solver.SolverNamespace method)@\spxentry{get\_objective\_function()}\spxextra{pygeox.solver.SolverNamespace method}}

\begin{fulllineitems}
\phantomsection\label{\detokenize{pygeox:pygeox.solver.SolverNamespace.get_objective_function}}
\pysigstartsignatures
\pysiglinewithargsret
{\sphinxbfcode{\sphinxupquote{get\_objective\_function}}}
{\sphinxparam{\DUrole{n}{variables}}\sphinxparamcomma \sphinxparam{\DUrole{n}{equations}}\sphinxparamcomma \sphinxparam{\DUrole{n}{alpha}\DUrole{o}{=}\DUrole{default_value}{1.0}}\sphinxparamcomma \sphinxparam{\DUrole{n}{beta}\DUrole{o}{=}\DUrole{default_value}{1.0}}\sphinxparamcomma \sphinxparam{\DUrole{n}{strict\_penalty}\DUrole{o}{=}\DUrole{default_value}{1.0}}\sphinxparamcomma \sphinxparam{\DUrole{n}{strict\_tol}\DUrole{o}{=}\DUrole{default_value}{0.0001}}\sphinxparamcomma \sphinxparam{\DUrole{n}{use\_numba}\DUrole{o}{=}\DUrole{default_value}{False}}}
{}
\pysigstopsignatures
\sphinxAtStartPar
Build the objective function for numerical optimization from constraints.

\sphinxAtStartPar
Supports both pure Python and Numba\sphinxhyphen{}jitted versions.
\begin{quote}\begin{description}
\sphinxlineitem{Parameters}\begin{itemize}
\item {} 
\sphinxAtStartPar
\sphinxstyleliteralstrong{\sphinxupquote{variables}} \textendash{} List of variables to optimize.

\item {} 
\sphinxAtStartPar
\sphinxstyleliteralstrong{\sphinxupquote{equations}} \textendash{} List of (description, sympy expression) constraints.

\item {} 
\sphinxAtStartPar
\sphinxstyleliteralstrong{\sphinxupquote{alpha}} (\sphinxstyleliteralemphasis{\sphinxupquote{float}}) \textendash{} Weight for equality constraints.

\item {} 
\sphinxAtStartPar
\sphinxstyleliteralstrong{\sphinxupquote{beta}} (\sphinxstyleliteralemphasis{\sphinxupquote{float}}) \textendash{} Weight for inequality constraints.

\item {} 
\sphinxAtStartPar
\sphinxstyleliteralstrong{\sphinxupquote{strict\_penalty}} (\sphinxstyleliteralemphasis{\sphinxupquote{float}}) \textendash{} Penalty for strict inequality violations.

\item {} 
\sphinxAtStartPar
\sphinxstyleliteralstrong{\sphinxupquote{strict\_tol}} (\sphinxstyleliteralemphasis{\sphinxupquote{float}}) \textendash{} Tolerance for strict inequalities.

\item {} 
\sphinxAtStartPar
\sphinxstyleliteralstrong{\sphinxupquote{use\_numba}} (\sphinxstyleliteralemphasis{\sphinxupquote{bool}}) \textendash{} If True, return a Numba\sphinxhyphen{}accelerated function.

\end{itemize}

\sphinxlineitem{Returns}
\sphinxAtStartPar
(objective function, variables, dictionary of sub\sphinxhyphen{}objectives)

\sphinxlineitem{Return type}
\sphinxAtStartPar
tuple

\end{description}\end{quote}

\end{fulllineitems}

\index{has\_any\_overlapping\_points() (pygeox.solver.SolverNamespace method)@\spxentry{has\_any\_overlapping\_points()}\spxextra{pygeox.solver.SolverNamespace method}}

\begin{fulllineitems}
\phantomsection\label{\detokenize{pygeox:pygeox.solver.SolverNamespace.has_any_overlapping_points}}
\pysigstartsignatures
\pysiglinewithargsret
{\sphinxbfcode{\sphinxupquote{has\_any\_overlapping\_points}}}
{\sphinxparam{\DUrole{n}{threshold}\DUrole{o}{=}\DUrole{default_value}{0.0001}}}
{}
\pysigstopsignatures
\end{fulllineitems}

\index{numerical() (pygeox.solver.SolverNamespace method)@\spxentry{numerical()}\spxextra{pygeox.solver.SolverNamespace method}}

\begin{fulllineitems}
\phantomsection\label{\detokenize{pygeox:pygeox.solver.SolverNamespace.numerical}}
\pysigstartsignatures
\pysiglinewithargsret
{\sphinxbfcode{\sphinxupquote{numerical}}}
{\sphinxparam{\DUrole{n}{method}\DUrole{o}{=}\DUrole{default_value}{\textquotesingle{}basinhoping\textquotesingle{}}}\sphinxparamcomma \sphinxparam{\DUrole{n}{seed}\DUrole{o}{=}\DUrole{default_value}{None}}\sphinxparamcomma \sphinxparam{\DUrole{n}{alpha}\DUrole{o}{=}\DUrole{default_value}{1.0}}\sphinxparamcomma \sphinxparam{\DUrole{n}{beta}\DUrole{o}{=}\DUrole{default_value}{1.0}}\sphinxparamcomma \sphinxparam{\DUrole{n}{strict\_penalty}\DUrole{o}{=}\DUrole{default_value}{1.0}}\sphinxparamcomma \sphinxparam{\DUrole{n}{strict\_tol}\DUrole{o}{=}\DUrole{default_value}{0.0001}}\sphinxparamcomma \sphinxparam{\DUrole{n}{distance\_penalty}\DUrole{o}{=}\DUrole{default_value}{1}}\sphinxparamcomma \sphinxparam{\DUrole{n}{min\_dist}\DUrole{o}{=}\DUrole{default_value}{None}}\sphinxparamcomma \sphinxparam{\DUrole{n}{use\_numba}\DUrole{o}{=}\DUrole{default_value}{True}}\sphinxparamcomma \sphinxparam{\DUrole{n}{save\_param\_history}\DUrole{o}{=}\DUrole{default_value}{False}}\sphinxparamcomma \sphinxparam{\DUrole{n}{verbose}\DUrole{o}{=}\DUrole{default_value}{True}}}
{}
\pysigstopsignatures
\sphinxAtStartPar
Solve the constraint system numerically using global or local optimization.

\sphinxAtStartPar
Supports methods like basinhopping, L\sphinxhyphen{}BFGS\sphinxhyphen{}B, and dual\_annealing.
\begin{quote}\begin{description}
\sphinxlineitem{Parameters}\begin{itemize}
\item {} 
\sphinxAtStartPar
\sphinxstyleliteralstrong{\sphinxupquote{method}} \textendash{} Optimization method name.

\item {} 
\sphinxAtStartPar
\sphinxstyleliteralstrong{\sphinxupquote{seed}} \textendash{} Random seed.

\item {} 
\sphinxAtStartPar
\sphinxstyleliteralstrong{\sphinxupquote{alpha}} \textendash{} Penalty weights and tolerances.

\item {} 
\sphinxAtStartPar
\sphinxstyleliteralstrong{\sphinxupquote{beta}} \textendash{} Penalty weights and tolerances.

\item {} 
\sphinxAtStartPar
\sphinxstyleliteralstrong{\sphinxupquote{strict\_penalty}} \textendash{} Penalty weights and tolerances.

\item {} 
\sphinxAtStartPar
\sphinxstyleliteralstrong{\sphinxupquote{strict\_tol}} \textendash{} Penalty weights and tolerances.

\item {} 
\sphinxAtStartPar
\sphinxstyleliteralstrong{\sphinxupquote{distance\_penalty}} (\sphinxstyleliteralemphasis{\sphinxupquote{float}}) \textendash{} Weight for penalty avoiding degenerate configurations.

\item {} 
\sphinxAtStartPar
\sphinxstyleliteralstrong{\sphinxupquote{min\_dist}} (\sphinxstyleliteralemphasis{\sphinxupquote{float}}\sphinxstyleliteralemphasis{\sphinxupquote{ | }}\sphinxstyleliteralemphasis{\sphinxupquote{None}}) \textendash{} Minimum allowed distance between points.

\item {} 
\sphinxAtStartPar
\sphinxstyleliteralstrong{\sphinxupquote{use\_numba}} (\sphinxstyleliteralemphasis{\sphinxupquote{bool}}) \textendash{} Use Numba\sphinxhyphen{}accelerated objective function.

\item {} 
\sphinxAtStartPar
\sphinxstyleliteralstrong{\sphinxupquote{save\_param\_history}} (\sphinxstyleliteralemphasis{\sphinxupquote{bool}}) \textendash{} Save parameter values during optimization.

\item {} 
\sphinxAtStartPar
\sphinxstyleliteralstrong{\sphinxupquote{verbose}} (\sphinxstyleliteralemphasis{\sphinxupquote{bool}}) \textendash{} Print detailed solver progress.

\end{itemize}

\sphinxlineitem{Returns}
\sphinxAtStartPar
Dictionary with all objects with scene and solved numerical values.

\sphinxlineitem{Return type}
\sphinxAtStartPar
dict

\sphinxlineitem{Raises}
\sphinxAtStartPar
\sphinxstyleliteralstrong{\sphinxupquote{ValueError}} \textendash{} If optimization fails.

\end{description}\end{quote}

\end{fulllineitems}

\index{numerical\_analysis() (pygeox.solver.SolverNamespace method)@\spxentry{numerical\_analysis()}\spxextra{pygeox.solver.SolverNamespace method}}

\begin{fulllineitems}
\phantomsection\label{\detokenize{pygeox:pygeox.solver.SolverNamespace.numerical_analysis}}
\pysigstartsignatures
\pysiglinewithargsret
{\sphinxbfcode{\sphinxupquote{numerical\_analysis}}}
{\sphinxparam{\DUrole{n}{x}\DUrole{o}{=}\DUrole{default_value}{None}}\sphinxparamcomma \sphinxparam{\DUrole{n}{verbose}\DUrole{o}{=}\DUrole{default_value}{True}}}
{}
\pysigstopsignatures
\sphinxAtStartPar
Evaluate and report the quality of a numerical solution.

\sphinxAtStartPar
Prints objective value, satisfaction of equality, inequality, strict,
and not\sphinxhyphen{}equal constraints, and warns about overlapping points.
\begin{quote}\begin{description}
\sphinxlineitem{Parameters}\begin{itemize}
\item {} 
\sphinxAtStartPar
\sphinxstyleliteralstrong{\sphinxupquote{x}} \textendash{} Optional solution vector to evaluate.

\item {} 
\sphinxAtStartPar
\sphinxstyleliteralstrong{\sphinxupquote{verbose}} \textendash{} Whether to print detailed report.

\end{itemize}

\sphinxlineitem{Returns}
\sphinxAtStartPar
Analysis summary including objective and constraint values.

\sphinxlineitem{Return type}
\sphinxAtStartPar
dict

\end{description}\end{quote}

\end{fulllineitems}

\index{pick\_solution() (pygeox.solver.SolverNamespace method)@\spxentry{pick\_solution()}\spxextra{pygeox.solver.SolverNamespace method}}

\begin{fulllineitems}
\phantomsection\label{\detokenize{pygeox:pygeox.solver.SolverNamespace.pick_solution}}
\pysigstartsignatures
\pysiglinewithargsret
{\sphinxbfcode{\sphinxupquote{pick\_solution}}}
{\sphinxparam{\DUrole{n}{index}\DUrole{o}{=}\DUrole{default_value}{None}}\sphinxparamcomma \sphinxparam{\DUrole{n}{ind\_var\_values}\DUrole{o}{=}\DUrole{default_value}{None}}}
{}
\pysigstopsignatures
\sphinxAtStartPar
Select and apply one solution among multiple symbolic solutions.
\begin{quote}\begin{description}
\sphinxlineitem{Parameters}\begin{itemize}
\item {} 
\sphinxAtStartPar
\sphinxstyleliteralstrong{\sphinxupquote{index}} (\sphinxstyleliteralemphasis{\sphinxupquote{int}}) \textendash{} Index of the solution to pick (random if None).

\item {} 
\sphinxAtStartPar
\sphinxstyleliteralstrong{\sphinxupquote{ind\_var\_values}} (\sphinxstyleliteralemphasis{\sphinxupquote{dict}}) \textendash{} Optional fixed values for independent variables.

\end{itemize}

\sphinxlineitem{Raises}
\sphinxAtStartPar
\sphinxstyleliteralstrong{\sphinxupquote{ValueError}} \textendash{} If solution contains invalid values or fails inequalities.

\end{description}\end{quote}

\end{fulllineitems}

\index{simplify\_constraints() (pygeox.solver.SolverNamespace method)@\spxentry{simplify\_constraints()}\spxextra{pygeox.solver.SolverNamespace method}}

\begin{fulllineitems}
\phantomsection\label{\detokenize{pygeox:pygeox.solver.SolverNamespace.simplify_constraints}}
\pysigstartsignatures
\pysiglinewithargsret
{\sphinxbfcode{\sphinxupquote{simplify\_constraints}}}
{\sphinxparam{\DUrole{n}{max\_linear}\DUrole{o}{=}\DUrole{default_value}{0}}}
{}
\pysigstopsignatures
\sphinxAtStartPar
Simplify the constraint system by applying known variable mappings, removing
trivial equalities, and optionally performing cheap linear eliminations.

\sphinxAtStartPar
This reduces the number of independent variables and constraints, aiding solver efficiency.
\begin{quote}\begin{description}
\sphinxlineitem{Parameters}
\sphinxAtStartPar
\sphinxstyleliteralstrong{\sphinxupquote{max\_linear}} (\sphinxstyleliteralemphasis{\sphinxupquote{int}}) \textendash{} Maximum number of additional cheap linear eliminations to attempt.

\sphinxlineitem{Returns}
\sphinxAtStartPar
\begin{itemize}
\item {} 
\sphinxAtStartPar
indep\_vars (list{[}sympy.Symbol{]}): Independent variables after simplification.

\item {} 
\sphinxAtStartPar
reduced\_eqs (list{[}sympy.Expr{]}): Simplified constraint expressions (lhs \sphinxhyphen{} rhs).

\item {} 
\sphinxAtStartPar
mapping (dict{[}sympy.Symbol, sympy.Expr{]}): Mapping from dependent to independent variables.

\item {} 
\sphinxAtStartPar
red\_full\_mapping (dict{[}sympy.Symbol, sympy.Expr{]}): Final mapping of variables to expressions.

\end{itemize}


\sphinxlineitem{Return type}
\sphinxAtStartPar
tuple

\end{description}\end{quote}

\end{fulllineitems}

\index{substitute\_in\_obj() (pygeox.solver.SolverNamespace static method)@\spxentry{substitute\_in\_obj()}\spxextra{pygeox.solver.SolverNamespace static method}}

\begin{fulllineitems}
\phantomsection\label{\detokenize{pygeox:pygeox.solver.SolverNamespace.substitute_in_obj}}
\pysigstartsignatures
\pysiglinewithargsret
{\sphinxbfcode{\sphinxupquote{\DUrole{k}{static}\DUrole{w}{ }}}\sphinxbfcode{\sphinxupquote{substitute\_in\_obj}}}
{\sphinxparam{\DUrole{n}{obj}}\sphinxparamcomma \sphinxparam{\DUrole{n}{sol}}\sphinxparamcomma \sphinxparam{\DUrole{n}{numeric}\DUrole{o}{=}\DUrole{default_value}{False}}\sphinxparamcomma \sphinxparam{\DUrole{n}{visited}\DUrole{o}{=}\DUrole{default_value}{None}}}
{}
\pysigstopsignatures\begin{quote}\begin{description}
\sphinxlineitem{Parameters}
\sphinxAtStartPar
\sphinxstyleliteralstrong{\sphinxupquote{numeric}} (\sphinxstyleliteralemphasis{\sphinxupquote{bool}})

\end{description}\end{quote}

\end{fulllineitems}


\end{fulllineitems}



\chapter{Basic Objects (point, lines, circles, arcs)}
\label{\detokenize{pygeox:basic-objects-point-lines-circles-arcs}}\index{module@\spxentry{module}!pygeox.basic\_objects@\spxentry{pygeox.basic\_objects}}\index{pygeox.basic\_objects@\spxentry{pygeox.basic\_objects}!module@\spxentry{module}}\phantomsection\label{\detokenize{pygeox:module-pygeox.basic_objects}}
\sphinxAtStartPar
basic\_objects.py: Fundamental 2D geometric primitives with symbolic or numeric coordinates.

\sphinxAtStartPar
Defines Point, LineLike, LineSegment, Line, Ray, Circle, and Arc classes with
methods for geometric computations such as distance, angle, slope, and length.
\index{Angle (class in pygeox.basic\_objects)@\spxentry{Angle}\spxextra{class in pygeox.basic\_objects}}

\begin{fulllineitems}
\phantomsection\label{\detokenize{pygeox:pygeox.basic_objects.Angle}}
\pysigstartsignatures
\pysiglinewithargsret
{\sphinxbfcode{\sphinxupquote{\DUrole{k}{class}\DUrole{w}{ }}}\sphinxcode{\sphinxupquote{pygeox.basic\_objects.}}\sphinxbfcode{\sphinxupquote{Angle}}}
{\sphinxparam{\DUrole{n}{A}}\sphinxparamcomma \sphinxparam{\DUrole{n}{vertex}}\sphinxparamcomma \sphinxparam{\DUrole{n}{B}}\sphinxparamcomma \sphinxparam{\DUrole{n}{plot\_sign}\DUrole{o}{=}\DUrole{default_value}{True}}\sphinxparamcomma \sphinxparam{\DUrole{n}{plot\_text}\DUrole{o}{=}\DUrole{default_value}{False}}\sphinxparamcomma \sphinxparam{\DUrole{n}{name}\DUrole{o}{=}\DUrole{default_value}{None}}}
{}
\pysigstopsignatures
\sphinxAtStartPar
Bases: \sphinxcode{\sphinxupquote{object}}

\sphinxAtStartPar
Represents an angle defined by three points with a specified vertex.
\begin{quote}\begin{description}
\sphinxlineitem{Parameters}\begin{itemize}
\item {} 
\sphinxAtStartPar
\sphinxstyleliteralstrong{\sphinxupquote{A}} (\sphinxstyleliteralemphasis{\sphinxupquote{Point}})

\item {} 
\sphinxAtStartPar
\sphinxstyleliteralstrong{\sphinxupquote{vertex}} (\sphinxstyleliteralemphasis{\sphinxupquote{Point}})

\item {} 
\sphinxAtStartPar
\sphinxstyleliteralstrong{\sphinxupquote{B}} (\sphinxstyleliteralemphasis{\sphinxupquote{Point}})

\item {} 
\sphinxAtStartPar
\sphinxstyleliteralstrong{\sphinxupquote{plot\_sign}} (\sphinxstyleliteralemphasis{\sphinxupquote{bool}}\sphinxstyleliteralemphasis{\sphinxupquote{ | }}\sphinxstyleliteralemphasis{\sphinxupquote{None}})

\item {} 
\sphinxAtStartPar
\sphinxstyleliteralstrong{\sphinxupquote{plot\_text}} (\sphinxstyleliteralemphasis{\sphinxupquote{bool}}\sphinxstyleliteralemphasis{\sphinxupquote{ | }}\sphinxstyleliteralemphasis{\sphinxupquote{None}})

\item {} 
\sphinxAtStartPar
\sphinxstyleliteralstrong{\sphinxupquote{name}} (\sphinxstyleliteralemphasis{\sphinxupquote{str}}\sphinxstyleliteralemphasis{\sphinxupquote{ | }}\sphinxstyleliteralemphasis{\sphinxupquote{None}})

\end{itemize}

\end{description}\end{quote}
\index{A (pygeox.basic\_objects.Angle attribute)@\spxentry{A}\spxextra{pygeox.basic\_objects.Angle attribute}}

\begin{fulllineitems}
\phantomsection\label{\detokenize{pygeox:pygeox.basic_objects.Angle.A}}
\pysigstartsignatures
\pysigline
{\sphinxbfcode{\sphinxupquote{A}}}
\pysigstopsignatures
\sphinxAtStartPar
One endpoint of the angle.
\begin{quote}\begin{description}
\sphinxlineitem{Type}
\sphinxAtStartPar
Point

\end{description}\end{quote}

\end{fulllineitems}

\index{vertex (pygeox.basic\_objects.Angle attribute)@\spxentry{vertex}\spxextra{pygeox.basic\_objects.Angle attribute}}

\begin{fulllineitems}
\phantomsection\label{\detokenize{pygeox:pygeox.basic_objects.Angle.vertex}}
\pysigstartsignatures
\pysigline
{\sphinxbfcode{\sphinxupquote{vertex}}}
\pysigstopsignatures
\sphinxAtStartPar
The vertex point of the angle.
\begin{quote}\begin{description}
\sphinxlineitem{Type}
\sphinxAtStartPar
Point

\end{description}\end{quote}

\end{fulllineitems}

\index{B (pygeox.basic\_objects.Angle attribute)@\spxentry{B}\spxextra{pygeox.basic\_objects.Angle attribute}}

\begin{fulllineitems}
\phantomsection\label{\detokenize{pygeox:pygeox.basic_objects.Angle.B}}
\pysigstartsignatures
\pysigline
{\sphinxbfcode{\sphinxupquote{B}}}
\pysigstopsignatures
\sphinxAtStartPar
The other endpoint of the angle.
\begin{quote}\begin{description}
\sphinxlineitem{Type}
\sphinxAtStartPar
Point

\end{description}\end{quote}

\end{fulllineitems}

\index{plot (pygeox.basic\_objects.Angle attribute)@\spxentry{plot}\spxextra{pygeox.basic\_objects.Angle attribute}}

\begin{fulllineitems}
\phantomsection\label{\detokenize{pygeox:pygeox.basic_objects.Angle.plot}}
\pysigstartsignatures
\pysigline
{\sphinxbfcode{\sphinxupquote{plot}}}
\pysigstopsignatures
\sphinxAtStartPar
Whether to plot the angle.
\begin{quote}\begin{description}
\sphinxlineitem{Type}
\sphinxAtStartPar
bool

\end{description}\end{quote}

\end{fulllineitems}

\index{name (pygeox.basic\_objects.Angle attribute)@\spxentry{name}\spxextra{pygeox.basic\_objects.Angle attribute}}

\begin{fulllineitems}
\phantomsection\label{\detokenize{pygeox:pygeox.basic_objects.Angle.name}}
\pysigstartsignatures
\pysigline
{\sphinxbfcode{\sphinxupquote{name}}}
\pysigstopsignatures
\sphinxAtStartPar
Name of the angle.
\begin{quote}\begin{description}
\sphinxlineitem{Type}
\sphinxAtStartPar
str

\end{description}\end{quote}

\end{fulllineitems}

\begin{description}
\sphinxlineitem{Properties:}
\sphinxAtStartPar
value: The angle measure at the vertex in degrees (or symbolic expression).

\end{description}
\index{value (pygeox.basic\_objects.Angle property)@\spxentry{value}\spxextra{pygeox.basic\_objects.Angle property}}

\begin{fulllineitems}
\phantomsection\label{\detokenize{pygeox:pygeox.basic_objects.Angle.value}}
\pysigstartsignatures
\pysigline
{\sphinxbfcode{\sphinxupquote{\DUrole{k}{property}\DUrole{w}{ }}}\sphinxbfcode{\sphinxupquote{value}}}
\pysigstopsignatures
\sphinxAtStartPar
Calculates and returns the angle value at the vertex formed by points A and B.
\begin{quote}\begin{description}
\sphinxlineitem{Returns}
\sphinxAtStartPar
The angle measure in degrees (or symbolic expression if symbolic).

\sphinxlineitem{Return type}
\sphinxAtStartPar
float or sympy.Expr

\end{description}\end{quote}

\end{fulllineitems}


\end{fulllineitems}

\index{BaseArc (class in pygeox.basic\_objects)@\spxentry{BaseArc}\spxextra{class in pygeox.basic\_objects}}

\begin{fulllineitems}
\phantomsection\label{\detokenize{pygeox:pygeox.basic_objects.BaseArc}}
\pysigstartsignatures
\pysiglinewithargsret
{\sphinxbfcode{\sphinxupquote{\DUrole{k}{class}\DUrole{w}{ }}}\sphinxcode{\sphinxupquote{pygeox.basic\_objects.}}\sphinxbfcode{\sphinxupquote{BaseArc}}}
{\sphinxparam{\DUrole{n}{center}}\sphinxparamcomma \sphinxparam{\DUrole{n}{start\_point}}\sphinxparamcomma \sphinxparam{\DUrole{n}{end\_point}}\sphinxparamcomma \sphinxparam{\DUrole{n}{name}\DUrole{o}{=}\DUrole{default_value}{None}}\sphinxparamcomma \sphinxparam{\DUrole{n}{arc\_type}\DUrole{o}{=}\DUrole{default_value}{\textquotesingle{}Arc\textquotesingle{}}}}
{}
\pysigstopsignatures
\sphinxAtStartPar
Bases: \sphinxcode{\sphinxupquote{object}}

\sphinxAtStartPar
Base class for arcs, defined by a center and two endpoints.
Not intended to be used directly.
\begin{quote}\begin{description}
\sphinxlineitem{Parameters}\begin{itemize}
\item {} 
\sphinxAtStartPar
\sphinxstyleliteralstrong{\sphinxupquote{center}} (\sphinxstyleliteralemphasis{\sphinxupquote{Point}})

\item {} 
\sphinxAtStartPar
\sphinxstyleliteralstrong{\sphinxupquote{start\_point}} (\sphinxstyleliteralemphasis{\sphinxupquote{Point}})

\item {} 
\sphinxAtStartPar
\sphinxstyleliteralstrong{\sphinxupquote{end\_point}} (\sphinxstyleliteralemphasis{\sphinxupquote{Point}})

\item {} 
\sphinxAtStartPar
\sphinxstyleliteralstrong{\sphinxupquote{name}} (\sphinxstyleliteralemphasis{\sphinxupquote{str}}\sphinxstyleliteralemphasis{\sphinxupquote{ | }}\sphinxstyleliteralemphasis{\sphinxupquote{None}})

\end{itemize}

\end{description}\end{quote}
\index{center (pygeox.basic\_objects.BaseArc attribute)@\spxentry{center}\spxextra{pygeox.basic\_objects.BaseArc attribute}}

\begin{fulllineitems}
\phantomsection\label{\detokenize{pygeox:pygeox.basic_objects.BaseArc.center}}
\pysigstartsignatures
\pysigline
{\sphinxbfcode{\sphinxupquote{center}}}
\pysigstopsignatures
\sphinxAtStartPar
The center of the arc’s circle.
\begin{quote}\begin{description}
\sphinxlineitem{Type}
\sphinxAtStartPar
Point

\end{description}\end{quote}

\end{fulllineitems}

\index{start\_point (pygeox.basic\_objects.BaseArc attribute)@\spxentry{start\_point}\spxextra{pygeox.basic\_objects.BaseArc attribute}}

\begin{fulllineitems}
\phantomsection\label{\detokenize{pygeox:pygeox.basic_objects.BaseArc.start_point}}
\pysigstartsignatures
\pysigline
{\sphinxbfcode{\sphinxupquote{start\_point}}}
\pysigstopsignatures
\sphinxAtStartPar
The start point of the arc.
\begin{quote}\begin{description}
\sphinxlineitem{Type}
\sphinxAtStartPar
Point

\end{description}\end{quote}

\end{fulllineitems}

\index{end\_point (pygeox.basic\_objects.BaseArc attribute)@\spxentry{end\_point}\spxextra{pygeox.basic\_objects.BaseArc attribute}}

\begin{fulllineitems}
\phantomsection\label{\detokenize{pygeox:pygeox.basic_objects.BaseArc.end_point}}
\pysigstartsignatures
\pysigline
{\sphinxbfcode{\sphinxupquote{end\_point}}}
\pysigstopsignatures
\sphinxAtStartPar
The end point of the arc.
\begin{quote}\begin{description}
\sphinxlineitem{Type}
\sphinxAtStartPar
Point

\end{description}\end{quote}

\end{fulllineitems}

\index{name (pygeox.basic\_objects.BaseArc attribute)@\spxentry{name}\spxextra{pygeox.basic\_objects.BaseArc attribute}}

\begin{fulllineitems}
\phantomsection\label{\detokenize{pygeox:pygeox.basic_objects.BaseArc.name}}
\pysigstartsignatures
\pysigline
{\sphinxbfcode{\sphinxupquote{name}}}
\pysigstopsignatures
\sphinxAtStartPar
The name of the arc.
\begin{quote}\begin{description}
\sphinxlineitem{Type}
\sphinxAtStartPar
str

\end{description}\end{quote}

\end{fulllineitems}

\index{arc\_type (pygeox.basic\_objects.BaseArc attribute)@\spxentry{arc\_type}\spxextra{pygeox.basic\_objects.BaseArc attribute}}

\begin{fulllineitems}
\phantomsection\label{\detokenize{pygeox:pygeox.basic_objects.BaseArc.arc_type}}
\pysigstartsignatures
\pysigline
{\sphinxbfcode{\sphinxupquote{arc\_type}}}
\pysigstopsignatures
\sphinxAtStartPar
The type of arc (“Arc”, “MinorArc”, “MajorArc”).
\begin{quote}\begin{description}
\sphinxlineitem{Type}
\sphinxAtStartPar
str

\end{description}\end{quote}

\end{fulllineitems}

\index{start\_radius\_line (pygeox.basic\_objects.BaseArc attribute)@\spxentry{start\_radius\_line}\spxextra{pygeox.basic\_objects.BaseArc attribute}}

\begin{fulllineitems}
\phantomsection\label{\detokenize{pygeox:pygeox.basic_objects.BaseArc.start_radius_line}}
\pysigstartsignatures
\pysigline
{\sphinxbfcode{\sphinxupquote{start\_radius\_line}}}
\pysigstopsignatures
\sphinxAtStartPar
Line from center to start\_point.
\begin{quote}\begin{description}
\sphinxlineitem{Type}
\sphinxAtStartPar
Line

\end{description}\end{quote}

\end{fulllineitems}

\index{end\_radius\_line (pygeox.basic\_objects.BaseArc attribute)@\spxentry{end\_radius\_line}\spxextra{pygeox.basic\_objects.BaseArc attribute}}

\begin{fulllineitems}
\phantomsection\label{\detokenize{pygeox:pygeox.basic_objects.BaseArc.end_radius_line}}
\pysigstartsignatures
\pysigline
{\sphinxbfcode{\sphinxupquote{end\_radius\_line}}}
\pysigstopsignatures
\sphinxAtStartPar
Line from center to end\_point.
\begin{quote}\begin{description}
\sphinxlineitem{Type}
\sphinxAtStartPar
Line

\end{description}\end{quote}

\end{fulllineitems}

\index{radius (pygeox.basic\_objects.BaseArc property)@\spxentry{radius}\spxextra{pygeox.basic\_objects.BaseArc property}}

\begin{fulllineitems}
\phantomsection\label{\detokenize{pygeox:pygeox.basic_objects.BaseArc.radius}}
\pysigstartsignatures
\pysigline
{\sphinxbfcode{\sphinxupquote{\DUrole{k}{property}\DUrole{w}{ }}}\sphinxbfcode{\sphinxupquote{radius}}}
\pysigstopsignatures
\sphinxAtStartPar
Returns the radius of the arc (distance from center to start\_point).
\begin{quote}\begin{description}
\sphinxlineitem{Returns}
\sphinxAtStartPar
The radius.

\sphinxlineitem{Return type}
\sphinxAtStartPar
sympy.Expr

\end{description}\end{quote}

\end{fulllineitems}


\end{fulllineitems}

\index{Circle (class in pygeox.basic\_objects)@\spxentry{Circle}\spxextra{class in pygeox.basic\_objects}}

\begin{fulllineitems}
\phantomsection\label{\detokenize{pygeox:pygeox.basic_objects.Circle}}
\pysigstartsignatures
\pysiglinewithargsret
{\sphinxbfcode{\sphinxupquote{\DUrole{k}{class}\DUrole{w}{ }}}\sphinxcode{\sphinxupquote{pygeox.basic\_objects.}}\sphinxbfcode{\sphinxupquote{Circle}}}
{\sphinxparam{\DUrole{n}{center}}\sphinxparamcomma \sphinxparam{\DUrole{n}{radius}}\sphinxparamcomma \sphinxparam{\DUrole{n}{name}\DUrole{o}{=}\DUrole{default_value}{None}}}
{}
\pysigstopsignatures
\sphinxAtStartPar
Bases: \sphinxcode{\sphinxupquote{object}}

\sphinxAtStartPar
Represents a circle defined by its center Point and radius (symbolic or numeric).
\begin{quote}\begin{description}
\sphinxlineitem{Parameters}\begin{itemize}
\item {} 
\sphinxAtStartPar
\sphinxstyleliteralstrong{\sphinxupquote{center}} (\sphinxstyleliteralemphasis{\sphinxupquote{Point}})

\item {} 
\sphinxAtStartPar
\sphinxstyleliteralstrong{\sphinxupquote{radius}} (\sphinxstyleliteralemphasis{\sphinxupquote{int}}\sphinxstyleliteralemphasis{\sphinxupquote{ | }}\sphinxstyleliteralemphasis{\sphinxupquote{float}}\sphinxstyleliteralemphasis{\sphinxupquote{ | }}\sphinxstyleliteralemphasis{\sphinxupquote{Expr}}\sphinxstyleliteralemphasis{\sphinxupquote{ | }}\sphinxstyleliteralemphasis{\sphinxupquote{None}})

\item {} 
\sphinxAtStartPar
\sphinxstyleliteralstrong{\sphinxupquote{name}} (\sphinxstyleliteralemphasis{\sphinxupquote{str}}\sphinxstyleliteralemphasis{\sphinxupquote{ | }}\sphinxstyleliteralemphasis{\sphinxupquote{None}})

\end{itemize}

\end{description}\end{quote}
\index{center (pygeox.basic\_objects.Circle attribute)@\spxentry{center}\spxextra{pygeox.basic\_objects.Circle attribute}}

\begin{fulllineitems}
\phantomsection\label{\detokenize{pygeox:pygeox.basic_objects.Circle.center}}
\pysigstartsignatures
\pysigline
{\sphinxbfcode{\sphinxupquote{center}}}
\pysigstopsignatures
\sphinxAtStartPar
The center of the circle.
\begin{quote}\begin{description}
\sphinxlineitem{Type}
\sphinxAtStartPar
Point

\end{description}\end{quote}

\end{fulllineitems}

\index{radius (pygeox.basic\_objects.Circle attribute)@\spxentry{radius}\spxextra{pygeox.basic\_objects.Circle attribute}}

\begin{fulllineitems}
\phantomsection\label{\detokenize{pygeox:pygeox.basic_objects.Circle.radius}}
\pysigstartsignatures
\pysigline
{\sphinxbfcode{\sphinxupquote{radius}}}
\pysigstopsignatures
\sphinxAtStartPar
The radius (can be symbolic or numeric).
\begin{quote}\begin{description}
\sphinxlineitem{Type}
\sphinxAtStartPar
Coordinate

\end{description}\end{quote}

\end{fulllineitems}

\index{name (pygeox.basic\_objects.Circle attribute)@\spxentry{name}\spxextra{pygeox.basic\_objects.Circle attribute}}

\begin{fulllineitems}
\phantomsection\label{\detokenize{pygeox:pygeox.basic_objects.Circle.name}}
\pysigstartsignatures
\pysigline
{\sphinxbfcode{\sphinxupquote{name}}}
\pysigstopsignatures
\sphinxAtStartPar
The name of the circle.
\begin{quote}\begin{description}
\sphinxlineitem{Type}
\sphinxAtStartPar
str

\end{description}\end{quote}

\end{fulllineitems}

\index{area (pygeox.basic\_objects.Circle property)@\spxentry{area}\spxextra{pygeox.basic\_objects.Circle property}}

\begin{fulllineitems}
\phantomsection\label{\detokenize{pygeox:pygeox.basic_objects.Circle.area}}
\pysigstartsignatures
\pysigline
{\sphinxbfcode{\sphinxupquote{\DUrole{k}{property}\DUrole{w}{ }}}\sphinxbfcode{\sphinxupquote{area}}}
\pysigstopsignatures
\sphinxAtStartPar
Returns the area of the circle.
\begin{quote}\begin{description}
\sphinxlineitem{Returns}
\sphinxAtStartPar
The area.

\sphinxlineitem{Return type}
\sphinxAtStartPar
sympy.Expr

\end{description}\end{quote}

\end{fulllineitems}

\index{diameter (pygeox.basic\_objects.Circle property)@\spxentry{diameter}\spxextra{pygeox.basic\_objects.Circle property}}

\begin{fulllineitems}
\phantomsection\label{\detokenize{pygeox:pygeox.basic_objects.Circle.diameter}}
\pysigstartsignatures
\pysigline
{\sphinxbfcode{\sphinxupquote{\DUrole{k}{property}\DUrole{w}{ }}}\sphinxbfcode{\sphinxupquote{diameter}}}
\pysigstopsignatures
\sphinxAtStartPar
Returns the diameter (length) of the circle.
\begin{quote}\begin{description}
\sphinxlineitem{Returns}
\sphinxAtStartPar
The diameter (2 * radius).

\sphinxlineitem{Return type}
\sphinxAtStartPar
sympy.Expr

\end{description}\end{quote}

\end{fulllineitems}

\index{perimeter (pygeox.basic\_objects.Circle property)@\spxentry{perimeter}\spxextra{pygeox.basic\_objects.Circle property}}

\begin{fulllineitems}
\phantomsection\label{\detokenize{pygeox:pygeox.basic_objects.Circle.perimeter}}
\pysigstartsignatures
\pysigline
{\sphinxbfcode{\sphinxupquote{\DUrole{k}{property}\DUrole{w}{ }}}\sphinxbfcode{\sphinxupquote{perimeter}}}
\pysigstopsignatures
\sphinxAtStartPar
Returns the circumference of the circle.
\begin{quote}\begin{description}
\sphinxlineitem{Returns}
\sphinxAtStartPar
The circumference.

\sphinxlineitem{Return type}
\sphinxAtStartPar
sympy.Expr

\end{description}\end{quote}

\end{fulllineitems}


\end{fulllineitems}

\index{Line (class in pygeox.basic\_objects)@\spxentry{Line}\spxextra{class in pygeox.basic\_objects}}

\begin{fulllineitems}
\phantomsection\label{\detokenize{pygeox:pygeox.basic_objects.Line}}
\pysigstartsignatures
\pysiglinewithargsret
{\sphinxbfcode{\sphinxupquote{\DUrole{k}{class}\DUrole{w}{ }}}\sphinxcode{\sphinxupquote{pygeox.basic\_objects.}}\sphinxbfcode{\sphinxupquote{Line}}}
{\sphinxparam{\DUrole{n}{point1}}\sphinxparamcomma \sphinxparam{\DUrole{n}{point2}}\sphinxparamcomma \sphinxparam{\DUrole{n}{name}\DUrole{o}{=}\DUrole{default_value}{None}}}
{}
\pysigstopsignatures
\sphinxAtStartPar
Bases: \sphinxcode{\sphinxupquote{LineLike}}

\sphinxAtStartPar
Represents an infinite line defined by two Point objects.
\begin{quote}\begin{description}
\sphinxlineitem{Parameters}\begin{itemize}
\item {} 
\sphinxAtStartPar
\sphinxstyleliteralstrong{\sphinxupquote{point1}} (\sphinxstyleliteralemphasis{\sphinxupquote{Point}})

\item {} 
\sphinxAtStartPar
\sphinxstyleliteralstrong{\sphinxupquote{point2}} (\sphinxstyleliteralemphasis{\sphinxupquote{Point}})

\item {} 
\sphinxAtStartPar
\sphinxstyleliteralstrong{\sphinxupquote{name}} (\sphinxstyleliteralemphasis{\sphinxupquote{str}}\sphinxstyleliteralemphasis{\sphinxupquote{ | }}\sphinxstyleliteralemphasis{\sphinxupquote{None}})

\end{itemize}

\end{description}\end{quote}
\index{equation (pygeox.basic\_objects.Line property)@\spxentry{equation}\spxextra{pygeox.basic\_objects.Line property}}

\begin{fulllineitems}
\phantomsection\label{\detokenize{pygeox:pygeox.basic_objects.Line.equation}}
\pysigstartsignatures
\pysigline
{\sphinxbfcode{\sphinxupquote{\DUrole{k}{property}\DUrole{w}{ }}}\sphinxbfcode{\sphinxupquote{equation}}\sphinxbfcode{\sphinxupquote{\DUrole{p}{:}\DUrole{w}{ }Expr}}}
\pysigstopsignatures
\sphinxAtStartPar
Returns the equation of the line in the form Ax + By + C = 0.
The result is a SymPy expression that evaluates to zero for points on the line.
\begin{quote}\begin{description}
\sphinxlineitem{Returns}
\sphinxAtStartPar
The symbolic equation of the line.

\sphinxlineitem{Return type}
\sphinxAtStartPar
sympy.Expr

\end{description}\end{quote}

\end{fulllineitems}


\end{fulllineitems}

\index{LineLike (class in pygeox.basic\_objects)@\spxentry{LineLike}\spxextra{class in pygeox.basic\_objects}}

\begin{fulllineitems}
\phantomsection\label{\detokenize{pygeox:pygeox.basic_objects.LineLike}}
\pysigstartsignatures
\pysiglinewithargsret
{\sphinxbfcode{\sphinxupquote{\DUrole{k}{class}\DUrole{w}{ }}}\sphinxcode{\sphinxupquote{pygeox.basic\_objects.}}\sphinxbfcode{\sphinxupquote{LineLike}}}
{\sphinxparam{\DUrole{n}{point1}}\sphinxparamcomma \sphinxparam{\DUrole{n}{point2}}\sphinxparamcomma \sphinxparam{\DUrole{n}{name}\DUrole{o}{=}\DUrole{default_value}{None}}}
{}
\pysigstopsignatures
\sphinxAtStartPar
Bases: \sphinxcode{\sphinxupquote{object}}

\sphinxAtStartPar
Base class for geometric entities that are line\sphinxhyphen{}like (LineSegment, Line, Ray).
Defined by two Point objects.
\begin{quote}\begin{description}
\sphinxlineitem{Parameters}\begin{itemize}
\item {} 
\sphinxAtStartPar
\sphinxstyleliteralstrong{\sphinxupquote{point1}} (\sphinxstyleliteralemphasis{\sphinxupquote{Point}})

\item {} 
\sphinxAtStartPar
\sphinxstyleliteralstrong{\sphinxupquote{point2}} (\sphinxstyleliteralemphasis{\sphinxupquote{Point}})

\item {} 
\sphinxAtStartPar
\sphinxstyleliteralstrong{\sphinxupquote{name}} (\sphinxstyleliteralemphasis{\sphinxupquote{str}}\sphinxstyleliteralemphasis{\sphinxupquote{ | }}\sphinxstyleliteralemphasis{\sphinxupquote{None}})

\end{itemize}

\end{description}\end{quote}
\index{direction\_vector (pygeox.basic\_objects.LineLike property)@\spxentry{direction\_vector}\spxextra{pygeox.basic\_objects.LineLike property}}

\begin{fulllineitems}
\phantomsection\label{\detokenize{pygeox:pygeox.basic_objects.LineLike.direction_vector}}
\pysigstartsignatures
\pysigline
{\sphinxbfcode{\sphinxupquote{\DUrole{k}{property}\DUrole{w}{ }}}\sphinxbfcode{\sphinxupquote{direction\_vector}}\sphinxbfcode{\sphinxupquote{\DUrole{p}{:}\DUrole{w}{ }tuple\DUrole{p}{{[}}int\DUrole{w}{ }\DUrole{p}{|}\DUrole{w}{ }float\DUrole{w}{ }\DUrole{p}{|}\DUrole{w}{ }Expr\DUrole{p}{,}\DUrole{w}{ }int\DUrole{w}{ }\DUrole{p}{|}\DUrole{w}{ }float\DUrole{w}{ }\DUrole{p}{|}\DUrole{w}{ }Expr\DUrole{p}{{]}}}}}
\pysigstopsignatures
\sphinxAtStartPar
Returns the direction vector of the ray, from start\_point to direction\_point.
\begin{quote}\begin{description}
\sphinxlineitem{Returns}
\sphinxAtStartPar
The (dx, dy) vector components.

\sphinxlineitem{Return type}
\sphinxAtStartPar
tuple{[}sympy.Expr, sympy.Expr{]}

\end{description}\end{quote}

\end{fulllineitems}

\index{intercept (pygeox.basic\_objects.LineLike property)@\spxentry{intercept}\spxextra{pygeox.basic\_objects.LineLike property}}

\begin{fulllineitems}
\phantomsection\label{\detokenize{pygeox:pygeox.basic_objects.LineLike.intercept}}
\pysigstartsignatures
\pysigline
{\sphinxbfcode{\sphinxupquote{\DUrole{k}{property}\DUrole{w}{ }}}\sphinxbfcode{\sphinxupquote{intercept}}\sphinxbfcode{\sphinxupquote{\DUrole{p}{:}\DUrole{w}{ }int\DUrole{w}{ }\DUrole{p}{|}\DUrole{w}{ }float\DUrole{w}{ }\DUrole{p}{|}\DUrole{w}{ }Expr}}}
\pysigstopsignatures
\sphinxAtStartPar
Returns the y\sphinxhyphen{}intercept (c) of the infinite line defined by point1 and point2
(in the form y = mx + c).
\begin{quote}\begin{description}
\sphinxlineitem{Returns}
\sphinxAtStartPar
The y\sphinxhyphen{}intercept, or sympy.nan for vertical lines.

\sphinxlineitem{Return type}
\sphinxAtStartPar
sympy.Expr

\end{description}\end{quote}

\end{fulllineitems}

\index{slope (pygeox.basic\_objects.LineLike property)@\spxentry{slope}\spxextra{pygeox.basic\_objects.LineLike property}}

\begin{fulllineitems}
\phantomsection\label{\detokenize{pygeox:pygeox.basic_objects.LineLike.slope}}
\pysigstartsignatures
\pysigline
{\sphinxbfcode{\sphinxupquote{\DUrole{k}{property}\DUrole{w}{ }}}\sphinxbfcode{\sphinxupquote{slope}}\sphinxbfcode{\sphinxupquote{\DUrole{p}{:}\DUrole{w}{ }int\DUrole{w}{ }\DUrole{p}{|}\DUrole{w}{ }float\DUrole{w}{ }\DUrole{p}{|}\DUrole{w}{ }Expr}}}
\pysigstopsignatures
\sphinxAtStartPar
Returns the slope of the infinite line passing through point1 and point2.
\begin{quote}\begin{description}
\sphinxlineitem{Returns}
\sphinxAtStartPar
The slope, or sympy.oo (symbolic infinity) for vertical lines.

\sphinxlineitem{Return type}
\sphinxAtStartPar
sympy.Expr

\end{description}\end{quote}

\end{fulllineitems}


\end{fulllineitems}

\index{LineSegment (class in pygeox.basic\_objects)@\spxentry{LineSegment}\spxextra{class in pygeox.basic\_objects}}

\begin{fulllineitems}
\phantomsection\label{\detokenize{pygeox:pygeox.basic_objects.LineSegment}}
\pysigstartsignatures
\pysiglinewithargsret
{\sphinxbfcode{\sphinxupquote{\DUrole{k}{class}\DUrole{w}{ }}}\sphinxcode{\sphinxupquote{pygeox.basic\_objects.}}\sphinxbfcode{\sphinxupquote{LineSegment}}}
{\sphinxparam{\DUrole{n}{point1}}\sphinxparamcomma \sphinxparam{\DUrole{n}{point2}}\sphinxparamcomma \sphinxparam{\DUrole{n}{name}\DUrole{o}{=}\DUrole{default_value}{None}}}
{}
\pysigstopsignatures
\sphinxAtStartPar
Bases: \sphinxcode{\sphinxupquote{LineLike}}

\sphinxAtStartPar
Represents a line segment defined by two Point objects.
This class contains methods specific to a finite segment.
\begin{quote}\begin{description}
\sphinxlineitem{Parameters}\begin{itemize}
\item {} 
\sphinxAtStartPar
\sphinxstyleliteralstrong{\sphinxupquote{point1}} (\sphinxstyleliteralemphasis{\sphinxupquote{Point}})

\item {} 
\sphinxAtStartPar
\sphinxstyleliteralstrong{\sphinxupquote{point2}} (\sphinxstyleliteralemphasis{\sphinxupquote{Point}})

\item {} 
\sphinxAtStartPar
\sphinxstyleliteralstrong{\sphinxupquote{name}} (\sphinxstyleliteralemphasis{\sphinxupquote{str}}\sphinxstyleliteralemphasis{\sphinxupquote{ | }}\sphinxstyleliteralemphasis{\sphinxupquote{None}})

\end{itemize}

\end{description}\end{quote}
\index{angle() (pygeox.basic\_objects.LineSegment method)@\spxentry{angle()}\spxextra{pygeox.basic\_objects.LineSegment method}}

\begin{fulllineitems}
\phantomsection\label{\detokenize{pygeox:pygeox.basic_objects.LineSegment.angle}}
\pysigstartsignatures
\pysiglinewithargsret
{\sphinxbfcode{\sphinxupquote{angle}}}
{\sphinxparam{\DUrole{n}{other\_line\_segment}}\sphinxparamcomma \sphinxparam{\DUrole{n}{signed}\DUrole{o}{=}\DUrole{default_value}{False}}}
{}
\pysigstopsignatures
\sphinxAtStartPar
Returns the angle between this line segment and another line segment
at their common intersection point (vertex).
\begin{quote}\begin{description}
\sphinxlineitem{Parameters}\begin{itemize}
\item {} 
\sphinxAtStartPar
\sphinxstyleliteralstrong{\sphinxupquote{other\_line\_segment}} (\sphinxstyleliteralemphasis{\sphinxupquote{LineSegment}}) \textendash{} The other line segment.

\item {} 
\sphinxAtStartPar
\sphinxstyleliteralstrong{\sphinxupquote{signed}} (\sphinxstyleliteralemphasis{\sphinxupquote{bool}}) \textendash{} If True, returns the signed angle.

\end{itemize}

\sphinxlineitem{Returns}
\sphinxAtStartPar
The angle in degrees.

\sphinxlineitem{Return type}
\sphinxAtStartPar
sympy.Expr

\sphinxlineitem{Raises}
\sphinxAtStartPar
\sphinxstyleliteralstrong{\sphinxupquote{ValueError}} \textendash{} If line segments do not meet at a single common vertex.

\end{description}\end{quote}

\end{fulllineitems}

\index{length (pygeox.basic\_objects.LineSegment property)@\spxentry{length}\spxextra{pygeox.basic\_objects.LineSegment property}}

\begin{fulllineitems}
\phantomsection\label{\detokenize{pygeox:pygeox.basic_objects.LineSegment.length}}
\pysigstartsignatures
\pysigline
{\sphinxbfcode{\sphinxupquote{\DUrole{k}{property}\DUrole{w}{ }}}\sphinxbfcode{\sphinxupquote{length}}\sphinxbfcode{\sphinxupquote{\DUrole{p}{:}\DUrole{w}{ }Expr}}}
\pysigstopsignatures
\sphinxAtStartPar
Returns the length of the line segment.
\begin{quote}\begin{description}
\sphinxlineitem{Returns}
\sphinxAtStartPar
The distance between point1 and point2.

\sphinxlineitem{Return type}
\sphinxAtStartPar
sympy.Expr

\end{description}\end{quote}

\end{fulllineitems}

\index{midpoint (pygeox.basic\_objects.LineSegment property)@\spxentry{midpoint}\spxextra{pygeox.basic\_objects.LineSegment property}}

\begin{fulllineitems}
\phantomsection\label{\detokenize{pygeox:pygeox.basic_objects.LineSegment.midpoint}}
\pysigstartsignatures
\pysigline
{\sphinxbfcode{\sphinxupquote{\DUrole{k}{property}\DUrole{w}{ }}}\sphinxbfcode{\sphinxupquote{midpoint}}\sphinxbfcode{\sphinxupquote{\DUrole{p}{:}\DUrole{w}{ }Point}}}
\pysigstopsignatures
\sphinxAtStartPar
Returns the midpoint of the line segment.
\begin{quote}\begin{description}
\sphinxlineitem{Returns}
\sphinxAtStartPar
The midpoint as a new Point object.

\sphinxlineitem{Return type}
\sphinxAtStartPar
Point

\end{description}\end{quote}

\end{fulllineitems}


\end{fulllineitems}

\index{MajorArc (class in pygeox.basic\_objects)@\spxentry{MajorArc}\spxextra{class in pygeox.basic\_objects}}

\begin{fulllineitems}
\phantomsection\label{\detokenize{pygeox:pygeox.basic_objects.MajorArc}}
\pysigstartsignatures
\pysiglinewithargsret
{\sphinxbfcode{\sphinxupquote{\DUrole{k}{class}\DUrole{w}{ }}}\sphinxcode{\sphinxupquote{pygeox.basic\_objects.}}\sphinxbfcode{\sphinxupquote{MajorArc}}}
{\sphinxparam{\DUrole{n}{center}}\sphinxparamcomma \sphinxparam{\DUrole{n}{start\_point}}\sphinxparamcomma \sphinxparam{\DUrole{n}{end\_point}}\sphinxparamcomma \sphinxparam{\DUrole{n}{name}\DUrole{o}{=}\DUrole{default_value}{None}}}
{}
\pysigstopsignatures
\sphinxAtStartPar
Bases: \sphinxcode{\sphinxupquote{BaseArc}}

\sphinxAtStartPar
Represents the major arc (larger arc) between two points on a circle.
\begin{quote}\begin{description}
\sphinxlineitem{Parameters}\begin{itemize}
\item {} 
\sphinxAtStartPar
\sphinxstyleliteralstrong{\sphinxupquote{center}} (\sphinxstyleliteralemphasis{\sphinxupquote{Point}})

\item {} 
\sphinxAtStartPar
\sphinxstyleliteralstrong{\sphinxupquote{start\_point}} (\sphinxstyleliteralemphasis{\sphinxupquote{Point}})

\item {} 
\sphinxAtStartPar
\sphinxstyleliteralstrong{\sphinxupquote{end\_point}} (\sphinxstyleliteralemphasis{\sphinxupquote{Point}})

\item {} 
\sphinxAtStartPar
\sphinxstyleliteralstrong{\sphinxupquote{name}} (\sphinxstyleliteralemphasis{\sphinxupquote{str}}\sphinxstyleliteralemphasis{\sphinxupquote{ | }}\sphinxstyleliteralemphasis{\sphinxupquote{None}})

\end{itemize}

\end{description}\end{quote}
\index{area (pygeox.basic\_objects.MajorArc property)@\spxentry{area}\spxextra{pygeox.basic\_objects.MajorArc property}}

\begin{fulllineitems}
\phantomsection\label{\detokenize{pygeox:pygeox.basic_objects.MajorArc.area}}
\pysigstartsignatures
\pysigline
{\sphinxbfcode{\sphinxupquote{\DUrole{k}{property}\DUrole{w}{ }}}\sphinxbfcode{\sphinxupquote{area}}}
\pysigstopsignatures
\sphinxAtStartPar
Returns the area of the sector defined by the major arc.
\begin{quote}\begin{description}
\sphinxlineitem{Returns}
\sphinxAtStartPar
The sector area.

\sphinxlineitem{Return type}
\sphinxAtStartPar
sympy.Expr

\end{description}\end{quote}

\end{fulllineitems}

\index{central\_angle (pygeox.basic\_objects.MajorArc property)@\spxentry{central\_angle}\spxextra{pygeox.basic\_objects.MajorArc property}}

\begin{fulllineitems}
\phantomsection\label{\detokenize{pygeox:pygeox.basic_objects.MajorArc.central_angle}}
\pysigstartsignatures
\pysigline
{\sphinxbfcode{\sphinxupquote{\DUrole{k}{property}\DUrole{w}{ }}}\sphinxbfcode{\sphinxupquote{central\_angle}}}
\pysigstopsignatures
\sphinxAtStartPar
Returns the central angle of the major arc in degrees (180 to 360).
\begin{quote}\begin{description}
\sphinxlineitem{Returns}
\sphinxAtStartPar
The angle in degrees.

\sphinxlineitem{Return type}
\sphinxAtStartPar
sympy.Expr

\end{description}\end{quote}

\end{fulllineitems}

\index{inscribed\_angle (pygeox.basic\_objects.MajorArc property)@\spxentry{inscribed\_angle}\spxextra{pygeox.basic\_objects.MajorArc property}}

\begin{fulllineitems}
\phantomsection\label{\detokenize{pygeox:pygeox.basic_objects.MajorArc.inscribed_angle}}
\pysigstartsignatures
\pysigline
{\sphinxbfcode{\sphinxupquote{\DUrole{k}{property}\DUrole{w}{ }}}\sphinxbfcode{\sphinxupquote{inscribed\_angle}}}
\pysigstopsignatures
\sphinxAtStartPar
Returns the inscribed angle subtended by the arc in degrees.
\begin{quote}\begin{description}
\sphinxlineitem{Returns}
\sphinxAtStartPar
The inscribed angle in degrees.

\sphinxlineitem{Return type}
\sphinxAtStartPar
sympy.Expr

\end{description}\end{quote}

\end{fulllineitems}

\index{length (pygeox.basic\_objects.MajorArc property)@\spxentry{length}\spxextra{pygeox.basic\_objects.MajorArc property}}

\begin{fulllineitems}
\phantomsection\label{\detokenize{pygeox:pygeox.basic_objects.MajorArc.length}}
\pysigstartsignatures
\pysigline
{\sphinxbfcode{\sphinxupquote{\DUrole{k}{property}\DUrole{w}{ }}}\sphinxbfcode{\sphinxupquote{length}}}
\pysigstopsignatures
\sphinxAtStartPar
Returns a non\sphinxhyphen{}standard ‘length’ of the major arc, calculated as radius * (angle in degrees).

\sphinxAtStartPar
Note: Standard arc length requires the angle in radians (s = r * theta\_radians).
\begin{quote}\begin{description}
\sphinxlineitem{Returns}
\sphinxAtStartPar
The arc length.

\sphinxlineitem{Return type}
\sphinxAtStartPar
sympy.Expr

\end{description}\end{quote}

\end{fulllineitems}

\index{midpoint (pygeox.basic\_objects.MajorArc property)@\spxentry{midpoint}\spxextra{pygeox.basic\_objects.MajorArc property}}

\begin{fulllineitems}
\phantomsection\label{\detokenize{pygeox:pygeox.basic_objects.MajorArc.midpoint}}
\pysigstartsignatures
\pysigline
{\sphinxbfcode{\sphinxupquote{\DUrole{k}{property}\DUrole{w}{ }}}\sphinxbfcode{\sphinxupquote{midpoint}}\sphinxbfcode{\sphinxupquote{\DUrole{p}{:}\DUrole{w}{ }Point}}}
\pysigstopsignatures
\sphinxAtStartPar
Calculates and returns the midpoint of the major arc.
\begin{quote}\begin{description}
\sphinxlineitem{Parameters}\begin{itemize}
\item {} 
\sphinxAtStartPar
\sphinxstyleliteralstrong{\sphinxupquote{arc}} (\sphinxstyleliteralemphasis{\sphinxupquote{MajorArc}}) \textendash{} The MajorArc instance.

\item {} 
\sphinxAtStartPar
\sphinxstyleliteralstrong{\sphinxupquote{name}} (\sphinxstyleliteralemphasis{\sphinxupquote{str}}\sphinxstyleliteralemphasis{\sphinxupquote{, }}\sphinxstyleliteralemphasis{\sphinxupquote{optional}}) \textendash{} Name for the midpoint. If None, auto\sphinxhyphen{}generated.

\end{itemize}

\sphinxlineitem{Returns}
\sphinxAtStartPar
The midpoint as a new Point object.

\sphinxlineitem{Return type}
\sphinxAtStartPar
Point

\end{description}\end{quote}

\end{fulllineitems}

\index{mirror\_across\_line() (pygeox.basic\_objects.MajorArc method)@\spxentry{mirror\_across\_line()}\spxextra{pygeox.basic\_objects.MajorArc method}}

\begin{fulllineitems}
\phantomsection\label{\detokenize{pygeox:pygeox.basic_objects.MajorArc.mirror_across_line}}
\pysigstartsignatures
\pysiglinewithargsret
{\sphinxbfcode{\sphinxupquote{mirror\_across\_line}}}
{\sphinxparam{\DUrole{n}{axis\_line}}\sphinxparamcomma \sphinxparam{\DUrole{n}{name}\DUrole{o}{=}\DUrole{default_value}{None}}}
{}
\pysigstopsignatures\begin{quote}\begin{description}
\sphinxlineitem{Parameters}
\sphinxAtStartPar
\sphinxstyleliteralstrong{\sphinxupquote{axis\_line}} (\sphinxstyleliteralemphasis{\sphinxupquote{Line}})

\sphinxlineitem{Return type}
\sphinxAtStartPar
\sphinxstyleemphasis{MinorArc}

\end{description}\end{quote}

\end{fulllineitems}


\end{fulllineitems}

\index{MinorArc (class in pygeox.basic\_objects)@\spxentry{MinorArc}\spxextra{class in pygeox.basic\_objects}}

\begin{fulllineitems}
\phantomsection\label{\detokenize{pygeox:pygeox.basic_objects.MinorArc}}
\pysigstartsignatures
\pysiglinewithargsret
{\sphinxbfcode{\sphinxupquote{\DUrole{k}{class}\DUrole{w}{ }}}\sphinxcode{\sphinxupquote{pygeox.basic\_objects.}}\sphinxbfcode{\sphinxupquote{MinorArc}}}
{\sphinxparam{\DUrole{n}{center}}\sphinxparamcomma \sphinxparam{\DUrole{n}{start\_point}}\sphinxparamcomma \sphinxparam{\DUrole{n}{end\_point}}\sphinxparamcomma \sphinxparam{\DUrole{n}{name}\DUrole{o}{=}\DUrole{default_value}{None}}}
{}
\pysigstopsignatures
\sphinxAtStartPar
Bases: \sphinxcode{\sphinxupquote{BaseArc}}

\sphinxAtStartPar
Represents the minor arc (smaller arc) between two points on a circle.
\begin{quote}\begin{description}
\sphinxlineitem{Parameters}\begin{itemize}
\item {} 
\sphinxAtStartPar
\sphinxstyleliteralstrong{\sphinxupquote{center}} (\sphinxstyleliteralemphasis{\sphinxupquote{Point}})

\item {} 
\sphinxAtStartPar
\sphinxstyleliteralstrong{\sphinxupquote{start\_point}} (\sphinxstyleliteralemphasis{\sphinxupquote{Point}})

\item {} 
\sphinxAtStartPar
\sphinxstyleliteralstrong{\sphinxupquote{end\_point}} (\sphinxstyleliteralemphasis{\sphinxupquote{Point}})

\item {} 
\sphinxAtStartPar
\sphinxstyleliteralstrong{\sphinxupquote{name}} (\sphinxstyleliteralemphasis{\sphinxupquote{str}}\sphinxstyleliteralemphasis{\sphinxupquote{ | }}\sphinxstyleliteralemphasis{\sphinxupquote{None}})

\end{itemize}

\end{description}\end{quote}
\index{area (pygeox.basic\_objects.MinorArc property)@\spxentry{area}\spxextra{pygeox.basic\_objects.MinorArc property}}

\begin{fulllineitems}
\phantomsection\label{\detokenize{pygeox:pygeox.basic_objects.MinorArc.area}}
\pysigstartsignatures
\pysigline
{\sphinxbfcode{\sphinxupquote{\DUrole{k}{property}\DUrole{w}{ }}}\sphinxbfcode{\sphinxupquote{area}}}
\pysigstopsignatures
\sphinxAtStartPar
Returns the area of the sector defined by the minor arc.
\begin{quote}\begin{description}
\sphinxlineitem{Returns}
\sphinxAtStartPar
The sector area.

\sphinxlineitem{Return type}
\sphinxAtStartPar
sympy.Expr

\end{description}\end{quote}

\end{fulllineitems}

\index{central\_angle (pygeox.basic\_objects.MinorArc property)@\spxentry{central\_angle}\spxextra{pygeox.basic\_objects.MinorArc property}}

\begin{fulllineitems}
\phantomsection\label{\detokenize{pygeox:pygeox.basic_objects.MinorArc.central_angle}}
\pysigstartsignatures
\pysigline
{\sphinxbfcode{\sphinxupquote{\DUrole{k}{property}\DUrole{w}{ }}}\sphinxbfcode{\sphinxupquote{central\_angle}}}
\pysigstopsignatures
\sphinxAtStartPar
Returns the central angle of the minor arc in degrees (0 to 180).
\begin{quote}\begin{description}
\sphinxlineitem{Returns}
\sphinxAtStartPar
The angle in degrees.

\sphinxlineitem{Return type}
\sphinxAtStartPar
sympy.Expr

\end{description}\end{quote}

\end{fulllineitems}

\index{inscribed\_angle (pygeox.basic\_objects.MinorArc property)@\spxentry{inscribed\_angle}\spxextra{pygeox.basic\_objects.MinorArc property}}

\begin{fulllineitems}
\phantomsection\label{\detokenize{pygeox:pygeox.basic_objects.MinorArc.inscribed_angle}}
\pysigstartsignatures
\pysigline
{\sphinxbfcode{\sphinxupquote{\DUrole{k}{property}\DUrole{w}{ }}}\sphinxbfcode{\sphinxupquote{inscribed\_angle}}}
\pysigstopsignatures
\sphinxAtStartPar
Returns the inscribed angle subtended by the arc in degrees.
\begin{quote}\begin{description}
\sphinxlineitem{Returns}
\sphinxAtStartPar
The inscribed angle in degrees.

\sphinxlineitem{Return type}
\sphinxAtStartPar
sympy.Expr

\end{description}\end{quote}

\end{fulllineitems}

\index{length (pygeox.basic\_objects.MinorArc property)@\spxentry{length}\spxextra{pygeox.basic\_objects.MinorArc property}}

\begin{fulllineitems}
\phantomsection\label{\detokenize{pygeox:pygeox.basic_objects.MinorArc.length}}
\pysigstartsignatures
\pysigline
{\sphinxbfcode{\sphinxupquote{\DUrole{k}{property}\DUrole{w}{ }}}\sphinxbfcode{\sphinxupquote{length}}}
\pysigstopsignatures
\sphinxAtStartPar
Returns a non\sphinxhyphen{}standard ‘length’ of the minor arc, calculated as radius * (angle in degrees).

\sphinxAtStartPar
Note: Standard arc length requires the angle in radians (s = r * theta\_radians).
\begin{quote}\begin{description}
\sphinxlineitem{Returns}
\sphinxAtStartPar
The arc length.

\sphinxlineitem{Return type}
\sphinxAtStartPar
sympy.Expr

\end{description}\end{quote}

\end{fulllineitems}


\end{fulllineitems}

\index{Point (class in pygeox.basic\_objects)@\spxentry{Point}\spxextra{class in pygeox.basic\_objects}}

\begin{fulllineitems}
\phantomsection\label{\detokenize{pygeox:pygeox.basic_objects.Point}}
\pysigstartsignatures
\pysiglinewithargsret
{\sphinxbfcode{\sphinxupquote{\DUrole{k}{class}\DUrole{w}{ }}}\sphinxcode{\sphinxupquote{pygeox.basic\_objects.}}\sphinxbfcode{\sphinxupquote{Point}}}
{\sphinxparam{\DUrole{n}{name}}\sphinxparamcomma \sphinxparam{\DUrole{n}{x}}\sphinxparamcomma \sphinxparam{\DUrole{n}{y}}}
{}
\pysigstopsignatures
\sphinxAtStartPar
Bases: \sphinxcode{\sphinxupquote{object}}

\sphinxAtStartPar
Represents a 2D point with symbolic or numeric coordinates.
\begin{quote}\begin{description}
\sphinxlineitem{Parameters}\begin{itemize}
\item {} 
\sphinxAtStartPar
\sphinxstyleliteralstrong{\sphinxupquote{name}} (\sphinxstyleliteralemphasis{\sphinxupquote{str}})

\item {} 
\sphinxAtStartPar
\sphinxstyleliteralstrong{\sphinxupquote{x}} (\sphinxstyleliteralemphasis{\sphinxupquote{int}}\sphinxstyleliteralemphasis{\sphinxupquote{ | }}\sphinxstyleliteralemphasis{\sphinxupquote{float}}\sphinxstyleliteralemphasis{\sphinxupquote{ | }}\sphinxstyleliteralemphasis{\sphinxupquote{Expr}}\sphinxstyleliteralemphasis{\sphinxupquote{ | }}\sphinxstyleliteralemphasis{\sphinxupquote{None}})

\item {} 
\sphinxAtStartPar
\sphinxstyleliteralstrong{\sphinxupquote{y}} (\sphinxstyleliteralemphasis{\sphinxupquote{int}}\sphinxstyleliteralemphasis{\sphinxupquote{ | }}\sphinxstyleliteralemphasis{\sphinxupquote{float}}\sphinxstyleliteralemphasis{\sphinxupquote{ | }}\sphinxstyleliteralemphasis{\sphinxupquote{Expr}}\sphinxstyleliteralemphasis{\sphinxupquote{ | }}\sphinxstyleliteralemphasis{\sphinxupquote{None}})

\end{itemize}

\end{description}\end{quote}
\index{name (pygeox.basic\_objects.Point attribute)@\spxentry{name}\spxextra{pygeox.basic\_objects.Point attribute}}

\begin{fulllineitems}
\phantomsection\label{\detokenize{pygeox:pygeox.basic_objects.Point.name}}
\pysigstartsignatures
\pysigline
{\sphinxbfcode{\sphinxupquote{name}}}
\pysigstopsignatures
\sphinxAtStartPar
Point name.
\begin{quote}\begin{description}
\sphinxlineitem{Type}
\sphinxAtStartPar
str

\end{description}\end{quote}

\end{fulllineitems}

\index{x (pygeox.basic\_objects.Point attribute)@\spxentry{x}\spxextra{pygeox.basic\_objects.Point attribute}}

\begin{fulllineitems}
\phantomsection\label{\detokenize{pygeox:pygeox.basic_objects.Point.x}}
\pysigstartsignatures
\pysigline
{\sphinxbfcode{\sphinxupquote{x}}}
\pysigstopsignatures
\sphinxAtStartPar
X coordinate.
\begin{quote}\begin{description}
\sphinxlineitem{Type}
\sphinxAtStartPar
Coordinate

\end{description}\end{quote}

\end{fulllineitems}

\index{y (pygeox.basic\_objects.Point attribute)@\spxentry{y}\spxextra{pygeox.basic\_objects.Point attribute}}

\begin{fulllineitems}
\phantomsection\label{\detokenize{pygeox:pygeox.basic_objects.Point.y}}
\pysigstartsignatures
\pysigline
{\sphinxbfcode{\sphinxupquote{y}}}
\pysigstopsignatures
\sphinxAtStartPar
Y coordinate.
\begin{quote}\begin{description}
\sphinxlineitem{Type}
\sphinxAtStartPar
Coordinate

\end{description}\end{quote}

\end{fulllineitems}

\index{angle() (pygeox.basic\_objects.Point method)@\spxentry{angle()}\spxextra{pygeox.basic\_objects.Point method}}

\begin{fulllineitems}
\phantomsection\label{\detokenize{pygeox:pygeox.basic_objects.Point.angle}}
\pysigstartsignatures
\pysiglinewithargsret
{\sphinxbfcode{\sphinxupquote{angle}}}
{\sphinxparam{\DUrole{n}{point\_b}}\sphinxparamcomma \sphinxparam{\DUrole{n}{point\_c}}\sphinxparamcomma \sphinxparam{\DUrole{n}{signed}\DUrole{o}{=}\DUrole{default_value}{False}}}
{}
\pysigstopsignatures
\sphinxAtStartPar
Compute angle at self between vectors to point\_b and point\_c in degrees.
\begin{quote}\begin{description}
\sphinxlineitem{Parameters}\begin{itemize}
\item {} 
\sphinxAtStartPar
\sphinxstyleliteralstrong{\sphinxupquote{point\_b}} (\sphinxstyleliteralemphasis{\sphinxupquote{Point}}) \textendash{} First point defining angle arm.

\item {} 
\sphinxAtStartPar
\sphinxstyleliteralstrong{\sphinxupquote{point\_c}} (\sphinxstyleliteralemphasis{\sphinxupquote{Point}}) \textendash{} Second point defining angle arm.

\item {} 
\sphinxAtStartPar
\sphinxstyleliteralstrong{\sphinxupquote{signed}} (\sphinxstyleliteralemphasis{\sphinxupquote{bool}}) \textendash{} If True, return signed angle (CCW positive).

\end{itemize}

\sphinxlineitem{Returns}
\sphinxAtStartPar
Angle in degrees.

\sphinxlineitem{Return type}
\sphinxAtStartPar
sympy.Expr

\end{description}\end{quote}

\end{fulllineitems}

\index{distance() (pygeox.basic\_objects.Point method)@\spxentry{distance()}\spxextra{pygeox.basic\_objects.Point method}}

\begin{fulllineitems}
\phantomsection\label{\detokenize{pygeox:pygeox.basic_objects.Point.distance}}
\pysigstartsignatures
\pysiglinewithargsret
{\sphinxbfcode{\sphinxupquote{distance}}}
{\sphinxparam{\DUrole{n}{other\_point}}}
{}
\pysigstopsignatures
\sphinxAtStartPar
Compute the Euclidean distance to another point.
\begin{quote}\begin{description}
\sphinxlineitem{Parameters}
\sphinxAtStartPar
\sphinxstyleliteralstrong{\sphinxupquote{other\_point}} (\sphinxstyleliteralemphasis{\sphinxupquote{Point}}) \textendash{} The other point.

\sphinxlineitem{Returns}
\sphinxAtStartPar
The distance between self and other\_point.

\sphinxlineitem{Return type}
\sphinxAtStartPar
sympy.Expr

\end{description}\end{quote}

\end{fulllineitems}

\index{distance\_squared() (pygeox.basic\_objects.Point method)@\spxentry{distance\_squared()}\spxextra{pygeox.basic\_objects.Point method}}

\begin{fulllineitems}
\phantomsection\label{\detokenize{pygeox:pygeox.basic_objects.Point.distance_squared}}
\pysigstartsignatures
\pysiglinewithargsret
{\sphinxbfcode{\sphinxupquote{distance\_squared}}}
{\sphinxparam{\DUrole{n}{other\_point}}}
{}
\pysigstopsignatures
\sphinxAtStartPar
Compute the squared Euclidean distance to another point.
\begin{quote}\begin{description}
\sphinxlineitem{Parameters}
\sphinxAtStartPar
\sphinxstyleliteralstrong{\sphinxupquote{other\_point}} (\sphinxstyleliteralemphasis{\sphinxupquote{Point}}) \textendash{} The other point.

\sphinxlineitem{Returns}
\sphinxAtStartPar
The squared distance between self and other\_point.

\sphinxlineitem{Return type}
\sphinxAtStartPar
sympy.Expr

\end{description}\end{quote}

\end{fulllineitems}


\end{fulllineitems}

\index{Ray (class in pygeox.basic\_objects)@\spxentry{Ray}\spxextra{class in pygeox.basic\_objects}}

\begin{fulllineitems}
\phantomsection\label{\detokenize{pygeox:pygeox.basic_objects.Ray}}
\pysigstartsignatures
\pysiglinewithargsret
{\sphinxbfcode{\sphinxupquote{\DUrole{k}{class}\DUrole{w}{ }}}\sphinxcode{\sphinxupquote{pygeox.basic\_objects.}}\sphinxbfcode{\sphinxupquote{Ray}}}
{\sphinxparam{\DUrole{n}{start\_point}}\sphinxparamcomma \sphinxparam{\DUrole{n}{direction\_point}}\sphinxparamcomma \sphinxparam{\DUrole{n}{name}\DUrole{o}{=}\DUrole{default_value}{None}}}
{}
\pysigstopsignatures
\sphinxAtStartPar
Bases: \sphinxcode{\sphinxupquote{LineLike}}

\sphinxAtStartPar
Represents a ray defined by a starting Point and a second Point indicating its direction.
The ray extends infinitely from the start\_point through the direction\_point.
\begin{quote}\begin{description}
\sphinxlineitem{Parameters}\begin{itemize}
\item {} 
\sphinxAtStartPar
\sphinxstyleliteralstrong{\sphinxupquote{start\_point}} (\sphinxstyleliteralemphasis{\sphinxupquote{Point}})

\item {} 
\sphinxAtStartPar
\sphinxstyleliteralstrong{\sphinxupquote{direction\_point}} (\sphinxstyleliteralemphasis{\sphinxupquote{Point}})

\item {} 
\sphinxAtStartPar
\sphinxstyleliteralstrong{\sphinxupquote{name}} (\sphinxstyleliteralemphasis{\sphinxupquote{str}}\sphinxstyleliteralemphasis{\sphinxupquote{ | }}\sphinxstyleliteralemphasis{\sphinxupquote{None}})

\end{itemize}

\end{description}\end{quote}
\index{length (pygeox.basic\_objects.Ray property)@\spxentry{length}\spxextra{pygeox.basic\_objects.Ray property}}

\begin{fulllineitems}
\phantomsection\label{\detokenize{pygeox:pygeox.basic_objects.Ray.length}}
\pysigstartsignatures
\pysigline
{\sphinxbfcode{\sphinxupquote{\DUrole{k}{property}\DUrole{w}{ }}}\sphinxbfcode{\sphinxupquote{length}}\sphinxbfcode{\sphinxupquote{\DUrole{p}{:}\DUrole{w}{ }oo}}}
\pysigstopsignatures
\sphinxAtStartPar
Returns the length of the ray, which is infinite.
\begin{quote}\begin{description}
\sphinxlineitem{Returns}
\sphinxAtStartPar
Symbolic infinity.

\sphinxlineitem{Return type}
\sphinxAtStartPar
sympy.oo

\end{description}\end{quote}

\end{fulllineitems}


\end{fulllineitems}



\chapter{Polygons}
\label{\detokenize{pygeox:module-pygeox.custom_objects}}\label{\detokenize{pygeox:polygons}}\index{module@\spxentry{module}!pygeox.custom\_objects@\spxentry{pygeox.custom\_objects}}\index{pygeox.custom\_objects@\spxentry{pygeox.custom\_objects}!module@\spxentry{module}}
\sphinxAtStartPar
custom\_objects.py: Classes for polygons and specialized geometric shapes within a geometric scene.

\sphinxAtStartPar
Defines general Polygon and its subclasses including CyclicPolygon, RegularPolygon,
and various specific polygons like Triangles, Quadrilaterals, and their special cases
(e.g., EquilateralTriangle, Rectangle, Square).

\sphinxAtStartPar
These classes encapsulate geometric properties, constraints, and relationships,
enabling symbolic geometry modeling, constraint solving, and scene management.
\index{AcuteTriangle (class in pygeox.custom\_objects)@\spxentry{AcuteTriangle}\spxextra{class in pygeox.custom\_objects}}

\begin{fulllineitems}
\phantomsection\label{\detokenize{pygeox:pygeox.custom_objects.AcuteTriangle}}
\pysigstartsignatures
\pysiglinewithargsret
{\sphinxbfcode{\sphinxupquote{\DUrole{k}{class}\DUrole{w}{ }}}\sphinxcode{\sphinxupquote{pygeox.custom\_objects.}}\sphinxbfcode{\sphinxupquote{AcuteTriangle}}}
{\sphinxparam{\DUrole{n}{geoscene}}\sphinxparamcomma \sphinxparam{\DUrole{n}{p1}}\sphinxparamcomma \sphinxparam{\DUrole{n}{p2}}\sphinxparamcomma \sphinxparam{\DUrole{n}{p3}}\sphinxparamcomma \sphinxparam{\DUrole{n}{name}\DUrole{o}{=}\DUrole{default_value}{None}}\sphinxparamcomma \sphinxparam{\DUrole{n}{add\_to\_scene\_as\_type}\DUrole{o}{=}\DUrole{default_value}{\textquotesingle{}AcuteTriangle\textquotesingle{}}}}
{}
\pysigstopsignatures
\sphinxAtStartPar
Bases: \sphinxcode{\sphinxupquote{Triangle}}

\sphinxAtStartPar
Triangle with all angles \textless{} 90 degrees.
\begin{quote}\begin{description}
\sphinxlineitem{Parameters}\begin{itemize}
\item {} 
\sphinxAtStartPar
\sphinxstyleliteralstrong{\sphinxupquote{p1}} (\sphinxstyleliteralemphasis{\sphinxupquote{Point}})

\item {} 
\sphinxAtStartPar
\sphinxstyleliteralstrong{\sphinxupquote{p2}} (\sphinxstyleliteralemphasis{\sphinxupquote{Point}})

\item {} 
\sphinxAtStartPar
\sphinxstyleliteralstrong{\sphinxupquote{p3}} (\sphinxstyleliteralemphasis{\sphinxupquote{Point}})

\item {} 
\sphinxAtStartPar
\sphinxstyleliteralstrong{\sphinxupquote{name}} (\sphinxstyleliteralemphasis{\sphinxupquote{str}}\sphinxstyleliteralemphasis{\sphinxupquote{ | }}\sphinxstyleliteralemphasis{\sphinxupquote{None}})

\item {} 
\sphinxAtStartPar
\sphinxstyleliteralstrong{\sphinxupquote{add\_to\_scene\_as\_type}} (\sphinxstyleliteralemphasis{\sphinxupquote{str}}\sphinxstyleliteralemphasis{\sphinxupquote{ | }}\sphinxstyleliteralemphasis{\sphinxupquote{None}})

\end{itemize}

\end{description}\end{quote}

\end{fulllineitems}

\index{CyclicPolygon (class in pygeox.custom\_objects)@\spxentry{CyclicPolygon}\spxextra{class in pygeox.custom\_objects}}

\begin{fulllineitems}
\phantomsection\label{\detokenize{pygeox:pygeox.custom_objects.CyclicPolygon}}
\pysigstartsignatures
\pysiglinewithargsret
{\sphinxbfcode{\sphinxupquote{\DUrole{k}{class}\DUrole{w}{ }}}\sphinxcode{\sphinxupquote{pygeox.custom\_objects.}}\sphinxbfcode{\sphinxupquote{CyclicPolygon}}}
{\sphinxparam{\DUrole{n}{geoscene}}\sphinxparamcomma \sphinxparam{\DUrole{o}{*}\DUrole{n}{points}}\sphinxparamcomma \sphinxparam{\DUrole{n}{name}\DUrole{o}{=}\DUrole{default_value}{None}}\sphinxparamcomma \sphinxparam{\DUrole{n}{add\_to\_scene\_as\_type}\DUrole{o}{=}\DUrole{default_value}{\textquotesingle{}CyclicPolygon\textquotesingle{}}}}
{}
\pysigstopsignatures
\sphinxAtStartPar
Bases: \sphinxcode{\sphinxupquote{Polygon}}

\sphinxAtStartPar
Polygon whose vertices all lie on a single circle (circumcircle).
Circumcenter and circumradius are always defined.
\begin{quote}\begin{description}
\sphinxlineitem{Parameters}\begin{itemize}
\item {} 
\sphinxAtStartPar
\sphinxstyleliteralstrong{\sphinxupquote{points}} (\sphinxstyleliteralemphasis{\sphinxupquote{Point}})

\item {} 
\sphinxAtStartPar
\sphinxstyleliteralstrong{\sphinxupquote{name}} (\sphinxstyleliteralemphasis{\sphinxupquote{str}}\sphinxstyleliteralemphasis{\sphinxupquote{ | }}\sphinxstyleliteralemphasis{\sphinxupquote{None}})

\item {} 
\sphinxAtStartPar
\sphinxstyleliteralstrong{\sphinxupquote{add\_to\_scene\_as\_type}} (\sphinxstyleliteralemphasis{\sphinxupquote{str}}\sphinxstyleliteralemphasis{\sphinxupquote{ | }}\sphinxstyleliteralemphasis{\sphinxupquote{None}})

\end{itemize}

\end{description}\end{quote}
\index{circumcenter (pygeox.custom\_objects.CyclicPolygon property)@\spxentry{circumcenter}\spxextra{pygeox.custom\_objects.CyclicPolygon property}}

\begin{fulllineitems}
\phantomsection\label{\detokenize{pygeox:pygeox.custom_objects.CyclicPolygon.circumcenter}}
\pysigstartsignatures
\pysigline
{\sphinxbfcode{\sphinxupquote{\DUrole{k}{property}\DUrole{w}{ }}}\sphinxbfcode{\sphinxupquote{circumcenter}}}
\pysigstopsignatures
\end{fulllineitems}

\index{circumradius (pygeox.custom\_objects.CyclicPolygon property)@\spxentry{circumradius}\spxextra{pygeox.custom\_objects.CyclicPolygon property}}

\begin{fulllineitems}
\phantomsection\label{\detokenize{pygeox:pygeox.custom_objects.CyclicPolygon.circumradius}}
\pysigstartsignatures
\pysigline
{\sphinxbfcode{\sphinxupquote{\DUrole{k}{property}\DUrole{w}{ }}}\sphinxbfcode{\sphinxupquote{circumradius}}}
\pysigstopsignatures
\end{fulllineitems}


\end{fulllineitems}

\index{EquilateralTriangle (class in pygeox.custom\_objects)@\spxentry{EquilateralTriangle}\spxextra{class in pygeox.custom\_objects}}

\begin{fulllineitems}
\phantomsection\label{\detokenize{pygeox:pygeox.custom_objects.EquilateralTriangle}}
\pysigstartsignatures
\pysiglinewithargsret
{\sphinxbfcode{\sphinxupquote{\DUrole{k}{class}\DUrole{w}{ }}}\sphinxcode{\sphinxupquote{pygeox.custom\_objects.}}\sphinxbfcode{\sphinxupquote{EquilateralTriangle}}}
{\sphinxparam{\DUrole{n}{geoscene}}\sphinxparamcomma \sphinxparam{\DUrole{n}{p1}}\sphinxparamcomma \sphinxparam{\DUrole{n}{p2}}\sphinxparamcomma \sphinxparam{\DUrole{n}{p3}}\sphinxparamcomma \sphinxparam{\DUrole{n}{name}\DUrole{o}{=}\DUrole{default_value}{None}}\sphinxparamcomma \sphinxparam{\DUrole{n}{add\_to\_scene\_as\_type}\DUrole{o}{=}\DUrole{default_value}{\textquotesingle{}EquilateralTriangle\textquotesingle{}}}}
{}
\pysigstopsignatures
\sphinxAtStartPar
Bases: \sphinxcode{\sphinxupquote{Triangle}}

\sphinxAtStartPar
Triangle with all sides equal and all angles 60 degrees.
\begin{quote}\begin{description}
\sphinxlineitem{Parameters}\begin{itemize}
\item {} 
\sphinxAtStartPar
\sphinxstyleliteralstrong{\sphinxupquote{p1}} (\sphinxstyleliteralemphasis{\sphinxupquote{Point}})

\item {} 
\sphinxAtStartPar
\sphinxstyleliteralstrong{\sphinxupquote{p2}} (\sphinxstyleliteralemphasis{\sphinxupquote{Point}})

\item {} 
\sphinxAtStartPar
\sphinxstyleliteralstrong{\sphinxupquote{p3}} (\sphinxstyleliteralemphasis{\sphinxupquote{Point}})

\item {} 
\sphinxAtStartPar
\sphinxstyleliteralstrong{\sphinxupquote{name}} (\sphinxstyleliteralemphasis{\sphinxupquote{str}}\sphinxstyleliteralemphasis{\sphinxupquote{ | }}\sphinxstyleliteralemphasis{\sphinxupquote{None}})

\item {} 
\sphinxAtStartPar
\sphinxstyleliteralstrong{\sphinxupquote{add\_to\_scene\_as\_type}} (\sphinxstyleliteralemphasis{\sphinxupquote{str}}\sphinxstyleliteralemphasis{\sphinxupquote{ | }}\sphinxstyleliteralemphasis{\sphinxupquote{None}})

\end{itemize}

\end{description}\end{quote}
\index{circumradius (pygeox.custom\_objects.EquilateralTriangle property)@\spxentry{circumradius}\spxextra{pygeox.custom\_objects.EquilateralTriangle property}}

\begin{fulllineitems}
\phantomsection\label{\detokenize{pygeox:pygeox.custom_objects.EquilateralTriangle.circumradius}}
\pysigstartsignatures
\pysigline
{\sphinxbfcode{\sphinxupquote{\DUrole{k}{property}\DUrole{w}{ }}}\sphinxbfcode{\sphinxupquote{circumradius}}}
\pysigstopsignatures
\end{fulllineitems}

\index{inradius (pygeox.custom\_objects.EquilateralTriangle property)@\spxentry{inradius}\spxextra{pygeox.custom\_objects.EquilateralTriangle property}}

\begin{fulllineitems}
\phantomsection\label{\detokenize{pygeox:pygeox.custom_objects.EquilateralTriangle.inradius}}
\pysigstartsignatures
\pysigline
{\sphinxbfcode{\sphinxupquote{\DUrole{k}{property}\DUrole{w}{ }}}\sphinxbfcode{\sphinxupquote{inradius}}}
\pysigstopsignatures
\end{fulllineitems}


\end{fulllineitems}

\index{IsoscelesTrapezoid (class in pygeox.custom\_objects)@\spxentry{IsoscelesTrapezoid}\spxextra{class in pygeox.custom\_objects}}

\begin{fulllineitems}
\phantomsection\label{\detokenize{pygeox:pygeox.custom_objects.IsoscelesTrapezoid}}
\pysigstartsignatures
\pysiglinewithargsret
{\sphinxbfcode{\sphinxupquote{\DUrole{k}{class}\DUrole{w}{ }}}\sphinxcode{\sphinxupquote{pygeox.custom\_objects.}}\sphinxbfcode{\sphinxupquote{IsoscelesTrapezoid}}}
{\sphinxparam{\DUrole{n}{geoscene}}\sphinxparamcomma \sphinxparam{\DUrole{n}{p1}}\sphinxparamcomma \sphinxparam{\DUrole{n}{p2}}\sphinxparamcomma \sphinxparam{\DUrole{n}{p3}}\sphinxparamcomma \sphinxparam{\DUrole{n}{p4}}\sphinxparamcomma \sphinxparam{\DUrole{n}{name}\DUrole{o}{=}\DUrole{default_value}{None}}\sphinxparamcomma \sphinxparam{\DUrole{n}{add\_to\_scene\_as\_type}\DUrole{o}{=}\DUrole{default_value}{\textquotesingle{}IsoscelesTrapezoid\textquotesingle{}}}\sphinxparamcomma \sphinxparam{\DUrole{n}{force\_convex}\DUrole{o}{=}\DUrole{default_value}{False}}}
{}
\pysigstopsignatures
\sphinxAtStartPar
Bases: \sphinxcode{\sphinxupquote{Trapezoid}}

\sphinxAtStartPar
Trapezoid with non\sphinxhyphen{}parallel sides equal in length. IsoscelesTrapezoid(A,B,C,D) means AB is parallel to CD and BC and DA have same lenght.
\begin{quote}\begin{description}
\sphinxlineitem{Parameters}
\sphinxAtStartPar
\sphinxstyleliteralstrong{\sphinxupquote{force\_convex}} (\sphinxstyleliteralemphasis{\sphinxupquote{bool}})

\end{description}\end{quote}

\end{fulllineitems}

\index{IsoscelesTriangle (class in pygeox.custom\_objects)@\spxentry{IsoscelesTriangle}\spxextra{class in pygeox.custom\_objects}}

\begin{fulllineitems}
\phantomsection\label{\detokenize{pygeox:pygeox.custom_objects.IsoscelesTriangle}}
\pysigstartsignatures
\pysiglinewithargsret
{\sphinxbfcode{\sphinxupquote{\DUrole{k}{class}\DUrole{w}{ }}}\sphinxcode{\sphinxupquote{pygeox.custom\_objects.}}\sphinxbfcode{\sphinxupquote{IsoscelesTriangle}}}
{\sphinxparam{\DUrole{n}{geoscene}}\sphinxparamcomma \sphinxparam{\DUrole{n}{p1}}\sphinxparamcomma \sphinxparam{\DUrole{n}{p2}}\sphinxparamcomma \sphinxparam{\DUrole{n}{p3}}\sphinxparamcomma \sphinxparam{\DUrole{n}{name}\DUrole{o}{=}\DUrole{default_value}{None}}\sphinxparamcomma \sphinxparam{\DUrole{n}{add\_to\_scene\_as\_type}\DUrole{o}{=}\DUrole{default_value}{\textquotesingle{}IsoscelesTriangle\textquotesingle{}}}}
{}
\pysigstopsignatures
\sphinxAtStartPar
Bases: \sphinxcode{\sphinxupquote{Triangle}}

\sphinxAtStartPar
Triangle with at least two equal sides. IsoscelesTriangle(A,B,C) means that AB=BC.
\begin{quote}\begin{description}
\sphinxlineitem{Parameters}\begin{itemize}
\item {} 
\sphinxAtStartPar
\sphinxstyleliteralstrong{\sphinxupquote{p1}} (\sphinxstyleliteralemphasis{\sphinxupquote{Point}})

\item {} 
\sphinxAtStartPar
\sphinxstyleliteralstrong{\sphinxupquote{p2}} (\sphinxstyleliteralemphasis{\sphinxupquote{Point}})

\item {} 
\sphinxAtStartPar
\sphinxstyleliteralstrong{\sphinxupquote{p3}} (\sphinxstyleliteralemphasis{\sphinxupquote{Point}})

\item {} 
\sphinxAtStartPar
\sphinxstyleliteralstrong{\sphinxupquote{name}} (\sphinxstyleliteralemphasis{\sphinxupquote{str}}\sphinxstyleliteralemphasis{\sphinxupquote{ | }}\sphinxstyleliteralemphasis{\sphinxupquote{None}})

\item {} 
\sphinxAtStartPar
\sphinxstyleliteralstrong{\sphinxupquote{add\_to\_scene\_as\_type}} (\sphinxstyleliteralemphasis{\sphinxupquote{str}}\sphinxstyleliteralemphasis{\sphinxupquote{ | }}\sphinxstyleliteralemphasis{\sphinxupquote{None}})

\end{itemize}

\end{description}\end{quote}

\end{fulllineitems}

\index{Kite (class in pygeox.custom\_objects)@\spxentry{Kite}\spxextra{class in pygeox.custom\_objects}}

\begin{fulllineitems}
\phantomsection\label{\detokenize{pygeox:pygeox.custom_objects.Kite}}
\pysigstartsignatures
\pysiglinewithargsret
{\sphinxbfcode{\sphinxupquote{\DUrole{k}{class}\DUrole{w}{ }}}\sphinxcode{\sphinxupquote{pygeox.custom\_objects.}}\sphinxbfcode{\sphinxupquote{Kite}}}
{\sphinxparam{\DUrole{n}{geoscene}}\sphinxparamcomma \sphinxparam{\DUrole{n}{p1}}\sphinxparamcomma \sphinxparam{\DUrole{n}{p2}}\sphinxparamcomma \sphinxparam{\DUrole{n}{p3}}\sphinxparamcomma \sphinxparam{\DUrole{n}{p4}}\sphinxparamcomma \sphinxparam{\DUrole{n}{name}\DUrole{o}{=}\DUrole{default_value}{None}}\sphinxparamcomma \sphinxparam{\DUrole{n}{add\_to\_scene\_as\_type}\DUrole{o}{=}\DUrole{default_value}{\textquotesingle{}Kite\textquotesingle{}}}\sphinxparamcomma \sphinxparam{\DUrole{n}{force\_convex}\DUrole{o}{=}\DUrole{default_value}{False}}}
{}
\pysigstopsignatures
\sphinxAtStartPar
Bases: \sphinxcode{\sphinxupquote{Quadrilateral}}

\sphinxAtStartPar
Quadrilateral with two distinct pairs of adjacent sides equal.
\begin{quote}\begin{description}
\sphinxlineitem{Parameters}
\sphinxAtStartPar
\sphinxstyleliteralstrong{\sphinxupquote{force\_convex}} (\sphinxstyleliteralemphasis{\sphinxupquote{bool}})

\end{description}\end{quote}
\index{area (pygeox.custom\_objects.Kite property)@\spxentry{area}\spxextra{pygeox.custom\_objects.Kite property}}

\begin{fulllineitems}
\phantomsection\label{\detokenize{pygeox:pygeox.custom_objects.Kite.area}}
\pysigstartsignatures
\pysigline
{\sphinxbfcode{\sphinxupquote{\DUrole{k}{property}\DUrole{w}{ }}}\sphinxbfcode{\sphinxupquote{area}}}
\pysigstopsignatures
\sphinxAtStartPar
Calculates the area of the polygon using the shoelace formula.
\begin{quote}\begin{description}
\sphinxlineitem{Returns}
\sphinxAtStartPar
The area. Units: square of input coordinates.

\sphinxlineitem{Return type}
\sphinxAtStartPar
sympy.Expr

\end{description}\end{quote}

\end{fulllineitems}


\end{fulllineitems}

\index{ObtuseTriangle (class in pygeox.custom\_objects)@\spxentry{ObtuseTriangle}\spxextra{class in pygeox.custom\_objects}}

\begin{fulllineitems}
\phantomsection\label{\detokenize{pygeox:pygeox.custom_objects.ObtuseTriangle}}
\pysigstartsignatures
\pysiglinewithargsret
{\sphinxbfcode{\sphinxupquote{\DUrole{k}{class}\DUrole{w}{ }}}\sphinxcode{\sphinxupquote{pygeox.custom\_objects.}}\sphinxbfcode{\sphinxupquote{ObtuseTriangle}}}
{\sphinxparam{\DUrole{n}{geoscene}}\sphinxparamcomma \sphinxparam{\DUrole{n}{p1}}\sphinxparamcomma \sphinxparam{\DUrole{n}{p2}}\sphinxparamcomma \sphinxparam{\DUrole{n}{p3}}\sphinxparamcomma \sphinxparam{\DUrole{n}{name}\DUrole{o}{=}\DUrole{default_value}{None}}\sphinxparamcomma \sphinxparam{\DUrole{n}{add\_to\_scene\_as\_type}\DUrole{o}{=}\DUrole{default_value}{\textquotesingle{}ObtuseTriangle\textquotesingle{}}}}
{}
\pysigstopsignatures
\sphinxAtStartPar
Bases: \sphinxcode{\sphinxupquote{Triangle}}

\sphinxAtStartPar
Triangle with one angle \textgreater{} 90 degrees.
\begin{quote}\begin{description}
\sphinxlineitem{Parameters}\begin{itemize}
\item {} 
\sphinxAtStartPar
\sphinxstyleliteralstrong{\sphinxupquote{p1}} (\sphinxstyleliteralemphasis{\sphinxupquote{Point}})

\item {} 
\sphinxAtStartPar
\sphinxstyleliteralstrong{\sphinxupquote{p2}} (\sphinxstyleliteralemphasis{\sphinxupquote{Point}})

\item {} 
\sphinxAtStartPar
\sphinxstyleliteralstrong{\sphinxupquote{p3}} (\sphinxstyleliteralemphasis{\sphinxupquote{Point}})

\item {} 
\sphinxAtStartPar
\sphinxstyleliteralstrong{\sphinxupquote{name}} (\sphinxstyleliteralemphasis{\sphinxupquote{str}}\sphinxstyleliteralemphasis{\sphinxupquote{ | }}\sphinxstyleliteralemphasis{\sphinxupquote{None}})

\item {} 
\sphinxAtStartPar
\sphinxstyleliteralstrong{\sphinxupquote{add\_to\_scene\_as\_type}} (\sphinxstyleliteralemphasis{\sphinxupquote{str}}\sphinxstyleliteralemphasis{\sphinxupquote{ | }}\sphinxstyleliteralemphasis{\sphinxupquote{None}})

\end{itemize}

\end{description}\end{quote}

\end{fulllineitems}

\index{Parallelogram (class in pygeox.custom\_objects)@\spxentry{Parallelogram}\spxextra{class in pygeox.custom\_objects}}

\begin{fulllineitems}
\phantomsection\label{\detokenize{pygeox:pygeox.custom_objects.Parallelogram}}
\pysigstartsignatures
\pysiglinewithargsret
{\sphinxbfcode{\sphinxupquote{\DUrole{k}{class}\DUrole{w}{ }}}\sphinxcode{\sphinxupquote{pygeox.custom\_objects.}}\sphinxbfcode{\sphinxupquote{Parallelogram}}}
{\sphinxparam{\DUrole{n}{geoscene}}\sphinxparamcomma \sphinxparam{\DUrole{n}{p1}}\sphinxparamcomma \sphinxparam{\DUrole{n}{p2}}\sphinxparamcomma \sphinxparam{\DUrole{n}{p3}}\sphinxparamcomma \sphinxparam{\DUrole{n}{p4}}\sphinxparamcomma \sphinxparam{\DUrole{n}{name}\DUrole{o}{=}\DUrole{default_value}{None}}\sphinxparamcomma \sphinxparam{\DUrole{n}{add\_to\_scene\_as\_type}\DUrole{o}{=}\DUrole{default_value}{\textquotesingle{}Parallelogram\textquotesingle{}}}\sphinxparamcomma \sphinxparam{\DUrole{n}{force\_convex}\DUrole{o}{=}\DUrole{default_value}{False}}}
{}
\pysigstopsignatures
\sphinxAtStartPar
Bases: \sphinxcode{\sphinxupquote{Quadrilateral}}

\sphinxAtStartPar
Quadrilateral with both pairs of opposite sides parallel.
\begin{quote}\begin{description}
\sphinxlineitem{Parameters}
\sphinxAtStartPar
\sphinxstyleliteralstrong{\sphinxupquote{force\_convex}} (\sphinxstyleliteralemphasis{\sphinxupquote{bool}})

\end{description}\end{quote}
\index{area (pygeox.custom\_objects.Parallelogram property)@\spxentry{area}\spxextra{pygeox.custom\_objects.Parallelogram property}}

\begin{fulllineitems}
\phantomsection\label{\detokenize{pygeox:pygeox.custom_objects.Parallelogram.area}}
\pysigstartsignatures
\pysigline
{\sphinxbfcode{\sphinxupquote{\DUrole{k}{property}\DUrole{w}{ }}}\sphinxbfcode{\sphinxupquote{area}}}
\pysigstopsignatures
\sphinxAtStartPar
Calculates the area of the polygon using the shoelace formula.
\begin{quote}\begin{description}
\sphinxlineitem{Returns}
\sphinxAtStartPar
The area. Units: square of input coordinates.

\sphinxlineitem{Return type}
\sphinxAtStartPar
sympy.Expr

\end{description}\end{quote}

\end{fulllineitems}


\end{fulllineitems}

\index{Polygon (class in pygeox.custom\_objects)@\spxentry{Polygon}\spxextra{class in pygeox.custom\_objects}}

\begin{fulllineitems}
\phantomsection\label{\detokenize{pygeox:pygeox.custom_objects.Polygon}}
\pysigstartsignatures
\pysiglinewithargsret
{\sphinxbfcode{\sphinxupquote{\DUrole{k}{class}\DUrole{w}{ }}}\sphinxcode{\sphinxupquote{pygeox.custom\_objects.}}\sphinxbfcode{\sphinxupquote{Polygon}}}
{\sphinxparam{\DUrole{n}{geoscene}}\sphinxparamcomma \sphinxparam{\DUrole{o}{*}\DUrole{n}{points}}\sphinxparamcomma \sphinxparam{\DUrole{n}{name}\DUrole{o}{=}\DUrole{default_value}{None}}\sphinxparamcomma \sphinxparam{\DUrole{n}{add\_to\_scene\_as\_type}\DUrole{o}{=}\DUrole{default_value}{\textquotesingle{}Polygon\textquotesingle{}}}\sphinxparamcomma \sphinxparam{\DUrole{n}{force\_convex}\DUrole{o}{=}\DUrole{default_value}{False}}}
{}
\pysigstopsignatures
\sphinxAtStartPar
Bases: \sphinxcode{\sphinxupquote{object}}

\sphinxAtStartPar
Represents a general polygon defined by an ordered list of Point objects.
\begin{quote}\begin{description}
\sphinxlineitem{Parameters}\begin{itemize}
\item {} 
\sphinxAtStartPar
\sphinxstyleliteralstrong{\sphinxupquote{geoscene}} \textendash{} The geometric scene/context to which this polygon belongs.

\item {} 
\sphinxAtStartPar
\sphinxstyleliteralstrong{\sphinxupquote{*points}} (\sphinxstyleliteralemphasis{\sphinxupquote{Point}}) \textendash{} The vertices of the polygon, in order.

\item {} 
\sphinxAtStartPar
\sphinxstyleliteralstrong{\sphinxupquote{name}} (\sphinxstyleliteralemphasis{\sphinxupquote{str}}\sphinxstyleliteralemphasis{\sphinxupquote{, }}\sphinxstyleliteralemphasis{\sphinxupquote{optional}}) \textendash{} Name for the polygon. If None, an automatic name is generated.

\item {} 
\sphinxAtStartPar
\sphinxstyleliteralstrong{\sphinxupquote{add\_to\_scene\_as\_type}} (\sphinxstyleliteralemphasis{\sphinxupquote{str}}\sphinxstyleliteralemphasis{\sphinxupquote{, }}\sphinxstyleliteralemphasis{\sphinxupquote{optional}}) \textendash{} Type string for scene registration. Default is ‘Polygon’.

\item {} 
\sphinxAtStartPar
\sphinxstyleliteralstrong{\sphinxupquote{force\_convex}} (\sphinxstyleliteralemphasis{\sphinxupquote{bool}}\sphinxstyleliteralemphasis{\sphinxupquote{, }}\sphinxstyleliteralemphasis{\sphinxupquote{optional}}) \textendash{} If True, enforces convexity and simplicity constraint (forcing also counterclockwise direction).

\end{itemize}

\sphinxlineitem{Returns}
\sphinxAtStartPar
A new Polygon object registered in the scene.

\sphinxlineitem{Return type}
\sphinxAtStartPar
Polygon

\end{description}\end{quote}
\index{area (pygeox.custom\_objects.Polygon property)@\spxentry{area}\spxextra{pygeox.custom\_objects.Polygon property}}

\begin{fulllineitems}
\phantomsection\label{\detokenize{pygeox:pygeox.custom_objects.Polygon.area}}
\pysigstartsignatures
\pysigline
{\sphinxbfcode{\sphinxupquote{\DUrole{k}{property}\DUrole{w}{ }}}\sphinxbfcode{\sphinxupquote{area}}}
\pysigstopsignatures
\sphinxAtStartPar
Calculates the area of the polygon using the shoelace formula.
\begin{quote}\begin{description}
\sphinxlineitem{Returns}
\sphinxAtStartPar
The area. Units: square of input coordinates.

\sphinxlineitem{Return type}
\sphinxAtStartPar
sympy.Expr

\end{description}\end{quote}

\end{fulllineitems}

\index{centroid (pygeox.custom\_objects.Polygon property)@\spxentry{centroid}\spxextra{pygeox.custom\_objects.Polygon property}}

\begin{fulllineitems}
\phantomsection\label{\detokenize{pygeox:pygeox.custom_objects.Polygon.centroid}}
\pysigstartsignatures
\pysigline
{\sphinxbfcode{\sphinxupquote{\DUrole{k}{property}\DUrole{w}{ }}}\sphinxbfcode{\sphinxupquote{centroid}}}
\pysigstopsignatures
\sphinxAtStartPar
Returns the vertex centroid (average of all vertices).
\begin{quote}\begin{description}
\sphinxlineitem{Returns}
\sphinxAtStartPar
The centroid as a Point object.

\sphinxlineitem{Return type}
\sphinxAtStartPar
Point

\end{description}\end{quote}

\end{fulllineitems}

\index{circumcircle (pygeox.custom\_objects.Polygon property)@\spxentry{circumcircle}\spxextra{pygeox.custom\_objects.Polygon property}}

\begin{fulllineitems}
\phantomsection\label{\detokenize{pygeox:pygeox.custom_objects.Polygon.circumcircle}}
\pysigstartsignatures
\pysigline
{\sphinxbfcode{\sphinxupquote{\DUrole{k}{property}\DUrole{w}{ }}}\sphinxbfcode{\sphinxupquote{circumcircle}}}
\pysigstopsignatures
\sphinxAtStartPar
Returns the circumcircle of the polygon if circumcenter and circumradius exist.
\begin{quote}\begin{description}
\sphinxlineitem{Returns}
\sphinxAtStartPar
The circumcircle.

\sphinxlineitem{Return type}
\sphinxAtStartPar
Circle

\sphinxlineitem{Raises}
\sphinxAtStartPar
\sphinxstyleliteralstrong{\sphinxupquote{NotImplementedError}} \textendash{} If circumcenter or circumradius is not defined for this polygon type.

\end{description}\end{quote}

\end{fulllineitems}

\index{incircle (pygeox.custom\_objects.Polygon property)@\spxentry{incircle}\spxextra{pygeox.custom\_objects.Polygon property}}

\begin{fulllineitems}
\phantomsection\label{\detokenize{pygeox:pygeox.custom_objects.Polygon.incircle}}
\pysigstartsignatures
\pysigline
{\sphinxbfcode{\sphinxupquote{\DUrole{k}{property}\DUrole{w}{ }}}\sphinxbfcode{\sphinxupquote{incircle}}}
\pysigstopsignatures
\sphinxAtStartPar
Returns the incircle of the polygon if incenter and inradius exist.
\begin{quote}\begin{description}
\sphinxlineitem{Returns}
\sphinxAtStartPar
The incircle.

\sphinxlineitem{Return type}
\sphinxAtStartPar
Circle

\sphinxlineitem{Raises}
\sphinxAtStartPar
\sphinxstyleliteralstrong{\sphinxupquote{NotImplementedError}} \textendash{} If incenter or inradius is not defined for this polygon type.

\end{description}\end{quote}

\end{fulllineitems}

\index{internal\_angle\_at\_point() (pygeox.custom\_objects.Polygon method)@\spxentry{internal\_angle\_at\_point()}\spxextra{pygeox.custom\_objects.Polygon method}}

\begin{fulllineitems}
\phantomsection\label{\detokenize{pygeox:pygeox.custom_objects.Polygon.internal_angle_at_point}}
\pysigstartsignatures
\pysiglinewithargsret
{\sphinxbfcode{\sphinxupquote{internal\_angle\_at\_point}}}
{\sphinxparam{\DUrole{n}{vertex\_point}}}
{}
\pysigstopsignatures
\sphinxAtStartPar
Returns the internal angle (in degrees) at a given vertex of the polygon,
correctly handling convex and concave cases.
\begin{quote}\begin{description}
\sphinxlineitem{Parameters}
\sphinxAtStartPar
\sphinxstyleliteralstrong{\sphinxupquote{vertex\_point}} (\sphinxstyleliteralemphasis{\sphinxupquote{Point}}) \textendash{} The vertex at which to compute the internal angle.

\sphinxlineitem{Returns}
\sphinxAtStartPar
Internal angle in degrees, in range (0, 360)

\sphinxlineitem{Return type}
\sphinxAtStartPar
sympy expression

\end{description}\end{quote}

\end{fulllineitems}

\index{internal\_angles (pygeox.custom\_objects.Polygon property)@\spxentry{internal\_angles}\spxextra{pygeox.custom\_objects.Polygon property}}

\begin{fulllineitems}
\phantomsection\label{\detokenize{pygeox:pygeox.custom_objects.Polygon.internal_angles}}
\pysigstartsignatures
\pysigline
{\sphinxbfcode{\sphinxupquote{\DUrole{k}{property}\DUrole{w}{ }}}\sphinxbfcode{\sphinxupquote{internal\_angles}}}
\pysigstopsignatures
\sphinxAtStartPar
Returns the internal angles at each vertex. Angles {[}DAB,ABC,BCD,..{]} for Polygon(A,B,C,D,…).
\begin{quote}\begin{description}
\sphinxlineitem{Returns}
\sphinxAtStartPar
The internal angles at each vertex in degrees.

\sphinxlineitem{Return type}
\sphinxAtStartPar
list{[}float{]}

\end{description}\end{quote}

\end{fulllineitems}

\index{num\_sides (pygeox.custom\_objects.Polygon property)@\spxentry{num\_sides}\spxextra{pygeox.custom\_objects.Polygon property}}

\begin{fulllineitems}
\phantomsection\label{\detokenize{pygeox:pygeox.custom_objects.Polygon.num_sides}}
\pysigstartsignatures
\pysigline
{\sphinxbfcode{\sphinxupquote{\DUrole{k}{property}\DUrole{w}{ }}}\sphinxbfcode{\sphinxupquote{num\_sides}}}
\pysigstopsignatures
\sphinxAtStartPar
Returns a string representation of the Polygon object.
\begin{quote}\begin{description}
\sphinxlineitem{Returns}
\sphinxAtStartPar
String representation.

\sphinxlineitem{Return type}
\sphinxAtStartPar
str

\end{description}\end{quote}

\end{fulllineitems}

\index{perimeter (pygeox.custom\_objects.Polygon property)@\spxentry{perimeter}\spxextra{pygeox.custom\_objects.Polygon property}}

\begin{fulllineitems}
\phantomsection\label{\detokenize{pygeox:pygeox.custom_objects.Polygon.perimeter}}
\pysigstartsignatures
\pysigline
{\sphinxbfcode{\sphinxupquote{\DUrole{k}{property}\DUrole{w}{ }}}\sphinxbfcode{\sphinxupquote{perimeter}}}
\pysigstopsignatures
\sphinxAtStartPar
Calculates the perimeter of the polygon by summing side lengths.
\begin{quote}\begin{description}
\sphinxlineitem{Returns}
\sphinxAtStartPar
The perimeter. Units: same as input coordinates.

\sphinxlineitem{Return type}
\sphinxAtStartPar
sympy.Expr

\end{description}\end{quote}

\end{fulllineitems}

\index{sides (pygeox.custom\_objects.Polygon property)@\spxentry{sides}\spxextra{pygeox.custom\_objects.Polygon property}}

\begin{fulllineitems}
\phantomsection\label{\detokenize{pygeox:pygeox.custom_objects.Polygon.sides}}
\pysigstartsignatures
\pysigline
{\sphinxbfcode{\sphinxupquote{\DUrole{k}{property}\DUrole{w}{ }}}\sphinxbfcode{\sphinxupquote{sides}}\sphinxbfcode{\sphinxupquote{\DUrole{p}{:}\DUrole{w}{ }List\DUrole{p}{{[}}LineSegment\DUrole{p}{{]}}}}}
\pysigstopsignatures
\sphinxAtStartPar
Returns the sides of the polygon as LijneSegment objects, in order. Sides {[}AB, BC, CD, …{]} for Polygon(A,B,C,D,…).
\begin{quote}\begin{description}
\sphinxlineitem{Returns}
\sphinxAtStartPar
The sides of the polygon.

\sphinxlineitem{Return type}
\sphinxAtStartPar
list{[}LijneSegment{]}

\end{description}\end{quote}

\end{fulllineitems}


\end{fulllineitems}

\index{Quadrilateral (class in pygeox.custom\_objects)@\spxentry{Quadrilateral}\spxextra{class in pygeox.custom\_objects}}

\begin{fulllineitems}
\phantomsection\label{\detokenize{pygeox:pygeox.custom_objects.Quadrilateral}}
\pysigstartsignatures
\pysiglinewithargsret
{\sphinxbfcode{\sphinxupquote{\DUrole{k}{class}\DUrole{w}{ }}}\sphinxcode{\sphinxupquote{pygeox.custom\_objects.}}\sphinxbfcode{\sphinxupquote{Quadrilateral}}}
{\sphinxparam{\DUrole{n}{geoscene}}\sphinxparamcomma \sphinxparam{\DUrole{n}{p1}}\sphinxparamcomma \sphinxparam{\DUrole{n}{p2}}\sphinxparamcomma \sphinxparam{\DUrole{n}{p3}}\sphinxparamcomma \sphinxparam{\DUrole{n}{p4}}\sphinxparamcomma \sphinxparam{\DUrole{n}{name}\DUrole{o}{=}\DUrole{default_value}{None}}\sphinxparamcomma \sphinxparam{\DUrole{n}{add\_to\_scene\_as\_type}\DUrole{o}{=}\DUrole{default_value}{\textquotesingle{}Quadrilateral\textquotesingle{}}}\sphinxparamcomma \sphinxparam{\DUrole{n}{force\_convex}\DUrole{o}{=}\DUrole{default_value}{False}}}
{}
\pysigstopsignatures
\sphinxAtStartPar
Bases: \sphinxcode{\sphinxupquote{Polygon}}

\sphinxAtStartPar
Represents a general quadrilateral defined by four Point objects.
Inherits area, perimeter, sides, and internal\_angles from Polygon.
\begin{quote}\begin{description}
\sphinxlineitem{Parameters}\begin{itemize}
\item {} 
\sphinxAtStartPar
\sphinxstyleliteralstrong{\sphinxupquote{p1}} (\sphinxstyleliteralemphasis{\sphinxupquote{Point}})

\item {} 
\sphinxAtStartPar
\sphinxstyleliteralstrong{\sphinxupquote{p2}} (\sphinxstyleliteralemphasis{\sphinxupquote{Point}})

\item {} 
\sphinxAtStartPar
\sphinxstyleliteralstrong{\sphinxupquote{p3}} (\sphinxstyleliteralemphasis{\sphinxupquote{Point}})

\item {} 
\sphinxAtStartPar
\sphinxstyleliteralstrong{\sphinxupquote{p4}} (\sphinxstyleliteralemphasis{\sphinxupquote{Point}})

\item {} 
\sphinxAtStartPar
\sphinxstyleliteralstrong{\sphinxupquote{name}} (\sphinxstyleliteralemphasis{\sphinxupquote{str}}\sphinxstyleliteralemphasis{\sphinxupquote{ | }}\sphinxstyleliteralemphasis{\sphinxupquote{None}})

\item {} 
\sphinxAtStartPar
\sphinxstyleliteralstrong{\sphinxupquote{add\_to\_scene\_as\_type}} (\sphinxstyleliteralemphasis{\sphinxupquote{str}}\sphinxstyleliteralemphasis{\sphinxupquote{ | }}\sphinxstyleliteralemphasis{\sphinxupquote{None}})

\item {} 
\sphinxAtStartPar
\sphinxstyleliteralstrong{\sphinxupquote{force\_convex}} (\sphinxstyleliteralemphasis{\sphinxupquote{bool}})

\end{itemize}

\end{description}\end{quote}
\index{diagonals (pygeox.custom\_objects.Quadrilateral property)@\spxentry{diagonals}\spxextra{pygeox.custom\_objects.Quadrilateral property}}

\begin{fulllineitems}
\phantomsection\label{\detokenize{pygeox:pygeox.custom_objects.Quadrilateral.diagonals}}
\pysigstartsignatures
\pysigline
{\sphinxbfcode{\sphinxupquote{\DUrole{k}{property}\DUrole{w}{ }}}\sphinxbfcode{\sphinxupquote{diagonals}}}
\pysigstopsignatures
\end{fulllineitems}

\index{midsegments (pygeox.custom\_objects.Quadrilateral property)@\spxentry{midsegments}\spxextra{pygeox.custom\_objects.Quadrilateral property}}

\begin{fulllineitems}
\phantomsection\label{\detokenize{pygeox:pygeox.custom_objects.Quadrilateral.midsegments}}
\pysigstartsignatures
\pysigline
{\sphinxbfcode{\sphinxupquote{\DUrole{k}{property}\DUrole{w}{ }}}\sphinxbfcode{\sphinxupquote{midsegments}}}
\pysigstopsignatures
\sphinxAtStartPar
Returns the two midsegments (lines connecting midpoints of opposite sides) of the quadrilateral.
\begin{quote}\begin{description}
\sphinxlineitem{Returns}
\sphinxAtStartPar
The two midsegments.

\sphinxlineitem{Return type}
\sphinxAtStartPar
list{[}LineSegment{]}

\end{description}\end{quote}

\end{fulllineitems}


\end{fulllineitems}

\index{Rectangle (class in pygeox.custom\_objects)@\spxentry{Rectangle}\spxextra{class in pygeox.custom\_objects}}

\begin{fulllineitems}
\phantomsection\label{\detokenize{pygeox:pygeox.custom_objects.Rectangle}}
\pysigstartsignatures
\pysiglinewithargsret
{\sphinxbfcode{\sphinxupquote{\DUrole{k}{class}\DUrole{w}{ }}}\sphinxcode{\sphinxupquote{pygeox.custom\_objects.}}\sphinxbfcode{\sphinxupquote{Rectangle}}}
{\sphinxparam{\DUrole{n}{geoscene}}\sphinxparamcomma \sphinxparam{\DUrole{n}{p1}}\sphinxparamcomma \sphinxparam{\DUrole{n}{p2}}\sphinxparamcomma \sphinxparam{\DUrole{n}{p3}}\sphinxparamcomma \sphinxparam{\DUrole{n}{p4}}\sphinxparamcomma \sphinxparam{\DUrole{n}{name}\DUrole{o}{=}\DUrole{default_value}{None}}\sphinxparamcomma \sphinxparam{\DUrole{n}{add\_to\_scene\_as\_type}\DUrole{o}{=}\DUrole{default_value}{\textquotesingle{}Rectangle\textquotesingle{}}}\sphinxparamcomma \sphinxparam{\DUrole{n}{force\_convex}\DUrole{o}{=}\DUrole{default_value}{False}}}
{}
\pysigstopsignatures
\sphinxAtStartPar
Bases: \sphinxcode{\sphinxupquote{Parallelogram}}

\sphinxAtStartPar
Parallelogram with all angles 90 degrees.
\begin{quote}\begin{description}
\sphinxlineitem{Parameters}
\sphinxAtStartPar
\sphinxstyleliteralstrong{\sphinxupquote{force\_convex}} (\sphinxstyleliteralemphasis{\sphinxupquote{bool}})

\end{description}\end{quote}
\index{area (pygeox.custom\_objects.Rectangle property)@\spxentry{area}\spxextra{pygeox.custom\_objects.Rectangle property}}

\begin{fulllineitems}
\phantomsection\label{\detokenize{pygeox:pygeox.custom_objects.Rectangle.area}}
\pysigstartsignatures
\pysigline
{\sphinxbfcode{\sphinxupquote{\DUrole{k}{property}\DUrole{w}{ }}}\sphinxbfcode{\sphinxupquote{area}}}
\pysigstopsignatures
\sphinxAtStartPar
Calculates the area of the polygon using the shoelace formula.
\begin{quote}\begin{description}
\sphinxlineitem{Returns}
\sphinxAtStartPar
The area. Units: square of input coordinates.

\sphinxlineitem{Return type}
\sphinxAtStartPar
sympy.Expr

\end{description}\end{quote}

\end{fulllineitems}

\index{circumcenter (pygeox.custom\_objects.Rectangle property)@\spxentry{circumcenter}\spxextra{pygeox.custom\_objects.Rectangle property}}

\begin{fulllineitems}
\phantomsection\label{\detokenize{pygeox:pygeox.custom_objects.Rectangle.circumcenter}}
\pysigstartsignatures
\pysigline
{\sphinxbfcode{\sphinxupquote{\DUrole{k}{property}\DUrole{w}{ }}}\sphinxbfcode{\sphinxupquote{circumcenter}}}
\pysigstopsignatures
\end{fulllineitems}

\index{circumradius (pygeox.custom\_objects.Rectangle property)@\spxentry{circumradius}\spxextra{pygeox.custom\_objects.Rectangle property}}

\begin{fulllineitems}
\phantomsection\label{\detokenize{pygeox:pygeox.custom_objects.Rectangle.circumradius}}
\pysigstartsignatures
\pysigline
{\sphinxbfcode{\sphinxupquote{\DUrole{k}{property}\DUrole{w}{ }}}\sphinxbfcode{\sphinxupquote{circumradius}}}
\pysigstopsignatures
\end{fulllineitems}

\index{diagonals (pygeox.custom\_objects.Rectangle property)@\spxentry{diagonals}\spxextra{pygeox.custom\_objects.Rectangle property}}

\begin{fulllineitems}
\phantomsection\label{\detokenize{pygeox:pygeox.custom_objects.Rectangle.diagonals}}
\pysigstartsignatures
\pysigline
{\sphinxbfcode{\sphinxupquote{\DUrole{k}{property}\DUrole{w}{ }}}\sphinxbfcode{\sphinxupquote{diagonals}}}
\pysigstopsignatures
\end{fulllineitems}

\index{perimeter (pygeox.custom\_objects.Rectangle property)@\spxentry{perimeter}\spxextra{pygeox.custom\_objects.Rectangle property}}

\begin{fulllineitems}
\phantomsection\label{\detokenize{pygeox:pygeox.custom_objects.Rectangle.perimeter}}
\pysigstartsignatures
\pysigline
{\sphinxbfcode{\sphinxupquote{\DUrole{k}{property}\DUrole{w}{ }}}\sphinxbfcode{\sphinxupquote{perimeter}}}
\pysigstopsignatures
\sphinxAtStartPar
Calculates the perimeter of the polygon by summing side lengths.
\begin{quote}\begin{description}
\sphinxlineitem{Returns}
\sphinxAtStartPar
The perimeter. Units: same as input coordinates.

\sphinxlineitem{Return type}
\sphinxAtStartPar
sympy.Expr

\end{description}\end{quote}

\end{fulllineitems}


\end{fulllineitems}

\index{RegularHeptagon (class in pygeox.custom\_objects)@\spxentry{RegularHeptagon}\spxextra{class in pygeox.custom\_objects}}

\begin{fulllineitems}
\phantomsection\label{\detokenize{pygeox:pygeox.custom_objects.RegularHeptagon}}
\pysigstartsignatures
\pysiglinewithargsret
{\sphinxbfcode{\sphinxupquote{\DUrole{k}{class}\DUrole{w}{ }}}\sphinxcode{\sphinxupquote{pygeox.custom\_objects.}}\sphinxbfcode{\sphinxupquote{RegularHeptagon}}}
{\sphinxparam{\DUrole{n}{geoscene}}\sphinxparamcomma \sphinxparam{\DUrole{o}{*}\DUrole{n}{points}}\sphinxparamcomma \sphinxparam{\DUrole{n}{name}\DUrole{o}{=}\DUrole{default_value}{None}}\sphinxparamcomma \sphinxparam{\DUrole{n}{add\_to\_scene\_as\_type}\DUrole{o}{=}\DUrole{default_value}{\textquotesingle{}RegularHeptagon\textquotesingle{}}}}
{}
\pysigstopsignatures
\sphinxAtStartPar
Bases: \sphinxcode{\sphinxupquote{RegularPolygon}}

\sphinxAtStartPar
A specific type of RegularPolygon with 7 sides.
\begin{quote}\begin{description}
\sphinxlineitem{Parameters}\begin{itemize}
\item {} 
\sphinxAtStartPar
\sphinxstyleliteralstrong{\sphinxupquote{points}} (\sphinxstyleliteralemphasis{\sphinxupquote{Point}})

\item {} 
\sphinxAtStartPar
\sphinxstyleliteralstrong{\sphinxupquote{name}} (\sphinxstyleliteralemphasis{\sphinxupquote{str}}\sphinxstyleliteralemphasis{\sphinxupquote{ | }}\sphinxstyleliteralemphasis{\sphinxupquote{None}})

\item {} 
\sphinxAtStartPar
\sphinxstyleliteralstrong{\sphinxupquote{add\_to\_scene\_as\_type}} (\sphinxstyleliteralemphasis{\sphinxupquote{str}}\sphinxstyleliteralemphasis{\sphinxupquote{ | }}\sphinxstyleliteralemphasis{\sphinxupquote{None}})

\end{itemize}

\end{description}\end{quote}

\end{fulllineitems}

\index{RegularHexagon (class in pygeox.custom\_objects)@\spxentry{RegularHexagon}\spxextra{class in pygeox.custom\_objects}}

\begin{fulllineitems}
\phantomsection\label{\detokenize{pygeox:pygeox.custom_objects.RegularHexagon}}
\pysigstartsignatures
\pysiglinewithargsret
{\sphinxbfcode{\sphinxupquote{\DUrole{k}{class}\DUrole{w}{ }}}\sphinxcode{\sphinxupquote{pygeox.custom\_objects.}}\sphinxbfcode{\sphinxupquote{RegularHexagon}}}
{\sphinxparam{\DUrole{n}{geoscene}}\sphinxparamcomma \sphinxparam{\DUrole{o}{*}\DUrole{n}{points}}\sphinxparamcomma \sphinxparam{\DUrole{n}{name}\DUrole{o}{=}\DUrole{default_value}{None}}\sphinxparamcomma \sphinxparam{\DUrole{n}{add\_to\_scene\_as\_type}\DUrole{o}{=}\DUrole{default_value}{\textquotesingle{}RegularHexagon\textquotesingle{}}}}
{}
\pysigstopsignatures
\sphinxAtStartPar
Bases: \sphinxcode{\sphinxupquote{RegularPolygon}}

\sphinxAtStartPar
A specific type of RegularPolygon with 6 sides.
\begin{quote}\begin{description}
\sphinxlineitem{Parameters}\begin{itemize}
\item {} 
\sphinxAtStartPar
\sphinxstyleliteralstrong{\sphinxupquote{points}} (\sphinxstyleliteralemphasis{\sphinxupquote{Point}})

\item {} 
\sphinxAtStartPar
\sphinxstyleliteralstrong{\sphinxupquote{name}} (\sphinxstyleliteralemphasis{\sphinxupquote{str}}\sphinxstyleliteralemphasis{\sphinxupquote{ | }}\sphinxstyleliteralemphasis{\sphinxupquote{None}})

\item {} 
\sphinxAtStartPar
\sphinxstyleliteralstrong{\sphinxupquote{add\_to\_scene\_as\_type}} (\sphinxstyleliteralemphasis{\sphinxupquote{str}}\sphinxstyleliteralemphasis{\sphinxupquote{ | }}\sphinxstyleliteralemphasis{\sphinxupquote{None}})

\end{itemize}

\end{description}\end{quote}

\end{fulllineitems}

\index{RegularOctagon (class in pygeox.custom\_objects)@\spxentry{RegularOctagon}\spxextra{class in pygeox.custom\_objects}}

\begin{fulllineitems}
\phantomsection\label{\detokenize{pygeox:pygeox.custom_objects.RegularOctagon}}
\pysigstartsignatures
\pysiglinewithargsret
{\sphinxbfcode{\sphinxupquote{\DUrole{k}{class}\DUrole{w}{ }}}\sphinxcode{\sphinxupquote{pygeox.custom\_objects.}}\sphinxbfcode{\sphinxupquote{RegularOctagon}}}
{\sphinxparam{\DUrole{n}{geoscene}}\sphinxparamcomma \sphinxparam{\DUrole{o}{*}\DUrole{n}{points}}\sphinxparamcomma \sphinxparam{\DUrole{n}{name}\DUrole{o}{=}\DUrole{default_value}{None}}\sphinxparamcomma \sphinxparam{\DUrole{n}{add\_to\_scene\_as\_type}\DUrole{o}{=}\DUrole{default_value}{\textquotesingle{}RegularOctagon\textquotesingle{}}}}
{}
\pysigstopsignatures
\sphinxAtStartPar
Bases: \sphinxcode{\sphinxupquote{RegularPolygon}}

\sphinxAtStartPar
A specific type of RegularPolygon with 8 sides.
\begin{quote}\begin{description}
\sphinxlineitem{Parameters}\begin{itemize}
\item {} 
\sphinxAtStartPar
\sphinxstyleliteralstrong{\sphinxupquote{points}} (\sphinxstyleliteralemphasis{\sphinxupquote{Point}})

\item {} 
\sphinxAtStartPar
\sphinxstyleliteralstrong{\sphinxupquote{name}} (\sphinxstyleliteralemphasis{\sphinxupquote{str}}\sphinxstyleliteralemphasis{\sphinxupquote{ | }}\sphinxstyleliteralemphasis{\sphinxupquote{None}})

\item {} 
\sphinxAtStartPar
\sphinxstyleliteralstrong{\sphinxupquote{add\_to\_scene\_as\_type}} (\sphinxstyleliteralemphasis{\sphinxupquote{str}}\sphinxstyleliteralemphasis{\sphinxupquote{ | }}\sphinxstyleliteralemphasis{\sphinxupquote{None}})

\end{itemize}

\end{description}\end{quote}

\end{fulllineitems}

\index{RegularPentagon (class in pygeox.custom\_objects)@\spxentry{RegularPentagon}\spxextra{class in pygeox.custom\_objects}}

\begin{fulllineitems}
\phantomsection\label{\detokenize{pygeox:pygeox.custom_objects.RegularPentagon}}
\pysigstartsignatures
\pysiglinewithargsret
{\sphinxbfcode{\sphinxupquote{\DUrole{k}{class}\DUrole{w}{ }}}\sphinxcode{\sphinxupquote{pygeox.custom\_objects.}}\sphinxbfcode{\sphinxupquote{RegularPentagon}}}
{\sphinxparam{\DUrole{n}{geoscene}}\sphinxparamcomma \sphinxparam{\DUrole{o}{*}\DUrole{n}{points}}\sphinxparamcomma \sphinxparam{\DUrole{n}{name}\DUrole{o}{=}\DUrole{default_value}{None}}\sphinxparamcomma \sphinxparam{\DUrole{n}{add\_to\_scene\_as\_type}\DUrole{o}{=}\DUrole{default_value}{\textquotesingle{}RegularPentagon\textquotesingle{}}}}
{}
\pysigstopsignatures
\sphinxAtStartPar
Bases: \sphinxcode{\sphinxupquote{RegularPolygon}}

\sphinxAtStartPar
A specific type of RegularPolygon with 5 sides.
\begin{quote}\begin{description}
\sphinxlineitem{Parameters}\begin{itemize}
\item {} 
\sphinxAtStartPar
\sphinxstyleliteralstrong{\sphinxupquote{points}} (\sphinxstyleliteralemphasis{\sphinxupquote{Point}})

\item {} 
\sphinxAtStartPar
\sphinxstyleliteralstrong{\sphinxupquote{name}} (\sphinxstyleliteralemphasis{\sphinxupquote{str}}\sphinxstyleliteralemphasis{\sphinxupquote{ | }}\sphinxstyleliteralemphasis{\sphinxupquote{None}})

\item {} 
\sphinxAtStartPar
\sphinxstyleliteralstrong{\sphinxupquote{add\_to\_scene\_as\_type}} (\sphinxstyleliteralemphasis{\sphinxupquote{str}}\sphinxstyleliteralemphasis{\sphinxupquote{ | }}\sphinxstyleliteralemphasis{\sphinxupquote{None}})

\end{itemize}

\end{description}\end{quote}

\end{fulllineitems}

\index{RegularPolygon (class in pygeox.custom\_objects)@\spxentry{RegularPolygon}\spxextra{class in pygeox.custom\_objects}}

\begin{fulllineitems}
\phantomsection\label{\detokenize{pygeox:pygeox.custom_objects.RegularPolygon}}
\pysigstartsignatures
\pysiglinewithargsret
{\sphinxbfcode{\sphinxupquote{\DUrole{k}{class}\DUrole{w}{ }}}\sphinxcode{\sphinxupquote{pygeox.custom\_objects.}}\sphinxbfcode{\sphinxupquote{RegularPolygon}}}
{\sphinxparam{\DUrole{n}{geoscene}}\sphinxparamcomma \sphinxparam{\DUrole{o}{*}\DUrole{n}{points}}\sphinxparamcomma \sphinxparam{\DUrole{n}{name}\DUrole{o}{=}\DUrole{default_value}{None}}\sphinxparamcomma \sphinxparam{\DUrole{n}{add\_to\_scene\_as\_type}\DUrole{o}{=}\DUrole{default_value}{\textquotesingle{}RegularPolygon\textquotesingle{}}}}
{}
\pysigstopsignatures
\sphinxAtStartPar
Bases: \sphinxcode{\sphinxupquote{Polygon}}

\sphinxAtStartPar
Represents a regular polygon (all sides and angles equal).
It imposes regularity constraints on the provided points.
Defines its own reduced parameter set (center, radius, orientation).
\begin{quote}\begin{description}
\sphinxlineitem{Parameters}\begin{itemize}
\item {} 
\sphinxAtStartPar
\sphinxstyleliteralstrong{\sphinxupquote{points}} (\sphinxstyleliteralemphasis{\sphinxupquote{Point}})

\item {} 
\sphinxAtStartPar
\sphinxstyleliteralstrong{\sphinxupquote{name}} (\sphinxstyleliteralemphasis{\sphinxupquote{str}}\sphinxstyleliteralemphasis{\sphinxupquote{ | }}\sphinxstyleliteralemphasis{\sphinxupquote{None}})

\item {} 
\sphinxAtStartPar
\sphinxstyleliteralstrong{\sphinxupquote{add\_to\_scene\_as\_type}} (\sphinxstyleliteralemphasis{\sphinxupquote{str}}\sphinxstyleliteralemphasis{\sphinxupquote{ | }}\sphinxstyleliteralemphasis{\sphinxupquote{None}})

\end{itemize}

\end{description}\end{quote}
\index{area (pygeox.custom\_objects.RegularPolygon property)@\spxentry{area}\spxextra{pygeox.custom\_objects.RegularPolygon property}}

\begin{fulllineitems}
\phantomsection\label{\detokenize{pygeox:pygeox.custom_objects.RegularPolygon.area}}
\pysigstartsignatures
\pysigline
{\sphinxbfcode{\sphinxupquote{\DUrole{k}{property}\DUrole{w}{ }}}\sphinxbfcode{\sphinxupquote{area}}}
\pysigstopsignatures
\sphinxAtStartPar
Calculates the area of the regular polygon using its properties.

\end{fulllineitems}

\index{center (pygeox.custom\_objects.RegularPolygon property)@\spxentry{center}\spxextra{pygeox.custom\_objects.RegularPolygon property}}

\begin{fulllineitems}
\phantomsection\label{\detokenize{pygeox:pygeox.custom_objects.RegularPolygon.center}}
\pysigstartsignatures
\pysigline
{\sphinxbfcode{\sphinxupquote{\DUrole{k}{property}\DUrole{w}{ }}}\sphinxbfcode{\sphinxupquote{center}}}
\pysigstopsignatures
\sphinxAtStartPar
Returns the center point of the regular polygon.

\end{fulllineitems}

\index{circumcenter (pygeox.custom\_objects.RegularPolygon property)@\spxentry{circumcenter}\spxextra{pygeox.custom\_objects.RegularPolygon property}}

\begin{fulllineitems}
\phantomsection\label{\detokenize{pygeox:pygeox.custom_objects.RegularPolygon.circumcenter}}
\pysigstartsignatures
\pysigline
{\sphinxbfcode{\sphinxupquote{\DUrole{k}{property}\DUrole{w}{ }}}\sphinxbfcode{\sphinxupquote{circumcenter}}}
\pysigstopsignatures
\sphinxAtStartPar
Returns:
Point: The circumcenter of the regular polygon, which coincides with the center for regular polygons.

\end{fulllineitems}

\index{circumradius (pygeox.custom\_objects.RegularPolygon property)@\spxentry{circumradius}\spxextra{pygeox.custom\_objects.RegularPolygon property}}

\begin{fulllineitems}
\phantomsection\label{\detokenize{pygeox:pygeox.custom_objects.RegularPolygon.circumradius}}
\pysigstartsignatures
\pysigline
{\sphinxbfcode{\sphinxupquote{\DUrole{k}{property}\DUrole{w}{ }}}\sphinxbfcode{\sphinxupquote{circumradius}}}
\pysigstopsignatures
\sphinxAtStartPar
Returns the circumradius of the regular polygon.

\end{fulllineitems}

\index{diagonal (pygeox.custom\_objects.RegularPolygon property)@\spxentry{diagonal}\spxextra{pygeox.custom\_objects.RegularPolygon property}}

\begin{fulllineitems}
\phantomsection\label{\detokenize{pygeox:pygeox.custom_objects.RegularPolygon.diagonal}}
\pysigstartsignatures
\pysigline
{\sphinxbfcode{\sphinxupquote{\DUrole{k}{property}\DUrole{w}{ }}}\sphinxbfcode{\sphinxupquote{diagonal}}}
\pysigstopsignatures
\sphinxAtStartPar
Longest diagonal of the regular polygon.

\end{fulllineitems}

\index{height (pygeox.custom\_objects.RegularPolygon property)@\spxentry{height}\spxextra{pygeox.custom\_objects.RegularPolygon property}}

\begin{fulllineitems}
\phantomsection\label{\detokenize{pygeox:pygeox.custom_objects.RegularPolygon.height}}
\pysigstartsignatures
\pysigline
{\sphinxbfcode{\sphinxupquote{\DUrole{k}{property}\DUrole{w}{ }}}\sphinxbfcode{\sphinxupquote{height}}}
\pysigstopsignatures
\sphinxAtStartPar
Height (distance from one side to the opposite vertex) of the regular polygon.

\end{fulllineitems}

\index{incenter (pygeox.custom\_objects.RegularPolygon property)@\spxentry{incenter}\spxextra{pygeox.custom\_objects.RegularPolygon property}}

\begin{fulllineitems}
\phantomsection\label{\detokenize{pygeox:pygeox.custom_objects.RegularPolygon.incenter}}
\pysigstartsignatures
\pysigline
{\sphinxbfcode{\sphinxupquote{\DUrole{k}{property}\DUrole{w}{ }}}\sphinxbfcode{\sphinxupquote{incenter}}}
\pysigstopsignatures
\sphinxAtStartPar
Returns:
Point: The incenter of the regular polygon, which coincides with the center for regular polygons.

\end{fulllineitems}

\index{inradius (pygeox.custom\_objects.RegularPolygon property)@\spxentry{inradius}\spxextra{pygeox.custom\_objects.RegularPolygon property}}

\begin{fulllineitems}
\phantomsection\label{\detokenize{pygeox:pygeox.custom_objects.RegularPolygon.inradius}}
\pysigstartsignatures
\pysigline
{\sphinxbfcode{\sphinxupquote{\DUrole{k}{property}\DUrole{w}{ }}}\sphinxbfcode{\sphinxupquote{inradius}}}
\pysigstopsignatures
\sphinxAtStartPar
Inradius (apothem) of the regular polygon.

\end{fulllineitems}

\index{internal\_angle (pygeox.custom\_objects.RegularPolygon property)@\spxentry{internal\_angle}\spxextra{pygeox.custom\_objects.RegularPolygon property}}

\begin{fulllineitems}
\phantomsection\label{\detokenize{pygeox:pygeox.custom_objects.RegularPolygon.internal_angle}}
\pysigstartsignatures
\pysigline
{\sphinxbfcode{\sphinxupquote{\DUrole{k}{property}\DUrole{w}{ }}}\sphinxbfcode{\sphinxupquote{internal\_angle}}}
\pysigstopsignatures
\sphinxAtStartPar
Returns:
float: The internal angle (in degrees) of the regular polygon.

\end{fulllineitems}

\index{orientation (pygeox.custom\_objects.RegularPolygon property)@\spxentry{orientation}\spxextra{pygeox.custom\_objects.RegularPolygon property}}

\begin{fulllineitems}
\phantomsection\label{\detokenize{pygeox:pygeox.custom_objects.RegularPolygon.orientation}}
\pysigstartsignatures
\pysigline
{\sphinxbfcode{\sphinxupquote{\DUrole{k}{property}\DUrole{w}{ }}}\sphinxbfcode{\sphinxupquote{orientation}}}
\pysigstopsignatures
\sphinxAtStartPar
Returns the orientation (rotation) of the regular polygon in radians.

\end{fulllineitems}

\index{perimeter (pygeox.custom\_objects.RegularPolygon property)@\spxentry{perimeter}\spxextra{pygeox.custom\_objects.RegularPolygon property}}

\begin{fulllineitems}
\phantomsection\label{\detokenize{pygeox:pygeox.custom_objects.RegularPolygon.perimeter}}
\pysigstartsignatures
\pysigline
{\sphinxbfcode{\sphinxupquote{\DUrole{k}{property}\DUrole{w}{ }}}\sphinxbfcode{\sphinxupquote{perimeter}}}
\pysigstopsignatures
\sphinxAtStartPar
Calculates the perimeter of the regular polygon using its properties.

\end{fulllineitems}

\index{side\_length (pygeox.custom\_objects.RegularPolygon property)@\spxentry{side\_length}\spxextra{pygeox.custom\_objects.RegularPolygon property}}

\begin{fulllineitems}
\phantomsection\label{\detokenize{pygeox:pygeox.custom_objects.RegularPolygon.side_length}}
\pysigstartsignatures
\pysigline
{\sphinxbfcode{\sphinxupquote{\DUrole{k}{property}\DUrole{w}{ }}}\sphinxbfcode{\sphinxupquote{side\_length}}}
\pysigstopsignatures
\sphinxAtStartPar
Side length of the regular polygon.

\end{fulllineitems}

\index{width (pygeox.custom\_objects.RegularPolygon property)@\spxentry{width}\spxextra{pygeox.custom\_objects.RegularPolygon property}}

\begin{fulllineitems}
\phantomsection\label{\detokenize{pygeox:pygeox.custom_objects.RegularPolygon.width}}
\pysigstartsignatures
\pysigline
{\sphinxbfcode{\sphinxupquote{\DUrole{k}{property}\DUrole{w}{ }}}\sphinxbfcode{\sphinxupquote{width}}}
\pysigstopsignatures
\sphinxAtStartPar
Width (distance between two farthest separated points) of the regular polygon.

\end{fulllineitems}


\end{fulllineitems}

\index{Rhomboid (class in pygeox.custom\_objects)@\spxentry{Rhomboid}\spxextra{class in pygeox.custom\_objects}}

\begin{fulllineitems}
\phantomsection\label{\detokenize{pygeox:pygeox.custom_objects.Rhomboid}}
\pysigstartsignatures
\pysiglinewithargsret
{\sphinxbfcode{\sphinxupquote{\DUrole{k}{class}\DUrole{w}{ }}}\sphinxcode{\sphinxupquote{pygeox.custom\_objects.}}\sphinxbfcode{\sphinxupquote{Rhomboid}}}
{\sphinxparam{\DUrole{n}{geoscene}}\sphinxparamcomma \sphinxparam{\DUrole{n}{p1}}\sphinxparamcomma \sphinxparam{\DUrole{n}{p2}}\sphinxparamcomma \sphinxparam{\DUrole{n}{p3}}\sphinxparamcomma \sphinxparam{\DUrole{n}{p4}}\sphinxparamcomma \sphinxparam{\DUrole{n}{name}\DUrole{o}{=}\DUrole{default_value}{None}}\sphinxparamcomma \sphinxparam{\DUrole{n}{add\_to\_scene\_as\_type}\DUrole{o}{=}\DUrole{default_value}{\textquotesingle{}Rhomboid\textquotesingle{}}}\sphinxparamcomma \sphinxparam{\DUrole{n}{force\_convex}\DUrole{o}{=}\DUrole{default_value}{False}}}
{}
\pysigstopsignatures
\sphinxAtStartPar
Bases: \sphinxcode{\sphinxupquote{Parallelogram}}

\sphinxAtStartPar
Parallelogram with adjacent sides of unequal lengths and angles not right.
\begin{quote}\begin{description}
\sphinxlineitem{Parameters}
\sphinxAtStartPar
\sphinxstyleliteralstrong{\sphinxupquote{force\_convex}} (\sphinxstyleliteralemphasis{\sphinxupquote{bool}})

\end{description}\end{quote}

\end{fulllineitems}

\index{Rhombus (class in pygeox.custom\_objects)@\spxentry{Rhombus}\spxextra{class in pygeox.custom\_objects}}

\begin{fulllineitems}
\phantomsection\label{\detokenize{pygeox:pygeox.custom_objects.Rhombus}}
\pysigstartsignatures
\pysiglinewithargsret
{\sphinxbfcode{\sphinxupquote{\DUrole{k}{class}\DUrole{w}{ }}}\sphinxcode{\sphinxupquote{pygeox.custom\_objects.}}\sphinxbfcode{\sphinxupquote{Rhombus}}}
{\sphinxparam{\DUrole{n}{geoscene}}\sphinxparamcomma \sphinxparam{\DUrole{n}{p1}}\sphinxparamcomma \sphinxparam{\DUrole{n}{p2}}\sphinxparamcomma \sphinxparam{\DUrole{n}{p3}}\sphinxparamcomma \sphinxparam{\DUrole{n}{p4}}\sphinxparamcomma \sphinxparam{\DUrole{n}{name}\DUrole{o}{=}\DUrole{default_value}{None}}\sphinxparamcomma \sphinxparam{\DUrole{n}{add\_to\_scene\_as\_type}\DUrole{o}{=}\DUrole{default_value}{\textquotesingle{}Rhombus\textquotesingle{}}}\sphinxparamcomma \sphinxparam{\DUrole{n}{force\_convex}\DUrole{o}{=}\DUrole{default_value}{False}}}
{}
\pysigstopsignatures
\sphinxAtStartPar
Bases: \sphinxcode{\sphinxupquote{Parallelogram}}

\sphinxAtStartPar
Parallelogram with all sides equal.
\begin{quote}\begin{description}
\sphinxlineitem{Parameters}
\sphinxAtStartPar
\sphinxstyleliteralstrong{\sphinxupquote{force\_convex}} (\sphinxstyleliteralemphasis{\sphinxupquote{bool}})

\end{description}\end{quote}
\index{area (pygeox.custom\_objects.Rhombus property)@\spxentry{area}\spxextra{pygeox.custom\_objects.Rhombus property}}

\begin{fulllineitems}
\phantomsection\label{\detokenize{pygeox:pygeox.custom_objects.Rhombus.area}}
\pysigstartsignatures
\pysigline
{\sphinxbfcode{\sphinxupquote{\DUrole{k}{property}\DUrole{w}{ }}}\sphinxbfcode{\sphinxupquote{area}}}
\pysigstopsignatures
\sphinxAtStartPar
Calculates the area of the polygon using the shoelace formula.
\begin{quote}\begin{description}
\sphinxlineitem{Returns}
\sphinxAtStartPar
The area. Units: square of input coordinates.

\sphinxlineitem{Return type}
\sphinxAtStartPar
sympy.Expr

\end{description}\end{quote}

\end{fulllineitems}

\index{incenter (pygeox.custom\_objects.Rhombus property)@\spxentry{incenter}\spxextra{pygeox.custom\_objects.Rhombus property}}

\begin{fulllineitems}
\phantomsection\label{\detokenize{pygeox:pygeox.custom_objects.Rhombus.incenter}}
\pysigstartsignatures
\pysigline
{\sphinxbfcode{\sphinxupquote{\DUrole{k}{property}\DUrole{w}{ }}}\sphinxbfcode{\sphinxupquote{incenter}}}
\pysigstopsignatures
\end{fulllineitems}

\index{inradius (pygeox.custom\_objects.Rhombus property)@\spxentry{inradius}\spxextra{pygeox.custom\_objects.Rhombus property}}

\begin{fulllineitems}
\phantomsection\label{\detokenize{pygeox:pygeox.custom_objects.Rhombus.inradius}}
\pysigstartsignatures
\pysigline
{\sphinxbfcode{\sphinxupquote{\DUrole{k}{property}\DUrole{w}{ }}}\sphinxbfcode{\sphinxupquote{inradius}}}
\pysigstopsignatures
\end{fulllineitems}


\end{fulllineitems}

\index{RightIsoscelesTriangle (class in pygeox.custom\_objects)@\spxentry{RightIsoscelesTriangle}\spxextra{class in pygeox.custom\_objects}}

\begin{fulllineitems}
\phantomsection\label{\detokenize{pygeox:pygeox.custom_objects.RightIsoscelesTriangle}}
\pysigstartsignatures
\pysiglinewithargsret
{\sphinxbfcode{\sphinxupquote{\DUrole{k}{class}\DUrole{w}{ }}}\sphinxcode{\sphinxupquote{pygeox.custom\_objects.}}\sphinxbfcode{\sphinxupquote{RightIsoscelesTriangle}}}
{\sphinxparam{\DUrole{n}{geoscene}}\sphinxparamcomma \sphinxparam{\DUrole{n}{p1}}\sphinxparamcomma \sphinxparam{\DUrole{n}{p2}}\sphinxparamcomma \sphinxparam{\DUrole{n}{p3}}\sphinxparamcomma \sphinxparam{\DUrole{n}{name}\DUrole{o}{=}\DUrole{default_value}{None}}\sphinxparamcomma \sphinxparam{\DUrole{n}{add\_to\_scene\_as\_type}\DUrole{o}{=}\DUrole{default_value}{\textquotesingle{}RightIsoscelesTriangle\textquotesingle{}}}}
{}
\pysigstopsignatures
\sphinxAtStartPar
Bases: \sphinxcode{\sphinxupquote{RightTriangle}}, \sphinxcode{\sphinxupquote{IsoscelesTriangle}}

\sphinxAtStartPar
Triangle with a right angle and two equal sides. RightIsoscelesTriangle(A,B,C) means that AB=BC and B is the right angle.
\begin{quote}\begin{description}
\sphinxlineitem{Parameters}\begin{itemize}
\item {} 
\sphinxAtStartPar
\sphinxstyleliteralstrong{\sphinxupquote{p1}} (\sphinxstyleliteralemphasis{\sphinxupquote{Point}})

\item {} 
\sphinxAtStartPar
\sphinxstyleliteralstrong{\sphinxupquote{p2}} (\sphinxstyleliteralemphasis{\sphinxupquote{Point}})

\item {} 
\sphinxAtStartPar
\sphinxstyleliteralstrong{\sphinxupquote{p3}} (\sphinxstyleliteralemphasis{\sphinxupquote{Point}})

\item {} 
\sphinxAtStartPar
\sphinxstyleliteralstrong{\sphinxupquote{name}} (\sphinxstyleliteralemphasis{\sphinxupquote{str}}\sphinxstyleliteralemphasis{\sphinxupquote{ | }}\sphinxstyleliteralemphasis{\sphinxupquote{None}})

\item {} 
\sphinxAtStartPar
\sphinxstyleliteralstrong{\sphinxupquote{add\_to\_scene\_as\_type}} (\sphinxstyleliteralemphasis{\sphinxupquote{str}}\sphinxstyleliteralemphasis{\sphinxupquote{ | }}\sphinxstyleliteralemphasis{\sphinxupquote{None}})

\end{itemize}

\end{description}\end{quote}
\index{circumcenter (pygeox.custom\_objects.RightIsoscelesTriangle property)@\spxentry{circumcenter}\spxextra{pygeox.custom\_objects.RightIsoscelesTriangle property}}

\begin{fulllineitems}
\phantomsection\label{\detokenize{pygeox:pygeox.custom_objects.RightIsoscelesTriangle.circumcenter}}
\pysigstartsignatures
\pysigline
{\sphinxbfcode{\sphinxupquote{\DUrole{k}{property}\DUrole{w}{ }}}\sphinxbfcode{\sphinxupquote{circumcenter}}}
\pysigstopsignatures
\end{fulllineitems}

\index{circumradius (pygeox.custom\_objects.RightIsoscelesTriangle property)@\spxentry{circumradius}\spxextra{pygeox.custom\_objects.RightIsoscelesTriangle property}}

\begin{fulllineitems}
\phantomsection\label{\detokenize{pygeox:pygeox.custom_objects.RightIsoscelesTriangle.circumradius}}
\pysigstartsignatures
\pysigline
{\sphinxbfcode{\sphinxupquote{\DUrole{k}{property}\DUrole{w}{ }}}\sphinxbfcode{\sphinxupquote{circumradius}}}
\pysigstopsignatures
\end{fulllineitems}


\end{fulllineitems}

\index{RightTrapezoid (class in pygeox.custom\_objects)@\spxentry{RightTrapezoid}\spxextra{class in pygeox.custom\_objects}}

\begin{fulllineitems}
\phantomsection\label{\detokenize{pygeox:pygeox.custom_objects.RightTrapezoid}}
\pysigstartsignatures
\pysiglinewithargsret
{\sphinxbfcode{\sphinxupquote{\DUrole{k}{class}\DUrole{w}{ }}}\sphinxcode{\sphinxupquote{pygeox.custom\_objects.}}\sphinxbfcode{\sphinxupquote{RightTrapezoid}}}
{\sphinxparam{\DUrole{n}{geoscene}}\sphinxparamcomma \sphinxparam{\DUrole{n}{p1}}\sphinxparamcomma \sphinxparam{\DUrole{n}{p2}}\sphinxparamcomma \sphinxparam{\DUrole{n}{p3}}\sphinxparamcomma \sphinxparam{\DUrole{n}{p4}}\sphinxparamcomma \sphinxparam{\DUrole{n}{name}\DUrole{o}{=}\DUrole{default_value}{None}}\sphinxparamcomma \sphinxparam{\DUrole{n}{add\_to\_scene\_as\_type}\DUrole{o}{=}\DUrole{default_value}{\textquotesingle{}IsoscelesTrapezoid\textquotesingle{}}}\sphinxparamcomma \sphinxparam{\DUrole{n}{force\_convex}\DUrole{o}{=}\DUrole{default_value}{False}}}
{}
\pysigstopsignatures
\sphinxAtStartPar
Bases: \sphinxcode{\sphinxupquote{Trapezoid}}

\sphinxAtStartPar
Trapezoid with non\sphinxhyphen{}parallel sides equal in length. RightTrapezoid(A,B,C,D) means AB is parallel to CD and angle A and D are right angles.
\begin{quote}\begin{description}
\sphinxlineitem{Parameters}
\sphinxAtStartPar
\sphinxstyleliteralstrong{\sphinxupquote{force\_convex}} (\sphinxstyleliteralemphasis{\sphinxupquote{bool}})

\end{description}\end{quote}

\end{fulllineitems}

\index{RightTriangle (class in pygeox.custom\_objects)@\spxentry{RightTriangle}\spxextra{class in pygeox.custom\_objects}}

\begin{fulllineitems}
\phantomsection\label{\detokenize{pygeox:pygeox.custom_objects.RightTriangle}}
\pysigstartsignatures
\pysiglinewithargsret
{\sphinxbfcode{\sphinxupquote{\DUrole{k}{class}\DUrole{w}{ }}}\sphinxcode{\sphinxupquote{pygeox.custom\_objects.}}\sphinxbfcode{\sphinxupquote{RightTriangle}}}
{\sphinxparam{\DUrole{n}{geoscene}}\sphinxparamcomma \sphinxparam{\DUrole{n}{p1}}\sphinxparamcomma \sphinxparam{\DUrole{n}{p2}}\sphinxparamcomma \sphinxparam{\DUrole{n}{p3}}\sphinxparamcomma \sphinxparam{\DUrole{n}{name}\DUrole{o}{=}\DUrole{default_value}{None}}\sphinxparamcomma \sphinxparam{\DUrole{n}{add\_to\_scene\_as\_type}\DUrole{o}{=}\DUrole{default_value}{\textquotesingle{}RightTriangle\textquotesingle{}}}}
{}
\pysigstopsignatures
\sphinxAtStartPar
Bases: \sphinxcode{\sphinxupquote{Triangle}}

\sphinxAtStartPar
Triangle with one right angle. RightTriangle(A,B,C) means that B is the right angle.
\begin{quote}\begin{description}
\sphinxlineitem{Parameters}\begin{itemize}
\item {} 
\sphinxAtStartPar
\sphinxstyleliteralstrong{\sphinxupquote{p1}} (\sphinxstyleliteralemphasis{\sphinxupquote{Point}})

\item {} 
\sphinxAtStartPar
\sphinxstyleliteralstrong{\sphinxupquote{p2}} (\sphinxstyleliteralemphasis{\sphinxupquote{Point}})

\item {} 
\sphinxAtStartPar
\sphinxstyleliteralstrong{\sphinxupquote{p3}} (\sphinxstyleliteralemphasis{\sphinxupquote{Point}})

\item {} 
\sphinxAtStartPar
\sphinxstyleliteralstrong{\sphinxupquote{name}} (\sphinxstyleliteralemphasis{\sphinxupquote{str}}\sphinxstyleliteralemphasis{\sphinxupquote{ | }}\sphinxstyleliteralemphasis{\sphinxupquote{None}})

\item {} 
\sphinxAtStartPar
\sphinxstyleliteralstrong{\sphinxupquote{add\_to\_scene\_as\_type}} (\sphinxstyleliteralemphasis{\sphinxupquote{str}}\sphinxstyleliteralemphasis{\sphinxupquote{ | }}\sphinxstyleliteralemphasis{\sphinxupquote{None}})

\end{itemize}

\end{description}\end{quote}

\end{fulllineitems}

\index{ScaleneTriangle (class in pygeox.custom\_objects)@\spxentry{ScaleneTriangle}\spxextra{class in pygeox.custom\_objects}}

\begin{fulllineitems}
\phantomsection\label{\detokenize{pygeox:pygeox.custom_objects.ScaleneTriangle}}
\pysigstartsignatures
\pysiglinewithargsret
{\sphinxbfcode{\sphinxupquote{\DUrole{k}{class}\DUrole{w}{ }}}\sphinxcode{\sphinxupquote{pygeox.custom\_objects.}}\sphinxbfcode{\sphinxupquote{ScaleneTriangle}}}
{\sphinxparam{\DUrole{n}{geoscene}}\sphinxparamcomma \sphinxparam{\DUrole{n}{p1}}\sphinxparamcomma \sphinxparam{\DUrole{n}{p2}}\sphinxparamcomma \sphinxparam{\DUrole{n}{p3}}\sphinxparamcomma \sphinxparam{\DUrole{n}{name}\DUrole{o}{=}\DUrole{default_value}{None}}\sphinxparamcomma \sphinxparam{\DUrole{n}{add\_to\_scene\_as\_type}\DUrole{o}{=}\DUrole{default_value}{\textquotesingle{}ScaleneTriangle\textquotesingle{}}}}
{}
\pysigstopsignatures
\sphinxAtStartPar
Bases: \sphinxcode{\sphinxupquote{Triangle}}

\sphinxAtStartPar
Triangle with all sides of different lengths.
\begin{quote}\begin{description}
\sphinxlineitem{Parameters}\begin{itemize}
\item {} 
\sphinxAtStartPar
\sphinxstyleliteralstrong{\sphinxupquote{p1}} (\sphinxstyleliteralemphasis{\sphinxupquote{Point}})

\item {} 
\sphinxAtStartPar
\sphinxstyleliteralstrong{\sphinxupquote{p2}} (\sphinxstyleliteralemphasis{\sphinxupquote{Point}})

\item {} 
\sphinxAtStartPar
\sphinxstyleliteralstrong{\sphinxupquote{p3}} (\sphinxstyleliteralemphasis{\sphinxupquote{Point}})

\item {} 
\sphinxAtStartPar
\sphinxstyleliteralstrong{\sphinxupquote{name}} (\sphinxstyleliteralemphasis{\sphinxupquote{str}}\sphinxstyleliteralemphasis{\sphinxupquote{ | }}\sphinxstyleliteralemphasis{\sphinxupquote{None}})

\item {} 
\sphinxAtStartPar
\sphinxstyleliteralstrong{\sphinxupquote{add\_to\_scene\_as\_type}} (\sphinxstyleliteralemphasis{\sphinxupquote{str}}\sphinxstyleliteralemphasis{\sphinxupquote{ | }}\sphinxstyleliteralemphasis{\sphinxupquote{None}})

\end{itemize}

\end{description}\end{quote}

\end{fulllineitems}

\index{Square (class in pygeox.custom\_objects)@\spxentry{Square}\spxextra{class in pygeox.custom\_objects}}

\begin{fulllineitems}
\phantomsection\label{\detokenize{pygeox:pygeox.custom_objects.Square}}
\pysigstartsignatures
\pysiglinewithargsret
{\sphinxbfcode{\sphinxupquote{\DUrole{k}{class}\DUrole{w}{ }}}\sphinxcode{\sphinxupquote{pygeox.custom\_objects.}}\sphinxbfcode{\sphinxupquote{Square}}}
{\sphinxparam{\DUrole{n}{geoscene}}\sphinxparamcomma \sphinxparam{\DUrole{n}{p1}}\sphinxparamcomma \sphinxparam{\DUrole{n}{p2}}\sphinxparamcomma \sphinxparam{\DUrole{n}{p3}}\sphinxparamcomma \sphinxparam{\DUrole{n}{p4}}\sphinxparamcomma \sphinxparam{\DUrole{n}{name}\DUrole{o}{=}\DUrole{default_value}{None}}\sphinxparamcomma \sphinxparam{\DUrole{n}{add\_to\_scene\_as\_type}\DUrole{o}{=}\DUrole{default_value}{\textquotesingle{}Square\textquotesingle{}}}\sphinxparamcomma \sphinxparam{\DUrole{n}{force\_convex}\DUrole{o}{=}\DUrole{default_value}{False}}}
{}
\pysigstopsignatures
\sphinxAtStartPar
Bases: \sphinxcode{\sphinxupquote{Rectangle}}

\sphinxAtStartPar
Rectangle with all sides equal.
\begin{quote}\begin{description}
\sphinxlineitem{Parameters}
\sphinxAtStartPar
\sphinxstyleliteralstrong{\sphinxupquote{force\_convex}} (\sphinxstyleliteralemphasis{\sphinxupquote{bool}})

\end{description}\end{quote}
\index{area (pygeox.custom\_objects.Square property)@\spxentry{area}\spxextra{pygeox.custom\_objects.Square property}}

\begin{fulllineitems}
\phantomsection\label{\detokenize{pygeox:pygeox.custom_objects.Square.area}}
\pysigstartsignatures
\pysigline
{\sphinxbfcode{\sphinxupquote{\DUrole{k}{property}\DUrole{w}{ }}}\sphinxbfcode{\sphinxupquote{area}}}
\pysigstopsignatures
\sphinxAtStartPar
Calculates the area of the polygon using the shoelace formula.
\begin{quote}\begin{description}
\sphinxlineitem{Returns}
\sphinxAtStartPar
The area. Units: square of input coordinates.

\sphinxlineitem{Return type}
\sphinxAtStartPar
sympy.Expr

\end{description}\end{quote}

\end{fulllineitems}

\index{circumcenter (pygeox.custom\_objects.Square property)@\spxentry{circumcenter}\spxextra{pygeox.custom\_objects.Square property}}

\begin{fulllineitems}
\phantomsection\label{\detokenize{pygeox:pygeox.custom_objects.Square.circumcenter}}
\pysigstartsignatures
\pysigline
{\sphinxbfcode{\sphinxupquote{\DUrole{k}{property}\DUrole{w}{ }}}\sphinxbfcode{\sphinxupquote{circumcenter}}}
\pysigstopsignatures
\end{fulllineitems}

\index{circumradius (pygeox.custom\_objects.Square property)@\spxentry{circumradius}\spxextra{pygeox.custom\_objects.Square property}}

\begin{fulllineitems}
\phantomsection\label{\detokenize{pygeox:pygeox.custom_objects.Square.circumradius}}
\pysigstartsignatures
\pysigline
{\sphinxbfcode{\sphinxupquote{\DUrole{k}{property}\DUrole{w}{ }}}\sphinxbfcode{\sphinxupquote{circumradius}}}
\pysigstopsignatures
\end{fulllineitems}

\index{incenter (pygeox.custom\_objects.Square property)@\spxentry{incenter}\spxextra{pygeox.custom\_objects.Square property}}

\begin{fulllineitems}
\phantomsection\label{\detokenize{pygeox:pygeox.custom_objects.Square.incenter}}
\pysigstartsignatures
\pysigline
{\sphinxbfcode{\sphinxupquote{\DUrole{k}{property}\DUrole{w}{ }}}\sphinxbfcode{\sphinxupquote{incenter}}}
\pysigstopsignatures
\end{fulllineitems}

\index{inradius (pygeox.custom\_objects.Square property)@\spxentry{inradius}\spxextra{pygeox.custom\_objects.Square property}}

\begin{fulllineitems}
\phantomsection\label{\detokenize{pygeox:pygeox.custom_objects.Square.inradius}}
\pysigstartsignatures
\pysigline
{\sphinxbfcode{\sphinxupquote{\DUrole{k}{property}\DUrole{w}{ }}}\sphinxbfcode{\sphinxupquote{inradius}}}
\pysigstopsignatures
\end{fulllineitems}

\index{perimeter (pygeox.custom\_objects.Square property)@\spxentry{perimeter}\spxextra{pygeox.custom\_objects.Square property}}

\begin{fulllineitems}
\phantomsection\label{\detokenize{pygeox:pygeox.custom_objects.Square.perimeter}}
\pysigstartsignatures
\pysigline
{\sphinxbfcode{\sphinxupquote{\DUrole{k}{property}\DUrole{w}{ }}}\sphinxbfcode{\sphinxupquote{perimeter}}}
\pysigstopsignatures
\sphinxAtStartPar
Calculates the perimeter of the polygon by summing side lengths.
\begin{quote}\begin{description}
\sphinxlineitem{Returns}
\sphinxAtStartPar
The perimeter. Units: same as input coordinates.

\sphinxlineitem{Return type}
\sphinxAtStartPar
sympy.Expr

\end{description}\end{quote}

\end{fulllineitems}


\end{fulllineitems}

\index{Trapezoid (class in pygeox.custom\_objects)@\spxentry{Trapezoid}\spxextra{class in pygeox.custom\_objects}}

\begin{fulllineitems}
\phantomsection\label{\detokenize{pygeox:pygeox.custom_objects.Trapezoid}}
\pysigstartsignatures
\pysiglinewithargsret
{\sphinxbfcode{\sphinxupquote{\DUrole{k}{class}\DUrole{w}{ }}}\sphinxcode{\sphinxupquote{pygeox.custom\_objects.}}\sphinxbfcode{\sphinxupquote{Trapezoid}}}
{\sphinxparam{\DUrole{n}{geoscene}}\sphinxparamcomma \sphinxparam{\DUrole{n}{p1}}\sphinxparamcomma \sphinxparam{\DUrole{n}{p2}}\sphinxparamcomma \sphinxparam{\DUrole{n}{p3}}\sphinxparamcomma \sphinxparam{\DUrole{n}{p4}}\sphinxparamcomma \sphinxparam{\DUrole{n}{name}\DUrole{o}{=}\DUrole{default_value}{None}}\sphinxparamcomma \sphinxparam{\DUrole{n}{add\_to\_scene\_as\_type}\DUrole{o}{=}\DUrole{default_value}{\textquotesingle{}Trapezoid\textquotesingle{}}}\sphinxparamcomma \sphinxparam{\DUrole{n}{force\_convex}\DUrole{o}{=}\DUrole{default_value}{False}}}
{}
\pysigstopsignatures
\sphinxAtStartPar
Bases: \sphinxcode{\sphinxupquote{Quadrilateral}}

\sphinxAtStartPar
Quadrilateral with at least one pair of parallel sides. Trapezoid(A,B,C,D) means AB is parallel to CD.
\begin{quote}\begin{description}
\sphinxlineitem{Parameters}
\sphinxAtStartPar
\sphinxstyleliteralstrong{\sphinxupquote{force\_convex}} (\sphinxstyleliteralemphasis{\sphinxupquote{bool}})

\end{description}\end{quote}

\end{fulllineitems}

\index{Triangle (class in pygeox.custom\_objects)@\spxentry{Triangle}\spxextra{class in pygeox.custom\_objects}}

\begin{fulllineitems}
\phantomsection\label{\detokenize{pygeox:pygeox.custom_objects.Triangle}}
\pysigstartsignatures
\pysiglinewithargsret
{\sphinxbfcode{\sphinxupquote{\DUrole{k}{class}\DUrole{w}{ }}}\sphinxcode{\sphinxupquote{pygeox.custom\_objects.}}\sphinxbfcode{\sphinxupquote{Triangle}}}
{\sphinxparam{\DUrole{n}{geoscene}}\sphinxparamcomma \sphinxparam{\DUrole{n}{p1}}\sphinxparamcomma \sphinxparam{\DUrole{n}{p2}}\sphinxparamcomma \sphinxparam{\DUrole{n}{p3}}\sphinxparamcomma \sphinxparam{\DUrole{n}{name}\DUrole{o}{=}\DUrole{default_value}{None}}\sphinxparamcomma \sphinxparam{\DUrole{n}{add\_to\_scene\_as\_type}\DUrole{o}{=}\DUrole{default_value}{\textquotesingle{}Triangle\textquotesingle{}}}}
{}
\pysigstopsignatures
\sphinxAtStartPar
Bases: \sphinxcode{\sphinxupquote{Polygon}}

\sphinxAtStartPar
Represents a general triangle defined by three Point objects.
\begin{quote}\begin{description}
\sphinxlineitem{Parameters}\begin{itemize}
\item {} 
\sphinxAtStartPar
\sphinxstyleliteralstrong{\sphinxupquote{p1}} (\sphinxstyleliteralemphasis{\sphinxupquote{Point}})

\item {} 
\sphinxAtStartPar
\sphinxstyleliteralstrong{\sphinxupquote{p2}} (\sphinxstyleliteralemphasis{\sphinxupquote{Point}})

\item {} 
\sphinxAtStartPar
\sphinxstyleliteralstrong{\sphinxupquote{p3}} (\sphinxstyleliteralemphasis{\sphinxupquote{Point}})

\item {} 
\sphinxAtStartPar
\sphinxstyleliteralstrong{\sphinxupquote{name}} (\sphinxstyleliteralemphasis{\sphinxupquote{str}}\sphinxstyleliteralemphasis{\sphinxupquote{ | }}\sphinxstyleliteralemphasis{\sphinxupquote{None}})

\item {} 
\sphinxAtStartPar
\sphinxstyleliteralstrong{\sphinxupquote{add\_to\_scene\_as\_type}} (\sphinxstyleliteralemphasis{\sphinxupquote{str}}\sphinxstyleliteralemphasis{\sphinxupquote{ | }}\sphinxstyleliteralemphasis{\sphinxupquote{None}})

\end{itemize}

\end{description}\end{quote}
\index{altitudes (pygeox.custom\_objects.Triangle property)@\spxentry{altitudes}\spxextra{pygeox.custom\_objects.Triangle property}}

\begin{fulllineitems}
\phantomsection\label{\detokenize{pygeox:pygeox.custom_objects.Triangle.altitudes}}
\pysigstartsignatures
\pysigline
{\sphinxbfcode{\sphinxupquote{\DUrole{k}{property}\DUrole{w}{ }}}\sphinxbfcode{\sphinxupquote{altitudes}}\sphinxbfcode{\sphinxupquote{\DUrole{p}{:}\DUrole{w}{ }list}}}
\pysigstopsignatures
\sphinxAtStartPar
Returns:
list{[}LineSegment{]}: The three altitudes of the triangle, each from a vertex perpendicular to the opposite side.

\end{fulllineitems}

\index{angle\_bisectors (pygeox.custom\_objects.Triangle property)@\spxentry{angle\_bisectors}\spxextra{pygeox.custom\_objects.Triangle property}}

\begin{fulllineitems}
\phantomsection\label{\detokenize{pygeox:pygeox.custom_objects.Triangle.angle_bisectors}}
\pysigstartsignatures
\pysigline
{\sphinxbfcode{\sphinxupquote{\DUrole{k}{property}\DUrole{w}{ }}}\sphinxbfcode{\sphinxupquote{angle\_bisectors}}\sphinxbfcode{\sphinxupquote{\DUrole{p}{:}\DUrole{w}{ }list}}}
\pysigstopsignatures
\sphinxAtStartPar
Returns:
list{[}LineSegment{]}: The three angle bisectors of the triangle, each bisecting one angle and ending at the opposite side.

\end{fulllineitems}

\index{area (pygeox.custom\_objects.Triangle property)@\spxentry{area}\spxextra{pygeox.custom\_objects.Triangle property}}

\begin{fulllineitems}
\phantomsection\label{\detokenize{pygeox:pygeox.custom_objects.Triangle.area}}
\pysigstartsignatures
\pysigline
{\sphinxbfcode{\sphinxupquote{\DUrole{k}{property}\DUrole{w}{ }}}\sphinxbfcode{\sphinxupquote{area}}}
\pysigstopsignatures
\sphinxAtStartPar
Calculates the area of the triangle using Heron’s formula.

\end{fulllineitems}

\index{circumcenter (pygeox.custom\_objects.Triangle property)@\spxentry{circumcenter}\spxextra{pygeox.custom\_objects.Triangle property}}

\begin{fulllineitems}
\phantomsection\label{\detokenize{pygeox:pygeox.custom_objects.Triangle.circumcenter}}
\pysigstartsignatures
\pysigline
{\sphinxbfcode{\sphinxupquote{\DUrole{k}{property}\DUrole{w}{ }}}\sphinxbfcode{\sphinxupquote{circumcenter}}}
\pysigstopsignatures
\end{fulllineitems}

\index{circumradius (pygeox.custom\_objects.Triangle property)@\spxentry{circumradius}\spxextra{pygeox.custom\_objects.Triangle property}}

\begin{fulllineitems}
\phantomsection\label{\detokenize{pygeox:pygeox.custom_objects.Triangle.circumradius}}
\pysigstartsignatures
\pysigline
{\sphinxbfcode{\sphinxupquote{\DUrole{k}{property}\DUrole{w}{ }}}\sphinxbfcode{\sphinxupquote{circumradius}}}
\pysigstopsignatures
\end{fulllineitems}

\index{heights (pygeox.custom\_objects.Triangle property)@\spxentry{heights}\spxextra{pygeox.custom\_objects.Triangle property}}

\begin{fulllineitems}
\phantomsection\label{\detokenize{pygeox:pygeox.custom_objects.Triangle.heights}}
\pysigstartsignatures
\pysigline
{\sphinxbfcode{\sphinxupquote{\DUrole{k}{property}\DUrole{w}{ }}}\sphinxbfcode{\sphinxupquote{heights}}}
\pysigstopsignatures
\sphinxAtStartPar
Returns the heights from each vertex to the opposite side.

\end{fulllineitems}

\index{incenter (pygeox.custom\_objects.Triangle property)@\spxentry{incenter}\spxextra{pygeox.custom\_objects.Triangle property}}

\begin{fulllineitems}
\phantomsection\label{\detokenize{pygeox:pygeox.custom_objects.Triangle.incenter}}
\pysigstartsignatures
\pysigline
{\sphinxbfcode{\sphinxupquote{\DUrole{k}{property}\DUrole{w}{ }}}\sphinxbfcode{\sphinxupquote{incenter}}}
\pysigstopsignatures
\end{fulllineitems}

\index{inradius (pygeox.custom\_objects.Triangle property)@\spxentry{inradius}\spxextra{pygeox.custom\_objects.Triangle property}}

\begin{fulllineitems}
\phantomsection\label{\detokenize{pygeox:pygeox.custom_objects.Triangle.inradius}}
\pysigstartsignatures
\pysigline
{\sphinxbfcode{\sphinxupquote{\DUrole{k}{property}\DUrole{w}{ }}}\sphinxbfcode{\sphinxupquote{inradius}}}
\pysigstopsignatures
\end{fulllineitems}

\index{medians (pygeox.custom\_objects.Triangle property)@\spxentry{medians}\spxextra{pygeox.custom\_objects.Triangle property}}

\begin{fulllineitems}
\phantomsection\label{\detokenize{pygeox:pygeox.custom_objects.Triangle.medians}}
\pysigstartsignatures
\pysigline
{\sphinxbfcode{\sphinxupquote{\DUrole{k}{property}\DUrole{w}{ }}}\sphinxbfcode{\sphinxupquote{medians}}\sphinxbfcode{\sphinxupquote{\DUrole{p}{:}\DUrole{w}{ }list}}}
\pysigstopsignatures
\sphinxAtStartPar
Returns:
list{[}LineSegment{]}: The three medians of the triangle, each from a vertex to the midpoint of the opposite side.

\end{fulllineitems}

\index{midsegments (pygeox.custom\_objects.Triangle property)@\spxentry{midsegments}\spxextra{pygeox.custom\_objects.Triangle property}}

\begin{fulllineitems}
\phantomsection\label{\detokenize{pygeox:pygeox.custom_objects.Triangle.midsegments}}
\pysigstartsignatures
\pysigline
{\sphinxbfcode{\sphinxupquote{\DUrole{k}{property}\DUrole{w}{ }}}\sphinxbfcode{\sphinxupquote{midsegments}}\sphinxbfcode{\sphinxupquote{\DUrole{p}{:}\DUrole{w}{ }list}}}
\pysigstopsignatures
\sphinxAtStartPar
Returns:
list{[}LineSegment{]}: The three midsegments of the triangle, each connecting the midpoints of two sides.

\end{fulllineitems}

\index{orthocenter (pygeox.custom\_objects.Triangle property)@\spxentry{orthocenter}\spxextra{pygeox.custom\_objects.Triangle property}}

\begin{fulllineitems}
\phantomsection\label{\detokenize{pygeox:pygeox.custom_objects.Triangle.orthocenter}}
\pysigstartsignatures
\pysigline
{\sphinxbfcode{\sphinxupquote{\DUrole{k}{property}\DUrole{w}{ }}}\sphinxbfcode{\sphinxupquote{orthocenter}}}
\pysigstopsignatures
\end{fulllineitems}


\end{fulllineitems}



\chapter{Animation (Under construction)}
\label{\detokenize{pygeox:animation-under-construction}}\index{module@\spxentry{module}!pygeox.animate@\spxentry{pygeox.animate}}\index{pygeox.animate@\spxentry{pygeox.animate}!module@\spxentry{module}}\index{AnimateNamespace (class in pygeox.animate)@\spxentry{AnimateNamespace}\spxextra{class in pygeox.animate}}\phantomsection\label{\detokenize{pygeox:module-pygeox.animate}}

\begin{fulllineitems}
\phantomsection\label{\detokenize{pygeox:pygeox.animate.AnimateNamespace}}
\pysigstartsignatures
\pysiglinewithargsret
{\sphinxbfcode{\sphinxupquote{\DUrole{k}{class}\DUrole{w}{ }}}\sphinxcode{\sphinxupquote{pygeox.animate.}}\sphinxbfcode{\sphinxupquote{AnimateNamespace}}}
{\sphinxparam{\DUrole{n}{scene}}}
{}
\pysigstopsignatures
\sphinxAtStartPar
Bases: \sphinxcode{\sphinxupquote{object}}

\sphinxAtStartPar
Namespace to manage animation parameters and generate animations
by varying parameters within a geometric scene.

\sphinxAtStartPar
Provides functionality to add animation parameters, run animations
by solving constraints at discrete parameter values, and export
animations as GIFs.
\index{add\_parameter() (pygeox.animate.AnimateNamespace method)@\spxentry{add\_parameter()}\spxextra{pygeox.animate.AnimateNamespace method}}

\begin{fulllineitems}
\phantomsection\label{\detokenize{pygeox:pygeox.animate.AnimateNamespace.add_parameter}}
\pysigstartsignatures
\pysiglinewithargsret
{\sphinxbfcode{\sphinxupquote{add\_parameter}}}
{\sphinxparam{\DUrole{n}{name}}\sphinxparamcomma \sphinxparam{\DUrole{n}{total\_time}\DUrole{o}{=}\DUrole{default_value}{1}}}
{}
\pysigstopsignatures
\sphinxAtStartPar
Add a new symbolic animation parameter to the scene.
\begin{quote}\begin{description}
\sphinxlineitem{Parameters}\begin{itemize}
\item {} 
\sphinxAtStartPar
\sphinxstyleliteralstrong{\sphinxupquote{name}} (\sphinxstyleliteralemphasis{\sphinxupquote{str}}) \textendash{} The name of the animation parameter.

\item {} 
\sphinxAtStartPar
\sphinxstyleliteralstrong{\sphinxupquote{total\_time}} (\sphinxstyleliteralemphasis{\sphinxupquote{int}}\sphinxstyleliteralemphasis{\sphinxupquote{ | }}\sphinxstyleliteralemphasis{\sphinxupquote{float}}\sphinxstyleliteralemphasis{\sphinxupquote{ | }}\sphinxstyleliteralemphasis{\sphinxupquote{Expr}}\sphinxstyleliteralemphasis{\sphinxupquote{ | }}\sphinxstyleliteralemphasis{\sphinxupquote{Symbol}})

\end{itemize}

\sphinxlineitem{Returns}
\sphinxAtStartPar
The sympy.Symbol representing the animation parameter.

\end{description}\end{quote}

\end{fulllineitems}

\index{eval\_piecewise\_at\_zero() (pygeox.animate.AnimateNamespace method)@\spxentry{eval\_piecewise\_at\_zero()}\spxextra{pygeox.animate.AnimateNamespace method}}

\begin{fulllineitems}
\phantomsection\label{\detokenize{pygeox:pygeox.animate.AnimateNamespace.eval_piecewise_at_zero}}
\pysigstartsignatures
\pysiglinewithargsret
{\sphinxbfcode{\sphinxupquote{eval\_piecewise\_at\_zero}}}
{\sphinxparam{\DUrole{n}{expr}}\sphinxparamcomma \sphinxparam{\DUrole{n}{anim\_syms}}}
{}
\pysigstopsignatures
\sphinxAtStartPar
If \sphinxtitleref{expr} contains Piecewise, pick the first branch whose condition is:
\sphinxhyphen{} literally True, or
\sphinxhyphen{} an And(…, alpha \textgreater{}= 0, …) (ignoring other bounds), or
\sphinxhyphen{} becomes True after substituting alpha \sphinxhyphen{}\textgreater{} 0.
Otherwise (no Piecewise), return expr unchanged.

\end{fulllineitems}

\index{get\_time\_params() (pygeox.animate.AnimateNamespace method)@\spxentry{get\_time\_params()}\spxextra{pygeox.animate.AnimateNamespace method}}

\begin{fulllineitems}
\phantomsection\label{\detokenize{pygeox:pygeox.animate.AnimateNamespace.get_time_params}}
\pysigstartsignatures
\pysiglinewithargsret
{\sphinxbfcode{\sphinxupquote{get\_time\_params}}}
{}
{}
\pysigstopsignatures
\sphinxAtStartPar
Retrieve all animation parameters currently registered in the scene.
\begin{quote}\begin{description}
\sphinxlineitem{Returns}
\sphinxAtStartPar
List of sympy.Symbol animation parameters.

\end{description}\end{quote}

\end{fulllineitems}

\index{gif() (pygeox.animate.AnimateNamespace method)@\spxentry{gif()}\spxextra{pygeox.animate.AnimateNamespace method}}

\begin{fulllineitems}
\phantomsection\label{\detokenize{pygeox:pygeox.animate.AnimateNamespace.gif}}
\pysigstartsignatures
\pysiglinewithargsret
{\sphinxbfcode{\sphinxupquote{gif}}}
{\sphinxparam{\DUrole{n}{save\_loc}\DUrole{o}{=}\DUrole{default_value}{\textquotesingle{}animation.gif\textquotesingle{}}}}
{}
\pysigstopsignatures
\sphinxAtStartPar
Generate and save an animated GIF from the current list of solved scenes.
\begin{quote}\begin{description}
\sphinxlineitem{Parameters}\begin{itemize}
\item {} 
\sphinxAtStartPar
\sphinxstyleliteralstrong{\sphinxupquote{save\_loc}} (\sphinxstyleliteralemphasis{\sphinxupquote{str}}) \textendash{} File path to save the GIF.

\item {} 
\sphinxAtStartPar
\sphinxstyleliteralstrong{\sphinxupquote{total\_time}} \textendash{} Total duration of the GIF in seconds.

\end{itemize}

\sphinxlineitem{Raises}
\sphinxAtStartPar
\sphinxstyleliteralstrong{\sphinxupquote{ValueError}} \textendash{} If there are no frames to generate.

\sphinxlineitem{Return type}
\sphinxAtStartPar
None

\end{description}\end{quote}

\end{fulllineitems}

\index{point\_sweep\_arc() (pygeox.animate.AnimateNamespace method)@\spxentry{point\_sweep\_arc()}\spxextra{pygeox.animate.AnimateNamespace method}}

\begin{fulllineitems}
\phantomsection\label{\detokenize{pygeox:pygeox.animate.AnimateNamespace.point_sweep_arc}}
\pysigstartsignatures
\pysiglinewithargsret
{\sphinxbfcode{\sphinxupquote{point\_sweep\_arc}}}
{\sphinxparam{\DUrole{n}{point}}\sphinxparamcomma \sphinxparam{\DUrole{n}{arc}}\sphinxparamcomma \sphinxparam{\DUrole{n}{counterclockwise}\DUrole{o}{=}\DUrole{default_value}{True}}\sphinxparamcomma \sphinxparam{\DUrole{n}{total\_time}\DUrole{o}{=}\DUrole{default_value}{2}}\sphinxparamcomma \sphinxparam{\DUrole{n}{speed}\DUrole{o}{=}\DUrole{default_value}{None}}}
{}
\pysigstopsignatures
\sphinxAtStartPar
Animate a point moving along an arc by directly constraining its x and y coordinates.
Point moves from arc.start\_point to arc.end\_point over the arc length.
\begin{quote}\begin{description}
\sphinxlineitem{Parameters}\begin{itemize}
\item {} 
\sphinxAtStartPar
\sphinxstyleliteralstrong{\sphinxupquote{point}} (\sphinxstyleliteralemphasis{\sphinxupquote{Point}})

\item {} 
\sphinxAtStartPar
\sphinxstyleliteralstrong{\sphinxupquote{arc}} (\sphinxstyleliteralemphasis{\sphinxupquote{MinorArc}}\sphinxstyleliteralemphasis{\sphinxupquote{ | }}\sphinxstyleliteralemphasis{\sphinxupquote{MajorArc}})

\item {} 
\sphinxAtStartPar
\sphinxstyleliteralstrong{\sphinxupquote{counterclockwise}} (\sphinxstyleliteralemphasis{\sphinxupquote{bool}})

\item {} 
\sphinxAtStartPar
\sphinxstyleliteralstrong{\sphinxupquote{total\_time}} (\sphinxstyleliteralemphasis{\sphinxupquote{int}}\sphinxstyleliteralemphasis{\sphinxupquote{ | }}\sphinxstyleliteralemphasis{\sphinxupquote{float}}\sphinxstyleliteralemphasis{\sphinxupquote{ | }}\sphinxstyleliteralemphasis{\sphinxupquote{Expr}}\sphinxstyleliteralemphasis{\sphinxupquote{ | }}\sphinxstyleliteralemphasis{\sphinxupquote{Symbol}}\sphinxstyleliteralemphasis{\sphinxupquote{ | }}\sphinxstyleliteralemphasis{\sphinxupquote{None}})

\item {} 
\sphinxAtStartPar
\sphinxstyleliteralstrong{\sphinxupquote{speed}} (\sphinxstyleliteralemphasis{\sphinxupquote{int}}\sphinxstyleliteralemphasis{\sphinxupquote{ | }}\sphinxstyleliteralemphasis{\sphinxupquote{float}}\sphinxstyleliteralemphasis{\sphinxupquote{ | }}\sphinxstyleliteralemphasis{\sphinxupquote{Expr}}\sphinxstyleliteralemphasis{\sphinxupquote{ | }}\sphinxstyleliteralemphasis{\sphinxupquote{Symbol}}\sphinxstyleliteralemphasis{\sphinxupquote{ | }}\sphinxstyleliteralemphasis{\sphinxupquote{None}})

\end{itemize}

\end{description}\end{quote}

\end{fulllineitems}

\index{point\_sweep\_circle() (pygeox.animate.AnimateNamespace method)@\spxentry{point\_sweep\_circle()}\spxextra{pygeox.animate.AnimateNamespace method}}

\begin{fulllineitems}
\phantomsection\label{\detokenize{pygeox:pygeox.animate.AnimateNamespace.point_sweep_circle}}
\pysigstartsignatures
\pysiglinewithargsret
{\sphinxbfcode{\sphinxupquote{point\_sweep\_circle}}}
{\sphinxparam{\DUrole{n}{point}}\sphinxparamcomma \sphinxparam{\DUrole{n}{circle}}\sphinxparamcomma \sphinxparam{\DUrole{n}{counterclockwise}\DUrole{o}{=}\DUrole{default_value}{True}}\sphinxparamcomma \sphinxparam{\DUrole{n}{total\_time}\DUrole{o}{=}\DUrole{default_value}{2}}\sphinxparamcomma \sphinxparam{\DUrole{n}{speed}\DUrole{o}{=}\DUrole{default_value}{None}}}
{}
\pysigstopsignatures
\sphinxAtStartPar
Animate a point moving around a circle by directly constraining its x and y coordinates.
Point starts at the top of the circle and moves counterclockwise by default.
\begin{quote}\begin{description}
\sphinxlineitem{Parameters}\begin{itemize}
\item {} 
\sphinxAtStartPar
\sphinxstyleliteralstrong{\sphinxupquote{point}} (\sphinxstyleliteralemphasis{\sphinxupquote{Point}})

\item {} 
\sphinxAtStartPar
\sphinxstyleliteralstrong{\sphinxupquote{circle}} (\sphinxstyleliteralemphasis{\sphinxupquote{Circle}})

\item {} 
\sphinxAtStartPar
\sphinxstyleliteralstrong{\sphinxupquote{counterclockwise}} (\sphinxstyleliteralemphasis{\sphinxupquote{bool}})

\item {} 
\sphinxAtStartPar
\sphinxstyleliteralstrong{\sphinxupquote{total\_time}} (\sphinxstyleliteralemphasis{\sphinxupquote{int}}\sphinxstyleliteralemphasis{\sphinxupquote{ | }}\sphinxstyleliteralemphasis{\sphinxupquote{float}}\sphinxstyleliteralemphasis{\sphinxupquote{ | }}\sphinxstyleliteralemphasis{\sphinxupquote{Expr}}\sphinxstyleliteralemphasis{\sphinxupquote{ | }}\sphinxstyleliteralemphasis{\sphinxupquote{Symbol}}\sphinxstyleliteralemphasis{\sphinxupquote{ | }}\sphinxstyleliteralemphasis{\sphinxupquote{None}})

\item {} 
\sphinxAtStartPar
\sphinxstyleliteralstrong{\sphinxupquote{speed}} (\sphinxstyleliteralemphasis{\sphinxupquote{int}}\sphinxstyleliteralemphasis{\sphinxupquote{ | }}\sphinxstyleliteralemphasis{\sphinxupquote{float}}\sphinxstyleliteralemphasis{\sphinxupquote{ | }}\sphinxstyleliteralemphasis{\sphinxupquote{Expr}}\sphinxstyleliteralemphasis{\sphinxupquote{ | }}\sphinxstyleliteralemphasis{\sphinxupquote{Symbol}}\sphinxstyleliteralemphasis{\sphinxupquote{ | }}\sphinxstyleliteralemphasis{\sphinxupquote{None}})

\end{itemize}

\end{description}\end{quote}

\end{fulllineitems}

\index{point\_sweep\_line() (pygeox.animate.AnimateNamespace method)@\spxentry{point\_sweep\_line()}\spxextra{pygeox.animate.AnimateNamespace method}}

\begin{fulllineitems}
\phantomsection\label{\detokenize{pygeox:pygeox.animate.AnimateNamespace.point_sweep_line}}
\pysigstartsignatures
\pysiglinewithargsret
{\sphinxbfcode{\sphinxupquote{point\_sweep\_line}}}
{\sphinxparam{\DUrole{n}{point}}\sphinxparamcomma \sphinxparam{\DUrole{n}{start\_point}}\sphinxparamcomma \sphinxparam{\DUrole{n}{end\_point}}\sphinxparamcomma \sphinxparam{\DUrole{n}{total\_time}\DUrole{o}{=}\DUrole{default_value}{2}}\sphinxparamcomma \sphinxparam{\DUrole{n}{speed}\DUrole{o}{=}\DUrole{default_value}{None}}}
{}
\pysigstopsignatures
\sphinxAtStartPar
Animate a point moving along a line segment by directly setting its x and y coordinates parametrically.
\begin{quote}\begin{description}
\sphinxlineitem{Parameters}\begin{itemize}
\item {} 
\sphinxAtStartPar
\sphinxstyleliteralstrong{\sphinxupquote{point}} (\sphinxstyleliteralemphasis{\sphinxupquote{Point}})

\item {} 
\sphinxAtStartPar
\sphinxstyleliteralstrong{\sphinxupquote{start\_point}} (\sphinxstyleliteralemphasis{\sphinxupquote{Point}})

\item {} 
\sphinxAtStartPar
\sphinxstyleliteralstrong{\sphinxupquote{end\_point}} (\sphinxstyleliteralemphasis{\sphinxupquote{Point}})

\item {} 
\sphinxAtStartPar
\sphinxstyleliteralstrong{\sphinxupquote{total\_time}} (\sphinxstyleliteralemphasis{\sphinxupquote{int}}\sphinxstyleliteralemphasis{\sphinxupquote{ | }}\sphinxstyleliteralemphasis{\sphinxupquote{float}}\sphinxstyleliteralemphasis{\sphinxupquote{ | }}\sphinxstyleliteralemphasis{\sphinxupquote{Expr}}\sphinxstyleliteralemphasis{\sphinxupquote{ | }}\sphinxstyleliteralemphasis{\sphinxupquote{Symbol}}\sphinxstyleliteralemphasis{\sphinxupquote{ | }}\sphinxstyleliteralemphasis{\sphinxupquote{None}})

\item {} 
\sphinxAtStartPar
\sphinxstyleliteralstrong{\sphinxupquote{speed}} (\sphinxstyleliteralemphasis{\sphinxupquote{int}}\sphinxstyleliteralemphasis{\sphinxupquote{ | }}\sphinxstyleliteralemphasis{\sphinxupquote{float}}\sphinxstyleliteralemphasis{\sphinxupquote{ | }}\sphinxstyleliteralemphasis{\sphinxupquote{Expr}}\sphinxstyleliteralemphasis{\sphinxupquote{ | }}\sphinxstyleliteralemphasis{\sphinxupquote{Symbol}}\sphinxstyleliteralemphasis{\sphinxupquote{ | }}\sphinxstyleliteralemphasis{\sphinxupquote{None}})

\end{itemize}

\end{description}\end{quote}

\end{fulllineitems}

\index{point\_sweep\_polyline() (pygeox.animate.AnimateNamespace method)@\spxentry{point\_sweep\_polyline()}\spxextra{pygeox.animate.AnimateNamespace method}}

\begin{fulllineitems}
\phantomsection\label{\detokenize{pygeox:pygeox.animate.AnimateNamespace.point_sweep_polyline}}
\pysigstartsignatures
\pysiglinewithargsret
{\sphinxbfcode{\sphinxupquote{point\_sweep\_polyline}}}
{\sphinxparam{\DUrole{n}{point}}\sphinxparamcomma \sphinxparam{\DUrole{n}{path\_points}}\sphinxparamcomma \sphinxparam{\DUrole{n}{total\_time}\DUrole{o}{=}\DUrole{default_value}{5}}\sphinxparamcomma \sphinxparam{\DUrole{n}{speed}\DUrole{o}{=}\DUrole{default_value}{None}}}
{}
\pysigstopsignatures
\sphinxAtStartPar
Animates a point moving smoothly along a polyline defined by a sequence of points.

\sphinxAtStartPar
A single parameter \sphinxtitleref{time\_param} in {[}0, 1{]} controls the point’s position.
As \sphinxtitleref{time\_param} increases from 0 to 1, the point traverses each segment of the polyline
in proportion to segment length. A single sympy.Piecewise constraint ensures that the point’s
coordinates match the appropriate interpolated position along the active segment.
\begin{quote}\begin{description}
\sphinxlineitem{Parameters}\begin{itemize}
\item {} 
\sphinxAtStartPar
\sphinxstyleliteralstrong{\sphinxupquote{point}} (\sphinxstyleliteralemphasis{\sphinxupquote{Point}}) \textendash{} The Point object to animate.

\item {} 
\sphinxAtStartPar
\sphinxstyleliteralstrong{\sphinxupquote{path\_points}} (\sphinxstyleliteralemphasis{\sphinxupquote{List}}\sphinxstyleliteralemphasis{\sphinxupquote{{[}}}\sphinxstyleliteralemphasis{\sphinxupquote{Point}}\sphinxstyleliteralemphasis{\sphinxupquote{{]}}}) \textendash{} Ordered list of points defining the polyline path.

\item {} 
\sphinxAtStartPar
\sphinxstyleliteralstrong{\sphinxupquote{total\_time}} (\sphinxstyleliteralemphasis{\sphinxupquote{int}}\sphinxstyleliteralemphasis{\sphinxupquote{ | }}\sphinxstyleliteralemphasis{\sphinxupquote{float}}\sphinxstyleliteralemphasis{\sphinxupquote{ | }}\sphinxstyleliteralemphasis{\sphinxupquote{Expr}}\sphinxstyleliteralemphasis{\sphinxupquote{ | }}\sphinxstyleliteralemphasis{\sphinxupquote{Symbol}}\sphinxstyleliteralemphasis{\sphinxupquote{ | }}\sphinxstyleliteralemphasis{\sphinxupquote{None}})

\item {} 
\sphinxAtStartPar
\sphinxstyleliteralstrong{\sphinxupquote{speed}} (\sphinxstyleliteralemphasis{\sphinxupquote{List}}\sphinxstyleliteralemphasis{\sphinxupquote{{[}}}\sphinxstyleliteralemphasis{\sphinxupquote{int}}\sphinxstyleliteralemphasis{\sphinxupquote{ | }}\sphinxstyleliteralemphasis{\sphinxupquote{float}}\sphinxstyleliteralemphasis{\sphinxupquote{ | }}\sphinxstyleliteralemphasis{\sphinxupquote{Expr}}\sphinxstyleliteralemphasis{\sphinxupquote{ | }}\sphinxstyleliteralemphasis{\sphinxupquote{Symbol}}\sphinxstyleliteralemphasis{\sphinxupquote{{]} }}\sphinxstyleliteralemphasis{\sphinxupquote{| }}\sphinxstyleliteralemphasis{\sphinxupquote{None}})

\end{itemize}

\end{description}\end{quote}

\end{fulllineitems}

\index{run() (pygeox.animate.AnimateNamespace method)@\spxentry{run()}\spxextra{pygeox.animate.AnimateNamespace method}}

\begin{fulllineitems}
\phantomsection\label{\detokenize{pygeox:pygeox.animate.AnimateNamespace.run}}
\pysigstartsignatures
\pysiglinewithargsret
{\sphinxbfcode{\sphinxupquote{run}}}
{\sphinxparam{\DUrole{n}{moving\_points}}\sphinxparamcomma \sphinxparam{\DUrole{n}{frames}\DUrole{o}{=}\DUrole{default_value}{20}}\sphinxparamcomma \sphinxparam{\DUrole{n}{opt\_method}\DUrole{o}{=}\DUrole{default_value}{\textquotesingle{}dual\_annealing\textquotesingle{}}}\sphinxparamcomma \sphinxparam{\DUrole{n}{distance\_penalty}\DUrole{o}{=}\DUrole{default_value}{1}}\sphinxparamcomma \sphinxparam{\DUrole{n}{min\_dist}\DUrole{o}{=}\DUrole{default_value}{None}}\sphinxparamcomma \sphinxparam{\DUrole{n}{run\_type}\DUrole{o}{=}\DUrole{default_value}{\textquotesingle{}independent\textquotesingle{}}}}
{}
\pysigstopsignatures
\sphinxAtStartPar
Generate a sequence of solved scenes by interpolating animation parameters
from 0 to 1, effectively creating frames of the animation.

\sphinxAtStartPar
Points not listed in \sphinxtitleref{moving\_points} are fixed to their initial positions.
\begin{quote}\begin{description}
\sphinxlineitem{Parameters}\begin{itemize}
\item {} 
\sphinxAtStartPar
\sphinxstyleliteralstrong{\sphinxupquote{moving\_points}} (\sphinxstyleliteralemphasis{\sphinxupquote{List}}\sphinxstyleliteralemphasis{\sphinxupquote{{[}}}\sphinxstyleliteralemphasis{\sphinxupquote{Point}}\sphinxstyleliteralemphasis{\sphinxupquote{{]}}}) \textendash{} A list of Point objects that are allowed to move during animation.
Points not in this list will remain fixed.

\item {} 
\sphinxAtStartPar
\sphinxstyleliteralstrong{\sphinxupquote{frames}} (\sphinxstyleliteralemphasis{\sphinxupquote{int}}) \textendash{} The number of frames to generate in the animation.

\item {} 
\sphinxAtStartPar
\sphinxstyleliteralstrong{\sphinxupquote{opt\_method}} (\sphinxstyleliteralemphasis{\sphinxupquote{str}}) \textendash{} The numerical solver method to use for constraint solving.
Refer to \sphinxtitleref{scipy.optimize.minimize} for available options.

\item {} 
\sphinxAtStartPar
\sphinxstyleliteralstrong{\sphinxupquote{distance\_penalty}} (\sphinxstyleliteralemphasis{\sphinxupquote{float}}\sphinxstyleliteralemphasis{\sphinxupquote{ | }}\sphinxstyleliteralemphasis{\sphinxupquote{None}}) \textendash{} The penalty weight to avoid degenerate cases related to
area and point distances. This is used in the first frame
and as a penalty constraint. It is removed from the moving points.

\item {} 
\sphinxAtStartPar
\sphinxstyleliteralstrong{\sphinxupquote{min\_dist}} (\sphinxstyleliteralemphasis{\sphinxupquote{float}}\sphinxstyleliteralemphasis{\sphinxupquote{ | }}\sphinxstyleliteralemphasis{\sphinxupquote{None}}) \textendash{} The minimum allowed distance between points. This is only used for
the first frame.

\item {} 
\sphinxAtStartPar
\sphinxstyleliteralstrong{\sphinxupquote{run\_type}} (\sphinxstyleliteralemphasis{\sphinxupquote{Literal}}\sphinxstyleliteralemphasis{\sphinxupquote{{[}}}\sphinxstyleliteralemphasis{\sphinxupquote{\textquotesingle{}independent\textquotesingle{}}}\sphinxstyleliteralemphasis{\sphinxupquote{, }}\sphinxstyleliteralemphasis{\sphinxupquote{\textquotesingle{}together\textquotesingle{}}}\sphinxstyleliteralemphasis{\sphinxupquote{{]}}})

\end{itemize}

\sphinxlineitem{Raises}
\sphinxAtStartPar
\sphinxstyleliteralstrong{\sphinxupquote{ValueError}} \textendash{} If no animation parameters are defined in the scene.

\end{description}\end{quote}

\end{fulllineitems}


\end{fulllineitems}

\end{quote}


\renewcommand{\indexname}{Python Module Index}
\begin{sphinxtheindex}
\let\bigletter\sphinxstyleindexlettergroup
\bigletter{p}
\item\relax\sphinxstyleindexentry{pygeox.constraint}\sphinxstyleindexpageref{pygeox:\detokenize{module-pygeox.constraint}}
\end{sphinxtheindex}

\renewcommand{\indexname}{Index}
\printindex
\end{document}